\documentclass[letterpaper,twocolumn,openany,nodeprecatedcode,hidelinks]{dndbook}
\usepackage[english]{babel}
\usepackage[utf8]{inputenc}
\usepackage{pdfpages}
\usepackage{graphicx}
\usepackage{wrapfig}
\usepackage{color}
\usepackage{tikz}
\usepackage{lipsum}
\usepackage{hyperref}
\usepackage{titling}
\usepackage{pagecolor}
\usepackage{afterpage}
\usepackage{titletoc}
\usepackage{everyshi}
\usepackage{changepage}
\usepackage{geometry}
\usepackage{xstring}
\usepackage{afterpage}
\usepackage{multicol}

\usepackage{tabularray}
\usepackage{amsmath}
\usepackage{amsfonts}
\usepackage{enumitem}
\usepackage{tipa}
\usepackage{tablefootnote}
\usepackage{linguex}
\usepackage{expex}

\usepackage{makecell}

\usepackage{float}

\renewcommand\glt{} % removes extra vertical space from translation line
\newcommand{\redcolor}[1]{\color{red}#1\normalcolor}

% \catcode`_=\active
% \def_#1{\ifmmode\sb{#1}\else$\sb{#1}$\fi}


\usetikzlibrary{calc}

% \usepackage[unicodeparmaite]{tengwarscript}
% \usepackage{tengwarscript}
% \usepackage[annataritalic]{tengwarscript}
% \pdfmapfile{=tengwarscript.map}


\usepackage{fontspec}
% \newfontface{\tovian}{Ancient-Tovian.ttf}
\newfontface{\tovian}{Tovian.ttf}
\newfontface{\tengwar}{tngani.ttf}
\newcommand*{\red}[1]{\color{red}#1\normalcolor}

\def\vc#1{$\vcenter{\vspace{-30pt}\hbox{\vspace*{-50pt}#1}}$}

\newenvironment{nopagebreak}
  {\par\nobreak\vfil\penalty0\vfilneg
   \vtop\bgroup}
  {\par\xdef\tpd{\the\prevdepth}\egroup
   \prevdepth=\tpd}

\begin{document}
% \setmainfont{Ancient-Tovian.ttf}

% \tengwarannataritalic[1.5]



% IDEAS:
% Speak vs stand means performance versus factual
% to tree: stay in once place long time (good thing)


% Fethrdriel
\title{
    % \Huge\Tformen\Tthuule\TTdotbelow\Troomen\Tanna\Tando\TTdot\Ttelco\TTdot\Tlambe\TTdotbelow\\
    % {\fontsize{60pt}{30pt}\tovian ««ishmahluorniara»»}\\ % ˈiʃmahɫuorɳiara
    {\fontsize{80pt}{30pt}\tovian «««fethrydriel»»»}\\ % ˈapɸeθjdrieɫ
    % {\Huge\tengwar fethrydriel}\\
    \vspace{1in}
    \large \textit{Beginner's Guide to Tovian} \\
    % \Tando\Tsilme\TTacute\Ttinco\TTdoubler\TTthreedots\Tyanta
     % \Tquesse\TTacute\Ttinco\TTdoubler\TTthreedots\\
     % \Tquesse\TTacute\Ttinco\TTdoubler\TTthreedots\\
    % \huge{Tovis}\\
    % \begin{tikzpicture}[remember picture]\\%
        % \node (main) {\includegraphics[width=7in]{img/tovis2.png}};\\%
    % \end{tikzpicture}
    % \vspace{-1in}
    }

% \author{\large Jackson Dean}
\date{}

\newcommand{\continue}[0]{
% \textit{continued on the next page...}
\newpage
}

\newcommand\blfootnote[1]{%
  \begingroup
  \renewcommand\thefootnote{}\footnotetext{#1}%
  %\addtocounter{footnote}{-1}% not needed
  \endgroup
}
\titlespacing*{\subsubsection}
{0pt}{0pt}{0pt}

\newcommand{\ra}[0]{*}
\newcommand{\linkplaceholder}[0]{\subsubsubsection{}}

\newcounter{tablectr}

\newcommand*{\itemrow}[1]{
    \stepcounter{tablectr}%
    \StrSubstitute{#1}{ }{}[\temp] % Replace spaces in the name with nothing
    \StrSubstitute{\temp}{'}{}[\templabel] % Replace apostrophe
    \StrSubstitute{\templabel}{,}{}[\templabel] % Replace apostrophe
    \label{\templabel} 
    \arabic{tablectr}\subsubsubsection{}\stepcounter{tablectr}-\arabic{tablectr} & \textbf{#1}
}

\newcommand{\rarity}[3]{%
    \small{\textit{
    \newline #1, #2 %
    \ifbool{#3}{(*)}%
    }}
}

\newcolumntype{B}{X}
\newcolumntype{s}{>{\hsize=0.20\hsize}X}





\newcommand{\multicolinterrupt}[2]{% Stuff to span both rows
\setcounter{tempcolnum}{\col@number}
\end{multicols}
#1%
\begin{multicols}{#2}
}

% \newcommand{\todo}[1]{\color{red}#1\normalcolor}
\newcommand{\todo}[1]{\textit{#1}}

\renewcommand{\chaptermark}[1]{\markboth{#1}{}}

\newcommand{\hanging}[0]{\hspace{0.5in}}


\setcounter{page}{1}
\maketitle



\newcommand{\CLan}[0]{la} % animate class
\newcommand{\CLin}[0]{lo} % inanimate class
\newcommand{\CLab}[0]{le} % abstract class

% cases
\newcommand{\CAnom}[0]{i} % nominative case (default, often omitted)
\newcommand{\CAacc}[0]{yu} % accusative case
\newcommand{\CAgen}[0]{hi} % genitive case
\newcommand{\CAdat}[0]{mi} % dative case
\newcommand{\CAloc}[0]{ti} % locative case (place)
\newcommand{\CAtem}[0]{tu} % temporal case (time)
\newcommand{\CAins}[0]{si} % instrumental case
\newcommand{\CApur}[0]{lhu} % purpose case
\newcommand{\CApri}[0]{ni} % privitive case
\newcommand{\CAcom}[0]{yi} % comlative case
\newcommand{\CAexc}[0]{fi} % excessive case
\newcommand{\CAcomp}[0]{asl} % comparative case
\newcommand{\CAade}[0]{asl} % adessive case
\newcommand{\CAall}[0]{su} % allative case
\newcommand{\CAess}[0]{pu} % allative case


\newcommand{\CAabl}[0]{di} 

\newcommand{\casestable}[0]{
\subsection*{Noun Classes}
\begin{itemize}
    \item \textbf{Animate (living beings, includes plants)}
    \begin{itemize}
        \item \textbf{Prefix:} \textit{\CLan} 
        \item \textbf{Example:} \textit{\CLan-driel} (person)
    \end{itemize}
    \item \textbf{Inanimate (objects and places)}
    \begin{itemize}
        \item \textbf{Prefix:} \textit{\CLin-}
        \item \textbf{Example:} \textit{\CLin-pelor} (city)
    \end{itemize}
    \item \textbf{Abstract (concepts and ideas)}
    \begin{itemize}
        \item \textbf{Prefix:} \textit{\CLab-}
        \item \textbf{Example:} \textit{\CLab-ara} (wisdom)
    \end{itemize}
\end{itemize}

\subsection*{Noun Cases}
\begin{itemize}
    \item \textbf{Nominative (subject)}
    \begin{itemize}
        \item \textbf{Prefix:} \textit{\CAnom-} (default case, often omitted)
        \item \textbf{Example:} \textit{\CAnom-\CLan-driel} (the person as the subject)
    \end{itemize}
    \item \textbf{Accusative (direct object)}
    \begin{itemize}
        \item \textbf{Prefix:} \textit{\CAacc-}
        \item \textbf{Example:} \textit{\CAacc-\CLan-driel} (the person as the direct object)
        \item \textbf{Sentence}: \textit{\CAacc-\CLin-tel fa mrila \CAnom-\CLan-driel} (the person sees the lake)
    \end{itemize}
    \item \textbf{Genitive (possession)}
    \begin{itemize}
        \item \textbf{Prefix:} \textit{\CAgen-}
        \item \textbf{Example:} \textit{\CAgen-\CLan-driel} (of the person)
    \end{itemize}
    \item \textbf{Dative (indirect object)}
    \begin{itemize}
        \item \textbf{Prefix:} \textit{\CAdat-}
        \item \textbf{Example:} \textit{\CAdat-\CLan-driel} (to/for the person)
    \end{itemize}
    \item \textbf{Locative (location)}
    \begin{itemize}
        \item \textbf{Prefix:} \textit{\CAloc-}
        \item \textbf{Example:} \textit{\CAloc-\CLin-pelor} (in the city)
    \end{itemize}
        \item \textbf{Temporal (time)}
    \begin{itemize}
        \item \textbf{Prefix:} \textit{\CAtem-}
        \item \textbf{Example:} \textit{\CAtem-\CLab-loe} (at two o'clock)
    \end{itemize}
    \item \textbf{Instrumental (means or tool)}
    \begin{itemize}
        \item \textbf{Prefix:} \textit{\CAins-}
        \item \textbf{Example:} \textit{\CAins-\CLin-rith} (with the air)
    \end{itemize}
        \item \textbf{Purpose (goal or reason)}
    \begin{itemize}
        \item \textbf{Prefix:} \textit{\CApur-}
        \item \textbf{Example:} \textit{\CApur-\CLab-aray} (for the purpose of wisdom)
    \end{itemize}

    \item \textbf{Ablative (away)}
    \begin{itemize}
        \item \textbf{Prefix:} \textit{\CAabl-}
        \item \textbf{Example:} \textit{\CAabl-\CLab-loe} (before 2)
    \end{itemize}

     \item \textbf{Comlative (together)}
    \begin{itemize}
        \item \textbf{Prefix:} \textit{\CAcom-}
        \item \textbf{Example:} \textit{\CAcom-\CLan-ta} (together with him/her/them)

             \item \textbf{Allative (towards)}
    \begin{itemize}
        \item \textbf{Prefix:} \textit{\CAall-}
        \item \textbf{Example:} \textit{\CAall-\CLan-tal} (towards the east)
    \end{itemize}
        

     \item \textbf{Essive (time or state)}
    \begin{itemize}
        \item \textbf{Prefix:} \textit{\CAess-}
        \item \textbf{Example:} \textit{\CAess-\CLan-ninadriel} (as a child)
    \end{itemize}


    % \end{itemize}
       % \item \textbf{Excessive (too)}
    % \begin{itemize}
        % \item \textbf{Prefix:} \textit{\CAreq-}
        % \item \textbf{Example:} \textit{\CAreq-\CLab-ara} (requires wisdom)
    % \end{itemize}
           % \item \textbf{Relative (relative clause)}
    % \begin{itemize}
        % \item \textbf{Prefix:} \textit{\CAreq-}
        % \item \textbf{Example:} \textit{\CAreq-\CLab-ara} (requires wisdom)
    % \end{itemize}
\end{itemize}


\subsection*{Combining Classes and Cases}
To use these prefixes effectively, the structure would be Case + Class + Noun.
\begin{itemize}
    \item \textbf{Example:} \textit{\CAacc-\Clan-driel} (the animate person as the direct object)
\end{itemize}

\subsection*{Examples}
\begin{itemize}
    \item \textbf{The person (subject)}: \newline \textit{\CLan-driel} or \textit{\CAnom-\CLan-driel}
    \item \textbf{The person (direct object)}: \newline \textit{\CAacc-\CLan-driel}
    \item \textbf{Of the person (possession)}: \newline \textit{\CAgen-\CLan-driel}
    \item \textbf{To/for the person (indirect object)}:\newline  \textit{\CAdat-\CLan-driel}
    \item \textbf{In the city (location)}:\newline  \textit{\CAloc-\CLin-pelor}
    \item \textbf{With the air (instrumental)}:\newline  \textit{\CAins-\CLin-rith}
        \item \textbf{With the air (comlative)}:\newline  \textit{\CAcom-\CLin-rith}
\end{itemize}

}



\onecolumn

\section*{\centering Alphabet}

\medskip

\begin{tabular}{l r}
{\large
\begin{tblr}{|Q[c,m] Q[l,m] Q[c,m]|}
\hline
\textbf{Romanization} & \textbf{Pronunciation} & \textbf{Symbol} \\
\hline
A& {\huge /a/}&{\fontsize{60pt}{10pt}\tovian a} \\
\hline
B& {\huge /b/}&{\fontsize{60pt}{10pt}\tovian b}\vspace{-0pt}\\
\hline
D& {\huge /d/}&{\fontsize{60pt}{10pt}\tovian d} \\ % could be \textrtaild 
\hline
E& {\huge /e/}&{\fontsize{60pt}{10pt}\tovian e} \vspace{-30pt}\\
\hline

F& {\huge /f/}&{\fontsize{60pt}{10pt}\tovian f} \vspace{-45pt}\\
\hline
G& {\huge /g/}&{\fontsize{60pt}{10pt}\tovian g}\\
\hline
H & {\huge /h/ }&{\fontsize{60pt}{10pt}\tovian h}\\
\hline
I& {\huge /i/}&{\fontsize{60pt}{10pt}\tovian i}\\
\hline
K& {\huge /k/}&{\fontsize{60pt}{10pt}\tovian k}\vspace{-10pt}\\

\hline
LH & {\huge /\textbeltl/ }&{\fontsize{60pt}{10pt}\tovian lh}\\\hline % I really wanted this to be \textbeltl, but could also be textsubring{l} or 
L & {\huge /l/ }&{\fontsize{60pt}{10pt}\tovian l}\\
\hline
M & {\huge /m/ }&{\fontsize{60pt}{10pt}\tovian m}\vspace{-20pt}\\
\hline
N& {\huge /n/}&{\fontsize{60pt}{10pt}\tovian n}\vspace{-20pt}\\
\hline
O& {\huge /o/}&{\fontsize{60pt}{10pt}\tovian o} \vspace{-20pt}\\
\hline
\end{tblr}
}
&
{\large
\begin{tblr}{|Q[c,m] Q[c,m] Q[c,m]|}
\hline
\textbf{Romanization} &\textbf{Pronunciation} & \textbf{Symbol} \\
\hline
P& {\huge/p/}&{\fontsize{60pt}{10pt}\tovian p}\\
\hline
R& {\huge /r/}&{\fontsize{60pt}{10pt}\tovian r}\vspace{-35pt}\\
\hline
S& {\huge /s/}&{\fontsize{60pt}{10pt}\tovian s} \vspace{-30pt}\\
\hline
T& {\huge /\textsubbridge{t}/}&{\fontsize{60pt}{10pt}\tovian t}\\
\hline
U & {\huge /u/  }&{\fontsize{60pt}{10pt}\tovian u} \vspace{-20pt} \\
\hline
V& {\huge /v/}&{\fontsize{60pt}{10pt}\tovian v} \vspace{-25pt}\\
\hline
W& {\huge /\textturnw/}&{\fontsize{60pt}{10pt}\tovian w} \\
\hline
X& {\huge /ks/}&{\fontsize{60pt}{10pt}\tovian x}\vspace{-30pt}\\
\hline
Y & {\huge /j/ }&{\fontsize{60pt}{10pt}\tovian y} \vspace{-20pt}\\
\hline
\:EA& {\huge /e.a/ }&{\fontsize{60pt}{10pt}\tovian ea}\\
\hline
\:SS & {\huge /\textipa{s:}/ }&{\fontsize{60pt}{10pt}\tovian ss}\vspace{-20pt}\\
\hline
SH & {\huge /{\textesh}/ }&{\fontsize{60pt}{10pt}\tovian sh} \\
\hline
TH& {\huge /\texttheta/} %ð?
&  {\fontsize{60pt}{10pt}\tovian th} \\
\hline

\end{tblr}
}

\end{tabular}

\newpage

\subsection*{vowels}

\hspace{1pt}

{\huge
\begin{table}[h]
    \centering
    \begin{tblr}{Q[c,m] Q[c,m] Q[c,m]}
         Romanized Vowel & Pronunciation & Glyph\\
         \hline
         {\huge A}& {\huge /a/}&{\fontsize{40pt}{10pt}\tovian a} \\\hline
         {\huge A}& {\huge /\textsubumlaut{a}/}&{\fontsize{40pt}{10pt}\tovian ah} \\\hline
         {\huge E}& {\huge /e/}&{\fontsize{40pt}{10pt}\tovian e}\\
         \hline
         {\huge I}& {\huge /i/}&{\fontsize{40pt}{10pt}\tovian i}\\
         \hline
         {\huge O}& {\huge /o/}&{\fontsize{40pt}{10pt}\tovian o}\\
         \hline
         {\huge U} & {\huge /u/}&{\fontsize{40pt}{10pt}\tovian u}\\
         
    \end{tblr}
    \label{tab:my_label}
\end{table}
}

\medskip
\begin{figure}[h]
    \centering

\begin{tikzpicture}[vowel/.style={fill, circle, inner sep=0pt, text height=1.25ex}]
\coordinate (hf) at (0,3); % the high front vertex
\coordinate (hb) at (4,3); % the high back vertex
\coordinate (lb) at (4,0); % the low back vertex
\coordinate (lf) at (2,0); % the low front vertex

\draw (hf) -- (hb) -- (lb) -- (lf) -- cycle; % draws the trapezoid

\draw ($(hf)!1/3!(lf)$) -- ($(hb)!1/3!(lb)$); % the high-mid line
\draw ($(hf)!2/3!(lf)$) -- ($(hb)!2/3!(lb)$); % the low-mid line
\draw ($(hf)!0.5!(hb)$) -- ($(lf)!0.5!(lb)$); % the center line

\node[vowel,label={[label distance=-1pt]below right:{\fontsize{30pt}{0pt}\tovian i}}] at (0.6,2.8) {}; 
\node[vowel,label={[label distance=-3pt]below left:{\fontsize{30pt}{0pt}\tovian o}}] at (3.7,1.6) {};
\node[vowel,label={[label distance=-3pt]below left:{\fontsize{30pt}{0pt}\tovian u}}] at (3.6,2.8) {};
\node[vowel,label={[label distance=-3pt]below right:{\fontsize{30pt}{0pt}\tovian e}}] at (1.2,1.6) {};
\node[vowel,label={[label distance=-3pt]above left:{\fontsize{30pt}{0pt}\tovian a}}] at (3.25,0.2) {};
\end{tikzpicture}
\end{figure}



% example oradriel: 
    % \textipa{'ɔrad͡ʒriɛlɛ}
% hlumelore ʎ̤uːmɛloːɹeɪ

\twocolumn

\section*{Consonant and Vowel Inventory}
\begin{itemize}
    \item \textbf{Consonants}: \textipa{/b/, /d/, f, 
    % /\textphi/,
    /g/, /h/, /k/, /\textbeltl/, /l/, /m/, /\ng/, /n/, /p/, /r/, /s/, /\textsubbridge{t}/, /v/, /w/, /j/, /\textesh/, /\texttheta/}
    \item \textbf{Vowels}: \textipa{/a/, /\textsubumlaut{a}/, /e/, /i/, /o/, /u/}
    \item \textbf{Dipthongs}: \textipa{/ai/, /uj/}
\end{itemize} 


% \subsection*{Rules for /f/ and /\textphi/}

% \textbf{/f/ (Voiceless Labiodental Fricative)}:

% At the beginning of words:
% \begin{itemize}
% \item \textipa{/falor/} (leader)
% \item \textipa{/fethr/} (speaker)
% \end{itemize}
% Between vowels:
% \begin{itemize}
% \item /afa/ (...)
% \item /ujfae/ (...)
% \end{itemize}

% When adjacent to voiced consonants:
% \begin{itemize}
% \item ...
% \end{itemize}
% \textbf{/\textphi/ (Voiceless Bilabial Fricative)}:

% At the end of words or plurals:
% \begin{itemize}
% \item \textipa{/te\textphi/} (fortress)
% \item \textipa{/te\textphi e/} (fortresses)
% \end{itemize}


% In newer or borrowed words to indicate foreign or technical terms:
% \begin{itemize}
% \item ... (machine, from gnomish)
% \end{itemize}

% \subsection*{Rules for /l/ and /\textbeltl/}

% ...

\section*{Syllable Structure}
\begin{itemize}
    \item \textbf{Basic Structure}: Each syllable should follow the general structure of (C)(C)V(C)(C), where:
    \begin{itemize}
        \item \textbf{C} = Consonant
        \item \textbf{V} = Vowel
    \end{itemize}
    \item \textbf{Consonant Clusters}:
    \begin{itemize}
        \item At the beginning of a syllable, allow a maximum of two consonants.
        \item At the end of a syllable, allow a single consonant or specific consonant clusters.
    \end{itemize}
\end{itemize}

\section*{Syllable Patterns}
\redcolor{TODO: CONFIRM WITH DICTIONARY}
\begin{enumerate}
    \item \textbf{V (Vowel Only)}: Rare but can occur in specific cases, especially with prefix and suffix modifications.
    \begin{itemize}
        \item \textipa{/a/, /e/, /i/}
    \end{itemize}
    \item \textbf{CV (Consonant + Vowel)}: The most common syllable pattern.
    \begin{itemize}
        \item Example: \textipa{/la/, /me/, /ti/}
    \end{itemize}
    \item \textbf{CVC (Consonant + Vowel + Consonant)}: Also common, adding a consonant to the end.
    \begin{itemize}
        \item Example: \textipa{/lan/, /mujl/, /tir/}
    \end{itemize}
    \item \textbf{CCV (Consonant Cluster + Vowel)}: Allowed at the beginning of words.
    \begin{itemize}
        \item Example: \textipa{/tla/, /fre/, /tru/}
    \end{itemize}
     \item \textbf{CCVC (Consonant Cluster + Vowel + Consonant)}: Complex but permissible.
    \begin{itemize}
        \item Example: \textipa{/tl\textsubumlaut{a}n/, /frel/, /trus/}
    \end{itemize}
    \item \textbf{CVCC (Consonant + Vowel + Consonant Cluster)}: Possible only at the end of words or as whole words, and only including /\texttheta/ or /\textesh/.
    \begin{itemize}
        \item Example: \textipa{/fe\texttheta r/, /pa\texttheta r/, /lor\textesh/}
    \end{itemize}
\end{enumerate}

\section*{Stress Patterns}
\begin{enumerate}
    \item \textbf{Primary Stress}: On the penultimate syllable.
    \begin{itemize}
        \item Example: /la-\textbf{dri}-el/, /lo-\textbf{pe}-lor/
    \end{itemize}
    \item \textbf{Secondary Stress}: On the third to last syllable if the word has more than three.
    \begin{itemize}
        \item Example: /ti-\textbf{la}-dri-el/, /ti-\textbf{lo}-pe-lor/, /mel-or-\textbf{e}-ra-e/
    \end{itemize}
\end{enumerate}

\section*{Diphthongs and Vowel Clusters}
Clearly enunciate most vowel clusters to maintain the phonetic distinctiveness and breathy quality.
\begin{itemize}
\item \textbf{Diphthongs}: Diphthongs are combinations of two vowel sounds within the same syllable, gliding smoothly from one vowel to the other. The primary diphthongs in are:
\begin{itemize}
\item \textipa{/ai/}: Pronounced as in "eye."
\item \textipa{/uj/}: Pronounced as in "buoy."
\end{itemize}
\item \textbf{Vowel Clusters}: When two vowels appear together but do not form a diphthong, they must be clearly enunciated to distinguish them as separate sounds.
\begin{itemize}
\item Example: \textipa{/ea/} similar to "idea", where both vowels are pronounced distinctly.
\end{itemize}
\end{itemize}

\section*{Examples}
\begin{description}
        \item /la-\textbf{dri}-el ya-\textbf{lo}-pe-lor \textbf{lo}-re \textbf{em}-ril \textbf{si}-ne\textesh/ 
        \hanging The person quickly sees the city.
\end{description}





Lefethr is an \textit{agglutanative} language, meaning morphemes do not change form


\section*{Plurals}
\textit{-e} indicates plurality. Verbs do not agree with number. Singular nouns therefore never end in \textit{-e}.

\begin{center}
\begin{tabular}{|c c|}
     \hline

    \textit{lothr} &  river \\
     \textit{lothre} & rivers \\

     \hline
     \hline

     \textit{nor} & sea \\

     \textit{nore} & seas \\
     \hline
     
\end{tabular}
\end{center}

\section*{Verbs}

Verbs are constructed around a consonant root that denotes its aspect. This root is combined with a suffix to indicate tense. The tense suffixes are \textit{-e} for past,\textit{ -a} for present, and \textit{-o} for future. Optional prefixes can be added to denote mood and voice. This final conjugated auxiliary verb is placed in front of the infinitive. The verb that indicates the action is never conjugated, only the auxiliary verb changes form to show aspect, tense, mood, and voice.

\begin{table}[H]
    \centering
    \begin{tabular}{|c|c|}
        \hline
        \textbf{Tense} & \textbf{Auxiliary Verb Suffix} \\
        \hline
        Past & -e \\
        Present & -a \\
        Future & -o \\
        \hline
    \end{tabular}
    % \caption{Person Suffixes in Tovian}
    \label{tab:verb-tense}
\end{table}


For example, to say "the person is speaking", start with the verb ``to speak": \textit{fethr}. Use the continuous aspect consonant\textit{sh} (from the verb \textit{shotl}, ``to flow") to form the auxiliary verb \textit{sha}. The auxilary verb goes before the infinitive, resulting in \textit{sha fethr}, ``is speaking". \textit{Driel} is ``the person", resulting in the final sentence: \textit{driel sha fethr}, ``the person is speaking". A very literal translation could be ``The person flows speaking".

\begin{table}[h!]
    \centering
    \subsubsection{Aspects with present tense examples}
    \begin{tabular}{c|c|c|c}
    \hline
         \textbf{Aspect}& \makecell{\textbf{Root}\\\textbf{Verb}}& \makecell{\textbf{Aux.}\\\textbf{root}} & \makecell{\textbf{Pres. tense}\\\textbf{Example}} \\
\hline
         Simple&\makecell{\textit{fethr} \\ speak}  & f- & \makecell{\textit{fa mrila}\\``sees"} \\ \hline
         \red{Imperfective}&\makecell{\textit{thraf} \\ wonder}  & th- & \makecell{\textit{tha mrila}\\ ``is seeing"} \\ \hline
         Perfect&\makecell{\textit{lhara} \\ know}  & lh- & \makecell{\textit{lha mrila} \\ ``has seen"} \\ \hline
         Near&\makecell{\textit{tlun} \\ share}& tl- & \makecell{\textit{tla mrila} \\ ``is about to see" }\\ \hline
         Immediate&\makecell{\textit{kesh} \\ burn}& k- & \makecell{\textit{ka mrila} \\ ``is seeing right now" }\\ \hline
         Habitual&\makecell{\textit{mel} \\ live}  & m- & \makecell{\textit{ma mrila} \\``sees regularly" }\\ \hline
         Progressive&\makecell{\textit{sethr} \\ run}& s- & \makecell{\textit{sa mrila} \\``is seeing now" }\\ \hline
         Continuous&\makecell{\textit{shotl} \\ flow}& sh- & \makecell{\textit{sha mrila} \\ ``is seeing" }\\ \hline
         Iterative&\makecell{\textit{nalh} \\ think}& n- & \makecell{\textit{na mrila} \\ ``sees repeatedly" }\\ \hline
         Inceptive&\makecell{\textit{yalor} \\ begin}& y- & \makecell{\textit{ya mrila} \\ ``begins  seeing" }\\ \hline
         Cessative&\makecell{\textit{panor} \\ end}& p- & \makecell{\textit{pa mrila} \\ ``stops seeing" }\\ \hline
         Remote&\makecell{\textit{huysh} \\ throw}  & h- & \makecell{\textit{ha mrila} \\``sees (distant)" \footnotemark[1]} \\
     \end{tabular}
     \\
     \medskip
     \raggedright
     {\footnotesize \footnotemark[1] \raggedleft In past and future tenses, the remote aspect refers to actions that took place long ago or far in the future, but the meaning in present tense is more nuanced. Elves tend to use the rare present remote aspect as an insult, for example \textit{wele ha fethr}: ``they speak but it is irrelevant"}
    \label{tab:aspects}
\end{table}

\medskip


\subsection*{Mood and Voice}

% \subsection*{All verb forms}
\begin{tabular}{rl}
&\textbf{Mood}\\
\textbf{unmarked}& $\rightarrow$ indicative \\
\textbf{wi-}& $\rightarrow$ subjunctive \\
\textbf{se-}& $\rightarrow$ imperative \\
\textbf{ko-} & $\rightarrow$ conditional \\
\textbf{xo-}& $\rightarrow$ counterfactual \\
\textbf{hle-}& $\rightarrow$ optative \\
\textbf{tlo-}& $\rightarrow$ obligative \\
\textbf{ef-}& $\rightarrow$ necessative \\
\textbf{sho-}& $\rightarrow$ potential \\



&\textbf{Voices}\\
\textbf{unmarked}& $\rightarrow$ active \\
\textbf{te-}& $\rightarrow$ reflexive \\
\textbf{pa-}& $\rightarrow$ passive \\
\textbf{mo-}& $\rightarrow$ middle \\
\textbf{ke-}& $\rightarrow$ causative \\
\textbf{ra-}& $\rightarrow$ reciprocal \\

% &\textbf{Other (not combined)}\\
% \textbf{li-}& $\rightarrow$ gerund \\

\end{tabular}

\medskip

Mood and voice can be combined by adding the mood prefix before the voice prefix. For example, adding the obligative mood marker \textit{tlo-} and the reflexive voice marker \textit{te-} to the habitual aspect auxilary root stem \textit{m-} with the present tense conjugation \textit{-a}, the final auxilary verb is \textit{tlotema}. 

\ex
\begingl
\gla idriel tlotema lorsel //
\glb NOM-person OBLIG-REFL-HAB.PRES wash-INF //
\glft "She should habitually wash herself." //
\endgl
\xe



\subsection{Unused Combinations}
Tovian is very flexible when it comes to combining different tenses, aspects, moods, and voices. However, certain verb conjugations result in combinations that are unlikely or never used due to their impractical or nonsensical nature. For example, the Past Imperative is very rare because commands most often refer to actions in the present or future.



\section{Copulas}
Copulas are formed by combining verbs and nouns to describe states and qualities. For example the word "mel" (life/living) can be used as a copula: "i-la-driel fa mela" means "The person is alive," or equally "The person lives" This method applies to various attributes by changing the verb root and noun. For instance, "driel fe araya" translates to "The person was wise." 

 \section{Comparative Constructions}

    Tovian uses comparative constructions to express excessiveness, for example: "bigger than appropriate." The ablative case is often used for this comparison: "the tree is too big" could be expressed as \textit{ilo-glar \CAabl-le-velar fa mota} ("The tree is big away-from-worthiness" or simply "The tree is bigger than what is worthy/right").

\section{Adjectives}

Tovian does not use a separate class of adjectives. Instead, all descriptive qualities are expressed through nouns. These attributive nouns follow the noun they modify in a fixed noun–noun structure. For example, the phrase \textit{driel aray} translates literally as “person wisdom,” conveying the meaning “the wise person.” To convey "people wisdom" or "the wisdom of people," one would switch the order to: \textit{aray driel}. This structure does not need inflection or agreement between adjectives and nouns, for example "the wise people" would simply be \textit{driele aray}. In cases requiring emphasis or clarification, a copular construction can be used instead: \textit{i-la-driel fa araya} (“The person is wise”). Many commonly used descriptive terms—such as colors, emotional states, or physical traits—exist as standalone nouns and follow the same syntactic pattern. This noun-only system reinforces the conceptual unity of properties and entities, treating attributes not as modifiers but as qualities with independent substance.


\section{Agreement}

Suffixes are added to verbs to indicate the person of the subject.

\begin{table}[H]
    \centering
    \begin{tabular}{|c|c|}
        \hline
        \textbf{Person} & \textbf{Suffix} \\
        \hline
        1st person & i \\
        2nd person & o \\
        3rd person & a \\
        \hline
    \end{tabular}
    % \caption{Person Suffixes in Tovian}
    \label{tab:suffixes}
\end{table}


\begin{table}[H]
    \centering
    \begin{tabular}{|c|c|c|}
        \hline
        \textbf{Phrase} & \textbf{Translation} & \textbf{Person} \\
        \hline
        ilane sa nalh\textbf{i} & "We think" & 1st \\
        ilana sa nalh\textbf{i} & "I think" & 1st \\
        ilawa sa nalh\textbf{o} & "You think" & 2nd \\
        ilawe sa nalh\textbf{o} & "You (plural) think" & 2nd \\
        ilata sa nalh\textbf{a} & "He/she/they think" & 3rd \\
        ilate sa nalh\textbf{a} & "They (plural) think" & 3rd \\
        \hline
    \end{tabular}
    % \caption{Agreement in Tovian}
    \label{tab:agreement}
\end{table}

%  Verb stems are added to nouns and always end in ``lore". Different tenses, aspects, moods, and voices can be applied by adding prefixes. The prefixes in the order they should be applied are as follows:\footnote{\noindent Complex verb forms are saved for formal communication.}

% \subsection{Simple Verb Forms}

% \begin{tabular}{rl}
% &\textbf{Tense}\\
% \textbf{} &$\rightarrow$$ present tense \\
% \textbf{vi-} &$\rightarrow$$ future tense \\
% \textbf{in-} &$\rightarrow$$ past tense \\

% &\textbf{Aspect}\\
% \textbf{de-}& $\rightarrow$ perfect \\
% \textbf{me-} & $\rightarrow$ participle \\
% \textbf{shi-}& $\rightarrow$ continuous \\

% &\textbf{Mood}\\
% \textbf{ko-} & $\rightarrow$ conditional \\
% \textbf{se-}& $\rightarrow$ imperative \\
% \textbf{wi-}& $\rightarrow$ subjunctive \\

% &\textbf{Voices}\\
% \textbf{te-}& $\rightarrow$ reflexive \\
% \textbf{pe-}& $\rightarrow$ passive \\

% &\textbf{Other (not combined)}\\
% \textbf{li-}& $\rightarrow$ gerund \\


% \end{tabular}
% % 55 combinations


% \subsection*{All verb forms}
% \begin{tabular}{rl}
% &\textbf{Tense}\\
% \textbf{} &$\rightarrow$$ present tense \\
% \textbf{vi-} &$\rightarrow$$ future tense \\
% \textbf{in-} &$\rightarrow$$ past tense \\
% &\textbf{Aspect}\\
% \textbf{de-}& $\rightarrow$ perfect \\
% \textbf{me-} & $\rightarrow$ participle \\
% \textbf{re-} & $\rightarrow$ habitual \\
% \textbf{ro-}& $\rightarrow$ iterative \\
% \textbf{ve-} & $\rightarrow$ immediate \\
% \textbf{shi-}& $\rightarrow$ continuous \\
% \textbf{wo-}& $\rightarrow$ remote \\
% \textbf{to-}& $\rightarrow$ near \\
% \textbf{yi-}& $\rightarrow$ inceptive \\
% \textbf{hlo-}& $\rightarrow$ cessative \\

% &\textbf{Mood}\\
% \textbf{ko-} & $\rightarrow$ conditional \\
% \textbf{se-}& $\rightarrow$ imperative \\
% \textbf{wi-}& $\rightarrow$ subjunctive \\
% \textbf{xo-}& $\rightarrow$ counterfactual \\

% &\textbf{Voices}\\
% \textbf{te-}& $\rightarrow$ reflexive \\
% \textbf{pa-}& $\rightarrow$ passive \\
% \textbf{mo-}& $\rightarrow$ middle \\
% \textbf{ke-}& $\rightarrow$ causative \\
% \textbf{ra-}& $\rightarrow$ reciprocal \\

% &\textbf{Other (not combined)}\\
% \textbf{li-}& $\rightarrow$ gerund \\


% \end{tabular}
% % 601 combinations

% \medskip

% \subsection*{Verb Conjugations}
% Using "speak" as an example: \textit{fethr}

% \subsubsection*{Tense}
% \begin{itemize}
%     \item \textbf{Present Tense:} \textit{} \\
%           \hanging Speak: \textit{lore-fethr}
          
%     \item \textbf{Future Tense:} \textit{vi-} \\
%           \hanging Will speak: \textit{vilore-fethr}
          
%     \item \textbf{Past Tense:} \textit{in-} \\
%           \hanging Spoke: \textit{inlore-fethr}
% \end{itemize}

% \subsubsection*{Aspect}
% \begin{itemize}
%     \item \textbf{Perfect:} \textit{de-} \\
%           \hanging Has spoken: \textit{delore-fethr}
          
%     \item \textbf{Continuous:} \textit{shi-} \\
%           \hanging Is speaking: \textit{shilore-fethr}
          
%     \item \textbf{Habitual:} \textit{re-} \\
%           \hanging Speaks regularly: \textit{relore-fethr}
          
%     \item \textbf{Immediate:} \textit{ve-} \\
%           \hanging Is about to speak: \textit{velore-fethr}
          
%     \item \textbf{Participle:} \textit{ma-} \\
%           \hanging Speaking (as a participle): \textit{malore-fethr}
% \end{itemize}

% \subsubsection*{Mood}
% \begin{itemize}
%     \item \textbf{Conditional:} \textit{ko-} \\
%           \hanging Would speak: \textit{kolore-fethr}
          
%     \item \textbf{Imperative:} \textit{se-} \\
%           \hanging Speak!: \textit{selore-fethr}
          
%     \item \textbf{Subjunctive:} \textit{wi-} \\
%           \hanging May speak: \textit{wilore-fethr}
          
%     \item \textbf{Counterfactual:} \textit{xo-} \\
%           \hanging Would have spoken: \textit{xolore-fethr}
% \end{itemize}

% \subsubsection*{Voice}
% \begin{itemize}
%     \item \textbf{Reflexive:} \textit{te-} \\
%           \hanging Speaks to oneself: \textit{telore-fethr}
          
%     \item \textbf{Passive:} \textit{pa-} \\
%           \hanging Is spoken: \textit{palore-fethr}
          
%     \item \textbf{Middle:} \textit{mo-} \\
%           \hanging Speaks for oneself: \textit{molore-fethr}
          
%     \item \textbf{Causative:} \textit{ke-} \\
%           \hanging Makes someone speak: \textit{kelore-fethr}
          
%     \item \textbf{Reciprocal:} \textit{ra-} \\
%           \hanging They speak to each other: \textit{ralore-fethr}
% \end{itemize}

% \subsubsection*{Other}
% \begin{itemize}
%     \item \textbf{Gerund:} \textit{li-} \\
%           \hanging Speaking (as a noun): \textit{lilore-fethr}
% \end{itemize}

% \subsubsection*{Combined Forms}
% \begin{itemize}
%     \item \textbf{Past Perfect:} \textit{de- + in-} \\
%           \hanging Had spoken: \textit{deinlore-fethr}
          
%     \item \textbf{Future Continuous:} \textit{shi- + vi-} \\
%           \hanging Will be speaking: \textit{shivilore-fethr}
          
%     \item \textbf{Conditional Habitual:} \textit{re- + ko-} \\
%           \hanging Would speak regularly: \textit{rekolore-fethr}
          
%     \item \textbf{Immediate Imperative:} \textit{ve- + se-} \\
%           \hanging Speak immediately!: \textit{veselore-fethr}
          
%     \item \textbf{Subjunctive Participle:} \textit{wi- + ma-} \\
%           \hanging Speaking in a hypothetical or wished-for sense: \textit{wimalore-fethr}
% \end{itemize}


% \subsection*{To be and to be like}
% All verbs with the exception of two follow the above structure. To express that the subject of the sentence is equivalent to the object, use ``'as" as a postfix on the object. To express similarity, use ``'asl" instead. For example:

% \begin{itemize}
%   \item \textbf{To be:} \\
%         \textit{Tel \CAacc-\CLin-helor'as.} (The lake is glass.) \\
%         - \textit{Tel} (subject: the lake) \\
%         - \textit{\CAacc-\CLin-helor'as} (object: glass with ``as" indicating "is")

%  \item \textbf{To be like:} \\
%         \textit{Tel \CAacc-\CLin-helor'asl.} (The lake is like glass.) \\
%         - \textit{Tel} (subject: the lake) \\
%         - \textit{helor'as} (object: glass with ``asl'" indicating "is like")

% \end{itemize}





\section*{Definiteness, Classes, and Cases}

Speakers usually on context to convey definiteness or indefiniteness.
In situations where there is ambiguity or in very formal communication, there are different prefixes:

\begin{table}[H]
    \centering
    \begin{tabular}{ccc}
         &  Subject& Non-subject\\
         Definite&  \textbf{i}-& \textbf{a}-\\
         Indefinite&  \textbf{e}-& \textbf{o}-\\
    \end{tabular}
\end{table}

{
\hspace{-30pt}
\begin{tabular}{l l}
\hline
    The sun shines in the sky  &   \textit{\textbf{i}-\CLin-efol fa ana \textbf{a}\CAloc-rithil}\\
    A sun shines in the sky   &  \textit{\textbf{e}-\CLin-efol fa ana \textbf{a}\CAloc-rithil}\\

        The sun shines in a sky  &   \textit{\textbf{i}-\CLin-efol fa ana \textbf{o}\CAloc-rithil}\\
    A sun shines in a sky   &  \textit{\textbf{e}-\CLin-efol fa ana \textbf{o}\CAloc-rithil}\\
\hline
     
\end{tabular}
}



\medskip

Nouns are categorized into three distinct \textbf{classes}, each identified by a unique prefix: \textit{\CLan-} for animate beings, \textit{\CLin-} for inanimate objects and places, and \textit{\CLab-} for abstract concepts and ideas. Nouns can also be inflected for grammatical \textbf{cases} to indicate their role within a sentence. Classes are not always used when the meaning is obvious, for example ``the person sees" would be written formally as \textit{\CLan-driel fa mrila}, but in informal communication, the class might be dropped to create \textit{driel fa mrila}, since the subject is obviously not a dead person (\textit{\CLin-driel}) or the abstract concept of people (\textit{\CLab-drie}l). When both cases and classes are required, cases are added before classes. For example \textit{\CLin-nume \textbf{\CApur}-\textbf{\CLab}-aray }is ``the paths for wisdom", where \textbf{\CApur} represents ``for", \textbf{\CLab} represents ``abstract concept of", and \textbf{ara} is ``wisdom." In formal or educational contexts, hyphens (romanization: ``-", Tovian: ``\vc{\fontsize{50pt}{-200pt}\tovian\textasciitilde-\textasciitilde}") are used to separate prefixes from their nouns. For example, in this document, a sentence could be written as \textit{\CAnom-\CLan-driel \CAacc-\CLin-tel fa mrila} (the person sees the lake), but in practice one might use \textit{driel \CAacc\CLin tel fa mrila}, dropping the hyphens, the nomative case, and the animate class.



\casestable

\subsection*{Pronouns}
\begin{tabular}{rl}
\textbf{na} & $\rightarrow$ me, I \\
\textbf{hina} & $\rightarrow$ my, mine \\
\textbf{ne} & $\rightarrow$ us or we \\
\textbf{hine} & $\rightarrow$ ours \\
\textbf{wa} & $\rightarrow$ you or your \\
\textbf{hiwa} & $\rightarrow$ yours \\
\textbf{we} & $\rightarrow$ you or (plural) \\
\textbf{hiwe} & $\rightarrow$ yours (plural) \\
\textbf{lata} & $\rightarrow$ them/she/he (animate)\\
\textbf{lahita} & $\rightarrow$ theirs/her/his \\
\textbf{late} & $\rightarrow$ them (plural, animate) \\
\textbf{hahite} & $\rightarrow$ theirs (plural) \\

\textbf{loda} & $\rightarrow$ it (inanimate)\\
\textbf{lohita} & $\rightarrow$ theirs \\
\textbf{lode} & $\rightarrow$ them (inanimate) \\
\textbf{lohite} & $\rightarrow$ theirs (plural, inanimate) \\

\end{tabular}



\subsection*{Word Order}
Sentences historically used a Subject-Object-Verb (SOV) word order. In this structure, the subject of the sentence comes first, followed by the object, and then the verb. For instance, "The person lives in the city" would be structured as "\CLan-driel \CAloc-\CLin-pelor fa mela", where "\CLan-driel" (the person) is the subject, "\CAloc-\CLin-pelor" (the city) is the direct (locative) object, and "fa mela" (lives) is the verb.


For emphasis, modern Tovian allows flexibility in word order. Elements can be fronted to the beginning of the sentence to add emphasis or focus. For example, to emphasize the city being lived in by the person, one might say "\CAloc-\CLin-pelor \CLan-driel fa mela," placing the object "\CAloc-\CLin-pelor" (the city) at the beginning of the sentence. This doesn't introduce ambiguity only because of the locative marker (\CAloc-) on the object of the sentence, which indicates it is never the subject despite coming first in word-order.

When changing word order, it is important to leave adjectives (remember, they are unmarked nouns) directly after the noun they modify. For example, "the person lives in the big city" would be either "\CLan-driel \CAloc-\CLin-pelor mot fa mela" or "\CAloc-\CLin-pelor mot \CLan-driel fa mela." Separating "mot" (big) and "pelor" (city) changes the meaning of the sentence. For example: "\CAloc-\CLin-pelor \CLan-driel mot fa mela" would be "the big person lives in the city." 


\subsection*{Questions}
\subsubsection{Yes/No questions}
To turn a declarative statement into a yes or no question, \textit{lhan} is added to the front of the sentence.

\textbf{Example: } \textit{lhan ila-we yalo-ta fa mrila?}

``Do you see it?"

\textbf{Response:} \textit{ila-na yalo-ta fa mrila}

``I see it"

In response to a question, declarative statements are sometimes simplified to just the conjugated auxilary + infinative, for example \textit{na yalo-ta fa mrila} could be shortened to \textit{fa mrila}, literally: ``sees", but the subject and object are inferred to be the same as the question since they are left out of the response. 

\medskip

\subsubsection{WH- questions}
 To form a non-yes/no question, add the question marker \textit{lhan} to the beginning of the sentence with a case-marking prefix to specify what question is being asked. 
 Examples:
     \begin{tabular}{c c}
        \textbf{ti}-lhan ta fe fethra & ``Where did he speak" \\
        \textbf{to}-lhan ta fe fethra & ``when did he speak?" \\
        \textbf{si}-lhan ta fe fethra & ``how did he speak?" \\
        \textbf{ya}-lhan ta fe fethra & ``to what/whom did he speak?" \\
          & 
     \end{tabular}

\subsubsection{Example conversation}

\noindent A: \textit{lhan driel se kwilora tolon?} % was the person dancing yesterday? 

\noindent B: \textit{ilata se nikwilora tolon, shan ilata lho kwilora ditom hitirlhon. lhan ilawa yata fo mrilo?} % she was not dancing yesterday, but she will have danced by the end of the week. Will you see her?

\noindent A: \textit{na yata fo mrili, tolhan ta fo kwilora?} % I am going to see her, what time will she dance?

\noindent B: \textit{tolerep tohethn} % At five o'clock tomorrow

{\tovian
\noindent lhan driel se kwilora tolon

\noindent ilata se nikwilora tolon, shan ilata lho kwilora ditom hitirlhon. lhan ilawa yata fo mrilo

\noindent na yata fo mrili, tolhan ta fo kwilora

\noindent torep tohethn
}


\ex
\begingl
\gla lhan driel se kwilora to-lon? //
\glb Q person PROG.PAST dance.1SG LOC-yesterday // 
\glft "Was the person dancing yesterday?" //
\endgl
\xe

\ex
\begingl
\gla ilata se nikwilora tolon, shan ilata lho kwilora ditom hitirlhon. lhan ilawa yata fo mrilo? //
\glb SBJ.ANIM.3SG PROG.PAST NEG-dance.1SG LOC-yesterday, but 3SG PERF dance.1SG ABL-end GEN-week. Q 2SG ACC-3SG FUT see.2SG? //
\glft "She was not dancing yesterday, but she will have danced by the end of the week. Will you see her?" //
\endgl
\xe

\ex
\begingl
\gla ila-na ya-ta fo mrili, to-lhan ta fo kwilora? //
\glb 1SG ACC-3SG FUT see.1SG, when-Q 3SG FUT dance.1SG? //
\glft "I will see her, what time will she dance?" //
\endgl
\xe
Sentence D

\ex
\begingl
\gla tole-rep to-hethn//
\glb LOC-five LOC-tomorrow//
\glft "At five o'clock tomorrow." //
\endgl
\xe


\newpage


% \section{Conjunctions}

% \subsection{Coordinating}
% \begin{itemize}
%     \item{\makebox[2cm][l]{And}  \_\_\_\_\_\_\_\_}
%     \item{\makebox[2cm][l]{But}  \_\_\_\_\_\_\_\_\_}
%     \item{\makebox[2cm][l]{So} \_\_\_\_\_\_\_\_}
% \end{itemize}

% \subsection{Subordinating}
% \begin{itemize}
%     \item{\makebox[2cm][l]{Because} \_\_\_\_\_\_\_\_}
%     \item{\makebox[2cm][l]{If} \_\_\_\_\_\_\_\_}
%     \item{\makebox[2cm][l]{Although} \_\_\_\_\_\_\_\_}
%     \item{\makebox[2cm][l]{When} \_\_\_\_\_\_\_\_}
%     \item{\makebox[2cm][l]{That} \_\_\_\_\_\_\_\_}
% \end{itemize}

% \subsection{Correlative}
% \begin{itemize}
%     \item{\makebox[2cm][l]{Either ... or} \_\_\_\_\_\_\_\_}
%     \item{\makebox[2cm][l]{Neither ... nor} \_\_\_\_\_\_\_\_}
% \end{itemize}



% 
% \section*{Vocabulary of Tovian}
% {\large \red{NEEDS TO BE UPDATED WITH NEW DICTIONARY}}


% \subsection*{Prefixes and Suffixes}
% \begin{tabular}{rl}
% \textbf{a-} & $\rightarrow$ of \\
% \textbf{y-} & $\rightarrow$ from \\
% \textbf{u-} & $\rightarrow$ with \\
% \textbf{hlu-} & $\rightarrow$ for \\
% \textbf{sho-} & $\rightarrow$ can (prefix) \\
% \textbf{ef-} & $\rightarrow$ requires (prefix) \\
% \textbf{ea-} & $\rightarrow$ noun to verb (prefix) \\
% \textbf{e} & $\rightarrow$ plural (postfix) \\
% \textbf{ni'} & $\rightarrow$ not/without (prefix) \\
% \textbf{nit'} & $\rightarrow$ opposite (prefix) \\
% \textbf{ana'} & $\rightarrow$ that which (prefix) \\
% \textbf{na'} & $\rightarrow$ which (prefix) \\
% \textbf{'as} & $\rightarrow$ to be (postfix to pronoun/noun) \\
% \textbf{'als} & $\rightarrow$ to be like (postfix to pronoun/noun) \\
% \textbf{thil} & $\rightarrow$ place or realm (postfix) \\
% \textbf{'nin} & $\rightarrow$ diminutive (postfix) \\
% \textbf{lum} & $\rightarrow$ good/desirable (postfix) \\
% \textbf{mah'} & $\rightarrow$ through (prefix).

% \end{tabular}

% \subsection*{Conjunctions, Prepositions, and Adverbs}
% \begin{tabular}{rl}
% \textbf{sa} & $\rightarrow$ but/however \\
% \textbf{tes} & $\rightarrow$ in order to \\
% \textbf{el} & $\rightarrow$ in \\
% \textbf{ta} & $\rightarrow$ while \\
% \textbf{sor} & $\rightarrow$ near/by \\
% \textbf{heth} & $\rightarrow$ ahead/front \\
% \textbf{hyul} & $\rightarrow$ behind/back \\
% \textbf{mah} & $\rightarrow$ through \\
% \textbf{wa} & $\rightarrow$ or \\
% \end{tabular}


% \subsection*{Natural Elements and Geography}
% \begin{tabular}{rl}
% \textbf{erathil} & $\rightarrow$ desert \\
% \end{tabular}

% \subsection*{Societal and Cultural Concepts}
% \begin{tabular}{rl}
% \textbf{nalhtlun} & $\rightarrow$ council \\
% \textbf{pelor} & $\rightarrow$ city \\
% \textbf{driel} & $\rightarrow$ person \\
% \textbf{ni'driel} & $\rightarrow$ non-person \\
% \textbf{nan} & $\rightarrow$ ancient \\
% \textbf{ni'nan} & $\rightarrow$ youth (not ancient) \\
% \textbf{ni'nan'nin} & $\rightarrow$ child (not ancient + diminutive) \\
% \end{tabular}

% \subsection*{Conceptual and Abstract Terms}
% \begin{tabular}{rl}
% \textbf{ara} & $\rightarrow$ wisdom \\
% \textbf{cana} & $\rightarrow$ change \\
% \textbf{thivm} & $\rightarrow$ important \\
% \textbf{haru} & $\rightarrow$ truth \\
% \textbf{kar} & $\rightarrow$ balance \\
% \textbf{bral} & $\rightarrow$ darkness \\
% \textbf{ahn} & $\rightarrow$ light \\
% \textbf{velar} & $\rightarrow$ worthy \\
% \end{tabular}

% \subsection*{Compound Words and Phrases}
% \begin{tabular}{r|l}
% \textbf{tkeshsul} & "under foot" \\
%  & = floor/ground \\
% \textbf{ana'leth} & "that which tastes" \\
%  & = food \\
% \textbf{ana'leth adrix} & meat \\
%  & = that which tastes of blood \\
% \textbf{shoiss} & bird \\
%  & = can fly \\
% \textbf{shoiss'mah'lothre} & river fish \\
%  & = can fly through rivers \\
% \textbf{shoiss'mah'nor} & sea fish \\
%  & = can fly through sea \\
% \textbf{leth'haru'sel} & true taste of water \\
% \textbf{marth'sul'melore} & forest air \\
%  & = clean air underneath \\
% \textbf{ruy'nin} & small feather \\
%  & = term of endearment \\
% \textbf{velar'vehkuilor} & worthy death by combat \\
% \textbf{num'hlu'ara} & path for wisdom \\
% \textbf{alu'sel'sir} & dread \\
% & = Deep water above\\
% \textbf{sura'lothr} & daily life \\
%  & = The routine river \\
% \textbf{yul'silytel} & the song created from the Lake \\
% \textbf{tkeshsil} & standing \\
%  & = over foot \\
% \textbf{kharànane} & old ways \\
%  & = balance of ancients \\
% \textbf{beth'har} & pain \\
%  & = feels alert \\
% \textbf{beth'asc} & fire \\
%  & = feels hot \\
% \textbf{eätul'har} & interesting \\
%  & = makes alert \\
% \textbf{heth'nume} & future \\
%  & = paths ahead \\
% \textbf{hyul'nume} & past \\
%  & = paths behind \\
% \textbf{rith'driel} & breath \\
%  & = person air \\
% \textbf{rith'driel'yul} & whistle \\
%  & = person air song \\
% \end{tabular}





% \subsection*{Numbers}
% \begin{tabular}{rl}
% \textbf{lo} & $\rightarrow$ one \\
% \textbf{loe} & $\rightarrow$ two \\
% \textbf{fem} & $\rightarrow$ three \\
% \textbf{tev} & $\rightarrow$ four \\
% \textbf{rep} & $\rightarrow$ five \\
% \textbf{tep} & $\rightarrow$ six \\
% \textbf{pol} & $\rightarrow$ seven \\
% \textbf{ath} & $\rightarrow$ eight \\
% \textbf{ki} & $\rightarrow$ nine \\
% \textbf{tir} & $\rightarrow$ ten \\
% \textbf{bey} & $\rightarrow$ eleven \\
% \textbf{hliya} & $\rightarrow$ twelve \\
% \textbf{hliya-lo} & $\rightarrow$ thirteen \\
% \textbf{hliya-loe} & $\rightarrow$ fourteen \\
% \textbf{hliya-fem} & $\rightarrow$ fifteen \\
% \textbf{hliya-ath} & $\rightarrow$ twenty \\
% \textbf{loe-hliyae} & $\rightarrow$ twenty four \\
% \textbf{loe-hliyae-lo} & $\rightarrow$ twenty five \\
% \textbf{loe-hliya-loe} & $\rightarrow$ twenty six \\
% \textbf{fem-hliyae} & $\rightarrow$ thirty six \\
% \textbf{tev-hliyae} & $\rightarrow$ forty eight \\
% \textbf{hliya-hliyae} & $\rightarrow$ one hundred forty four\\
% \textbf{emiyr} & $\rightarrow$ one thousand and eight\\


% \end{tabular}

% \newpage

% \subsection*{Root Words}
% \begin{tabular}{r|l}
% % \textbf{trom} & $\rightarrow$ no \\
% % \textbf{te} & $\rightarrow$ yes \\
% \textbf{lor} & $\rightarrow$ do\\
% \textbf{or} & $\rightarrow$ group\\
% \textbf{driel} & $\rightarrow$ person\\
% \textbf{lem} & $\rightarrow$ sun\\
% \textbf{elan} & $\rightarrow$ hand\\
% \textbf{bral} & $\rightarrow$ darkness \\
% \textbf{ahn} & $\rightarrow$ light \\
% \textbf{sor} & $\rightarrow$ closeness \\
% \textbf{heth} & $\rightarrow$ front \\
% \textbf{hul} & $\rightarrow$ back \\
% \textbf{melor} & $\rightarrow$ plant\\
% \textbf{glar} & $\rightarrow$ forest \\
% \textbf{era} & $\rightarrow$ sand \\
% \textbf{lothr} & $\rightarrow$ river \\
% \textbf{nor} & $\rightarrow$ sea \\
% \textbf{sel} & $\rightarrow$ water \\
% \textbf{rith} & $\rightarrow$ air \\
% \textbf{sul} & $\rightarrow$ earth \\
% \textbf{ask} & $\rightarrow$ hot \\
% \textbf{tkshu} & $\rightarrow$ cold \\
% \textbf{pathr} & $\rightarrow$ weight\\
% \textbf{har} & $\rightarrow$ pain\\
% \textbf{???} & $\rightarrow$ size\\

% \textbf{green} & $\rightarrow$ \\
% \textbf{red} & $\rightarrow$ \\
% \textbf{blue} & $\rightarrow$dora\\
% \textbf{???} & $\rightarrow$ \\
% \textbf{???} & $\rightarrow$ \\


% \textbf{kana} & $\rightarrow$ change \\
% \textbf{thul} & $\rightarrow$ top \\
% \textbf{sul} & $\rightarrow$ bottom \\
% \textbf{het} & $\rightarrow$ left \\
% \textbf{lhoy} & $\rightarrow$ right \\
% \textbf{lhel} & $\rightarrow$ center\\
% \textbf{lae} & $\rightarrow$ north \\
% \textbf{tal} & $\rightarrow$ east \\
% \textbf{ria} & $\rightarrow$ south \\
% \textbf{fil} & $\rightarrow$ west \\

% \textbf{ara} & $\rightarrow$ wisdom \\
% \textbf{velar} & $\rightarrow$ worth \\
% \textbf{thivem} & $\rightarrow$ importance \\
% \textbf{haru} & $\rightarrow$ correct/truth \\
% \textbf{kar} & $\rightarrow$ balance\\
% \textbf{quil} & $\rightarrow$ movement\\
% \textbf{tlesh} & $\rightarrow$ speed\\
% \textbf{lis} & $\rightarrow$ possession\\
% \textbf{lor} & $\rightarrow$ action\\
% \textbf{} & $\rightarrow$ name\\
% \textbf{tom} & $\rightarrow$ end\\

% \hline

% \textbf{hethnume} & $\rightarrow$ future\\
% \textbf{hulnume} & $\rightarrow$ past\\
% \textbf{lhelnume} & $\rightarrow$ present\\

% \textbf{lhonum} & $\rightarrow$ day (period)\\
% \textbf{tirlhon} & $\rightarrow$ week (period)\\
% \textbf{lheln} & $\rightarrow$ today\\
% \textbf{hethn} & $\rightarrow$ tomorrow\\
% \textbf{huln} & $\rightarrow$ yesterday\\
% \textbf{lorsel} & $\rightarrow$ wash\\
% \textbf{quilhum} & $\rightarrow$ dance\\

% \end{tabular}

% 
% \newpage
\onecolumn
\section*{Example Poem}
\textbf{Original Poem:}


\begin{quote}
{\fontsize{30pt}{10pt}\tovian
eth drixadrielas\\
na el lothre ea-iss\\
Emah glaranane\\
peloresul lea-kuil \\
analeth hlumelore
}
\end{quote}



\begin{quote}
Eth drixádriel'as\\
Ti-lothre lore-iss \\
Ti-glaranane\\
Peloresul lilore-quil \\
Ana'leth hlumelore
\end{quote}

\textbf{Translation:}
\begin{quote}
``I am the blood of elves,\\
Which swims in rivers,\\
Through ancient forests,\\
Moving under cities,\\
That which sustains the plants."
\end{quote}

This verse expresses an person's connection to their heritage, their harmony with nature, and their journey through both ancient and urban landscapes, ultimately acknowledging their role in the cycle of life as nourishment for flora.


% 
\section*{Example Sentences [WIP]}


\begin{tabular}{l|l c}
\textbf{Sentence} & \textbf{Tovian Translation} & &{\fontsize{20pt}{10pt}\tovian }\\ 

The sun shines. & lem lore-arthul &{\fontsize{20pt}{10pt}\tovian lem lore-arthul}\\
 
The sun is shining. & lem lore-arthul &{\fontsize{20pt}{10pt}\tovian lem lore-arthul}\\
 
The sun shone. & lem inlore-arthul &{\fontsize{20pt}{10pt}\tovian lem inlore-arthul}\\
 
The sun will shine. & lem vilore-arthul&{\fontsize{20pt}{10pt}\tovian lem vilore-arthul}\\
 
The sun has been shining. & lem deshilore-arthul &{\fontsize{20pt}{10pt}\tovian lem deshilore-arthul}\\
 
The sun is shining again. & &{\fontsize{20pt}{10pt}\tovian }\\
 
The sun will shine tomorrow. & &{\fontsize{20pt}{10pt}\tovian }\\
 
The sun shines brightly. & &{\fontsize{20pt}{10pt}\tovian }\\
 
The bright sun shines. & &{\fontsize{20pt}{10pt}\tovian }\\
 
The sun is rising now. & &{\fontsize{20pt}{10pt}\tovian }\\
 
All the people shouted. & &{\fontsize{20pt}{10pt}\tovian }\\
 
Some of the people shouted. & &{\fontsize{20pt}{10pt}\tovian }\\
 
Many of the people shouted twice. & &{\fontsize{20pt}{10pt}\tovian }\\
 
Happy people often shout. & &{\fontsize{20pt}{10pt}\tovian }\\
 
The kitten jumped up. & &{\fontsize{20pt}{10pt}\tovian }\\
 
The kitten jumped onto the table. & &{\fontsize{20pt}{10pt}\tovian }\\
 
My little kitten walked away. & &{\fontsize{20pt}{10pt}\tovian }\\
 
It's raining. & &{\fontsize{20pt}{10pt}\tovian }\\
 
The rain came down. & &{\fontsize{20pt}{10pt}\tovian }\\
 
The kitten is playing in the rain. & &{\fontsize{20pt}{10pt}\tovian }\\
 
The rain has stopped. & &{\fontsize{20pt}{10pt}\tovian }\\
 
Soon the rain will stop. & &{\fontsize{20pt}{10pt}\tovian }\\
 
I hope the rain stops soon. & &{\fontsize{20pt}{10pt}\tovian }\\
 
Once wild animals lived here. & &{\fontsize{20pt}{10pt}\tovian }\\
 
Slowly she looked around. & &{\fontsize{20pt}{10pt}\tovian }\\
 
Go away! & &{\fontsize{20pt}{10pt}\tovian }\\
 
Let's go! & &{\fontsize{20pt}{10pt}\tovian }\\
 
You should go. & &{\fontsize{20pt}{10pt}\tovian }\\
 
I will be happy to go. & &{\fontsize{20pt}{10pt}\tovian }\\
 
He will arrive soon. & &{\fontsize{20pt}{10pt}\tovian }\\
 
The baby's ball has rolled away. & &{\fontsize{20pt}{10pt}\tovian }\\
 
The two boys are working together. & &{\fontsize{20pt}{10pt}\tovian }\\
 
This mist will probably clear away. & &{\fontsize{20pt}{10pt}\tovian }\\
 
Lovely flowers are growing everywhere. & &{\fontsize{20pt}{10pt}\tovian }\\
 
We should eat more slowly. & &{\fontsize{20pt}{10pt}\tovian }\\
 
You have come too soon. & &{\fontsize{20pt}{10pt}\tovian }\\
 
You must write more neatly. & &{\fontsize{20pt}{10pt}\tovian }\\
 
Directly opposite stands a wonderful palace. & &{\fontsize{20pt}{10pt}\tovian }\\
 
Henry's dog is lost. & &{\fontsize{20pt}{10pt}\tovian }\\
 
My cat is black. & &{\fontsize{20pt}{10pt}\tovian }\\
 
The little girl's doll is broken. & &{\fontsize{20pt}{10pt}\tovian }\\
 
I usually sleep soundly. & &{\fontsize{20pt}{10pt}\tovian }\\
 
The children ran after Jack. & &{\fontsize{20pt}{10pt}\tovian }\\
 
I can play after school. & &{\fontsize{20pt}{10pt}\tovian }\\
 
We went to the village for a visit. & &{\fontsize{20pt}{10pt}\tovian }\\
 
We arrived at the river. & &{\fontsize{20pt}{10pt}\tovian }\\
 
I have been waiting for you. & &{\fontsize{20pt}{10pt}\tovian }\\
 
The campers sat around the fire. & &{\fontsize{20pt}{10pt}\tovian }\\
 
A little girl with a kitten sat near me. & &{\fontsize{20pt}{10pt}\tovian }\\
 
\end{tabular}

\begin{tabular}{l|l l}
\textbf{Sentence} & \textbf{Tovian Translation} \\ 
The child waited at the door for her father. & &{\fontsize{20pt}{10pt}\tovian }\\
 
Yesterday the oldest girl in the village lost her kitten. & &{\fontsize{20pt}{10pt}\tovian }\\
 
Were you born in this village? & &{\fontsize{20pt}{10pt}\tovian }\\
 

Can your brother dance well? & &{\fontsize{20pt}{10pt}\tovian }\\
 
Did the man leave? & &{\fontsize{20pt}{10pt}\tovian }\\
 
Is your sister coming for you? & &{\fontsize{20pt}{10pt}\tovian }\\
 
Can you come tomorrow? & &{\fontsize{20pt}{10pt}\tovian }\\
Have the neighbors gone away for the winter? & &{\fontsize{20pt}{10pt}\tovian }\\
 
Does the robin sing in the rain? & &{\fontsize{20pt}{10pt}\tovian }\\
 
Are you going with us to the concert? & &{\fontsize{20pt}{10pt}\tovian }\\
 
Have you ever travelled in the jungle? & &{\fontsize{20pt}{10pt}\tovian }\\
 
We sailed down the river for several miles. & &{\fontsize{20pt}{10pt}\tovian }\\
 
Everybody knows about hunting. & &{\fontsize{20pt}{10pt}\tovian }\\
 
On a Sunny morning after the solstice we \\ \indent started for the mountains. & &{\fontsize{20pt}{10pt}\tovian }\\
 
Tom laughed at the monkey's tricks. & &{\fontsize{20pt}{10pt}\tovian }\\
 
An old man with a walking stick stood beside the fence. & &{\fontsize{20pt}{10pt}\tovian }\\
 
The squirrel's nest was hidden by drooping boughs. & &{\fontsize{20pt}{10pt}\tovian }\\
 
The little seeds waited patiently under the snow for \\ \indent the warm spring sun. & &{\fontsize{20pt}{10pt}\tovian }\\
 
 
Many little girls with wreaths of flowers on their heads \\ \indent danced around the bonfire. & &{\fontsize{20pt}{10pt}\tovian }\\

 
The cover of the basket fell to the floor. & &{\fontsize{20pt}{10pt}\tovian }\\
 
The first boy in the line stopped at the entrance. & &{\fontsize{20pt}{10pt}\tovian }\\
On the top of the hill in a little hut lived a wise old woman. & &{\fontsize{20pt}{10pt}\tovian }\\
 
During our residence in the country we often \\ \indent walked in the pastures. & &{\fontsize{20pt}{10pt}\tovian }\\
 
When will your guests from the city arrive? & &{\fontsize{20pt}{10pt}\tovian }\\
 
Near the mouth of the river, its course turns sharply \\ \indent towards the East. & &{\fontsize{20pt}{10pt}\tovian }\\
 
 
Between the two lofty mountains lay a fertile valley. & &{\fontsize{20pt}{10pt}\tovian }\\
 
Among the wheat grew tall red poppies. & &{\fontsize{20pt}{10pt}\tovian }\\
 
The strong roots of the oak trees were torn from the ground. & &{\fontsize{20pt}{10pt}\tovian }\\
 
The sun looked down through the branches \\ \indent upon the children at play. & &{\fontsize{20pt}{10pt}\tovian }\\
 
The west wind blew across my face like a friendly caress. & &{\fontsize{20pt}{10pt}\tovian }\\
 
The spool of thread rolled across the floor. & &{\fontsize{20pt}{10pt}\tovian }\\
 
A box of growing plants stood in the Window. & &{\fontsize{20pt}{10pt}\tovian }\\
 
I am very happy. & &{\fontsize{20pt}{10pt}\tovian }\\
 
These oranges are juicy. & &{\fontsize{20pt}{10pt}\tovian }\\
 
Sea water is salty. & &{\fontsize{20pt}{10pt}\tovian }\\
 
The streets are full of people. & &{\fontsize{20pt}{10pt}\tovian }\\
 
Sugar tastes sweet. & &{\fontsize{20pt}{10pt}\tovian }\\
 
The fire feels hot. & &{\fontsize{20pt}{10pt}\tovian }\\
 
The little girl seemed lonely. & &{\fontsize{20pt}{10pt}\tovian }\\
 
The little boy's father had once been a sailor. & &{\fontsize{20pt}{10pt}\tovian }\\
 
I have lost my blanket. & &{\fontsize{20pt}{10pt}\tovian }\\
 
A robin has built his nest in the apple tree. & &{\fontsize{20pt}{10pt}\tovian }\\
 
At noon we ate our lunch by the roadside. & &{\fontsize{20pt}{10pt}\tovian }\\
 
Mr. Jones made a knife for his little boy. & &{\fontsize{20pt}{10pt}\tovian }\\
 
Their voices sound very happy. & &{\fontsize{20pt}{10pt}\tovian }\\
 
Is today Monday? & &{\fontsize{20pt}{10pt}\tovian }\\
 
Have all the leaves fallen from the tree? & &{\fontsize{20pt}{10pt}\tovian }\\
 
Will you be ready on time? & &{\fontsize{20pt}{10pt}\tovian }\\
 
Will you send this message for me? & &{\fontsize{20pt}{10pt}\tovian }\\
 
Are you waiting for me? & &{\fontsize{20pt}{10pt}\tovian }\\
 

\end{tabular}

\begin{tabular}{l|l l}
Is this the first kitten of the litter? & &{\fontsize{20pt}{10pt}\tovian }\\
 
Are these shoes too big for you? & &{\fontsize{20pt}{10pt}\tovian }\\
 
How wide is the River? & &{\fontsize{20pt}{10pt}\tovian }\\
 
Listen. & &{\fontsize{20pt}{10pt}\tovian }\\
 
Sit here by me. & &{\fontsize{20pt}{10pt}\tovian }\\
 
Keep this secret until tomorrow. & &{\fontsize{20pt}{10pt}\tovian }\\
 
Come with us. & &{\fontsize{20pt}{10pt}\tovian }\\
 
Bring your friends with you. & &{\fontsize{20pt}{10pt}\tovian }\\
 
Be careful. & &{\fontsize{20pt}{10pt}\tovian }\\
 
Have some tea. & &{\fontsize{20pt}{10pt}\tovian }\\
 
Pip and his dog were great friends. & &\\
 
John and Elizabeth are brother and sister. & &\\
 
You and I will go together. & &\\
 
They opened all the doors and windows. & &\\
 
He is small, but strong. & &\\
 
Is this tree an oak or a maple? & &\\
 
Does the sky look blue or gray? & &\\
 
Come with your father or mother. & &\\
 
I am tired, but very happy. & &\\
 
He played a tune on his wonderful flute. & &\\
 
Toward the end of August the days grow much shorter. & &\\
 
A company of soldiers marched over the hill  &\\ \indent and across the meadow. & &\\
 
The first part of the story is very interesting. & &\\
 
The crow dropped some pebbles into the pitcher  &\\ \indent and raised the water to the brim. & &\\
 
The baby clapped her hands and laughed in glee. & &\\
 
Stop your game and be quiet. & &\\
 
The sound of the drums grew louder and louder. & &\\
 
Do you like summer or winter better? & &\\
 
That boy will have a wonderful trip. & &\\
 
They popped corn, and then sat around  &\\ \indent  the fire and ate it. & &\\
 
They won the first two games, but lost the last one. & &\\
 
Take this note, carry it to your  &\\ \indent mother; and wait for an answer. & &\\
 
I awoke early, dressed hastily, and went down to breakfast. & &\\
 
Aha! I have caught you! & &\\
 
This string is too short! & &\\
 
Oh, dear! the wind has blown my hat away! & &\\
 
Alas! that news is sad indeed! & &\\
 
Whew! that cold wind freezes my nose! & &\\
 
Are you warm enough now? & &\\
 
They heard the warning too late. & &\\
 
We are a brave people, and love our country. & &\\
 
All the children came except Mary. & &\\
 
Jack seized a handful of pebbles and threw  &\\ \indent  them into the lake. & &\\
 
This cottage stood on a low hill, at some  &\\ \indent  distance from the village. & &\\
 
On a fine summer evening, the two old people  &\\ \indent  were sitting  outside the door of their cottage. & &\\
 
We visited my uncle's village, the largest village in the world. & &\\
 
We learn something new each day. & &\\
 
The market begins five minutes earlier this week. & &\\
 
Did you find the distance too great? & &\\
 
Hurry, children. & &\\
 
Madam, I will obey your command. & &\\
 
Here under this tree they gave their guests a splendid feast. & &\\
 
 
 

\end{tabular}


\begin{tabular}{l|l l }
\textbf{Sentence} & \textbf{Tovian Translation} &\\ 
In winter I get up at night, and dress by yellow candlelight. & &\\
 
Tell the last part of that story again. & &\\
 
 
Our bird's name is Jacko. & &\\
 
The river knows the way to the sea. & &\\
 
The boat sails away, like a bird on the wing. & &\\
 
They looked cautiously about, but saw nothing. & &\\
 
The little house had three rooms, a sitting room, a bedroom, &\\ \indent and a tiny kitchen. & &\\
 
Be quick or you will be too late. & ersh'seloreas wa shol vilore-fe'ninesh & \\ 
&{\fontsize{20pt}{10pt}\tovian ersh'seloreas wa shol vilore-fe'ninesh}&\\
 
Will you go with us or wait here? & &\\
 
She was always, shabby, often ragged, and on cold  &\\ \indent days very uncomfortable. & &\\
 
Think first and then act. & &\\
I stood, a little mite of a girl, upon a chair by the window, &\\ \indent and watched the falling snowflakes. & &\\
 
Show the guests these shells, my son, and tell them  &\\ \indent their strange history. & &\\
Be satisfied with nothing but your best. & &\\

We consider them our faithful friends. & &\\
 
We will make this place our home. & &\\
 
The squirrels make their nests warm and snug with  &\\ \indent soft moss and leaves. & &\\
 
The little girl made the doll's dress herself. & &\\
 
I hurt myself. & &\\
 
She was talking to herself. & &\\
 
He proved himself trustworthy. & &\\
 
We could see ourselves in the water. & &\\
 
Do it yourself. & &\\
 
I feel ashamed of myself. & &\\
 
Sit here by yourself. & &\\
 
The dress of the little princess was embroidered with roses,  &\\ \indent the national flower of the Country. & &\\
 
They wore red caps, the symbol of liberty. & &\\
 
With him as our protector, we fear no danger. & &\\
 
All her finery, lace, ribbons, and feathers,  &\\ \indent  was packed away in a trunk. & &\\
 
Light he thought her, like a feather. & &\\
 
Every spring and fall our cousins pay us a long visit. & &\\
 
In our climate the grass remains green all winter. & &\\
 
The boy who brought the book has gone. & &\\
 
These are the flowers that you ordered. & &\\
 
I have lost the book that you gave me. & &\\
 
The fisherman who owned the boat now demanded payment. & &\\
 
Come when you are called. & &\\
 
I shall stay at home if it rains. & &\\
 
When he saw me, he stopped. & &\\
 
Do not laugh at me because I seem so absent minded. & &\\
 
I shall lend you the books that you need. & &\\
 
Come early next Monday if you can. & &\\
 

\end{tabular}


\begin{tabular}{l|l l}
\textbf{Sentence} & \textbf{Tovian Translation} &\\ 
 
If you come early, wait in the hall. & shol kolore-quilsor nesh, selore-dolr im'moryuth!&\\
 
I had a younger brother whose name was Antonio. & &\\
 
Gnomes are little men who live under the ground. & &\\
 
He is loved by everybody, because he has  &\\ \indent  a gentle disposition. & &\\
 
Hold the horse while I run and get my cap. & &\\
 
I have found the ring I lost. & &\\
 
Play and I will sing. & &\\
 
That is the funniest story I ever heard. & &\\
 
She is taller than her brother. & &\\
 
They are no wiser than we. & &\\
 
Light travels faster than sound. & &\\
 
We have more time than they. & &\\
 
She has more friends than enemies. & &\\
 
He was very poor, and with his wife and five children lived  &\\ \indent in a little low cabin of logs and stones. & &\\
 
When the wind blew, the traveler wrapped his mantle more  &\\ \indent closely around him. & &\\
 
I am sure that we can go. & &\\
 
We went back to the place where we saw the roses. & &\\

"This tree is fifty feet high," said the gardener. & &\\
 
I think that this train leaves five minutes earlier today. & &\\
 
My opinion is that the governor will grant him a pardon. & &\\
 
Why he has left the city is a mystery. & risa wel pelor delore-quilnisor pfetr'as  &\\ 
 
The house stands where three roads meet. & &\\
 
He has far more money than brains. & &\\
 
Evidently that gate is never opened, for the long grass and the &\\ \indent  great hemlocks grow close against it. & &\\
 
I met a little cottage girl; she was eight years old, she said. & &\\
&&\\

If he has wisdom (he might), he would speak truthfully. &
La-wel yale-ara indekolore-lis, la-wel kolore-fethr sile-haru &\\
&
{\fontsize{14pt}{10pt}\tovian la-wel yale-ara indekolore-lis la-wel kolore-fethr sile-haru}&
\\

If he had wisdom (he doesn't), he would speak truthfully. &
La-wel yale-ara indexolore-lis, la-wel xolore-fethr sile-haru &\\
&
{\fontsize{14pt}{10pt}\tovian la-wel yale-ara xolore-lis la-wel xolore-fethr sile-haru}&
\\

If he had wisdom (he might), would he speak truthfully? &
La-wel yale-ara indekolore-lis, la-wel kolore-fethr sile-haru? &\\
&
{\fontsize{14pt}{10pt}\tovian la-wel yale-ara indekolore-lis, la-wel kolore-fethr sile-haru}&
\\

\end{tabular}

% \twocolumn

\section*{Additional Vocabulary}

\begin{tabular}{r c c l}
\textbf{???} &&$\rightarrow$ hello \\
\textbf{???} &&$\rightarrow$ goodbye \\
\textbf{kului} &&$\rightarrow$ cinnamon \\
\textbf{helor} &&$\rightarrow$ glass \\
\textbf{dora} &&$\rightarrow$ blue \\
\textbf{fethr} &&$\rightarrow$ speaker \\
\textbf{drix} &&$\rightarrow$ blood \\
\textbf{iss} &&$\rightarrow$ fly/swim \\
\textbf{tes} &&$\rightarrow$ sky \\
\textbf{imalor} &&$\rightarrow$ experience \\
\textbf{emethl} &&$\rightarrow$ intelligence \\
\textbf{fer} &&$\rightarrow$ outdoors \\
\textbf{lom} &&$\rightarrow$ soft \\
\textbf{marth} &&$\rightarrow$ clean \\
\textbf{sura} &&$\rightarrow$ regular/routine \\
\textbf{thu} &&$\rightarrow$ high \\
\textbf{gol} &&$\rightarrow$ passive \\
\textbf{val} &&$\rightarrow$ skilled \\
\textbf{har} &&$\rightarrow$ alert \\
\textbf{curth} &&$\rightarrow$ expensive \\
\textbf{sepa} &&$\rightarrow$ fresh \\
\textbf{asc} &&$\rightarrow$ hot \\
\textbf{bem} &&$\rightarrow$ sweet \\
\textbf{bhal} &&$\rightarrow$ sour \\
\textbf{leth'har} &&$\rightarrow$ spicy \\
\textbf{fex} &&$\rightarrow$ bitter \\
\textbf{som} &&$\rightarrow$ umami \\
\textbf{alu} &&$\rightarrow$ deep \\
\textbf{tkesh} &&$\rightarrow$ foot \\
\textbf{leth} &&$\rightarrow$ taste \\
\textbf{melor} &&$\rightarrow$ plant \\
\textbf{tith} &&$\rightarrow$ quiet \\
\textbf{nin} &&$\rightarrow$ small (diminutive) \\
\textbf{num} &&$\rightarrow$ path \\
\textbf{yul} &&$\rightarrow$ song \\
\textbf{mril} &&$\rightarrow$ eye \\
\textbf{ruy} &&$\rightarrow$ feather \\
\textbf{pila} &&$\rightarrow$ love \\
\textbf{vehk} &&$\rightarrow$ kill \\
\textbf{ilor} &&$\rightarrow$ death \\
\textbf{quil} &&$\rightarrow$ movement \\
\textbf{???} &&$\rightarrow$ carry \\
\textbf{wen} &&$\rightarrow$ beauty \\
\textbf{asp} &&$\rightarrow$ sharp \\
\textbf{???} &&$\rightarrow$ hold \\
\textbf{lim} &&$\rightarrow$ star \\
\textbf{???} &&$\rightarrow$ within \\
\end{tabular}

\begin{tabular}{rl}
\textbf{???} & $\rightarrow$ vast \\
\textbf{elan} & $\rightarrow$ hand \\
\textbf{???} & $\rightarrow$ search \\
\textbf{melorerae} & $\rightarrow$ oasis \\
\textbf{???} & $\rightarrow$ dusk \\
\textbf{???} & $\rightarrow$ dawn \\
\textbf{iro} & $\rightarrow$ morning \\
\textbf{} & $\rightarrow$ afternoon \\
\textbf{tuml} & $\rightarrow$ night \\
\textbf{???} & $\rightarrow$ attack \\
\textbf{???} & $\rightarrow$ friend \\
\textbf{yush} & $\rightarrow$ block \\
\textbf{a} & $\rightarrow$ of \\
\textbf{ul} & $\rightarrow$ and \\
\textbf{???} & $\rightarrow$ to \\
\textbf{???} & $\rightarrow$ in \\
\textbf{???} & $\rightarrow$ that \\
\textbf{???} & $\rightarrow$ was \\
\textbf{???} & $\rightarrow$ being \\
\textbf{???} & $\rightarrow$ with \\
\textbf{???} & $\rightarrow$ is \\
\textbf{???} & $\rightarrow$ for \\
\textbf{???} & $\rightarrow$ as \\
\textbf{???} & $\rightarrow$ had \\
\textbf{???} & $\rightarrow$ you \\
\textbf{???} & $\rightarrow$ not \\
\textbf{???} & $\rightarrow$ on \\
\textbf{???} & $\rightarrow$ at \\
\textbf{???} & $\rightarrow$ by \\
\textbf{???} & $\rightarrow$ which \\
\textbf{lore-lis} & $\rightarrow$ have \\
\textbf{lis} & $\rightarrow$ possession\\
\textbf{wa} & $\rightarrow$ or \\
\textbf{???} & $\rightarrow$ from \\
\textbf{???} & $\rightarrow$ this \\
\textbf{???} & $\rightarrow$ but \\
\textbf{???} & $\rightarrow$ all \\
\textbf{nesh} & $\rightarrow$ early \\
\end{tabular}


\begin{tabular}{rl}
\textbf{???} & $\rightarrow$ were \\
\textbf{???} & $\rightarrow$ are \\
\textbf{???} & $\rightarrow$ one \\
\textbf{???} & $\rightarrow$ so \\
\textbf{???} & $\rightarrow$ an \\
\textbf{???} & $\rightarrow$ said \\
\textbf{???} & $\rightarrow$ they \\
\textbf{???} & $\rightarrow$ who \\
\textbf{???} & $\rightarrow$ would \\
\textbf{???} & $\rightarrow$ been \\
\textbf{???} & $\rightarrow$ will (N) \\
\textbf{???} & $\rightarrow$ when \\
\textbf{???} & $\rightarrow$ there \\
\textbf{???} & $\rightarrow$ if \\
\textbf{???} & $\rightarrow$ more \\
\textbf{???} & $\rightarrow$ out \\
\textbf{???} & $\rightarrow$ up \\
\textbf{???} & $\rightarrow$ into \\
\textbf{???} & $\rightarrow$ do \\
\textbf{???} & $\rightarrow$ any \\
\textbf{???} & $\rightarrow$ your \\
\textbf{???} & $\rightarrow$ what \\
\textbf{???} & $\rightarrow$ has \\
\textbf{or} & $\rightarrow$ person \\
\textbf{???} & $\rightarrow$ could \\
\textbf{???} & $\rightarrow$ other \\
\textbf{???} & $\rightarrow$ than \\
\textbf{???} & $\rightarrow$ our \\
\textbf{???} & $\rightarrow$ some \\
\textbf{???} & $\rightarrow$ very \\
\textbf{???} & $\rightarrow$ time \\
\textbf{???} & $\rightarrow$ upon \\
\textbf{???} & $\rightarrow$ about \\
\textbf{???} & $\rightarrow$ may \\
\textbf{???} & $\rightarrow$ its \\
\textbf{???} & $\rightarrow$ only \\
\textbf{im'num} & $\rightarrow$ now \\
\textbf{???} & $\rightarrow$ like \\
\textbf{???} & $\rightarrow$ then \\
\textbf{???} & $\rightarrow$ can \\
\textbf{???} & $\rightarrow$ should \\
\textbf{???} & $\rightarrow$ did \\
\textbf{???} & $\rightarrow$ such \\
\textbf{???} & $\rightarrow$ great \\
\textbf{???} & $\rightarrow$ before \\
\textbf{???} & $\rightarrow$ must \\
\textbf{???} & $\rightarrow$ two \\
\textbf{???} & $\rightarrow$ these \\
\textbf{???} & $\rightarrow$ see \\
\textbf{arash} & $\rightarrow$ knowledge \\
\textbf{???} & $\rightarrow$ over \\
\textbf{???} & $\rightarrow$ down \\
\textbf{???} & $\rightarrow$ after \\
\textbf{???} & $\rightarrow$ first \\
\textbf{???} & $\rightarrow$ good \\
\textbf{???} & $\rightarrow$ never \\
\end{tabular}


\begin{tabular}{rl}
\textbf{???} & $\rightarrow$ most \\
\textbf{???} & $\rightarrow$ old \\
\textbf{???} & $\rightarrow$ day \\
\textbf{???} & $\rightarrow$ where \\
\textbf{asor} & $\rightarrow$ here \\
\textbf{sor} & $\rightarrow$ close \\
\textbf{quilsor} & $\rightarrow$ to come \\
\textbf{???} & $\rightarrow$ way \\
\textbf{???} & $\rightarrow$ work \\
\textbf{???} & $\rightarrow$ life \\
\textbf{???} & $\rightarrow$ without \\
\textbf{quil} & $\rightarrow$ to go \\
\textbf{???} & $\rightarrow$ well \\
\textbf{ema} & $\rightarrow$ through \\
\textbf{???} & $\rightarrow$ long \\
\textbf{???} & $\rightarrow$ say \\
\textbf{???} & $\rightarrow$ might \\
\textbf{???} & $\rightarrow$ how \\
\textbf{???} & $\rightarrow$ too \\
\textbf{???} & $\rightarrow$ even \\
\textbf{???} & $\rightarrow$ again \\
\textbf{???} & $\rightarrow$ many \\
\textbf{???} & $\rightarrow$ back \\
\textbf{???} & $\rightarrow$ think \\
\textbf{???} & $\rightarrow$ every \\
\textbf{???} & $\rightarrow$ people \\
\textbf{???} & $\rightarrow$ went \\
\textbf{???} & $\rightarrow$ same \\
\textbf{???} & $\rightarrow$ last \\
\textbf{tason} & $\rightarrow$ worse \\
\textbf{áemethl} & $\rightarrow$ thought \\
\textbf{???} & $\rightarrow$ away \\
\textbf{???} & $\rightarrow$ under \\
\textbf{???} & $\rightarrow$ take \\
\textbf{???} & $\rightarrow$ found \\
\textbf{elan} & $\rightarrow$ hand \\
\textbf{elane} & $\rightarrow$ hands \\
\textbf{mrile} & $\rightarrow$ eyes \\
\textbf{???} & $\rightarrow$ still \\
\textbf{???} & $\rightarrow$ place \\
\textbf{???} & $\rightarrow$ while \\
\textbf{???} & $\rightarrow$ just \\
\textbf{???} & $\rightarrow$ also \\
\textbf{???} & $\rightarrow$ young \\
\textbf{???} & $\rightarrow$ yet \\
\textbf{???} & $\rightarrow$ though \\
\textbf{???} & $\rightarrow$ against \\
\textbf{???} & $\rightarrow$ things \\
\textbf{???} & $\rightarrow$ get \\
\textbf{???} & $\rightarrow$ ever \\
\textbf{???} & $\rightarrow$ give \\
\textbf{???} & $\rightarrow$ years \\
\textbf{???} & $\rightarrow$ off \\
\textbf{???} & $\rightarrow$ face \\
\end{tabular}


\begin{tabular}{rl}
\textbf{???} & $\rightarrow$ nothing \\
\textbf{???} & $\rightarrow$ right \\
\textbf{???} & $\rightarrow$ once \\
\textbf{???} & $\rightarrow$ left \\
\textbf{???} & $\rightarrow$ part \\
\textbf{???} & $\rightarrow$ saw \\
\textbf{mril} & $\rightarrow$ world \\
\textbf{lem} & $\rightarrow$ sun \\
\textbf{ea-arthul} & $\rightarrow$ shine \\
\textbf{arthul} & $\rightarrow$ bright \\
\textbf{nisil} & $\rightarrow$ break \\
\textbf{???} & $\rightarrow$ head \\
\textbf{???} & $\rightarrow$ took \\
\textbf{???} & $\rightarrow$ new \\
\textbf{simra} & $\rightarrow$ always \\
\textbf{???} & $\rightarrow$ put \\
\textbf{???} & $\rightarrow$ night \\
\textbf{???} & $\rightarrow$ each \\
\textbf{rethin} & $\rightarrow$ between \\
\textbf{???} & $\rightarrow$ tell \\
\textbf{???} & $\rightarrow$ mind \\
\textbf{arpel} & $\rightarrow$ heart \\
\textbf{???} & $\rightarrow$ few \\
\textbf{elan} & $\rightarrow$ because \\
\textbf{???} & $\rightarrow$ thing \\
\textbf{???} & $\rightarrow$ far \\
\end{tabular}


\begin{tabular}{rl}
\textbf{???} & $\rightarrow$ seemed \\
\textbf{???} & $\rightarrow$ look \\
\textbf{???} & $\rightarrow$ called \\
\textbf{???} & $\rightarrow$ whole \\
\textbf{???} & $\rightarrow$ set \\
\textbf{???} & $\rightarrow$ both \\
\textbf{???} & $\rightarrow$ got \\
\textbf{???} & $\rightarrow$ find \\
\textbf{sil} & $\rightarrow$ done \\
\textbf{???} & $\rightarrow$ heard \\
\textbf{???} & $\rightarrow$ name \\
\textbf{???} & $\rightarrow$ days \\
\textbf{???} & $\rightarrow$ told \\
\textbf{???} & $\rightarrow$ let \\
\textbf{edom} & $\rightarrow$ country \\
\textbf{???} & $\rightarrow$ asked \\
\textbf{???} & $\rightarrow$ going \\
\textbf{???} & $\rightarrow$ seen \\
\textbf{???} & $\rightarrow$ better \\
\textbf{toki} & $\rightarrow$ dog \\
\textbf{???} & $\rightarrow$ having \\
\textbf{???} & $\rightarrow$ home \\
\textbf{???} & $\rightarrow$ knew \\
\textbf{???} & $\rightarrow$ side \\
\textbf{???} & $\rightarrow$ something \\
\textbf{???} & $\rightarrow$ moment \\
\textbf{???} & $\rightarrow$ among \\
\textbf{???} & $\rightarrow$ course \\
\textbf{???} & $\rightarrow$ enough \\
\textbf{???} & $\rightarrow$ word \\
\textbf{???} & $\rightarrow$ soon \\
\textbf{bon} & $\rightarrow$ full \\
\textbf{???} & $\rightarrow$ end \\
\textbf{???} & $\rightarrow$ gave \\
\textbf{???} & $\rightarrow$ room \\
\textbf{???} & $\rightarrow$ almost \\
\textbf{???} & $\rightarrow$ want \\
\textbf{???} & $\rightarrow$ however \\
\end{tabular}


\begin{tabular}{rl}
\textbf{ahn} & $\rightarrow$ light \\
\textbf{???} & $\rightarrow$ quite \\
\textbf{???} & $\rightarrow$ bring \\
\textbf{???} & $\rightarrow$ given \\
\textbf{yuthil} & $\rightarrow$ door \\
\textbf{thul} & $\rightarrow$ best \\
\textbf{???} & $\rightarrow$ turned \\
\textbf{???} & $\rightarrow$ taken \\
\textbf{???} & $\rightarrow$ does \\
\textbf{???} & $\rightarrow$ use \\
\textbf{???} & $\rightarrow$ felt \\
\textbf{???} & $\rightarrow$ until \\
\textbf{???} & $\rightarrow$ since \\
\textbf{lathn} & $\rightarrow$ power \\
\textbf{aipte} & $\rightarrow$ themselves \\
\textbf{???} & $\rightarrow$ used \\
\textbf{???} & $\rightarrow$ rather \\
\textbf{???} & $\rightarrow$ began \\
\textbf{???} & $\rightarrow$ present \\
\textbf{pfetr} & $\rightarrow$ mystery \\
\textbf{hlufethr} & $\rightarrow$ voice \\
\textbf{???} & $\rightarrow$ others \\
\textbf{???} & $\rightarrow$ works \\
\textbf{???} & $\rightarrow$ less \\
\textbf{???} & $\rightarrow$ money \\
\textbf{???} & $\rightarrow$ next \\
\textbf{???} & $\rightarrow$ poor \\
\textbf{???} & $\rightarrow$ stood \\
\textbf{???} & $\rightarrow$ form \\
\textbf{???} & $\rightarrow$ within \\
\textbf{usol} & $\rightarrow$ together \\
\textbf{???} & $\rightarrow$ till \\
\textbf{???} & $\rightarrow$ large \\
\textbf{???} & $\rightarrow$ matter \\
\textbf{???} & $\rightarrow$ kind \\
\textbf{???} & $\rightarrow$ often \\
\textbf{???} & $\rightarrow$ year \\
\textbf{???} & $\rightarrow$ friend \\
\textbf{???} & $\rightarrow$ order \\
\textbf{???} & $\rightarrow$ round \\
\textbf{???} & $\rightarrow$ true \\
\textbf{???} & $\rightarrow$ anything \\
\textbf{???} & $\rightarrow$ keep \\
\end{tabular}


\begin{tabular}{rl}
\textbf{???} & $\rightarrow$ sent \\
\textbf{???} & $\rightarrow$ means \\
\textbf{???} & $\rightarrow$ believe \\
\textbf{yu'quil}  & $\rightarrow$ pass \\
\textbf{???} & $\rightarrow$ feet \\
\textbf{???} & $\rightarrow$ near \\
\textbf{???} & $\rightarrow$ public \\
\textbf{???} & $\rightarrow$ state \\
\textbf{???} & $\rightarrow$ hundred \\
\textbf{???} & $\rightarrow$ hope \\
\textbf{???} & $\rightarrow$ alone \\
\textbf{sir} & $\rightarrow$ above \\
\textbf{???} & $\rightarrow$ case \\
\textbf{???} & $\rightarrow$ high \\
\textbf{???} & $\rightarrow$ read \\
\textbf{???} & $\rightarrow$ received \\
\textbf{haru} & $\rightarrow$ fact \\
\textbf{???} & $\rightarrow$ gone \\
\textbf{???} & $\rightarrow$ girl \\
\textbf{pea-arash} & $\rightarrow$ known \\
\textbf{???} & $\rightarrow$ hear \\
\textbf{???} & $\rightarrow$ times \\
\textbf{???} & $\rightarrow$ least \\
\textbf{???} & $\rightarrow$ perhaps \\
\textbf{???} & $\rightarrow$ sure \\
\textbf{???} & $\rightarrow$ indeed \\
\textbf{???} & $\rightarrow$ open \\
\textbf{???} & $\rightarrow$ body \\
\textbf{aipt} & $\rightarrow$ itself \\
\textbf{???} & $\rightarrow$ along \\
\textbf{???} & $\rightarrow$ land \\
\textbf{???} & $\rightarrow$ return \\
\textbf{ea-quilnisor} & $\rightarrow$ leave (depart) \\
\textbf{ea-wushed} & $\rightarrow$ leave (abandon) \\
\textbf{wushed} & $\rightarrow$ abandonment \\
\textbf{ishma} & $\rightarrow$ guide \\
\textbf{rith} & $\rightarrow$ air \\
\textbf{anafer} & $\rightarrow$ nature \\
\textbf{???} & $\rightarrow$ answer \\
\textbf{???} & $\rightarrow$ either \\
\textbf{???} & $\rightarrow$ law \\
\textbf{???} & $\rightarrow$ help \\
\textbf{ea-quil'sul} & $\rightarrow$ lay \\
\end{tabular}


\begin{tabular}{rl}
\textbf{???} & $\rightarrow$ point \\
\textbf{ni'nan'nin} & $\rightarrow$ child \\
\textbf{???} & $\rightarrow$ letter \\
\textbf{???} & $\rightarrow$ wish \\
\textbf{???} & $\rightarrow$ cried \\
\textbf{ersh} & $\rightarrow$ fast \\
\textbf{???} & $\rightarrow$ number \\
\textbf{???} & $\rightarrow$ therefore \\
\textbf{???} & $\rightarrow$ hour \\
\textbf{???} & $\rightarrow$ held \\
\textbf{???} & $\rightarrow$ free \\
\textbf{keshl} & $\rightarrow$ war \\
\textbf{???} & $\rightarrow$ during \\
\textbf{???} & $\rightarrow$ several \\
\textbf{shissan} & $\rightarrow$ business \\
\textbf{???} & $\rightarrow$ whether \\
\textbf{???} & $\rightarrow$ manner \\
\textbf{???} & $\rightarrow$ second \\
\textbf{???} & $\rightarrow$ reason \\
\textbf{???} & $\rightarrow$ reply \\
\textbf{???} & $\rightarrow$ call \\
\textbf{???} & $\rightarrow$ general \\
\textbf{risa} & $\rightarrow$ why \\
\textbf{???} & $\rightarrow$ behind \\
\textbf{???} & $\rightarrow$ become \\
\textbf{???} & $\rightarrow$ lost \\
\textbf{???} & $\rightarrow$ forth \\
\textbf{emiyr} & $\rightarrow$ thousand \\
\textbf{???} & $\rightarrow$ family \\
\textbf{ea-beth} & $\rightarrow$ feel \\
\textbf{hluara} & $\rightarrow$ soul \\
\textbf{hluara} & $\rightarrow$ spirit \\
\textbf{asher} & $\rightarrow$ question \\
\textbf{???} & $\rightarrow$ care \\
\textbf{aharu} & $\rightarrow$ truth \\
\textbf{tkesh'sul} & $\rightarrow$ ground \\
\textbf{???} & $\rightarrow$ really \\
\textbf{???} & $\rightarrow$ rest \\
\textbf{???} & $\rightarrow$ mean \\
\textbf{???} & $\rightarrow$ different \\
\textbf{ea-sil} & $\rightarrow$ make \\
\textbf{???} & $\rightarrow$ possible \\
\textbf{???} & $\rightarrow$ fell \\
\textbf{yu} & $\rightarrow$ towards \\
\textbf{tasondriel} & $\rightarrow$ human \\
\textbf{aglardriel} & $\rightarrow$ wood-person \\
\textbf{or'keshl} & $\rightarrow$ orc \\
\textbf{???} & $\rightarrow$ kept \\
\textbf{tros} & $\rightarrow$ sad \\
\textbf{trossel} & $\rightarrow$ tear \\
\textbf{dren} & $\rightarrow$ flower \\
\textbf{tis} & $\rightarrow$ stone \\
\textbf{tisàhlúara} & $\rightarrow$ soul stone \\
\textbf{cyran} & $\rightarrow$ dapple \\
\textbf{lerth} & $\rightarrow$ tree \\
\textbf{merth} & $\rightarrow$ sound\\
\textbf{yulara}  & $\rightarrow$ name\\
\textbf{merthyara} & $\rightarrow$ prayer\\
\textbf{dorl} & $\rightarrow$ wait\\
\textbf{passage} & $\rightarrow$ moryuth\\
\textbf{hallway} & $\rightarrow$ moryuth\\
\end{tabular}


\iffalse
# See the documentation linked above for full details.

# These set up assimilations ("np" will be changed to "mp") and a few other
# minor things.  It is possible for assimilations to produce phonemes you
# don't want, which is why one of the filters below is: ŋ > n.
with: std-ipa-features std-assimilations coronal-metathesis

# This determines the sort order of the output, and also helps
# the 'with:' line settings understand what is going on.
letters: a b d e f g h i q l m n o p r s t u v w x y sh th

# THE ORDER OF PHONEME CLASSES MATTERS.  The very first phoneme
# will be picked much, much more often than the last.  The order
# here is approximately natural for a lot of languages.  You can
# change the character of a language a lot by shuffling which
# phonemes occur most often, which the least.
C = th sh r f d n s h l m p t g v q w x y b
D = f d n p t
V = a e i o u

# Macros are only a convenience to reduce typing, and do simple
# substitution only.  You can't use macros inside other macros.
$S = CVD?

# The first word shape will be picked most often, the last least
# often.
words: V?$S$S V?$S V?$S$S$S

# Rejections and filters use Javascript regular expressions.
# 'reject:' simply throws away a word.  Filter turns certain
# paterns into something else.
reject: wu yi w$ y$ h$ ʔʔ (p|t|k|ʔ)h
filter: nr > tr; mr > pr; ŋ > n

# Haplology - remove repeats.
reject: (..+)\1+
\fi



Potential words:

gashl

pathr






% {\centering Vowels
\begin{center}
    \begin{tabular}{cccccccc}
        & \textbf{Front} & & \textbf{Central} & & \textbf{Back} & & \\
        \textbf{Close} & i & &  & & u & &  \\
        \textbf{Close-mid} & e & &  & & o & &  \\
        \textbf{Mid} &  & &  & &  & &  \\
        \textbf{Open-mid} &  & &  & &  & &  \\
        \textbf{Open} & a & &  & &  & &  \\
    \end{tabular}
\end{center}

\medskip

{\huge WIP, unfinished/incorrect:}

\setlength{\tabcolsep}{2pt} % Adjust column separation (default is 6pt)
\begin{center}
{\fontsize{7pt}{10pt}
    \begin{tabular}{ccccccccccccccc}
        & & \textbf{Bilabial} & & \textbf{Labiodental} & & \textbf{Dental} & & \textbf{Alveolar} & & \textbf{Postalveolar} & & \textbf{Retroflex} & & \textbf{Palatal} & & \textbf{Velar} & & \textbf{Uvular} & & \textbf{Pharyngeal} & & \textbf{Glottal} \\
        \textbf{Plosive} & & p & &  & & \textsubbridge{t} & & & &  & & \textrtaild & & q & &  & &  & &  \\
        \textbf{Nasal} & & m & &  & &  & & n & &  & &  & & ŋ & &  & &  & &  \\
        \textbf{Trill} & &  & &  & &  & & r & &  & &  & &  & &  & &  & &  \\
        \textbf{Tap/Flap} & &  & &  & &  & &  & &  & &  & &  & &  & &  & &  \\
        \textbf{Fricative} & &  & &  & &  & & s & & \textesh & &\texttheta & & \textphi & &  & & h \\
        \textbf{Lateral Fricative} & &  & &  & &  & &  & &  & &  & &  & &  \\
        \textbf{Approximant} & &  & &  & &  & &  & &  & &  & & j & &  & &  \\
        \textbf{Lateral Approximant} & &  & &  & &  & & ɫ & &  & &  & &  & &  & &  & &  \\
    \end{tabular}
}

Other sounds:

\textturnw

\textltilde


\end{center}
}




\onecolumn

\section{Phonetic Inventory}

{\fontsize{8pt}{10pt}
\setlength{\tabcolsep}{8pt} % Adjust the column spacing here
\begin{center}
\begin{tabular}{|r|c|c|c|c|c|c|c|c|}
\hline
 & \textbf{Bilabial} & \textbf{Labiodental} & \textbf{Dental} & \textbf{Alveolar} & \textbf{Palatal} & \textbf{Postalveolar} & \textbf{Velar} & \textbf{Glottal} \\
\hline
\textbf{Plosive} & \textipa{p b} &  &  & \textipa{t d} & \textipa{\textbardotlessj} &  & \textipa{k g} &  \\
\hline
\textbf{Nasal} & \textipa{m} &  &  & \textipa{n} &  &  & \textipa{\ng} &  \\
\hline
\textbf{Trill} &  &  &  & \textipa{r} &  &  &  &  \\
\hline
\textbf{Fricative} & \textipa{\textphi} & \textipa{f v} & \textipa{\texttheta \red{\textipa{\dh}}} & \textipa{s \textcolor{red}{z}} &  & \textipa{\textesh \textcolor{red}{\textyogh}} &  & \textipa{h} \\
\hline
\textbf{Lateral Fricative} &  &  &  & \textipa{\textbeltl} &  &  &  &  \\
\hline
\textbf{Approximant} &  &  &  &  &\textipa{j} &  &  & \\
\hline
\textbf{Lateral Approximant} &  &  &  & \textipa{l} &  &  &  &  \\
\hline
\end{tabular}
\\
other: \textipa{w}. Sounds in \red{red} were in Ancient Lefethr but have since been lost.
\end{center}
}
\medskip
\medskip

\section{History of Sound Changes}
{
\setlength{\tabcolsep}{8pt} % Adjust the column spacing here
\begin{tabular}{ll}
\textbf{Year} & \textbf{Change}\\
\hline
1000 & Vowel loss between voiceless consonants in unstressed syllables \\& (p t f \texttheta\ s \textesh\ h k m n \ng) \\

\hline
2000 & Assimilation: Voiceless stop between voiced sounds become voiced \\& \{p t k\} $\rightarrow$ \{ b d g\} \\

\hline
2300 & \textschwa\ lost \\

\hline
3000 & No voiceless stops in clusters, \\& e.g., ttelama $\rightarrow$ telama \\

\hline
3500 & \textbardotlessj\ $\rightarrow$ j \\

\hline
3500 & Rhotacism GsG $\rightarrow$ GrG and GʒG $\rightarrow$ GrG \\

\hline
3501 & Rhotacism VsV $\rightarrow$ VrV to VʒV $\rightarrow$ VrV \\

\hline
4500 & No stops after fricatives \\

\hline
4501 & No stops after liquids \\

\hline
4502 & No fricative clusters \\

\hline
4503 & No stops after glides \\

\hline
5000 & h is lost between vowel and at the end of words, \\&
e.g., drihel $\rightarrow$ diel\\

\hline
5500 & ai $\rightarrow$ i / \_C, aa $\rightarrow$ a\textlengthmark\rightarrow a, ei 
$\rightarrow$ e \\

\hline
6000 & Nasal assimilation: \\& md $\rightarrow$ nd, np $\rightarrow$ mp, nk $\rightarrow$ \ng g \\

\hline
6500 & No consonant reduplication\\

\hline
7500 & Word-initial vowel loss: \\& \#V $\rightarrow$ \# \O \\

\hline
8000 & Vowel loss \texttheta Vr $\rightarrow$ \texttheta r unless stressed \\

\hline
8500 & No stops after nasals: \\& \ng k $\rightarrow$ \ng, mp $\rightarrow$ m, \ng g $\rightarrow$ g, etc. \\

\hline
8750 & No stops after any sonorant \\

\hline
9500 & Word-final vowel loss, unless stressed: \\& \#V $\rightarrow$ \O \\

\hline
10000 & ae $\rightarrow$ a \\

\hline
11000 & No coda stops: \\& e.g., \{p, b\} $\rightarrow$ m /V\_ \\

\hline
12000 & \textipa{\textyogh} $\rightarrow$ \textesh \\

\hline
12001 & \textipa{\dh} $\rightarrow$ \texttheta \\

\hline
12003, No repeated vowels\\
\hline
    12004 & No word-final e\\
\hline
    13000 & No fricative clusters\\
\hline
    14000 & No fricatives after affricates\\
\hline
   14001 & No affricates after fricatives\\
\end{tabular}
}


\newpage

\section{Phrases}
good morning

wishe-lhano yi-aray




\newpage
\section{Dictionary}
\twocolumn
\vspace{15pt}
\begin{nopagebreak}
\noindent{\fontsize{20pt}{10pt}\textbf{keluloresyaleth} } \textit{1 of autumn/September 1} (N)\\
\noindent {\tovian \fontsize{20pt}{10pt} \textbf{keluloresyaleth} }\\
\noindent /keluloresy{\textprimstress}ale{\texttheta}/\\
\noindent lit. branch+grain\\
\noindent \textit{Starrye}\\


\noindent History:

\vspace{-0pt}
\hspace{40pt}
\begin{tabular}{ccc}
\textit{6000} & /kelumlloresyale{\dh}a/&$\rightarrow$ & \textit{6500} & /kelumloresyale{\dh}a/&$\rightarrow$ & \textit{7600} & /keluloresyale{\dh}a/&$\rightarrow$ & \textit{9500} & /keluloresyale{\dh}/&$\rightarrow$ & \textit{12002} & /keluloresyale{\texttheta}/& \\
\end{tabular}

\vspace{20pt}\hline

\end{nopagebreak}
\filbreak



\vspace{15pt}
\begin{nopagebreak}
\noindent{\fontsize{20pt}{10pt}\textbf{ysklumlor} } \textit{1 of spring/March 1} (N)\\
\noindent {\tovian \fontsize{20pt}{10pt} \textbf{ysklumlor} }\\
\noindent /yskl{\textprimstress}umlor/\\
\noindent lit. moon+vine\\
\noindent \textit{Lunevigne}\\


\noindent History:

\vspace{-0pt}
\hspace{40pt}
\begin{tabular}{ccc}
\textit{6000} & /aliyskluwmllore/&$\rightarrow$ & \textit{6300} & /alyskluwmllore/&$\rightarrow$ & \textit{6500} & /alyskluwmlore/&$\rightarrow$ & \textit{7500} & /lyskluwmlore/&$\rightarrow$ & \textit{7501} & /yskluwmlore/&$\rightarrow$ & \textit{7600} & /ysklumlore/&$\rightarrow$ & \textit{9500} & /ysklumlor/& \\
\end{tabular}

\vspace{20pt}\hline

\end{nopagebreak}
\filbreak



\vspace{15pt}
\begin{nopagebreak}
\noindent{\fontsize{20pt}{10pt}\textbf{mloanam} } \textit{1 of summer/June 1} (N)\\
\noindent {\tovian \fontsize{20pt}{10pt} \textbf{mloanam} }\\
\noindent /mlo{\textprimstress}anam/\\
\noindent lit. morning+bloom\\
\noindent \textit{Dawnbloom}\\


\noindent History:

\vspace{-0pt}
\hspace{40pt}
\begin{tabular}{ccc}
\textit{6000} & /ploanadrmpfede/&$\rightarrow$ & \textit{7600} & /ploanadmpfede/&$\rightarrow$ & \textit{8500} & /ploanadmfede/&$\rightarrow$ & \textit{9500} & /ploanadmfed/&$\rightarrow$ & \textit{11000} & /mloanadmfed/&$\rightarrow$ & \textit{12003} & /mloanadmd/&$\rightarrow$ & \textit{13001} & /mloanadm/&$\rightarrow$ & \textit{13004} & /mloanam/& \\
\end{tabular}

\vspace{20pt}\hline

\end{nopagebreak}
\filbreak



\vspace{15pt}
\begin{nopagebreak}
\noindent{\fontsize{20pt}{10pt}\textbf{lrithdrem} } \textit{1 of winter/december 1} (N)\\
\noindent {\tovian \fontsize{20pt}{10pt} \textbf{lrithdrem} }\\
\noindent /lr{\textprimstress}i{\texttheta}drem/\\
\noindent lit. breath+flower\\
\noindent \textit{Rithwildrem}\\


\noindent History:

\vspace{-0pt}
\hspace{40pt}
\begin{tabular}{ccc}
\textit{6000} & /kwlyri{\texttheta}dremp/&$\rightarrow$ & \textit{7501} & /wlyri{\texttheta}dremp/&$\rightarrow$ & \textit{7600} & /lri{\texttheta}dremp/&$\rightarrow$ & \textit{8750} & /lri{\texttheta}drem/& \\
\end{tabular}

\vspace{20pt}\hline

\end{nopagebreak}
\filbreak



\vspace{15pt}
\begin{nopagebreak}
\noindent{\fontsize{20pt}{10pt}\textbf{nanathuv} } \textit{10 of autumn/September 10} (N)\\
\noindent {\tovian \fontsize{20pt}{10pt} \textbf{nanathuv} }\\
\noindent /nan{\textprimstress}a{\texttheta}uv/\\
\noindent lit. star+sword\\
\noindent \textit{Stareblad}\\


\noindent History:

\vspace{-0pt}
\hspace{40pt}
\begin{tabular}{ccc}
\textit{6000} & /nnaana{\dh}n{\textbeltl}uv/&$\rightarrow$ & \textit{6400} & /nnaana{\dh}{\textbeltl}uv/&$\rightarrow$ & \textit{6500} & /naana{\dh}{\textbeltl}uv/&$\rightarrow$ & \textit{12002} & /naana{\texttheta}{\textbeltl}uv/&$\rightarrow$ & \textit{12005} & /nana{\texttheta}{\textbeltl}uv/&$\rightarrow$ & \textit{13000} & /nana{\texttheta}uv/& \\
\end{tabular}

\vspace{20pt}\hline

\end{nopagebreak}
\filbreak



\vspace{15pt}
\begin{nopagebreak}
\noindent{\fontsize{20pt}{10pt}\textbf{todawen} } \textit{10 of spring/ March 10} (N)\\
\noindent {\tovian \fontsize{20pt}{10pt} \textbf{todawen} }\\
\noindent /{\textsubbridge{t}}od{\textprimstress}awen/\\
\noindent lit. eternal+beauty\\
\noindent \textit{Eternavial}\\


\noindent History:

\vspace{-0pt}
\hspace{40pt}
\begin{tabular}{ccc}
\textit{6000} & /{\textsubbridge{t}}odawen/& \\
\end{tabular}

\vspace{20pt}\hline

\end{nopagebreak}
\filbreak



\vspace{15pt}
\begin{nopagebreak}
\noindent{\fontsize{20pt}{10pt}\textbf{yethilemnan} } \textit{10 of summer/June 10} (N)\\
\noindent {\tovian \fontsize{20pt}{10pt} \textbf{yethilemnan} }\\
\noindent /ye{\texttheta}il{\textprimstress}emnan/\\
\noindent lit. home+protection\\
\noindent \textit{Nestkeep}\\


\noindent History:

\vspace{-0pt}
\hspace{40pt}
\begin{tabular}{ccc}
\textit{6000} & /dryel{\texttheta}ilebnani/&$\rightarrow$ & \textit{7501} & /ryel{\texttheta}ilebnani/&$\rightarrow$ & \textit{7600} & /yel{\texttheta}ilebnani/&$\rightarrow$ & \textit{9500} & /yel{\texttheta}ilebnan/&$\rightarrow$ & \textit{11000} & /yel{\texttheta}ilemnan/&$\rightarrow$ & \textit{12502} & /ye{\texttheta}ilemnan/& \\
\end{tabular}

\vspace{20pt}\hline

\end{nopagebreak}
\filbreak



\vspace{15pt}
\begin{nopagebreak}
\noindent{\fontsize{20pt}{10pt}\textbf{lriskwil} } \textit{10 of winter/december 10} (N)\\
\noindent {\tovian \fontsize{20pt}{10pt} \textbf{lriskwil} }\\
\noindent /lr{\textprimstress}iskwil/\\
\noindent lit. breath+movement\\
\noindent \textit{Galespad}\\


\noindent History:

\vspace{-0pt}
\hspace{40pt}
\begin{tabular}{ccc}
\textit{6000} & /kwlyri{\texttheta}kwil/&$\rightarrow$ & \textit{7501} & /wlyri{\texttheta}kwil/&$\rightarrow$ & \textit{7600} & /lri{\texttheta}kwil/&$\rightarrow$ & \textit{10500} & /lriskwil/& \\
\end{tabular}

\vspace{20pt}\hline

\end{nopagebreak}
\filbreak



\vspace{15pt}
\begin{nopagebreak}
\noindent{\fontsize{20pt}{10pt}\textbf{???} } \textit{11 of autumn/September 11} (N)\\
\noindent {\tovian \fontsize{20pt}{10pt} \textbf{???} }\\
\noindent /???/\\
\noindent lit. ???\\
\noindent \textit{Ciantro}\\


\noindent History:

\vspace{-0pt}
\hspace{40pt}
\begin{tabular}{ccc}
\textit{6000} & /???/& \\
\end{tabular}

\vspace{20pt}\hline

\end{nopagebreak}
\filbreak



\vspace{15pt}
\begin{nopagebreak}
\noindent{\fontsize{20pt}{10pt}\textbf{mralm} } \textit{11 of spring/ March 11} (N)\\
\noindent {\tovian \fontsize{20pt}{10pt} \textbf{mralm} }\\
\noindent /mr{\textprimstress}alm/\\
\noindent lit. darkness+bloom\\
\noindent \textit{Shadowbulb}\\


\noindent History:

\vspace{-0pt}
\hspace{40pt}
\begin{tabular}{ccc}
\textit{6000} & /braldrmpfede/&$\rightarrow$ & \textit{7600} & /braldmpfede/&$\rightarrow$ & \textit{8500} & /braldmfede/&$\rightarrow$ & \textit{8750} & /bralmfede/&$\rightarrow$ & \textit{9500} & /bralmfed/&$\rightarrow$ & \textit{11000} & /mralmfed/&$\rightarrow$ & \textit{12003} & /mralmd/&$\rightarrow$ & \textit{13001} & /mralm/& \\
\end{tabular}

\vspace{20pt}\hline

\end{nopagebreak}
\filbreak



\vspace{15pt}
\begin{nopagebreak}
\noindent{\fontsize{20pt}{10pt}\textbf{thuvglar} } \textit{11 of summer/June 11} (N)\\
\noindent {\tovian \fontsize{20pt}{10pt} \textbf{thuvglar} }\\
\noindent /{\texttheta}{\textprimstress}uvglar/\\
\noindent lit. sword+woods\\
\noindent \textit{Vekhangor}\\


\noindent History:

\vspace{-0pt}
\hspace{40pt}
\begin{tabular}{ccc}
\textit{6000} & /{\dh}n{\textbeltl}uvglar/&$\rightarrow$ & \textit{6400} & /{\dh}{\textbeltl}uvglar/&$\rightarrow$ & \textit{12002} & /{\texttheta}{\textbeltl}uvglar/&$\rightarrow$ & \textit{13000} & /{\texttheta}uvglar/& \\
\end{tabular}

\vspace{20pt}\hline

\end{nopagebreak}
\filbreak



\vspace{15pt}
\begin{nopagebreak}
\noindent{\fontsize{20pt}{10pt}\textbf{fedrlad} } \textit{11 of winter/december 11} (N)\\
\noindent {\tovian \fontsize{20pt}{10pt} \textbf{fedrlad} }\\
\noindent /f{\textprimstress}edrlad/\\
\noindent lit. showing+jealousy\\
\noindent \textit{Narcissus}\\


\noindent History:

\vspace{-0pt}
\hspace{40pt}
\begin{tabular}{ccc}
\textit{6000} & /fedeylrladl/&$\rightarrow$ & \textit{6300} & /fedylrladl/&$\rightarrow$ & \textit{6503} & /fedyrlladl/&$\rightarrow$ & \textit{7600} & /fedrladl/&$\rightarrow$ & \textit{13001} & /fedrlad/& \\
\end{tabular}

\vspace{20pt}\hline

\end{nopagebreak}
\filbreak



\vspace{15pt}
\begin{nopagebreak}
\noindent{\fontsize{20pt}{10pt}\textbf{thilethan} } \textit{12 of autumn/September 12} (N)\\
\noindent {\tovian \fontsize{20pt}{10pt} \textbf{thilethan} }\\
\noindent /{\texttheta}il{\textprimstress}e{\texttheta}an/\\
\noindent lit. place+tool\\
\noindent \textit{Earchoke}\\


\noindent History:

\vspace{-0pt}
\hspace{40pt}
\begin{tabular}{ccc}
\textit{6000} & /{\texttheta}ile{\dh}an/&$\rightarrow$ & \textit{12002} & /{\texttheta}ile{\texttheta}an/& \\
\end{tabular}

\vspace{20pt}\hline

\end{nopagebreak}
\filbreak



\vspace{15pt}
\begin{nopagebreak}
\noindent{\fontsize{20pt}{10pt}\textbf{mloelhanalor} } \textit{12 of spring/ March 12} (N)\\
\noindent {\tovian \fontsize{20pt}{10pt} \textbf{mloelhanalor} }\\
\noindent /mloe{\textbeltl}an{\textprimstress}alor/\\
\noindent lit. rebirth+plant\\
\noindent \textit{Féenixflour}\\


\noindent History:

\vspace{-0pt}
\hspace{40pt}
\begin{tabular}{ccc}
\textit{6000} & /ploe{\textbeltl}aanamllore/&$\rightarrow$ & \textit{6500} & /ploe{\textbeltl}aanamlore/&$\rightarrow$ & \textit{7600} & /ploe{\textbeltl}aanalore/&$\rightarrow$ & \textit{9500} & /ploe{\textbeltl}aanalor/&$\rightarrow$ & \textit{11000} & /mloe{\textbeltl}aanalor/&$\rightarrow$ & \textit{12005} & /mloe{\textbeltl}analor/& \\
\end{tabular}

\vspace{20pt}\hline

\end{nopagebreak}
\filbreak



\vspace{15pt}
\begin{nopagebreak}
\noindent{\fontsize{20pt}{10pt}\textbf{nador} } \textit{12 of summer/June 12} (N)\\
\noindent {\tovian \fontsize{20pt}{10pt} \textbf{nador} }\\
\noindent /n{\textprimstress}ador/\\
\noindent lit. light+blue\\
\noindent \textit{Gleamwoad}\\


\noindent History:

\vspace{-0pt}
\hspace{40pt}
\begin{tabular}{ccc}
\textit{6000} & /anadora/&$\rightarrow$ & \textit{7500} & /nadora/&$\rightarrow$ & \textit{9500} & /nador/& \\
\end{tabular}

\vspace{20pt}\hline

\end{nopagebreak}
\filbreak



\vspace{15pt}
\begin{nopagebreak}
\noindent{\fontsize{20pt}{10pt}\textbf{rithaglar} } \textit{12 of winter/december 12} (N)\\
\noindent {\tovian \fontsize{20pt}{10pt} \textbf{rithaglar} }\\
\noindent /ri{\texttheta}{\textprimstress}aglar/\\
\noindent lit. air+tree\\
\noindent \textit{Windelm}\\


\noindent History:

\vspace{-0pt}
\hspace{40pt}
\begin{tabular}{ccc}
\textit{6000} & /ri{\texttheta}aglare/&$\rightarrow$ & \textit{9500} & /ri{\texttheta}aglar/& \\
\end{tabular}

\vspace{20pt}\hline

\end{nopagebreak}
\filbreak



\vspace{15pt}
\begin{nopagebreak}
\noindent{\fontsize{20pt}{10pt}\textbf{shithan} } \textit{13 of autumn/September 13} (N)\\
\noindent {\tovian \fontsize{20pt}{10pt} \textbf{shithan} }\\
\noindent /{\textesh}{\textprimstress}i{\texttheta}an/\\
\noindent lit. sharpness+tool\\
\noindent \textit{Spicesta}\\


\noindent History:

\vspace{-0pt}
\hspace{40pt}
\begin{tabular}{ccc}
\textit{6000} & /{\textyogh}i{\texttheta}{\dh}an/&$\rightarrow$ & \textit{12001} & /{\textesh}i{\texttheta}{\dh}an/&$\rightarrow$ & \textit{12002} & /{\textesh}i{\texttheta}{\texttheta}an/&$\rightarrow$ & \textit{12007} & /{\textesh}i{\texttheta}an/& \\
\end{tabular}

\vspace{20pt}\hline

\end{nopagebreak}
\filbreak



\vspace{15pt}
\begin{nopagebreak}
\noindent{\fontsize{20pt}{10pt}\textbf{tushudol} } \textit{13 of spring/ March 13} (N)\\
\noindent {\tovian \fontsize{20pt}{10pt} \textbf{tushudol} }\\
\noindent /{\textsubbridge{t}}u{\textesh}{\textprimstress}udol/\\
\noindent lit. cold+fruit\\
\noindent \textit{Víntergourd}\\


\noindent History:

\vspace{-0pt}
\hspace{40pt}
\begin{tabular}{ccc}
\textit{6000} & /{\textsubbridge{t}}u{\textesh}u{\textyogh}ola/&$\rightarrow$ & \textit{9500} & /{\textsubbridge{t}}u{\textesh}u{\textyogh}ol/&$\rightarrow$ & \textit{12001} & /{\textsubbridge{t}}u{\textesh}u{\textesh}ol/&$\rightarrow$ & \textit{12501} & /{\textsubbridge{t}}u{\textesh}udol/& \\
\end{tabular}

\vspace{20pt}\hline

\end{nopagebreak}
\filbreak



\vspace{15pt}
\begin{nopagebreak}
\noindent{\fontsize{20pt}{10pt}\textbf{mralsyaleth} } \textit{13 of summer/June 13} (N)\\
\noindent {\tovian \fontsize{20pt}{10pt} \textbf{mralsyaleth} }\\
\noindent /mralsy{\textprimstress}ale{\texttheta}/\\
\noindent lit. darkness+food\\
\noindent \textit{Nightmush}\\


\noindent History:

\vspace{-0pt}
\hspace{40pt}
\begin{tabular}{ccc}
\textit{6000} & /bralsyale{\dh}a/&$\rightarrow$ & \textit{9500} & /bralsyale{\dh}/&$\rightarrow$ & \textit{11000} & /mralsyale{\dh}/&$\rightarrow$ & \textit{12002} & /mralsyale{\texttheta}/& \\
\end{tabular}

\vspace{20pt}\hline

\end{nopagebreak}
\filbreak



\vspace{15pt}
\begin{nopagebreak}
\noindent{\fontsize{20pt}{10pt}\textbf{shotlan} } \textit{13 of winter/december 13} (N)\\
\noindent {\tovian \fontsize{20pt}{10pt} \textbf{shotlan} }\\
\noindent /{\textesh}{\textprimstress}o{\texttoptiebar{t\textbeltl}}an/\\
\noindent lit. flowing+tool\\
\noindent \textit{Mistifen}\\


\noindent History:

\vspace{-0pt}
\hspace{40pt}
\begin{tabular}{ccc}
\textit{6000} & /{\textesh}o{\texttoptiebar{t\textbeltl}}{\dh}an/&$\rightarrow$ & \textit{12002} & /{\textesh}o{\texttoptiebar{t\textbeltl}}{\texttheta}an/&$\rightarrow$ & \textit{14000} & /{\textesh}o{\texttoptiebar{t\textbeltl}}an/& \\
\end{tabular}

\vspace{20pt}\hline

\end{nopagebreak}
\filbreak



\vspace{15pt}
\begin{nopagebreak}
\noindent{\fontsize{20pt}{10pt}\textbf{tlurinin} } \textit{14 of autumn/September 14} (N)\\
\noindent {\tovian \fontsize{20pt}{10pt} \textbf{tlurinin} }\\
\noindent /{\texttoptiebar{t\textbeltl}}ur{\textprimstress}inin/\\
\noindent lit. mist+???\\
\noindent \textit{Mistlaven}\\


\noindent History:

\vspace{-0pt}
\hspace{40pt}
\begin{tabular}{ccc}
\textit{6000} & /i{\ng}g{\texttoptiebar{t\textbeltl}}unri{\texttheta}nina???/&$\rightarrow$ & \textit{7500} & /{\ng}g{\texttoptiebar{t\textbeltl}}unri{\texttheta}nina???/&$\rightarrow$ & \textit{7501} & /g{\texttoptiebar{t\textbeltl}}unri{\texttheta}nina???/&$\rightarrow$ & \textit{7600} & /g{\texttoptiebar{t\textbeltl}}uri{\texttheta}nina???/&$\rightarrow$ & \textit{8500} & /g{\texttoptiebar{t\textbeltl}}uri{\texttheta}nina/&$\rightarrow$ & \textit{9500} & /g{\texttoptiebar{t\textbeltl}}uri{\texttheta}nin/&$\rightarrow$ & \textit{11501} & /{\texttoptiebar{t\textbeltl}}uri{\texttheta}nin/&$\rightarrow$ & \textit{13002} & /{\texttoptiebar{t\textbeltl}}urinin/& \\
\end{tabular}

\vspace{20pt}\hline

\end{nopagebreak}
\filbreak



\vspace{15pt}
\begin{nopagebreak}
\noindent{\fontsize{20pt}{10pt}\textbf{menan} } \textit{14 of spring/ March 14} (N)\\
\noindent {\tovian \fontsize{20pt}{10pt} \textbf{menan} }\\
\noindent /m{\textprimstress}enan/\\
\noindent lit. silver+???\\
\noindent \textit{Silvreséda}\\


\noindent History:

\vspace{-0pt}
\hspace{40pt}
\begin{tabular}{ccc}
\textit{6000} & /{\dh}nhmernana???/&$\rightarrow$ & \textit{6400} & /{\dh}hmernana???/&$\rightarrow$ & \textit{7501} & /hmernana???/&$\rightarrow$ & \textit{7600} & /hmenana???/&$\rightarrow$ & \textit{8500} & /hmenana/&$\rightarrow$ & \textit{9500} & /hmenan/&$\rightarrow$ & \textit{13002} & /menan/& \\
\end{tabular}

\vspace{20pt}\hline

\end{nopagebreak}
\filbreak



\vspace{15pt}
\begin{nopagebreak}
\noindent{\fontsize{20pt}{10pt}\textbf{lriskwil} } \textit{14 of summer/June 14} (N)\\
\noindent {\tovian \fontsize{20pt}{10pt} \textbf{lriskwil} }\\
\noindent /lr{\textprimstress}iskwil/\\
\noindent lit. wind+movement\\
\noindent \textit{Rithwil}\\


\noindent History:

\vspace{-0pt}
\hspace{40pt}
\begin{tabular}{ccc}
\textit{6000} & /kwlyri{\texttheta}kwil/&$\rightarrow$ & \textit{7501} & /wlyri{\texttheta}kwil/&$\rightarrow$ & \textit{7600} & /lri{\texttheta}kwil/&$\rightarrow$ & \textit{10500} & /lriskwil/& \\
\end{tabular}

\vspace{20pt}\hline

\end{nopagebreak}
\filbreak



\vspace{15pt}
\begin{nopagebreak}
\noindent{\fontsize{20pt}{10pt}\textbf{fedetush} } \textit{14 of winter/december 14} (N)\\
\noindent {\tovian \fontsize{20pt}{10pt} \textbf{fedetush} }\\
\noindent /fed{\textprimstress}e{\textsubbridge{t}}u{\textesh}/\\
\noindent lit. showing+cold\\
\noindent \textit{Galblom}\\


\noindent History:

\vspace{-0pt}
\hspace{40pt}
\begin{tabular}{ccc}
\textit{6000} & /fede{\textsubbridge{t}}u{\textesh}u/&$\rightarrow$ & \textit{9500} & /fede{\textsubbridge{t}}u{\textesh}/& \\
\end{tabular}

\vspace{20pt}\hline

\end{nopagebreak}
\filbreak



\vspace{15pt}
\begin{nopagebreak}
\noindent{\fontsize{20pt}{10pt}\textbf{???} } \textit{15 of autumn/September 15} (N)\\
\noindent {\tovian \fontsize{20pt}{10pt} \textbf{???} }\\
\noindent /???/\\
\noindent lit. ???\\
\noindent \textit{Chamois}\\


\noindent History:

\vspace{-0pt}
\hspace{40pt}
\begin{tabular}{ccc}
\textit{6000} & /???/& \\
\end{tabular}

\vspace{20pt}\hline

\end{nopagebreak}
\filbreak



\vspace{15pt}
\begin{nopagebreak}
\noindent{\fontsize{20pt}{10pt}\textbf{rithilekel} } \textit{15 of spring/ March 15} (N)\\
\noindent {\tovian \fontsize{20pt}{10pt} \textbf{rithilekel} }\\
\noindent /ri{\texttheta}il{\textprimstress}ekel/\\
\noindent lit. sky+branch\\
\noindent \textit{Skyaspe}\\


\noindent History:

\vspace{-0pt}
\hspace{40pt}
\begin{tabular}{ccc}
\textit{6000} & /ri{\texttheta}ilekelu/&$\rightarrow$ & \textit{9500} & /ri{\texttheta}ilekel/& \\
\end{tabular}

\vspace{20pt}\hline

\end{nopagebreak}
\filbreak



\vspace{15pt}
\begin{nopagebreak}
\noindent{\fontsize{20pt}{10pt}\textbf{forengothr} } \textit{15 of summer/June 15} (N)\\
\noindent {\tovian \fontsize{20pt}{10pt} \textbf{forengothr} }\\
\noindent /for{\textprimstress}e{\ng}o{\texttheta}r/\\
\noindent lit. writing+time\\
\noindent \textit{Felor Day}\\


\noindent History:

\vspace{-0pt}
\hspace{40pt}
\begin{tabular}{ccc}
\textit{6000} & /flore{\ng}go{\texttheta}r/&$\rightarrow$ & \textit{8750} & /flore{\ng}o{\texttheta}r/&$\rightarrow$ & \textit{12500} & /fore{\ng}o{\texttheta}r/& \\
\end{tabular}

\vspace{20pt}\hline

\end{nopagebreak}
\filbreak



\vspace{15pt}
\begin{nopagebreak}
\noindent{\fontsize{20pt}{10pt}\textbf{lrithen} } \textit{15 of winter/december 15} (N)\\
\noindent {\tovian \fontsize{20pt}{10pt} \textbf{lrithen} }\\
\noindent /lr{\textprimstress}i{\texttheta}en/\\
\noindent lit. breath+dream\\
\noindent \textit{Windrith}\\


\noindent History:

\vspace{-0pt}
\hspace{40pt}
\begin{tabular}{ccc}
\textit{6000} & /kwlyri{\texttheta}{\textbeltl}ena/&$\rightarrow$ & \textit{7501} & /wlyri{\texttheta}{\textbeltl}ena/&$\rightarrow$ & \textit{7600} & /lri{\texttheta}{\textbeltl}ena/&$\rightarrow$ & \textit{9500} & /lri{\texttheta}{\textbeltl}en/&$\rightarrow$ & \textit{13000} & /lri{\texttheta}en/& \\
\end{tabular}

\vspace{20pt}\hline

\end{nopagebreak}
\filbreak



\vspace{15pt}
\begin{nopagebreak}
\noindent{\fontsize{20pt}{10pt}\textbf{rithol} } \textit{16 of autumn/September 16} (N)\\
\noindent {\tovian \fontsize{20pt}{10pt} \textbf{rithol} }\\
\noindent /r{\textprimstress}i{\texttheta}ol/\\
\noindent lit. air+sweet\\
\noindent \textit{Smokesweet}\\


\noindent History:

\vspace{-0pt}
\hspace{40pt}
\begin{tabular}{ccc}
\textit{6000} & /ri{\texttheta}hola/&$\rightarrow$ & \textit{9500} & /ri{\texttheta}hol/&$\rightarrow$ & \textit{13000} & /ri{\texttheta}ol/& \\
\end{tabular}

\vspace{20pt}\hline

\end{nopagebreak}
\filbreak



\vspace{15pt}
\begin{nopagebreak}
\noindent{\fontsize{20pt}{10pt}\textbf{mrawen} } \textit{16 of spring/ March 16} (N)\\
\noindent {\tovian \fontsize{20pt}{10pt} \textbf{mrawen} }\\
\noindent /mr{\textprimstress}awen/\\
\noindent lit. darkness+beauty\\
\noindent \textit{Nightbelle}\\


\noindent History:

\vspace{-0pt}
\hspace{40pt}
\begin{tabular}{ccc}
\textit{6000} & /bralwen/&$\rightarrow$ & \textit{7600} & /brawen/&$\rightarrow$ & \textit{11000} & /mrawen/& \\
\end{tabular}

\vspace{20pt}\hline

\end{nopagebreak}
\filbreak



\vspace{15pt}
\begin{nopagebreak}
\noindent{\fontsize{20pt}{10pt}\textbf{semlor} } \textit{16 of summer/June 16} (N)\\
\noindent {\tovian \fontsize{20pt}{10pt} \textbf{semlor} }\\
\noindent /s{\textprimstress}emlor/\\
\noindent lit. water+plant\\
\noindent \textit{Dewlettuce}\\


\noindent History:

\vspace{-0pt}
\hspace{40pt}
\begin{tabular}{ccc}
\textit{6000} & /zelmllore/&$\rightarrow$ & \textit{6500} & /zelmlore/&$\rightarrow$ & \textit{7600} & /zemlore/&$\rightarrow$ & \textit{9500} & /zemlor/&$\rightarrow$ & \textit{12000} & /semlor/& \\
\end{tabular}

\vspace{20pt}\hline

\end{nopagebreak}
\filbreak



\vspace{15pt}
\begin{nopagebreak}
\noindent{\fontsize{20pt}{10pt}\textbf{loreayeol} } \textit{16 of winter/december 16} (N)\\
\noindent {\tovian \fontsize{20pt}{10pt} \textbf{loreayeol} }\\
\noindent /loreay{\textprimstress}eol/\\
\noindent lit. dance+wisdom\\
\noindent \textit{Dansinath}\\


\noindent History:

\vspace{-0pt}
\hspace{40pt}
\begin{tabular}{ccc}
\textit{6000} & /kwllorearyeol/&$\rightarrow$ & \textit{6500} & /kwlorearyeol/&$\rightarrow$ & \textit{7501} & /wlorearyeol/&$\rightarrow$ & \textit{7600} & /loreayeol/& \\
\end{tabular}

\vspace{20pt}\hline

\end{nopagebreak}
\filbreak



\vspace{15pt}
\begin{nopagebreak}
\noindent{\fontsize{20pt}{10pt}\textbf{sholalhum} } \textit{17 of autumn/September 17} (N)\\
\noindent {\tovian \fontsize{20pt}{10pt} \textbf{sholalhum} }\\
\noindent /{\textesh}ol{\textprimstress}a{\textbeltl}um/\\
\noindent lit. fruit+goodness\\
\noindent \textit{Curranthum}\\


\noindent History:

\vspace{-0pt}
\hspace{40pt}
\begin{tabular}{ccc}
\textit{6000} & /{\textyogh}ola{\textbeltl}um/&$\rightarrow$ & \textit{12001} & /{\textesh}ola{\textbeltl}um/& \\
\end{tabular}

\vspace{20pt}\hline

\end{nopagebreak}
\filbreak



\vspace{15pt}
\begin{nopagebreak}
\noindent{\fontsize{20pt}{10pt}\textbf{tlafor} } \textit{17 of spring/ March 17} (N)\\
\noindent {\tovian \fontsize{20pt}{10pt} \textbf{tlafor} }\\
\noindent /{\texttoptiebar{t\textbeltl}}{\textprimstress}afor/\\
\noindent lit. peace+plant\\
\noindent \textit{Tlafen}\\


\noindent History:

\vspace{-0pt}
\hspace{40pt}
\begin{tabular}{ccc}
\textit{6000} & /{\texttoptiebar{t\textbeltl}}afmllore/&$\rightarrow$ & \textit{6500} & /{\texttoptiebar{t\textbeltl}}afmlore/&$\rightarrow$ & \textit{7600} & /{\texttoptiebar{t\textbeltl}}aflore/&$\rightarrow$ & \textit{9500} & /{\texttoptiebar{t\textbeltl}}aflor/&$\rightarrow$ & \textit{12500} & /{\texttoptiebar{t\textbeltl}}afor/& \\
\end{tabular}

\vspace{20pt}\hline

\end{nopagebreak}
\filbreak



\vspace{15pt}
\begin{nopagebreak}
\noindent{\fontsize{20pt}{10pt}\textbf{gaglar} } \textit{17 of summer/June 17} (N)\\
\noindent {\tovian \fontsize{20pt}{10pt} \textbf{gaglar} }\\
\noindent /g{\textprimstress}aglar/\\
\noindent lit. warmth+tree\\
\noindent \textit{Frostlarch}\\


\noindent History:

\vspace{-0pt}
\hspace{40pt}
\begin{tabular}{ccc}
\textit{6000} & /agaaglare/&$\rightarrow$ & \textit{7500} & /gaaglare/&$\rightarrow$ & \textit{9500} & /gaaglar/&$\rightarrow$ & \textit{12005} & /gaglar/& \\
\end{tabular}

\vspace{20pt}\hline

\end{nopagebreak}
\filbreak



\vspace{15pt}
\begin{nopagebreak}
\noindent{\fontsize{20pt}{10pt}\textbf{lanam} } \textit{17 of winter/december 17} (N)\\
\noindent {\tovian \fontsize{20pt}{10pt} \textbf{lanam} }\\
\noindent /l{\textprimstress}anam/\\
\noindent lit. twilight+bloom\\
\noindent \textit{Twilibloom}\\


\noindent History:

\vspace{-0pt}
\hspace{40pt}
\begin{tabular}{ccc}
\textit{6000} & /brlanadrmpfede/&$\rightarrow$ & \textit{7501} & /rlanadrmpfede/&$\rightarrow$ & \textit{7600} & /lanadmpfede/&$\rightarrow$ & \textit{8500} & /lanadmfede/&$\rightarrow$ & \textit{9500} & /lanadmfed/&$\rightarrow$ & \textit{12003} & /lanadmd/&$\rightarrow$ & \textit{13001} & /lanadm/&$\rightarrow$ & \textit{13004} & /lanam/& \\
\end{tabular}

\vspace{20pt}\hline

\end{nopagebreak}
\filbreak



\vspace{15pt}
\begin{nopagebreak}
\noindent{\fontsize{20pt}{10pt}\textbf{sholaglar} } \textit{18 of autumn/September 18} (N)\\
\noindent {\tovian \fontsize{20pt}{10pt} \textbf{sholaglar} }\\
\noindent /{\textesh}ol{\textprimstress}aglar/\\
\noindent lit. fruit+tree\\
\noindent \textit{Peawil}\\


\noindent History:

\vspace{-0pt}
\hspace{40pt}
\begin{tabular}{ccc}
\textit{6000} & /{\textyogh}olaaglare/&$\rightarrow$ & \textit{9500} & /{\textyogh}olaaglar/&$\rightarrow$ & \textit{12001} & /{\textesh}olaaglar/&$\rightarrow$ & \textit{12005} & /{\textesh}olaglar/& \\
\end{tabular}

\vspace{20pt}\hline

\end{nopagebreak}
\filbreak



\vspace{15pt}
\begin{nopagebreak}
\noindent{\fontsize{20pt}{10pt}\textbf{nanamel} } \textit{18 of spring/ March 18} (N)\\
\noindent {\tovian \fontsize{20pt}{10pt} \textbf{nanamel} }\\
\noindent /nan{\textprimstress}amel/\\
\noindent lit. star+life\\
\noindent \textit{Starzin}\\


\noindent History:

\vspace{-0pt}
\hspace{40pt}
\begin{tabular}{ccc}
\textit{6000} & /nnaanamel/&$\rightarrow$ & \textit{6500} & /naanamel/&$\rightarrow$ & \textit{12005} & /nanamel/& \\
\end{tabular}

\vspace{20pt}\hline

\end{nopagebreak}
\filbreak



\vspace{15pt}
\begin{nopagebreak}
\noindent{\fontsize{20pt}{10pt}\textbf{mrallhenadrem} } \textit{18 of summer/June 18} (N)\\
\noindent {\tovian \fontsize{20pt}{10pt} \textbf{mrallhenadrem} }\\
\noindent /mral{\textbeltl}en{\textprimstress}adrem/\\
\noindent lit. darkness+dream+flower\\
\noindent \textit{Darkbloom}\\


\noindent History:

\vspace{-0pt}
\hspace{40pt}
\begin{tabular}{ccc}
\textit{6000} & /bral{\textbeltl}enadremp/&$\rightarrow$ & \textit{8750} & /bral{\textbeltl}enadrem/&$\rightarrow$ & \textit{11000} & /mral{\textbeltl}enadrem/& \\
\end{tabular}

\vspace{20pt}\hline

\end{nopagebreak}
\filbreak



\vspace{15pt}
\begin{nopagebreak}
\noindent{\fontsize{20pt}{10pt}\textbf{nanalor} } \textit{18 of winter/december 18} (N)\\
\noindent {\tovian \fontsize{20pt}{10pt} \textbf{nanalor} }\\
\noindent /nan{\textprimstress}alor/\\
\noindent lit. star+plant\\
\noindent \textit{Starsweed}\\


\noindent History:

\vspace{-0pt}
\hspace{40pt}
\begin{tabular}{ccc}
\textit{6000} & /nnaanamllore/&$\rightarrow$ & \textit{6500} & /naanamlore/&$\rightarrow$ & \textit{7600} & /naanalore/&$\rightarrow$ & \textit{9500} & /naanalor/&$\rightarrow$ & \textit{12005} & /nanalor/& \\
\end{tabular}

\vspace{20pt}\hline

\end{nopagebreak}
\filbreak



\vspace{15pt}
\begin{nopagebreak}
\noindent{\fontsize{20pt}{10pt}\textbf{lshanin} } \textit{19 of autumn/September 19} (N)\\
\noindent {\tovian \fontsize{20pt}{10pt} \textbf{lshanin} }\\
\noindent /l{\textesh}{\textprimstress}anin/\\
\noindent lit. sun+berry\\
\noindent \textit{Sunberry}\\


\noindent History:

\vspace{-0pt}
\hspace{40pt}
\begin{tabular}{ccc}
\textit{6000} & /efl{\textyogh}lanina/&$\rightarrow$ & \textit{7500} & /fl{\textyogh}lanina/&$\rightarrow$ & \textit{7501} & /l{\textyogh}lanina/&$\rightarrow$ & \textit{9500} & /l{\textyogh}lanin/&$\rightarrow$ & \textit{12001} & /l{\textesh}lanin/&$\rightarrow$ & \textit{12500} & /l{\textesh}anin/& \\
\end{tabular}

\vspace{20pt}\hline

\end{nopagebreak}
\filbreak



\vspace{15pt}
\begin{nopagebreak}
\noindent{\fontsize{20pt}{10pt}\textbf{efem} } \textit{19 of spring/ March 19} (N)\\
\noindent {\tovian \fontsize{20pt}{10pt} \textbf{efem} }\\
\noindent /{\textprimstress}efem/\\
\noindent lit. sun+flower\\
\noindent \textit{Soldrem}\\


\noindent History:

\vspace{-0pt}
\hspace{40pt}
\begin{tabular}{ccc}
\textit{6000} & /efldremp/&$\rightarrow$ & \textit{8750} & /eflrem/&$\rightarrow$ & \textit{12500} & /efem/& \\
\end{tabular}

\vspace{20pt}\hline

\end{nopagebreak}
\filbreak



\vspace{15pt}
\begin{nopagebreak}
\noindent{\fontsize{20pt}{10pt}\textbf{tomomen} } \textit{19 of summer/June 19} (N)\\
\noindent {\tovian \fontsize{20pt}{10pt} \textbf{tomomen} }\\
\noindent /{\textsubbridge{t}}om{\textprimstress}omen/\\
\noindent lit. root+earth\\
\noindent \textit{Starroot}\\


\noindent History:

\vspace{-0pt}
\hspace{40pt}
\begin{tabular}{ccc}
\textit{6000} & /{\textsubbridge{t}}omomern/&$\rightarrow$ & \textit{7600} & /{\textsubbridge{t}}omomen/& \\
\end{tabular}

\vspace{20pt}\hline

\end{nopagebreak}
\filbreak



\vspace{15pt}
\begin{nopagebreak}
\noindent{\fontsize{20pt}{10pt}\textbf{ayswen} } \textit{19 of winter/december 19} (N)\\
\noindent {\tovian \fontsize{20pt}{10pt} \textbf{ayswen} }\\
\noindent /{\textprimstress}ayswen/\\
\noindent lit. moon+beauty\\
\noindent \textit{Mooncheris}\\


\noindent History:

\vspace{-0pt}
\hspace{40pt}
\begin{tabular}{ccc}
\textit{6000} & /aliyswen/&$\rightarrow$ & \textit{6300} & /alyswen/&$\rightarrow$ & \textit{7600} & /ayswen/& \\
\end{tabular}

\vspace{20pt}\hline

\end{nopagebreak}
\filbreak



\vspace{15pt}
\begin{nopagebreak}
\noindent{\fontsize{20pt}{10pt}\textbf{ysoresyaleth} } \textit{2 of autumn/September 2} (N)\\
\noindent {\tovian \fontsize{20pt}{10pt} \textbf{ysoresyaleth} }\\
\noindent /ysoresy{\textprimstress}ale{\texttheta}/\\
\noindent lit. moon+grain\\
\noindent \textit{Moonoat}\\


\noindent History:

\vspace{-0pt}
\hspace{40pt}
\begin{tabular}{ccc}
\textit{6000} & /aliysmlloresyale{\dh}a/&$\rightarrow$ & \textit{6300} & /alysmlloresyale{\dh}a/&$\rightarrow$ & \textit{6500} & /alysmloresyale{\dh}a/&$\rightarrow$ & \textit{7500} & /lysmloresyale{\dh}a/&$\rightarrow$ & \textit{7501} & /ysmloresyale{\dh}a/&$\rightarrow$ & \textit{7600} & /ysloresyale{\dh}a/&$\rightarrow$ & \textit{9500} & /ysloresyale{\dh}/&$\rightarrow$ & \textit{12002} & /ysloresyale{\texttheta}/&$\rightarrow$ & \textit{12500} & /ysoresyale{\texttheta}/& \\
\end{tabular}

\vspace{20pt}\hline

\end{nopagebreak}
\filbreak



\vspace{15pt}
\begin{nopagebreak}
\noindent{\fontsize{20pt}{10pt}\textbf{efith} } \textit{2 of spring/March 2} (N)\\
\noindent {\tovian \fontsize{20pt}{10pt} \textbf{efith} }\\
\noindent /{\textprimstress}efi{\texttheta}/\\
\noindent lit. sun+sharpness\\
\noindent \textit{Solspica}\\


\noindent History:

\vspace{-0pt}
\hspace{40pt}
\begin{tabular}{ccc}
\textit{6000} & /efl{\textyogh}i{\texttheta}/&$\rightarrow$ & \textit{12001} & /efl{\textesh}i{\texttheta}/&$\rightarrow$ & \textit{12500} & /ef{\textesh}i{\texttheta}/&$\rightarrow$ & \textit{13000} & /efi{\texttheta}/& \\
\end{tabular}

\vspace{20pt}\hline

\end{nopagebreak}
\filbreak



\vspace{15pt}
\begin{nopagebreak}
\noindent{\fontsize{20pt}{10pt}\textbf{rithileaglar} } \textit{2 of summer/June 2} (N)\\
\noindent {\tovian \fontsize{20pt}{10pt} \textbf{rithileaglar} }\\
\noindent /ri{\texttheta}ile{\textprimstress}aglar/\\
\noindent lit. sky+tree\\
\noindent \textit{Skybark}\\


\noindent History:

\vspace{-0pt}
\hspace{40pt}
\begin{tabular}{ccc}
\textit{6000} & /ri{\texttheta}ileaglare/&$\rightarrow$ & \textit{9500} & /ri{\texttheta}ileaglar/& \\
\end{tabular}

\vspace{20pt}\hline

\end{nopagebreak}
\filbreak



\vspace{15pt}
\begin{nopagebreak}
\noindent{\fontsize{20pt}{10pt}\textbf{dolshanin} } \textit{2 of winter/december 2} (N)\\
\noindent {\tovian \fontsize{20pt}{10pt} \textbf{dolshanin} }\\
\noindent /dol{\textesh}{\textprimstress}anin/\\
\noindent lit. death+berry\\
\noindent \textit{Galbery}\\


\noindent History:

\vspace{-0pt}
\hspace{40pt}
\begin{tabular}{ccc}
\textit{6000} & /dol{\textyogh}lanina/&$\rightarrow$ & \textit{9500} & /dol{\textyogh}lanin/&$\rightarrow$ & \textit{12001} & /dol{\textesh}lanin/&$\rightarrow$ & \textit{12500} & /dol{\textesh}anin/& \\
\end{tabular}

\vspace{20pt}\hline

\end{nopagebreak}
\filbreak



\vspace{15pt}
\begin{nopagebreak}
\noindent{\fontsize{20pt}{10pt}\textbf{tilethil} } \textit{20 of autumn/September 20} (N)\\
\noindent {\tovian \fontsize{20pt}{10pt} \textbf{tilethil} }\\
\noindent /{\textsubbridge{t}}il{\textprimstress}e{\texttheta}il/\\
\noindent lit. plain+place\\
\noindent \textit{Meadowrealm}\\


\noindent History:

\vspace{-0pt}
\hspace{40pt}
\begin{tabular}{ccc}
\textit{6000} & /{\textsubbridge{t}}n{\texttheta}ile{\texttheta}ile/&$\rightarrow$ & \textit{6400} & /{\textsubbridge{t}}{\texttheta}ile{\texttheta}ile/&$\rightarrow$ & \textit{9500} & /{\textsubbridge{t}}{\texttheta}ile{\texttheta}il/&$\rightarrow$ & \textit{10700} & /{\textsubbridge{t}}{\textsubbridge{t}}ile{\texttheta}il/&$\rightarrow$ & \textit{12004} & /{\textsubbridge{t}}ile{\texttheta}il/& \\
\end{tabular}

\vspace{20pt}\hline

\end{nopagebreak}
\filbreak



\vspace{15pt}
\begin{nopagebreak}
\noindent{\fontsize{20pt}{10pt}\textbf{nanashith} } \textit{20 of spring/March 20} (N)\\
\noindent {\tovian \fontsize{20pt}{10pt} \textbf{nanashith} }\\
\noindent /nan{\textprimstress}a{\textesh}i{\texttheta}/\\
\noindent lit. star+sharpness\\
\noindent \textit{Starpress}\\


\noindent History:

\vspace{-0pt}
\hspace{40pt}
\begin{tabular}{ccc}
\textit{6000} & /nnaana{\textyogh}i{\texttheta}/&$\rightarrow$ & \textit{6500} & /naana{\textyogh}i{\texttheta}/&$\rightarrow$ & \textit{12001} & /naana{\textesh}i{\texttheta}/&$\rightarrow$ & \textit{12005} & /nana{\textesh}i{\texttheta}/& \\
\end{tabular}

\vspace{20pt}\hline

\end{nopagebreak}
\filbreak



\vspace{15pt}
\begin{nopagebreak}
\noindent{\fontsize{20pt}{10pt}\textbf{holadyethil} } \textit{20 of summer/June 20} (N)\\
\noindent {\tovian \fontsize{20pt}{10pt} \textbf{holadyethil} }\\
\noindent /ħolady{\textprimstress}e{\texttheta}il/\\
\noindent lit. sweet+home\\
\noindent \textit{Honeyhall}\\


\noindent History:

\vspace{-0pt}
\hspace{40pt}
\begin{tabular}{ccc}
\textit{6000} & /holadryel{\texttheta}ile/&$\rightarrow$ & \textit{7600} & /holadyel{\texttheta}ile/&$\rightarrow$ & \textit{9500} & /holadyel{\texttheta}il/&$\rightarrow$ & \textit{12502} & /holadye{\texttheta}il/&$\rightarrow$ & \textit{14006} & /ħoladye{\texttheta}il/& \\
\end{tabular}

\vspace{20pt}\hline

\end{nopagebreak}
\filbreak



\vspace{15pt}
\begin{nopagebreak}
\noindent{\fontsize{20pt}{10pt}\textbf{lrithun} } \textit{20 of winter/december 20} (N)\\
\noindent {\tovian \fontsize{20pt}{10pt} \textbf{lrithun} }\\
\noindent /lr{\textprimstress}i{\texttheta}un/\\
\noindent lit. breath+joining\\
\noindent \textit{Windcord}\\


\noindent History:

\vspace{-0pt}
\hspace{40pt}
\begin{tabular}{ccc}
\textit{6000} & /kwlyri{\texttheta}{\texttoptiebar{t\textbeltl}}una/&$\rightarrow$ & \textit{7501} & /wlyri{\texttheta}{\texttoptiebar{t\textbeltl}}una/&$\rightarrow$ & \textit{7600} & /lri{\texttheta}{\texttoptiebar{t\textbeltl}}una/&$\rightarrow$ & \textit{9500} & /lri{\texttheta}{\texttoptiebar{t\textbeltl}}un/&$\rightarrow$ & \textit{14001} & /lri{\texttheta}un/& \\
\end{tabular}

\vspace{20pt}\hline

\end{nopagebreak}
\filbreak



\vspace{15pt}
\begin{nopagebreak}
\noindent{\fontsize{20pt}{10pt}\textbf{shush} } \textit{20 of winter/January 20} (N)\\
\noindent {\tovian \fontsize{20pt}{10pt} \textbf{shush} }\\
\noindent /{\textesh}{\textprimstress}u{\textesh}/\\
\noindent lit. ice+???\\
\noindent \textit{Icetole}\\


\noindent History:

\vspace{-0pt}
\hspace{40pt}
\begin{tabular}{ccc}
\textit{6000} & /{\textesh}u{\textesh}a???/&$\rightarrow$ & \textit{8500} & /{\textesh}u{\textesh}a/&$\rightarrow$ & \textit{9500} & /{\textesh}u{\textesh}/& \\
\end{tabular}

\vspace{20pt}\hline

\end{nopagebreak}
\filbreak



\vspace{15pt}
\begin{nopagebreak}
\noindent{\fontsize{20pt}{10pt}\textbf{riskel} } \textit{21 of autumn/September 21} (N)\\
\noindent {\tovian \fontsize{20pt}{10pt} \textbf{riskel} }\\
\noindent /r{\textprimstress}iskel/\\
\noindent lit. air+branch\\
\noindent \textit{Breezemint}\\


\noindent History:

\vspace{-0pt}
\hspace{40pt}
\begin{tabular}{ccc}
\textit{6000} & /ri{\texttheta}kelu/&$\rightarrow$ & \textit{9500} & /ri{\texttheta}kel/&$\rightarrow$ & \textit{10500} & /riskel/& \\
\end{tabular}

\vspace{20pt}\hline

\end{nopagebreak}
\filbreak



\vspace{15pt}
\begin{nopagebreak}
\noindent{\fontsize{20pt}{10pt}\textbf{dremlor} } \textit{21 of spring/March 21} (N)\\
\noindent {\tovian \fontsize{20pt}{10pt} \textbf{dremlor} }\\
\noindent /dr{\textprimstress}emlor/\\
\noindent lit. flower+plant\\
\noindent \textit{Dremhemp}\\


\noindent History:

\vspace{-0pt}
\hspace{40pt}
\begin{tabular}{ccc}
\textit{6000} & /drempmllore/&$\rightarrow$ & \textit{6500} & /drempmlore/&$\rightarrow$ & \textit{7600} & /dremplore/&$\rightarrow$ & \textit{8500} & /dremlore/&$\rightarrow$ & \textit{9500} & /dremlor/& \\
\end{tabular}

\vspace{20pt}\hline

\end{nopagebreak}
\filbreak



\vspace{15pt}
\begin{nopagebreak}
\noindent{\fontsize{20pt}{10pt}\textbf{meloremnan} } \textit{21 of summer/June 21} (N)\\
\noindent {\tovian \fontsize{20pt}{10pt} \textbf{meloremnan} }\\
\noindent /melor{\textprimstress}emnan/\\
\noindent lit. plant+protection\\
\noindent \textit{Seedguard}\\


\noindent History:

\vspace{-0pt}
\hspace{40pt}
\begin{tabular}{ccc}
\textit{6000} & /mllorebnani/&$\rightarrow$ & \textit{6500} & /mlorebnani/&$\rightarrow$ & \textit{7502} & /melorebnani/&$\rightarrow$ & \textit{9500} & /melorebnan/&$\rightarrow$ & \textit{11000} & /meloremnan/& \\
\end{tabular}

\vspace{20pt}\hline

\end{nopagebreak}
\filbreak



\vspace{15pt}
\begin{nopagebreak}
\noindent{\fontsize{20pt}{10pt}\textbf{lhilaglar} } \textit{21 of winter/december 21} (N)\\
\noindent {\tovian \fontsize{20pt}{10pt} \textbf{lhilaglar} }\\
\noindent /{\textbeltl}il{\textprimstress}aglar/\\
\noindent lit. night+tree\\
\noindent \textit{Nightroot}\\


\noindent History:

\vspace{-0pt}
\hspace{40pt}
\begin{tabular}{ccc}
\textit{6000} & /a{\textbeltl}inlaaglare/&$\rightarrow$ & \textit{7500} & /{\textbeltl}inlaaglare/&$\rightarrow$ & \textit{7600} & /{\textbeltl}ilaaglare/&$\rightarrow$ & \textit{9500} & /{\textbeltl}ilaaglar/&$\rightarrow$ & \textit{12005} & /{\textbeltl}ilaglar/& \\
\end{tabular}

\vspace{20pt}\hline

\end{nopagebreak}
\filbreak



\vspace{15pt}
\begin{nopagebreak}
\noindent{\fontsize{20pt}{10pt}\textbf{mraltlun} } \textit{22 of autumn/September 22} (N)\\
\noindent {\tovian \fontsize{20pt}{10pt} \textbf{mraltlun} }\\
\noindent /mr{\textprimstress}al{\texttoptiebar{t\textbeltl}}un/\\
\noindent lit. darkness+joining\\
\noindent \textit{Nightspice}\\


\noindent History:

\vspace{-0pt}
\hspace{40pt}
\begin{tabular}{ccc}
\textit{6000} & /bral{\texttoptiebar{t\textbeltl}}una/&$\rightarrow$ & \textit{9500} & /bral{\texttoptiebar{t\textbeltl}}un/&$\rightarrow$ & \textit{11000} & /mral{\texttoptiebar{t\textbeltl}}un/& \\
\end{tabular}

\vspace{20pt}\hline

\end{nopagebreak}
\filbreak



\vspace{15pt}
\begin{nopagebreak}
\noindent{\fontsize{20pt}{10pt}\textbf{rithiledrem} } \textit{22 of spring/March 22} (N)\\
\noindent {\tovian \fontsize{20pt}{10pt} \textbf{rithiledrem} }\\
\noindent /ri{\texttheta}il{\textprimstress}edrem/\\
\noindent lit. sky+flower\\
\noindent \textit{Skypech}\\


\noindent History:

\vspace{-0pt}
\hspace{40pt}
\begin{tabular}{ccc}
\textit{6000} & /ri{\texttheta}iledremp/&$\rightarrow$ & \textit{8750} & /ri{\texttheta}iledrem/& \\
\end{tabular}

\vspace{20pt}\hline

\end{nopagebreak}
\filbreak



\vspace{15pt}
\begin{nopagebreak}
\noindent{\fontsize{20pt}{10pt}\textbf{shotlathang} } \textit{22 of summer/June 22} (N)\\
\noindent {\tovian \fontsize{20pt}{10pt} \textbf{shotlathang} }\\
\noindent /{\textesh}o{\texttoptiebar{t\textbeltl}}{\textprimstress}a{\texttheta}a{\ng}/\\
\noindent lit. flowing+leaf\\
\noindent \textit{Mistleaf}\\


\noindent History:

\vspace{-0pt}
\hspace{40pt}
\begin{tabular}{ccc}
\textit{6000} & /{\textesh}o{\texttoptiebar{t\textbeltl}}a{\dh}a{\ng}ge/&$\rightarrow$ & \textit{8750} & /{\textesh}o{\texttoptiebar{t\textbeltl}}a{\dh}a{\ng}e/&$\rightarrow$ & \textit{9500} & /{\textesh}o{\texttoptiebar{t\textbeltl}}a{\dh}a{\ng}/&$\rightarrow$ & \textit{12002} & /{\textesh}o{\texttoptiebar{t\textbeltl}}a{\texttheta}a{\ng}/& \\
\end{tabular}

\vspace{20pt}\hline

\end{nopagebreak}
\filbreak



\vspace{15pt}
\begin{nopagebreak}
\noindent{\fontsize{20pt}{10pt}\textbf{lan} } \textit{22 of winter/december 22} (N)\\
\noindent {\tovian \fontsize{20pt}{10pt} \textbf{lan} }\\
\noindent /l{\textprimstress}an/\\
\noindent lit. twilight+???\\
\noindent \textit{Duspars}\\


\noindent History:

\vspace{-0pt}
\hspace{40pt}
\begin{tabular}{ccc}
\textit{6000} & /brlana???/&$\rightarrow$ & \textit{7501} & /rlana???/&$\rightarrow$ & \textit{7600} & /lana???/&$\rightarrow$ & \textit{8500} & /lana/&$\rightarrow$ & \textit{9500} & /lan/& \\
\end{tabular}

\vspace{20pt}\hline

\end{nopagebreak}
\filbreak



\vspace{15pt}
\begin{nopagebreak}
\noindent{\fontsize{20pt}{10pt}\textbf{kelushol} } \textit{23 of autumn/September 23} (N)\\
\noindent {\tovian \fontsize{20pt}{10pt} \textbf{kelushol} }\\
\noindent /kel{\textprimstress}u{\textesh}ol/\\
\noindent lit. branch+fruit\\
\noindent \textit{Starbean}\\


\noindent History:

\vspace{-0pt}
\hspace{40pt}
\begin{tabular}{ccc}
\textit{6000} & /kelu{\textyogh}ola/&$\rightarrow$ & \textit{9500} & /kelu{\textyogh}ol/&$\rightarrow$ & \textit{12001} & /kelu{\textesh}ol/& \\
\end{tabular}

\vspace{20pt}\hline

\end{nopagebreak}
\filbreak



\vspace{15pt}
\begin{nopagebreak}
\noindent{\fontsize{20pt}{10pt}\textbf{shasith} } \textit{23 of spring/March 23} (N)\\
\noindent {\tovian \fontsize{20pt}{10pt} \textbf{shasith} }\\
\noindent /{\textesh}{\textprimstress}asi{\texttheta}/\\
\noindent lit. ice+ship\\
\noindent \textit{Frostnave}\\


\noindent History:

\vspace{-0pt}
\hspace{40pt}
\begin{tabular}{ccc}
\textit{6000} & /{\textesh}u{\textesh}azli{\texttheta}a/&$\rightarrow$ & \textit{9200} & /{\textesh}{\textschwa}{\textesh}azli{\texttheta}a/&$\rightarrow$ & \textit{9500} & /{\textesh}{\textschwa}{\textesh}azli{\texttheta}/&$\rightarrow$ & \textit{11990} & /{\textesh}{\textesh}azli{\texttheta}/&$\rightarrow$ & \textit{12000} & /{\textesh}{\textesh}asli{\texttheta}/&$\rightarrow$ & \textit{12007} & /{\textesh}asli{\texttheta}/&$\rightarrow$ & \textit{12500} & /{\textesh}asi{\texttheta}/& \\
\end{tabular}

\vspace{20pt}\hline

\end{nopagebreak}
\filbreak



\vspace{15pt}
\begin{nopagebreak}
\noindent{\fontsize{20pt}{10pt}\textbf{melaglar} } \textit{23 of summer/June 23} (N)\\
\noindent {\tovian \fontsize{20pt}{10pt} \textbf{melaglar} }\\
\noindent /mel{\textprimstress}aglar/\\
\noindent lit. life+tree\\
\noindent \textit{Heartree}\\


\noindent History:

\vspace{-0pt}
\hspace{40pt}
\begin{tabular}{ccc}
\textit{6000} & /melaglare/&$\rightarrow$ & \textit{9500} & /melaglar/& \\
\end{tabular}

\vspace{20pt}\hline

\end{nopagebreak}
\filbreak



\vspace{15pt}
\begin{nopagebreak}
\noindent{\fontsize{20pt}{10pt}\textbf{lanasyaleth} } \textit{23 of winter/december 23} (N)\\
\noindent {\tovian \fontsize{20pt}{10pt} \textbf{lanasyaleth} }\\
\noindent /lanasy{\textprimstress}ale{\texttheta}/\\
\noindent lit. twilight+food\\
\noindent \textit{Twilispoon}\\


\noindent History:

\vspace{-0pt}
\hspace{40pt}
\begin{tabular}{ccc}
\textit{6000} & /brlanasyale{\dh}a/&$\rightarrow$ & \textit{7501} & /rlanasyale{\dh}a/&$\rightarrow$ & \textit{7600} & /lanasyale{\dh}a/&$\rightarrow$ & \textit{9500} & /lanasyale{\dh}/&$\rightarrow$ & \textit{12002} & /lanasyale{\texttheta}/& \\
\end{tabular}

\vspace{20pt}\hline

\end{nopagebreak}
\filbreak



\vspace{15pt}
\begin{nopagebreak}
\noindent{\fontsize{20pt}{10pt}\textbf{mloanamen} } \textit{24 of autumn/September 24} (N)\\
\noindent {\tovian \fontsize{20pt}{10pt} \textbf{mloanamen} }\\
\noindent /mloan{\textprimstress}amen/\\
\noindent lit. morning+earth\\
\noindent \textit{Dawnborage}\\


\noindent History:

\vspace{-0pt}
\hspace{40pt}
\begin{tabular}{ccc}
\textit{6000} & /ploanamern/&$\rightarrow$ & \textit{7600} & /ploanamen/&$\rightarrow$ & \textit{11000} & /mloanamen/& \\
\end{tabular}

\vspace{20pt}\hline

\end{nopagebreak}
\filbreak



\vspace{15pt}
\begin{nopagebreak}
\noindent{\fontsize{20pt}{10pt}\textbf{nanadrem} } \textit{24 of spring/March 24} (N)\\
\noindent {\tovian \fontsize{20pt}{10pt} \textbf{nanadrem} }\\
\noindent /nan{\textprimstress}adrem/\\
\noindent lit. star+flower\\
\noindent \textit{Starlilis}\\


\noindent History:

\vspace{-0pt}
\hspace{40pt}
\begin{tabular}{ccc}
\textit{6000} & /nnaanadremp/&$\rightarrow$ & \textit{6500} & /naanadremp/&$\rightarrow$ & \textit{8750} & /naanadrem/&$\rightarrow$ & \textit{12005} & /nanadrem/& \\
\end{tabular}

\vspace{20pt}\hline

\end{nopagebreak}
\filbreak



\vspace{15pt}
\begin{nopagebreak}
\noindent{\fontsize{20pt}{10pt}\textbf{tleshor} } \textit{24 of summer/June 24} (N)\\
\noindent {\tovian \fontsize{20pt}{10pt} \textbf{tleshor} }\\
\noindent /{\texttoptiebar{t\textbeltl}}{\textprimstress}e{\textesh}or/\\
\noindent lit. speed+plant\\
\noindent \textit{Swiftgreen}\\


\noindent History:

\vspace{-0pt}
\hspace{40pt}
\begin{tabular}{ccc}
\textit{6000} & /{\texttoptiebar{t\textbeltl}}e{\textyogh}mllore/&$\rightarrow$ & \textit{6500} & /{\texttoptiebar{t\textbeltl}}e{\textyogh}mlore/&$\rightarrow$ & \textit{7600} & /{\texttoptiebar{t\textbeltl}}e{\textyogh}lore/&$\rightarrow$ & \textit{9500} & /{\texttoptiebar{t\textbeltl}}e{\textyogh}lor/&$\rightarrow$ & \textit{12001} & /{\texttoptiebar{t\textbeltl}}e{\textesh}lor/&$\rightarrow$ & \textit{12500} & /{\texttoptiebar{t\textbeltl}}e{\textesh}or/& \\
\end{tabular}

\vspace{20pt}\hline

\end{nopagebreak}
\filbreak



\vspace{15pt}
\begin{nopagebreak}
\noindent{\fontsize{20pt}{10pt}\textbf{lhonrem} } \textit{24 of winter/december 24} (N)\\
\noindent {\tovian \fontsize{20pt}{10pt} \textbf{lhonrem} }\\
\noindent /{\textbeltl}{\textprimstress}onrem/\\
\noindent lit. day+flower\\
\noindent \textit{Daydaisy}\\


\noindent History:

\vspace{-0pt}
\hspace{40pt}
\begin{tabular}{ccc}
\textit{6000} & /{\textbeltl}onmdremp/&$\rightarrow$ & \textit{6000} & /{\textbeltl}onndremp/&$\rightarrow$ & \textit{6500} & /{\textbeltl}ondremp/&$\rightarrow$ & \textit{8500} & /{\textbeltl}onremp/&$\rightarrow$ & \textit{8750} & /{\textbeltl}onrem/& \\
\end{tabular}

\vspace{20pt}\hline

\end{nopagebreak}
\filbreak



\vspace{15pt}
\begin{nopagebreak}
\noindent{\fontsize{20pt}{10pt}\textbf{???} } \textit{25 of autumn/September 25} (N)\\
\noindent {\tovian \fontsize{20pt}{10pt} \textbf{???} }\\
\noindent /???/\\
\noindent lit. ???\\
\noindent \textit{Skuinea}\\


\noindent History:

\vspace{-0pt}
\hspace{40pt}
\begin{tabular}{ccc}
\textit{6000} & /???/& \\
\end{tabular}

\vspace{20pt}\hline

\end{nopagebreak}
\filbreak



\vspace{15pt}
\begin{nopagebreak}
\noindent{\fontsize{20pt}{10pt}\textbf{lrithan} } \textit{25 of spring/March 25} (N)\\
\noindent {\tovian \fontsize{20pt}{10pt} \textbf{lrithan} }\\
\noindent /lr{\textprimstress}i{\texttheta}an/\\
\noindent lit. storm+animal\\
\noindent \textit{Stormbull}\\


\noindent History:

\vspace{-0pt}
\hspace{40pt}
\begin{tabular}{ccc}
\textit{6000} & /kwlyri{\texttheta}l{\textsubbridge{t}}lamn/&$\rightarrow$ & \textit{6401} & /kwlyri{\texttheta}l{\texttoptiebar{t\textbeltl}}amn/&$\rightarrow$ & \textit{7501} & /wlyri{\texttheta}l{\texttoptiebar{t\textbeltl}}amn/&$\rightarrow$ & \textit{7600} & /lri{\texttheta}l{\texttoptiebar{t\textbeltl}}an/&$\rightarrow$ & \textit{12500} & /lri{\texttheta}{\texttoptiebar{t\textbeltl}}an/&$\rightarrow$ & \textit{14001} & /lri{\texttheta}an/& \\
\end{tabular}

\vspace{20pt}\hline

\end{nopagebreak}
\filbreak



\vspace{15pt}
\begin{nopagebreak}
\noindent{\fontsize{20pt}{10pt}\textbf{rithileay} } \textit{25 of summer/June 25} (N)\\
\noindent {\tovian \fontsize{20pt}{10pt} \textbf{rithileay} }\\
\noindent /ri{\texttheta}il{\textprimstress}eay/\\
\noindent lit. sky+feather\\
\noindent \textit{Skyfeather}\\


\noindent History:

\vspace{-0pt}
\hspace{40pt}
\begin{tabular}{ccc}
\textit{6000} & /ri{\texttheta}ilearya/&$\rightarrow$ & \textit{7600} & /ri{\texttheta}ileaya/&$\rightarrow$ & \textit{9500} & /ri{\texttheta}ileay/& \\
\end{tabular}

\vspace{20pt}\hline

\end{nopagebreak}
\filbreak



\vspace{15pt}
\begin{nopagebreak}
\noindent{\fontsize{20pt}{10pt}\textbf{nan} } \textit{25 of winter/december 25} (N)\\
\noindent {\tovian \fontsize{20pt}{10pt} \textbf{nan} }\\
\noindent /n{\textprimstress}an/\\
\noindent lit. star+???\\
\noindent \textit{Starfis}\\


\noindent History:

\vspace{-0pt}
\hspace{40pt}
\begin{tabular}{ccc}
\textit{6000} & /nnaana???/&$\rightarrow$ & \textit{6500} & /naana???/&$\rightarrow$ & \textit{8500} & /naana/&$\rightarrow$ & \textit{9500} & /naan/&$\rightarrow$ & \textit{12005} & /nan/& \\
\end{tabular}

\vspace{20pt}\hline

\end{nopagebreak}
\filbreak



\vspace{15pt}
\begin{nopagebreak}
\noindent{\fontsize{20pt}{10pt}\textbf{melorelor} } \textit{26 of autumn/September 26} (N)\\
\noindent {\tovian \fontsize{20pt}{10pt} \textbf{melorelor} }\\
\noindent /melor{\textprimstress}elor/\\
\noindent lit. plant+plant\\
\noindent \textit{Hayherb}\\


\noindent History:

\vspace{-0pt}
\hspace{40pt}
\begin{tabular}{ccc}
\textit{6000} & /mlloremllore/&$\rightarrow$ & \textit{6500} & /mloremlore/&$\rightarrow$ & \textit{7502} & /meloremlore/&$\rightarrow$ & \textit{7600} & /melorelore/&$\rightarrow$ & \textit{9500} & /melorelor/& \\
\end{tabular}

\vspace{20pt}\hline

\end{nopagebreak}
\filbreak



\vspace{15pt}
\begin{nopagebreak}
\noindent{\fontsize{20pt}{10pt}\textbf{mramlorelhen} } \textit{26 of spring/March 26} (N)\\
\noindent {\tovian \fontsize{20pt}{10pt} \textbf{mramlorelhen} }\\
\noindent /mramlor{\textprimstress}e{\textbeltl}en/\\
\noindent lit. darkness+plant+dream\\
\noindent \textit{Nightshade}\\


\noindent History:

\vspace{-0pt}
\hspace{40pt}
\begin{tabular}{ccc}
\textit{6000} & /bralmllore{\textbeltl}ena/&$\rightarrow$ & \textit{6500} & /bralmlore{\textbeltl}ena/&$\rightarrow$ & \textit{7600} & /bramlore{\textbeltl}ena/&$\rightarrow$ & \textit{9500} & /bramlore{\textbeltl}en/&$\rightarrow$ & \textit{11000} & /mramlore{\textbeltl}en/& \\
\end{tabular}

\vspace{20pt}\hline

\end{nopagebreak}
\filbreak



\vspace{15pt}
\begin{nopagebreak}
\noindent{\fontsize{20pt}{10pt}\textbf{dolithor} } \textit{26 of summer/June 26} (N)\\
\noindent {\tovian \fontsize{20pt}{10pt} \textbf{dolithor} }\\
\noindent /dol{\textprimstress}i{\texttheta}or/\\
\noindent lit. death+smell\\
\noindent \textit{Duskscent}\\


\noindent History:

\vspace{-0pt}
\hspace{40pt}
\begin{tabular}{ccc}
\textit{6000} & /dolri{\texttheta}lore/&$\rightarrow$ & \textit{6503} & /dorli{\texttheta}lore/&$\rightarrow$ & \textit{7600} & /doli{\texttheta}lore/&$\rightarrow$ & \textit{9500} & /doli{\texttheta}lor/&$\rightarrow$ & \textit{12500} & /doli{\texttheta}or/& \\
\end{tabular}

\vspace{20pt}\hline

\end{nopagebreak}
\filbreak



\vspace{15pt}
\begin{nopagebreak}
\noindent{\fontsize{20pt}{10pt}\textbf{shotlodrem} } \textit{26 of winter/december 26} (N)\\
\noindent {\tovian \fontsize{20pt}{10pt} \textbf{shotlodrem} }\\
\noindent /{\textesh}o{\texttoptiebar{t\textbeltl}}{\textprimstress}odrem/\\
\noindent lit. flowing+flower\\
\noindent \textit{Mistdande}\\


\noindent History:

\vspace{-0pt}
\hspace{40pt}
\begin{tabular}{ccc}
\textit{6000} & /{\textesh}o{\texttoptiebar{t\textbeltl}}dremp/&$\rightarrow$ & \textit{8750} & /{\textesh}o{\texttoptiebar{t\textbeltl}}drem/&$\rightarrow$ & \textit{9300} & /{\textesh}o{\texttoptiebar{t\textbeltl}}odrem/& \\
\end{tabular}

\vspace{20pt}\hline

\end{nopagebreak}
\filbreak



\vspace{15pt}
\begin{nopagebreak}
\noindent{\fontsize{20pt}{10pt}\textbf{thullhen} } \textit{27 of autumn/September 27} (N)\\
\noindent {\tovian \fontsize{20pt}{10pt} \textbf{thullhen} }\\
\noindent /{\texttheta}{\textprimstress}ul{\textbeltl}en/\\
\noindent lit. top+dream\\
\noindent \textit{Sarlic}\\


\noindent History:

\vspace{-0pt}
\hspace{40pt}
\begin{tabular}{ccc}
\textit{6000} & /{\dh}ul{\textbeltl}ena/&$\rightarrow$ & \textit{9500} & /{\dh}ul{\textbeltl}en/&$\rightarrow$ & \textit{12002} & /{\texttheta}ul{\textbeltl}en/& \\
\end{tabular}

\vspace{20pt}\hline

\end{nopagebreak}
\filbreak



\vspace{15pt}
\begin{nopagebreak}
\noindent{\fontsize{20pt}{10pt}\textbf{gadrem} } \textit{27 of spring/March 27} (N)\\
\noindent {\tovian \fontsize{20pt}{10pt} \textbf{gadrem} }\\
\noindent /g{\textprimstress}adrem/\\
\noindent lit. hotness+flower\\
\noindent \textit{Ascbloom}\\


\noindent History:

\vspace{-0pt}
\hspace{40pt}
\begin{tabular}{ccc}
\textit{6000} & /agadremp/&$\rightarrow$ & \textit{7500} & /gadremp/&$\rightarrow$ & \textit{8750} & /gadrem/& \\
\end{tabular}

\vspace{20pt}\hline

\end{nopagebreak}
\filbreak



\vspace{15pt}
\begin{nopagebreak}
\noindent{\fontsize{20pt}{10pt}\textbf{lrithdrem} } \textit{27 of summer/June 27} (N)\\
\noindent {\tovian \fontsize{20pt}{10pt} \textbf{lrithdrem} }\\
\noindent /lr{\textprimstress}i{\texttheta}drem/\\
\noindent lit. wind+flower\\
\noindent \textit{Windflower}\\


\noindent History:

\vspace{-0pt}
\hspace{40pt}
\begin{tabular}{ccc}
\textit{6000} & /kwlyri{\texttheta}dremp/&$\rightarrow$ & \textit{7501} & /wlyri{\texttheta}dremp/&$\rightarrow$ & \textit{7600} & /lri{\texttheta}dremp/&$\rightarrow$ & \textit{8750} & /lri{\texttheta}drem/& \\
\end{tabular}

\vspace{20pt}\hline

\end{nopagebreak}
\filbreak



\vspace{15pt}
\begin{nopagebreak}
\noindent{\fontsize{20pt}{10pt}\textbf{doraklrith} } \textit{27 of winter/december 27} (N)\\
\noindent {\tovian \fontsize{20pt}{10pt} \textbf{doraklrith} }\\
\noindent /dor{\textprimstress}aklri{\texttheta}/\\
\noindent lit. blue+breath\\
\noindent \textit{Windwoad}\\


\noindent History:

\vspace{-0pt}
\hspace{40pt}
\begin{tabular}{ccc}
\textit{6000} & /dorakwlyri{\texttheta}/&$\rightarrow$ & \textit{7600} & /doraklri{\texttheta}/& \\
\end{tabular}

\vspace{20pt}\hline

\end{nopagebreak}
\filbreak



\vspace{15pt}
\begin{nopagebreak}
\noindent{\fontsize{20pt}{10pt}\textbf{???} } \textit{28 of autumn/September 28} (N)\\
\noindent {\tovian \fontsize{20pt}{10pt} \textbf{???} }\\
\noindent /???/\\
\noindent lit. ???\\
\noindent \textit{Whivetch}\\


\noindent History:

\vspace{-0pt}
\hspace{40pt}
\begin{tabular}{ccc}
\textit{6000} & /???/& \\
\end{tabular}

\vspace{20pt}\hline

\end{nopagebreak}
\filbreak



\vspace{15pt}
\begin{nopagebreak}
\noindent{\fontsize{20pt}{10pt}\textbf{lshol} } \textit{28 of spring/March 28} (N)\\
\noindent {\tovian \fontsize{20pt}{10pt} \textbf{lshol} }\\
\noindent /l{\textesh}{\textprimstress}ol/\\
\noindent lit. sun+fruit\\
\noindent \textit{Sunfruit}\\


\noindent History:

\vspace{-0pt}
\hspace{40pt}
\begin{tabular}{ccc}
\textit{6000} & /efl{\textyogh}ola/&$\rightarrow$ & \textit{7500} & /fl{\textyogh}ola/&$\rightarrow$ & \textit{7501} & /l{\textyogh}ola/&$\rightarrow$ & \textit{9500} & /l{\textyogh}ol/&$\rightarrow$ & \textit{12001} & /l{\textesh}ol/& \\
\end{tabular}

\vspace{20pt}\hline

\end{nopagebreak}
\filbreak



\vspace{15pt}
\begin{nopagebreak}
\noindent{\fontsize{20pt}{10pt}\textbf{ithrem} } \textit{28 of summer/June 28} (N)\\
\noindent {\tovian \fontsize{20pt}{10pt} \textbf{ithrem} }\\
\noindent /{\textprimstress}i{\texttheta}rem/\\
\noindent lit. thought+flower\\
\noindent \textit{Thoughtbloom}\\


\noindent History:

\vspace{-0pt}
\hspace{40pt}
\begin{tabular}{ccc}
\textit{6000} & /i{\texttheta}rdremp/&$\rightarrow$ & \textit{8750} & /i{\texttheta}rrem/&$\rightarrow$ & \textit{12007} & /i{\texttheta}rem/& \\
\end{tabular}

\vspace{20pt}\hline

\end{nopagebreak}
\filbreak



\vspace{15pt}
\begin{nopagebreak}
\noindent{\fontsize{20pt}{10pt}\textbf{klumloremlan} } \textit{28 of winter/december 28} (N)\\
\noindent {\tovian \fontsize{20pt}{10pt} \textbf{klumloremlan} }\\
\noindent /klumlor{\textprimstress}emlan/\\
\noindent lit. vine+twilight\\
\noindent \textit{Twilifern}\\


\noindent History:

\vspace{-0pt}
\hspace{40pt}
\begin{tabular}{ccc}
\textit{6000} & /kluwmllorebrlana/&$\rightarrow$ & \textit{6500} & /kluwmlorebrlana/&$\rightarrow$ & \textit{7600} & /klumloreblana/&$\rightarrow$ & \textit{9500} & /klumloreblan/&$\rightarrow$ & \textit{11000} & /klumloremlan/& \\
\end{tabular}

\vspace{20pt}\hline

\end{nopagebreak}
\filbreak



\vspace{15pt}
\begin{nopagebreak}
\noindent{\fontsize{20pt}{10pt}\textbf{vareloresyaleth} } \textit{29 of autumn/September 29} (N)\\
\noindent {\tovian \fontsize{20pt}{10pt} \textbf{vareloresyaleth} }\\
\noindent /vareloresy{\textprimstress}ale{\texttheta}/\\
\noindent lit. value+grain\\
\noindent \textit{Goldwheat}\\


\noindent History:

\vspace{-0pt}
\hspace{40pt}
\begin{tabular}{ccc}
\textit{6000} & /vlaremlloresyale{\dh}a/&$\rightarrow$ & \textit{6500} & /vlaremloresyale{\dh}a/&$\rightarrow$ & \textit{7600} & /vlareloresyale{\dh}a/&$\rightarrow$ & \textit{9500} & /vlareloresyale{\dh}/&$\rightarrow$ & \textit{12002} & /vlareloresyale{\texttheta}/&$\rightarrow$ & \textit{12500} & /vareloresyale{\texttheta}/& \\
\end{tabular}

\vspace{20pt}\hline

\end{nopagebreak}
\filbreak



\vspace{15pt}
\begin{nopagebreak}
\noindent{\fontsize{20pt}{10pt}\textbf{nanaloresyaleth} } \textit{29 of spring/March 29} (N)\\
\noindent {\tovian \fontsize{20pt}{10pt} \textbf{nanaloresyaleth} }\\
\noindent /nanaloresy{\textprimstress}ale{\texttheta}/\\
\noindent lit. star+grain\\
\noindent \textit{Starbarley}\\


\noindent History:

\vspace{-0pt}
\hspace{40pt}
\begin{tabular}{ccc}
\textit{6000} & /nnaanamlloresyale{\dh}a/&$\rightarrow$ & \textit{6500} & /naanamloresyale{\dh}a/&$\rightarrow$ & \textit{7600} & /naanaloresyale{\dh}a/&$\rightarrow$ & \textit{9500} & /naanaloresyale{\dh}/&$\rightarrow$ & \textit{12002} & /naanaloresyale{\texttheta}/&$\rightarrow$ & \textit{12005} & /nanaloresyale{\texttheta}/& \\
\end{tabular}

\vspace{20pt}\hline

\end{nopagebreak}
\filbreak



\vspace{15pt}
\begin{nopagebreak}
\noindent{\fontsize{20pt}{10pt}\textbf{mralshol} } \textit{29 of summer/June 29} (N)\\
\noindent {\tovian \fontsize{20pt}{10pt} \textbf{mralshol} }\\
\noindent /mr{\textprimstress}al{\textesh}ol/\\
\noindent lit. darkness+fruit\\
\noindent \textit{Nightberry}\\


\noindent History:

\vspace{-0pt}
\hspace{40pt}
\begin{tabular}{ccc}
\textit{6000} & /bral{\textyogh}ola/&$\rightarrow$ & \textit{9500} & /bral{\textyogh}ol/&$\rightarrow$ & \textit{11000} & /mral{\textyogh}ol/&$\rightarrow$ & \textit{12001} & /mral{\textesh}ol/& \\
\end{tabular}

\vspace{20pt}\hline

\end{nopagebreak}
\filbreak



\vspace{15pt}
\begin{nopagebreak}
\noindent{\fontsize{20pt}{10pt}\textbf{ayskesh} } \textit{29 of winter/december 29} (N)\\
\noindent {\tovian \fontsize{20pt}{10pt} \textbf{ayskesh} }\\
\noindent /{\textprimstress}ayske{\textesh}/\\
\noindent lit. moon+fire\\
\noindent \textit{Moonash}\\


\noindent History:

\vspace{-0pt}
\hspace{40pt}
\begin{tabular}{ccc}
\textit{6000} & /aliyske{\textesh}/&$\rightarrow$ & \textit{6300} & /alyske{\textesh}/&$\rightarrow$ & \textit{7600} & /ayske{\textesh}/& \\
\end{tabular}

\vspace{20pt}\hline

\end{nopagebreak}
\filbreak



\vspace{15pt}
\begin{nopagebreak}
\noindent{\fontsize{20pt}{10pt}\textbf{mralaysm} } \textit{3 of autumn/September 3} (N)\\
\noindent {\tovian \fontsize{20pt}{10pt} \textbf{mralaysm} }\\
\noindent /mr{\textprimstress}alaysm/\\
\noindent lit. darkness+moon+bloom\\
\noindent \textit{Nightbulb}\\


\noindent History:

\vspace{-0pt}
\hspace{40pt}
\begin{tabular}{ccc}
\textit{6000} & /bralaliysdrmpfede/&$\rightarrow$ & \textit{6300} & /bralalysdrmpfede/&$\rightarrow$ & \textit{7600} & /bralaysdmpfede/&$\rightarrow$ & \textit{8500} & /bralaysdmfede/&$\rightarrow$ & \textit{9500} & /bralaysdmfed/&$\rightarrow$ & \textit{11000} & /mralaysdmfed/&$\rightarrow$ & \textit{12003} & /mralaysdmd/&$\rightarrow$ & \textit{13001} & /mralaysdm/&$\rightarrow$ & \textit{13004} & /mralaysm/& \\
\end{tabular}

\vspace{20pt}\hline

\end{nopagebreak}
\filbreak



\vspace{15pt}
\begin{nopagebreak}
\noindent{\fontsize{20pt}{10pt}\textbf{thanshol} } \textit{3 of spring/March 3} (N)\\
\noindent {\tovian \fontsize{20pt}{10pt} \textbf{thanshol} }\\
\noindent /{\texttheta}{\textprimstress}an{\textesh}ol/\\
\noindent lit. tool+fruit\\
\noindent \textit{Stellnut}\\


\noindent History:

\vspace{-0pt}
\hspace{40pt}
\begin{tabular}{ccc}
\textit{6000} & /{\dh}an{\textyogh}ola/&$\rightarrow$ & \textit{9500} & /{\dh}an{\textyogh}ol/&$\rightarrow$ & \textit{12001} & /{\dh}an{\textesh}ol/&$\rightarrow$ & \textit{12002} & /{\texttheta}an{\textesh}ol/& \\
\end{tabular}

\vspace{20pt}\hline

\end{nopagebreak}
\filbreak



\vspace{15pt}
\begin{nopagebreak}
\noindent{\fontsize{20pt}{10pt}\textbf{aysan} } \textit{3 of summer/June 3} (N)\\
\noindent {\tovian \fontsize{20pt}{10pt} \textbf{aysan} }\\
\noindent /{\textprimstress}aysan/\\
\noindent lit. moon+tool\\
\noindent \textit{Moonspear}\\


\noindent History:

\vspace{-0pt}
\hspace{40pt}
\begin{tabular}{ccc}
\textit{6000} & /aliys{\dh}an/&$\rightarrow$ & \textit{6300} & /alys{\dh}an/&$\rightarrow$ & \textit{7600} & /ays{\dh}an/&$\rightarrow$ & \textit{12002} & /ays{\texttheta}an/&$\rightarrow$ & \textit{13000} & /aysan/& \\
\end{tabular}

\vspace{20pt}\hline

\end{nopagebreak}
\filbreak



\vspace{15pt}
\begin{nopagebreak}
\noindent{\fontsize{20pt}{10pt}\textbf{???} } \textit{3 of winter/december 3} (N)\\
\noindent {\tovian \fontsize{20pt}{10pt} \textbf{???} }\\
\noindent /???/\\
\noindent lit. ???\\
\noindent \textit{Brezgilly}\\


\noindent History:

\vspace{-0pt}
\hspace{40pt}
\begin{tabular}{ccc}
\textit{6000} & /???/& \\
\end{tabular}

\vspace{20pt}\hline

\end{nopagebreak}
\filbreak



\vspace{15pt}
\begin{nopagebreak}
\noindent{\fontsize{20pt}{10pt}\textbf{rithtesh} } \textit{30 of autumn/September 30} (N)\\
\noindent {\tovian \fontsize{20pt}{10pt} \textbf{rithtesh} }\\
\noindent /r{\textprimstress}i{\texttheta}{\textsubbridge{t}}e{\textesh}/\\
\noindent lit. air+sound\\
\noindent \textit{Windflute}\\


\noindent History:

\vspace{-0pt}
\hspace{40pt}
\begin{tabular}{ccc}
\textit{6000} & /ri{\texttheta}{\textsubbridge{t}}e{\textyogh}/&$\rightarrow$ & \textit{12001} & /ri{\texttheta}{\textsubbridge{t}}e{\textesh}/& \\
\end{tabular}

\vspace{20pt}\hline

\end{nopagebreak}
\filbreak



\vspace{15pt}
\begin{nopagebreak}
\noindent{\fontsize{20pt}{10pt}\textbf{shimilhom} } \textit{30 of spring/March 30} (N)\\
\noindent {\tovian \fontsize{20pt}{10pt} \textbf{shimilhom} }\\
\noindent /{\textesh}im{\textprimstress}i{\textbeltl}om/\\
\noindent lit. mystery+container\\
\noindent \textit{Mystbarrel}\\


\noindent History:

\vspace{-0pt}
\hspace{40pt}
\begin{tabular}{ccc}
\textit{6000} & /{\textesh}imi{\textbeltl}oma/&$\rightarrow$ & \textit{9500} & /{\textesh}imi{\textbeltl}om/& \\
\end{tabular}

\vspace{20pt}\hline

\end{nopagebreak}
\filbreak



\vspace{15pt}
\begin{nopagebreak}
\noindent{\fontsize{20pt}{10pt}\textbf{melor} } \textit{30 of summer/June 30} (N)\\
\noindent {\tovian \fontsize{20pt}{10pt} \textbf{melor} }\\
\noindent /m{\textprimstress}elor/\\
\noindent lit. plant+???\\
\noindent \textit{Sproutbles}\\


\noindent History:

\vspace{-0pt}
\hspace{40pt}
\begin{tabular}{ccc}
\textit{6000} & /mllore???/&$\rightarrow$ & \textit{6500} & /mlore???/&$\rightarrow$ & \textit{7502} & /melore???/&$\rightarrow$ & \textit{8500} & /melore/&$\rightarrow$ & \textit{9500} & /melor/& \\
\end{tabular}

\vspace{20pt}\hline

\end{nopagebreak}
\filbreak



\vspace{15pt}
\begin{nopagebreak}
\noindent{\fontsize{20pt}{10pt}\textbf{lrith} } \textit{30 of winter/december 30} (N)\\
\noindent {\tovian \fontsize{20pt}{10pt} \textbf{lrith} }\\
\noindent /lr{\textprimstress}i{\texttheta}/\\
\noindent lit. breath+???\\
\noindent \textit{Galedibble}\\


\noindent History:

\vspace{-0pt}
\hspace{40pt}
\begin{tabular}{ccc}
\textit{6000} & /kwlyri{\texttheta}???/&$\rightarrow$ & \textit{7501} & /wlyri{\texttheta}???/&$\rightarrow$ & \textit{7600} & /lri{\texttheta}???/&$\rightarrow$ & \textit{8500} & /lri{\texttheta}/& \\
\end{tabular}

\vspace{20pt}\hline

\end{nopagebreak}
\filbreak



\vspace{15pt}
\begin{nopagebreak}
\noindent{\fontsize{20pt}{10pt}\textbf{leshol} } \textit{31 of autumn/October 1} (N)\\
\noindent {\tovian \fontsize{20pt}{10pt} \textbf{leshol} }\\
\noindent /l{\textprimstress}e{\textesh}ol/\\
\noindent lit. harvest+fruit\\
\noindent \textit{Starplum}\\


\noindent History:

\vspace{-0pt}
\hspace{40pt}
\begin{tabular}{ccc}
\textit{6000} & /mllre{\textyogh}{\textyogh}ola/&$\rightarrow$ & \textit{6500} & /mlre{\textyogh}ola/&$\rightarrow$ & \textit{6503} & /mrle{\textyogh}ola/&$\rightarrow$ & \textit{7501} & /rle{\textyogh}ola/&$\rightarrow$ & \textit{7600} & /le{\textyogh}ola/&$\rightarrow$ & \textit{9500} & /le{\textyogh}ol/&$\rightarrow$ & \textit{12001} & /le{\textesh}ol/& \\
\end{tabular}

\vspace{20pt}\hline

\end{nopagebreak}
\filbreak



\vspace{15pt}
\begin{nopagebreak}
\noindent{\fontsize{20pt}{10pt}\textbf{rithileshol} } \textit{31 of spring/April 1} (N)\\
\noindent {\tovian \fontsize{20pt}{10pt} \textbf{rithileshol} }\\
\noindent /ri{\texttheta}il{\textprimstress}e{\textesh}ol/\\
\noindent lit. sky+fruit\\
\noindent \textit{Skyfruit}\\


\noindent History:

\vspace{-0pt}
\hspace{40pt}
\begin{tabular}{ccc}
\textit{6000} & /ri{\texttheta}ile{\textyogh}ola/&$\rightarrow$ & \textit{9500} & /ri{\texttheta}ile{\textyogh}ol/&$\rightarrow$ & \textit{12001} & /ri{\texttheta}ile{\textesh}ol/& \\
\end{tabular}

\vspace{20pt}\hline

\end{nopagebreak}
\filbreak



\vspace{15pt}
\begin{nopagebreak}
\noindent{\fontsize{20pt}{10pt}\textbf{lhenananaloresyaleth} } \textit{31 of summer/July 1} (N)\\
\noindent {\tovian \fontsize{20pt}{10pt} \textbf{lhenananaloresyaleth} }\\
\noindent /{\textbeltl}enananaloresy{\textprimstress}ale{\texttheta}/\\
\noindent lit. dream+star+grain\\
\noindent \textit{Startriticale}\\


\noindent History:

\vspace{-0pt}
\hspace{40pt}
\begin{tabular}{ccc}
\textit{6000} & /{\textbeltl}enannaanamlloresyale{\dh}a/&$\rightarrow$ & \textit{6500} & /{\textbeltl}enanaanamloresyale{\dh}a/&$\rightarrow$ & \textit{7600} & /{\textbeltl}enanaanaloresyale{\dh}a/&$\rightarrow$ & \textit{9500} & /{\textbeltl}enanaanaloresyale{\dh}/&$\rightarrow$ & \textit{12002} & /{\textbeltl}enanaanaloresyale{\texttheta}/&$\rightarrow$ & \textit{12005} & /{\textbeltl}enananaloresyale{\texttheta}/& \\
\end{tabular}

\vspace{20pt}\hline

\end{nopagebreak}
\filbreak



\vspace{15pt}
\begin{nopagebreak}
\noindent{\fontsize{20pt}{10pt}\textbf{???} } \textit{31 of winter/January 31} (N)\\
\noindent {\tovian \fontsize{20pt}{10pt} \textbf{???} }\\
\noindent /???/\\
\noindent lit. ???\\
\noindent \textit{Frohstlef}\\


\noindent History:

\vspace{-0pt}
\hspace{40pt}
\begin{tabular}{ccc}
\textit{6000} & /???/& \\
\end{tabular}

\vspace{20pt}\hline

\end{nopagebreak}
\filbreak



\vspace{15pt}
\begin{nopagebreak}
\noindent{\fontsize{20pt}{10pt}\textbf{klumloreloresyaleth} } \textit{32 of autumn/October 2} (N)\\
\noindent {\tovian \fontsize{20pt}{10pt} \textbf{klumloreloresyaleth} }\\
\noindent /klumloreloresy{\textprimstress}ale{\texttheta}/\\
\noindent lit. vine+grain\\
\noindent \textit{Duskmillet}\\


\noindent History:

\vspace{-0pt}
\hspace{40pt}
\begin{tabular}{ccc}
\textit{6000} & /kluwmlloremlloresyale{\dh}a/&$\rightarrow$ & \textit{6500} & /kluwmloremloresyale{\dh}a/&$\rightarrow$ & \textit{7600} & /klumloreloresyale{\dh}a/&$\rightarrow$ & \textit{9500} & /klumloreloresyale{\dh}/&$\rightarrow$ & \textit{12002} & /klumloreloresyale{\texttheta}/& \\
\end{tabular}

\vspace{20pt}\hline

\end{nopagebreak}
\filbreak



\vspace{15pt}
\begin{nopagebreak}
\noindent{\fontsize{20pt}{10pt}\textbf{aysith} } \textit{32 of spring/April 2} (N)\\
\noindent {\tovian \fontsize{20pt}{10pt} \textbf{aysith} }\\
\noindent /{\textprimstress}aysi{\texttheta}/\\
\noindent lit. moon+sharpness\\
\noindent \textit{Mooncrísp}\\


\noindent History:

\vspace{-0pt}
\hspace{40pt}
\begin{tabular}{ccc}
\textit{6000} & /aliys{\textyogh}i{\texttheta}/&$\rightarrow$ & \textit{6300} & /alys{\textyogh}i{\texttheta}/&$\rightarrow$ & \textit{7600} & /ays{\textyogh}i{\texttheta}/&$\rightarrow$ & \textit{12001} & /ays{\textesh}i{\texttheta}/&$\rightarrow$ & \textit{13000} & /aysi{\texttheta}/& \\
\end{tabular}

\vspace{20pt}\hline

\end{nopagebreak}
\filbreak



\vspace{15pt}
\begin{nopagebreak}
\noindent{\fontsize{20pt}{10pt}\textbf{yeomlor} } \textit{32 of summer/July 2} (N)\\
\noindent {\tovian \fontsize{20pt}{10pt} \textbf{yeomlor} }\\
\noindent /ye{\textprimstress}omlor/\\
\noindent lit. wisdom+plant\\
\noindent \textit{Moonmullen}\\


\noindent History:

\vspace{-0pt}
\hspace{40pt}
\begin{tabular}{ccc}
\textit{6000} & /aryeolmllore/&$\rightarrow$ & \textit{6500} & /aryeolmlore/&$\rightarrow$ & \textit{7500} & /ryeolmlore/&$\rightarrow$ & \textit{7600} & /yeomlore/&$\rightarrow$ & \textit{9500} & /yeomlor/& \\
\end{tabular}

\vspace{20pt}\hline

\end{nopagebreak}
\filbreak



\vspace{15pt}
\begin{nopagebreak}
\noindent{\fontsize{20pt}{10pt}\textbf{tushuthok} } \textit{32 of winter/January 32} (N)\\
\noindent {\tovian \fontsize{20pt}{10pt} \textbf{tushuthok} }\\
\noindent /{\textsubbridge{t}}u{\textesh}{\textprimstress}u{\texttheta}ok/\\
\noindent lit. cold+ground\\
\noindent \textit{Wintersnap}\\


\noindent History:

\vspace{-0pt}
\hspace{40pt}
\begin{tabular}{ccc}
\textit{6000} & /{\textsubbridge{t}}u{\textesh}u{\texttheta}ok{\textesh}/&$\rightarrow$ & \textit{13001} & /{\textsubbridge{t}}u{\textesh}u{\texttheta}ok/& \\
\end{tabular}

\vspace{20pt}\hline

\end{nopagebreak}
\filbreak



\vspace{15pt}
\begin{nopagebreak}
\noindent{\fontsize{20pt}{10pt}\textbf{???} } \textit{33 of autumn/October 3} (N)\\
\noindent {\tovian \fontsize{20pt}{10pt} \textbf{???} }\\
\noindent /???/\\
\noindent lit. ???\\
\noindent \textit{Munpuf}\\


\noindent History:

\vspace{-0pt}
\hspace{40pt}
\begin{tabular}{ccc}
\textit{6000} & /???/& \\
\end{tabular}

\vspace{20pt}\hline

\end{nopagebreak}
\filbreak



\vspace{15pt}
\begin{nopagebreak}
\noindent{\fontsize{20pt}{10pt}\textbf{lrithwen} } \textit{33 of spring/April 3} (N)\\
\noindent {\tovian \fontsize{20pt}{10pt} \textbf{lrithwen} }\\
\noindent /lr{\textprimstress}i{\texttheta}wen/\\
\noindent lit. storm+beauty\\
\noindent \textit{Thunderpearl}\\


\noindent History:

\vspace{-0pt}
\hspace{40pt}
\begin{tabular}{ccc}
\textit{6000} & /kwlyri{\texttheta}lwen/&$\rightarrow$ & \textit{7501} & /wlyri{\texttheta}lwen/&$\rightarrow$ & \textit{7600} & /lri{\texttheta}wen/& \\
\end{tabular}

\vspace{20pt}\hline

\end{nopagebreak}
\filbreak



\vspace{15pt}
\begin{nopagebreak}
\noindent{\fontsize{20pt}{10pt}\textbf{llhenashol} } \textit{33 of summer/July 3} (N)\\
\noindent {\tovian \fontsize{20pt}{10pt} \textbf{llhenashol} }\\
\noindent /l{\textbeltl}en{\textprimstress}a{\textesh}ol/\\
\noindent lit. sun+dream+fruit\\
\noindent \textit{Sanfruta}\\


\noindent History:

\vspace{-0pt}
\hspace{40pt}
\begin{tabular}{ccc}
\textit{6000} & /efl{\textbeltl}ena{\textyogh}ola/&$\rightarrow$ & \textit{7500} & /fl{\textbeltl}ena{\textyogh}ola/&$\rightarrow$ & \textit{7501} & /l{\textbeltl}ena{\textyogh}ola/&$\rightarrow$ & \textit{9500} & /l{\textbeltl}ena{\textyogh}ol/&$\rightarrow$ & \textit{12001} & /l{\textbeltl}ena{\textesh}ol/& \\
\end{tabular}

\vspace{20pt}\hline

\end{nopagebreak}
\filbreak



\vspace{15pt}
\begin{nopagebreak}
\noindent{\fontsize{20pt}{10pt}\textbf{mralrem} } \textit{33 of winter/January 33} (N)\\
\noindent {\tovian \fontsize{20pt}{10pt} \textbf{mralrem} }\\
\noindent /mr{\textprimstress}alrem/\\
\noindent lit. darkness+flower\\
\noindent \textit{Níghtbloom}\\


\noindent History:

\vspace{-0pt}
\hspace{40pt}
\begin{tabular}{ccc}
\textit{6000} & /braldremp/&$\rightarrow$ & \textit{8750} & /bralrem/&$\rightarrow$ & \textit{11000} & /mralrem/& \\
\end{tabular}

\vspace{20pt}\hline

\end{nopagebreak}
\filbreak



\vspace{15pt}
\begin{nopagebreak}
\noindent{\fontsize{20pt}{10pt}\textbf{lhilaloresyaleth} } \textit{34 of autumn/October 4} (N)\\
\noindent {\tovian \fontsize{20pt}{10pt} \textbf{lhilaloresyaleth} }\\
\noindent /{\textbeltl}ilaloresy{\textprimstress}ale{\texttheta}/\\
\noindent lit. night+grain\\
\noindent \textit{Twilbar}\\


\noindent History:

\vspace{-0pt}
\hspace{40pt}
\begin{tabular}{ccc}
\textit{6000} & /a{\textbeltl}inlamlloresyale{\dh}a/&$\rightarrow$ & \textit{6500} & /a{\textbeltl}inlamloresyale{\dh}a/&$\rightarrow$ & \textit{7500} & /{\textbeltl}inlamloresyale{\dh}a/&$\rightarrow$ & \textit{7600} & /{\textbeltl}ilaloresyale{\dh}a/&$\rightarrow$ & \textit{9500} & /{\textbeltl}ilaloresyale{\dh}/&$\rightarrow$ & \textit{12002} & /{\textbeltl}ilaloresyale{\texttheta}/& \\
\end{tabular}

\vspace{20pt}\hline

\end{nopagebreak}
\filbreak



\vspace{15pt}
\begin{nopagebreak}
\noindent{\fontsize{20pt}{10pt}\textbf{dlelaglar} } \textit{34 of spring/April 4} (N)\\
\noindent {\tovian \fontsize{20pt}{10pt} \textbf{dlelaglar} }\\
\noindent /dlel{\textprimstress}aglar/\\
\noindent lit. body+tree\\
\noindent \textit{Heartroot}\\


\noindent History:

\vspace{-0pt}
\hspace{40pt}
\begin{tabular}{ccc}
\textit{6000} & /adllenlaglare/&$\rightarrow$ & \textit{6500} & /adlenlaglare/&$\rightarrow$ & \textit{7500} & /dlenlaglare/&$\rightarrow$ & \textit{7600} & /dlelaglare/&$\rightarrow$ & \textit{9500} & /dlelaglar/& \\
\end{tabular}

\vspace{20pt}\hline

\end{nopagebreak}
\filbreak



\vspace{15pt}
\begin{nopagebreak}
\noindent{\fontsize{20pt}{10pt}\textbf{melorethanglar} } \textit{34 of summer/July 4} (N)\\
\noindent {\tovian \fontsize{20pt}{10pt} \textbf{melorethanglar} }\\
\noindent /melore{\texttheta}{\textprimstress}a{\ng}lar/\\
\noindent lit. plant+tool+woods\\
\noindent \textit{Melorethangor}\\


\noindent History:

\vspace{-0pt}
\hspace{40pt}
\begin{tabular}{ccc}
\textit{6000} & /mllore{\dh}anglar/&$\rightarrow$ & \textit{6000} & /mllore{\dh}a{\ng}glar/&$\rightarrow$ & \textit{6500} & /mlore{\dh}a{\ng}glar/&$\rightarrow$ & \textit{7502} & /melore{\dh}a{\ng}glar/&$\rightarrow$ & \textit{8500} & /melore{\dh}a{\ng}lar/&$\rightarrow$ & \textit{12002} & /melore{\texttheta}a{\ng}lar/& \\
\end{tabular}

\vspace{20pt}\hline

\end{nopagebreak}
\filbreak



\vspace{15pt}
\begin{nopagebreak}
\noindent{\fontsize{20pt}{10pt}\textbf{shashanin} } \textit{34 of winter/January 34} (N)\\
\noindent {\tovian \fontsize{20pt}{10pt} \textbf{shashanin} }\\
\noindent /{\textesh}a{\textesh}{\textprimstress}anin/\\
\noindent lit. ice+berry\\
\noindent \textit{Icepear}\\


\noindent History:

\vspace{-0pt}
\hspace{40pt}
\begin{tabular}{ccc}
\textit{6000} & /{\textesh}u{\textesh}a{\textyogh}lanina/&$\rightarrow$ & \textit{9200} & /{\textesh}{\textschwa}{\textesh}a{\textyogh}lanina/&$\rightarrow$ & \textit{9500} & /{\textesh}{\textschwa}{\textesh}a{\textyogh}lanin/&$\rightarrow$ & \textit{11990} & /{\textesh}{\textesh}a{\textyogh}lanin/&$\rightarrow$ & \textit{12001} & /{\textesh}{\textesh}a{\textesh}lanin/&$\rightarrow$ & \textit{12004} & /{\textesh}a{\textesh}lanin/&$\rightarrow$ & \textit{12500} & /{\textesh}a{\textesh}anin/& \\
\end{tabular}

\vspace{20pt}\hline

\end{nopagebreak}
\filbreak



\vspace{15pt}
\begin{nopagebreak}
\noindent{\fontsize{20pt}{10pt}\textbf{shotlan} } \textit{35 of autumn/October 5} (N)\\
\noindent {\tovian \fontsize{20pt}{10pt} \textbf{shotlan} }\\
\noindent /{\textesh}{\textprimstress}o{\texttoptiebar{t\textbeltl}}an/\\
\noindent lit. river+light\\
\noindent \textit{Rivshine}\\


\noindent History:

\vspace{-0pt}
\hspace{40pt}
\begin{tabular}{ccc}
\textit{6000} & /{\textesh}o{\texttoptiebar{t\textbeltl}}zlana/&$\rightarrow$ & \textit{9500} & /{\textesh}o{\texttoptiebar{t\textbeltl}}zlan/&$\rightarrow$ & \textit{12000} & /{\textesh}o{\texttoptiebar{t\textbeltl}}slan/&$\rightarrow$ & \textit{12500} & /{\textesh}o{\texttoptiebar{t\textbeltl}}san/&$\rightarrow$ & \textit{14000} & /{\textesh}o{\texttoptiebar{t\textbeltl}}an/& \\
\end{tabular}

\vspace{20pt}\hline

\end{nopagebreak}
\filbreak



\vspace{15pt}
\begin{nopagebreak}
\noindent{\fontsize{20pt}{10pt}\textbf{lrithay} } \textit{35 of spring/April 5} (N)\\
\noindent {\tovian \fontsize{20pt}{10pt} \textbf{lrithay} }\\
\noindent /lr{\textprimstress}i{\texttheta}ay/\\
\noindent lit. storm+feather\\
\noindent \textit{Stormféther}\\


\noindent History:

\vspace{-0pt}
\hspace{40pt}
\begin{tabular}{ccc}
\textit{6000} & /kwlyri{\texttheta}larya/&$\rightarrow$ & \textit{7501} & /wlyri{\texttheta}larya/&$\rightarrow$ & \textit{7600} & /lri{\texttheta}laya/&$\rightarrow$ & \textit{9500} & /lri{\texttheta}lay/&$\rightarrow$ & \textit{12500} & /lri{\texttheta}ay/& \\
\end{tabular}

\vspace{20pt}\hline

\end{nopagebreak}
\filbreak



\vspace{15pt}
\begin{nopagebreak}
\noindent{\fontsize{20pt}{10pt}\textbf{lhomtlan} } \textit{35 of summer/July 5} (N)\\
\noindent {\tovian \fontsize{20pt}{10pt} \textbf{lhomtlan} }\\
\noindent /{\textbeltl}{\textprimstress}om{\texttoptiebar{t\textbeltl}}an/\\
\noindent lit. day+animal\\
\noindent \textit{Dayram}\\


\noindent History:

\vspace{-0pt}
\hspace{40pt}
\begin{tabular}{ccc}
\textit{6000} & /{\textbeltl}onm{\textsubbridge{t}}lamn/&$\rightarrow$ & \textit{6401} & /{\textbeltl}onm{\texttoptiebar{t\textbeltl}}amn/&$\rightarrow$ & \textit{7600} & /{\textbeltl}om{\texttoptiebar{t\textbeltl}}an/& \\
\end{tabular}

\vspace{20pt}\hline

\end{nopagebreak}
\filbreak



\vspace{15pt}
\begin{nopagebreak}
\noindent{\fontsize{20pt}{10pt}\textbf{selshith} } \textit{35 of winter/January 35} (N)\\
\noindent {\tovian \fontsize{20pt}{10pt} \textbf{selshith} }\\
\noindent /s{\textprimstress}el{\textesh}i{\texttheta}/\\
\noindent lit. water+sharpness\\
\noindent \textit{Snowbristle}\\


\noindent History:

\vspace{-0pt}
\hspace{40pt}
\begin{tabular}{ccc}
\textit{6000} & /zel{\textyogh}i{\texttheta}/&$\rightarrow$ & \textit{12000} & /sel{\textyogh}i{\texttheta}/&$\rightarrow$ & \textit{12001} & /sel{\textesh}i{\texttheta}/& \\
\end{tabular}

\vspace{20pt}\hline

\end{nopagebreak}
\filbreak



\vspace{15pt}
\begin{nopagebreak}
\noindent{\fontsize{20pt}{10pt}\textbf{ysnal} } \textit{36 of autumn/October 6} (N)\\
\noindent {\tovian \fontsize{20pt}{10pt} \textbf{ysnal} }\\
\noindent /ysn{\textprimstress}al/\\
\noindent lit. moon+sharing\\
\noindent \textit{Tlunalis}\\


\noindent History:

\vspace{-0pt}
\hspace{40pt}
\begin{tabular}{ccc}
\textit{6000} & /aliys{\texttoptiebar{t\textbeltl}}naels/&$\rightarrow$ & \textit{6300} & /alys{\texttoptiebar{t\textbeltl}}naels/&$\rightarrow$ & \textit{7500} & /lys{\texttoptiebar{t\textbeltl}}naels/&$\rightarrow$ & \textit{7501} & /ys{\texttoptiebar{t\textbeltl}}naels/&$\rightarrow$ & \textit{10000} & /ys{\texttoptiebar{t\textbeltl}}nals/&$\rightarrow$ & \textit{13001} & /ys{\texttoptiebar{t\textbeltl}}nal/&$\rightarrow$ & \textit{14001} & /ysnal/& \\
\end{tabular}

\vspace{20pt}\hline

\end{nopagebreak}
\filbreak



\vspace{15pt}
\begin{nopagebreak}
\noindent{\fontsize{20pt}{10pt}\textbf{teshglar} } \textit{36 of spring/April 6} (N)\\
\noindent {\tovian \fontsize{20pt}{10pt} \textbf{teshglar} }\\
\noindent /{\textsubbridge{t}}{\textprimstress}e{\textesh}glar/\\
\noindent lit. sound+woods\\
\noindent \textit{Teshangor}\\


\noindent History:

\vspace{-0pt}
\hspace{40pt}
\begin{tabular}{ccc}
\textit{6000} & /{\textsubbridge{t}}e{\textyogh}glar/&$\rightarrow$ & \textit{12001} & /{\textsubbridge{t}}e{\textesh}glar/& \\
\end{tabular}

\vspace{20pt}\hline

\end{nopagebreak}
\filbreak



\vspace{15pt}
\begin{nopagebreak}
\noindent{\fontsize{20pt}{10pt}\textbf{mloanalor} } \textit{36 of summer/July 6} (N)\\
\noindent {\tovian \fontsize{20pt}{10pt} \textbf{mloanalor} }\\
\noindent /mloan{\textprimstress}alor/\\
\noindent lit. dawn+plant\\
\noindent \textit{Dawnfern}\\


\noindent History:

\vspace{-0pt}
\hspace{40pt}
\begin{tabular}{ccc}
\textit{6000} & /ploanamllore/&$\rightarrow$ & \textit{6500} & /ploanamlore/&$\rightarrow$ & \textit{7600} & /ploanalore/&$\rightarrow$ & \textit{9500} & /ploanalor/&$\rightarrow$ & \textit{11000} & /mloanalor/& \\
\end{tabular}

\vspace{20pt}\hline

\end{nopagebreak}
\filbreak



\vspace{15pt}
\begin{nopagebreak}
\noindent{\fontsize{20pt}{10pt}\textbf{klumloreshud} } \textit{36 of winter/January 36} (N)\\
\noindent {\tovian \fontsize{20pt}{10pt} \textbf{klumloreshud} }\\
\noindent /klumlor{\textprimstress}e{\textesh}ud/\\
\noindent lit. vine+ice\\
\noindent \textit{Frostfern}\\


\noindent History:

\vspace{-0pt}
\hspace{40pt}
\begin{tabular}{ccc}
\textit{6000} & /kluwmllore{\textesh}u{\textesh}a/&$\rightarrow$ & \textit{6500} & /kluwmlore{\textesh}u{\textesh}a/&$\rightarrow$ & \textit{7600} & /klumlore{\textesh}u{\textesh}a/&$\rightarrow$ & \textit{9500} & /klumlore{\textesh}u{\textesh}/&$\rightarrow$ & \textit{12501} & /klumlore{\textesh}ud/& \\
\end{tabular}

\vspace{20pt}\hline

\end{nopagebreak}
\filbreak



\vspace{15pt}
\begin{nopagebreak}
\noindent{\fontsize{20pt}{10pt}\textbf{hol} } \textit{37 of autumn/October 7} (N)\\
\noindent {\tovian \fontsize{20pt}{10pt} \textbf{hol} }\\
\noindent /ħ{\textprimstress}ol/\\
\noindent lit. sweet+???\\
\noindent \textit{Sweetspore}\\


\noindent History:

\vspace{-0pt}
\hspace{40pt}
\begin{tabular}{ccc}
\textit{6000} & /hola???/&$\rightarrow$ & \textit{8500} & /hola/&$\rightarrow$ & \textit{9500} & /hol/&$\rightarrow$ & \textit{14006} & /ħol/& \\
\end{tabular}

\vspace{20pt}\hline

\end{nopagebreak}
\filbreak



\vspace{15pt}
\begin{nopagebreak}
\noindent{\fontsize{20pt}{10pt}\textbf{dremshol} } \textit{37 of spring/April 7} (N)\\
\noindent {\tovian \fontsize{20pt}{10pt} \textbf{dremshol} }\\
\noindent /dr{\textprimstress}em{\textesh}ol/\\
\noindent lit. flower+fruit\\
\noindent \textit{Drémfig}\\


\noindent History:

\vspace{-0pt}
\hspace{40pt}
\begin{tabular}{ccc}
\textit{6000} & /dremp{\textyogh}ola/&$\rightarrow$ & \textit{8500} & /drem{\textyogh}ola/&$\rightarrow$ & \textit{9500} & /drem{\textyogh}ol/&$\rightarrow$ & \textit{12001} & /drem{\textesh}ol/& \\
\end{tabular}

\vspace{20pt}\hline

\end{nopagebreak}
\filbreak



\vspace{15pt}
\begin{nopagebreak}
\noindent{\fontsize{20pt}{10pt}\textbf{lhenamral} } \textit{37 of summer/July 7} (N)\\
\noindent {\tovian \fontsize{20pt}{10pt} \textbf{lhenamral} }\\
\noindent /{\textbeltl}en{\textprimstress}amral/\\
\noindent lit. dream+darkness\\
\noindent \textit{Nightwood}\\


\noindent History:

\vspace{-0pt}
\hspace{40pt}
\begin{tabular}{ccc}
\textit{6000} & /{\textbeltl}enabral/&$\rightarrow$ & \textit{11000} & /{\textbeltl}enamral/& \\
\end{tabular}

\vspace{20pt}\hline

\end{nopagebreak}
\filbreak



\vspace{15pt}
\begin{nopagebreak}
\noindent{\fontsize{20pt}{10pt}\textbf{aysel} } \textit{37 of winter/January 37} (N)\\
\noindent {\tovian \fontsize{20pt}{10pt} \textbf{aysel} }\\
\noindent /{\textprimstress}aysel/\\
\noindent lit. moon+center\\
\noindent \textit{Mooncrést}\\


\noindent History:

\vspace{-0pt}
\hspace{40pt}
\begin{tabular}{ccc}
\textit{6000} & /aliys{\textbeltl}el/&$\rightarrow$ & \textit{6300} & /alys{\textbeltl}el/&$\rightarrow$ & \textit{7600} & /ays{\textbeltl}el/&$\rightarrow$ & \textit{13000} & /aysel/& \\
\end{tabular}

\vspace{20pt}\hline

\end{nopagebreak}
\filbreak



\vspace{15pt}
\begin{nopagebreak}
\noindent{\fontsize{20pt}{10pt}\textbf{lhenadrem} } \textit{38 of autumn/October 8} (N)\\
\noindent {\tovian \fontsize{20pt}{10pt} \textbf{lhenadrem} }\\
\noindent /{\textbeltl}en{\textprimstress}adrem/\\
\noindent lit. dream+flower\\
\noindent \textit{Dremblossom}\\


\noindent History:

\vspace{-0pt}
\hspace{40pt}
\begin{tabular}{ccc}
\textit{6000} & /{\textbeltl}enadremp/&$\rightarrow$ & \textit{8750} & /{\textbeltl}enadrem/& \\
\end{tabular}

\vspace{20pt}\hline

\end{nopagebreak}
\filbreak



\vspace{15pt}
\begin{nopagebreak}
\noindent{\fontsize{20pt}{10pt}\textbf{mraladlel} } \textit{38 of spring/April 8} (N)\\
\noindent {\tovian \fontsize{20pt}{10pt} \textbf{mraladlel} }\\
\noindent /mral{\textprimstress}adlel/\\
\noindent lit. darkness+body\\
\noindent \textit{Níghtstalk}\\


\noindent History:

\vspace{-0pt}
\hspace{40pt}
\begin{tabular}{ccc}
\textit{6000} & /braladllenl/&$\rightarrow$ & \textit{6500} & /braladlenl/&$\rightarrow$ & \textit{7600} & /braladlel/&$\rightarrow$ & \textit{11000} & /mraladlel/& \\
\end{tabular}

\vspace{20pt}\hline

\end{nopagebreak}
\filbreak



\vspace{15pt}
\begin{nopagebreak}
\noindent{\fontsize{20pt}{10pt}\textbf{llhenalor} } \textit{38 of summer/July 8} (N)\\
\noindent {\tovian \fontsize{20pt}{10pt} \textbf{llhenalor} }\\
\noindent /l{\textbeltl}en{\textprimstress}alor/\\
\noindent lit. sun+dream+plant\\
\noindent \textit{Sunthistle}\\


\noindent History:

\vspace{-0pt}
\hspace{40pt}
\begin{tabular}{ccc}
\textit{6000} & /efl{\textbeltl}enamllore/&$\rightarrow$ & \textit{6500} & /efl{\textbeltl}enamlore/&$\rightarrow$ & \textit{7500} & /fl{\textbeltl}enamlore/&$\rightarrow$ & \textit{7501} & /l{\textbeltl}enamlore/&$\rightarrow$ & \textit{7600} & /l{\textbeltl}enalore/&$\rightarrow$ & \textit{9500} & /l{\textbeltl}enalor/& \\
\end{tabular}

\vspace{20pt}\hline

\end{nopagebreak}
\filbreak



\vspace{15pt}
\begin{nopagebreak}
\noindent{\fontsize{20pt}{10pt}\textbf{fol} } \textit{38 of winter/January 38} (N)\\
\noindent {\tovian \fontsize{20pt}{10pt} \textbf{fol} }\\
\noindent /f{\textprimstress}ol/\\
\noindent lit. sun+sweet\\
\noindent \textit{Sunnectar}\\


\noindent History:

\vspace{-0pt}
\hspace{40pt}
\begin{tabular}{ccc}
\textit{6000} & /eflhola/&$\rightarrow$ & \textit{6402} & /ef{\textbeltl}ola/&$\rightarrow$ & \textit{7500} & /f{\textbeltl}ola/&$\rightarrow$ & \textit{9500} & /f{\textbeltl}ol/&$\rightarrow$ & \textit{13000} & /fol/& \\
\end{tabular}

\vspace{20pt}\hline

\end{nopagebreak}
\filbreak



\vspace{15pt}
\begin{nopagebreak}
\noindent{\fontsize{20pt}{10pt}\textbf{mralaglar} } \textit{39 of autumn/October 9} (N)\\
\noindent {\tovian \fontsize{20pt}{10pt} \textbf{mralaglar} }\\
\noindent /mral{\textprimstress}aglar/\\
\noindent lit. darkness+tree\\
\noindent \textit{Shadroot}\\


\noindent History:

\vspace{-0pt}
\hspace{40pt}
\begin{tabular}{ccc}
\textit{6000} & /bralaglare/&$\rightarrow$ & \textit{9500} & /bralaglar/&$\rightarrow$ & \textit{11000} & /mralaglar/& \\
\end{tabular}

\vspace{20pt}\hline

\end{nopagebreak}
\filbreak



\vspace{15pt}
\begin{nopagebreak}
\noindent{\fontsize{20pt}{10pt}\textbf{yeolaglar} } \textit{39 of spring/April 9} (N)\\
\noindent {\tovian \fontsize{20pt}{10pt} \textbf{yeolaglar} }\\
\noindent /yeol{\textprimstress}aglar/\\
\noindent lit. truth+tree\\
\noindent \textit{Wísertre}\\


\noindent History:

\vspace{-0pt}
\hspace{40pt}
\begin{tabular}{ccc}
\textit{6000} & /aryeolaglare/&$\rightarrow$ & \textit{7500} & /ryeolaglare/&$\rightarrow$ & \textit{7600} & /yeolaglare/&$\rightarrow$ & \textit{9500} & /yeolaglar/& \\
\end{tabular}

\vspace{20pt}\hline

\end{nopagebreak}
\filbreak



\vspace{15pt}
\begin{nopagebreak}
\noindent{\fontsize{20pt}{10pt}\textbf{shaninalhil} } \textit{39 of summer/July 9} (N)\\
\noindent {\tovian \fontsize{20pt}{10pt} \textbf{shaninalhil} }\\
\noindent /{\textesh}anin{\textprimstress}a{\textbeltl}il/\\
\noindent lit. berry+night\\
\noindent \textit{Twilightberry}\\


\noindent History:

\vspace{-0pt}
\hspace{40pt}
\begin{tabular}{ccc}
\textit{6000} & /{\textyogh}laninaa{\textbeltl}inla/&$\rightarrow$ & \textit{7600} & /{\textyogh}laninaa{\textbeltl}ila/&$\rightarrow$ & \textit{9500} & /{\textyogh}laninaa{\textbeltl}il/&$\rightarrow$ & \textit{12001} & /{\textesh}laninaa{\textbeltl}il/&$\rightarrow$ & \textit{12005} & /{\textesh}lanina{\textbeltl}il/&$\rightarrow$ & \textit{12500} & /{\textesh}anina{\textbeltl}il/& \\
\end{tabular}

\vspace{20pt}\hline

\end{nopagebreak}
\filbreak



\vspace{15pt}
\begin{nopagebreak}
\noindent{\fontsize{20pt}{10pt}\textbf{lhenashanin} } \textit{39 of winter/January 39} (N)\\
\noindent {\tovian \fontsize{20pt}{10pt} \textbf{lhenashanin} }\\
\noindent /{\textbeltl}ena{\textesh}{\textprimstress}anin/\\
\noindent lit. spirit+berry\\
\noindent \textit{Spirbery}\\


\noindent History:

\vspace{-0pt}
\hspace{40pt}
\begin{tabular}{ccc}
\textit{6000} & /{\textbeltl}ena{\textyogh}lanina/&$\rightarrow$ & \textit{9500} & /{\textbeltl}ena{\textyogh}lanin/&$\rightarrow$ & \textit{12001} & /{\textbeltl}ena{\textesh}lanin/&$\rightarrow$ & \textit{12500} & /{\textbeltl}ena{\textesh}anin/& \\
\end{tabular}

\vspace{20pt}\hline

\end{nopagebreak}
\filbreak



\vspace{15pt}
\begin{nopagebreak}
\noindent{\fontsize{20pt}{10pt}\textbf{???} } \textit{4 of autumn/September 4} (N)\\
\noindent {\tovian \fontsize{20pt}{10pt} \textbf{???} }\\
\noindent /???/\\
\noindent lit. ???\\
\noindent \textit{Daeronica}\\


\noindent History:

\vspace{-0pt}
\hspace{40pt}
\begin{tabular}{ccc}
\textit{6000} & /???/& \\
\end{tabular}

\vspace{20pt}\hline

\end{nopagebreak}
\filbreak



\vspace{15pt}
\begin{nopagebreak}
\noindent{\fontsize{20pt}{10pt}\textbf{dremrem} } \textit{4 of spring/March 4} (N)\\
\noindent {\tovian \fontsize{20pt}{10pt} \textbf{dremrem} }\\
\noindent /dr{\textprimstress}emrem/\\
\noindent lit. flower+flower\\
\noindent \textit{Drembloom}\\


\noindent History:

\vspace{-0pt}
\hspace{40pt}
\begin{tabular}{ccc}
\textit{6000} & /drempdremp/&$\rightarrow$ & \textit{8500} & /dremremp/&$\rightarrow$ & \textit{8750} & /dremrem/& \\
\end{tabular}

\vspace{20pt}\hline

\end{nopagebreak}
\filbreak



\vspace{15pt}
\begin{nopagebreak}
\noindent{\fontsize{20pt}{10pt}\textbf{lel} } \textit{4 of summer/June 4} (N)\\
\noindent {\tovian \fontsize{20pt}{10pt} \textbf{lel} }\\
\noindent /l{\textprimstress}el/\\
\noindent lit. sun+branch\\
\noindent \textit{Sunpetal}\\


\noindent History:

\vspace{-0pt}
\hspace{40pt}
\begin{tabular}{ccc}
\textit{6000} & /eflkelu/&$\rightarrow$ & \textit{7500} & /flkelu/&$\rightarrow$ & \textit{7501} & /lkelu/&$\rightarrow$ & \textit{8750} & /lelu/&$\rightarrow$ & \textit{9500} & /lel/& \\
\end{tabular}

\vspace{20pt}\hline

\end{nopagebreak}
\filbreak



\vspace{15pt}
\begin{nopagebreak}
\noindent{\fontsize{20pt}{10pt}\textbf{???} } \textit{4 of winter/december 4} (N)\\
\noindent {\tovian \fontsize{20pt}{10pt} \textbf{???} }\\
\noindent /???/\\
\noindent lit. ???\\
\noindent \textit{Stormprive}\\


\noindent History:

\vspace{-0pt}
\hspace{40pt}
\begin{tabular}{ccc}
\textit{6000} & /???/& \\
\end{tabular}

\vspace{20pt}\hline

\end{nopagebreak}
\filbreak



\vspace{15pt}
\begin{nopagebreak}
\noindent{\fontsize{20pt}{10pt}\textbf{nanathan} } \textit{40 of autumn/October 10} (N)\\
\noindent {\tovian \fontsize{20pt}{10pt} \textbf{nanathan} }\\
\noindent /nan{\textprimstress}a{\texttheta}an/\\
\noindent lit. star+tool\\
\noindent \textit{Starladder}\\


\noindent History:

\vspace{-0pt}
\hspace{40pt}
\begin{tabular}{ccc}
\textit{6000} & /nnaana{\dh}an/&$\rightarrow$ & \textit{6500} & /naana{\dh}an/&$\rightarrow$ & \textit{12002} & /naana{\texttheta}an/&$\rightarrow$ & \textit{12005} & /nana{\texttheta}an/& \\
\end{tabular}

\vspace{20pt}\hline

\end{nopagebreak}
\filbreak



\vspace{15pt}
\begin{nopagebreak}
\noindent{\fontsize{20pt}{10pt}\textbf{rithilenum} } \textit{40 of spring} (N)\\
\noindent {\tovian \fontsize{20pt}{10pt} \textbf{rithilenum} }\\
\noindent /ri{\texttheta}il{\textprimstress}enum/\\
\noindent lit. sky+path\\
\noindent \textit{Cloudfurrow}\\


\noindent History:

\vspace{-0pt}
\hspace{40pt}
\begin{tabular}{ccc}
\textit{6000} & /ri{\texttheta}ilenum/& \\
\end{tabular}

\vspace{20pt}\hline

\end{nopagebreak}
\filbreak



\vspace{15pt}
\begin{nopagebreak}
\noindent{\fontsize{20pt}{10pt}\textbf{selnan} } \textit{40 of summer/July 10} (N)\\
\noindent {\tovian \fontsize{20pt}{10pt} \textbf{selnan} }\\
\noindent /s{\textprimstress}elnan/\\
\noindent lit. water+protection\\
\noindent \textit{Rainkeeper}\\


\noindent History:

\vspace{-0pt}
\hspace{40pt}
\begin{tabular}{ccc}
\textit{6000} & /zelbnani/&$\rightarrow$ & \textit{8750} & /zelnani/&$\rightarrow$ & \textit{9500} & /zelnan/&$\rightarrow$ & \textit{12000} & /selnan/& \\
\end{tabular}

\vspace{20pt}\hline

\end{nopagebreak}
\filbreak



\vspace{15pt}
\begin{nopagebreak}
\noindent{\fontsize{20pt}{10pt}\textbf{shaingtlurith} } \textit{40 of winter/January 40} (N)\\
\noindent {\tovian \fontsize{20pt}{10pt} \textbf{shaingtlurith} }\\
\noindent /{\textesh}ai{\ng}{\texttoptiebar{t\textbeltl}}{\textprimstress}uri{\texttheta}/\\
\noindent lit. ice+fog\\
\noindent \textit{Frostfoge}\\


\noindent History:

\vspace{-0pt}
\hspace{40pt}
\begin{tabular}{ccc}
\textit{6000} & /{\textesh}u{\textesh}ai{\ng}g{\texttoptiebar{t\textbeltl}}unri{\texttheta}/&$\rightarrow$ & \textit{7600} & /{\textesh}u{\textesh}ai{\ng}g{\texttoptiebar{t\textbeltl}}uri{\texttheta}/&$\rightarrow$ & \textit{8500} & /{\textesh}u{\textesh}ai{\ng}{\texttoptiebar{t\textbeltl}}uri{\texttheta}/&$\rightarrow$ & \textit{9200} & /{\textesh}{\textschwa}{\textesh}ai{\ng}{\texttoptiebar{t\textbeltl}}uri{\texttheta}/&$\rightarrow$ & \textit{11990} & /{\textesh}{\textesh}ai{\ng}{\texttoptiebar{t\textbeltl}}uri{\texttheta}/&$\rightarrow$ & \textit{12004} & /{\textesh}ai{\ng}{\texttoptiebar{t\textbeltl}}uri{\texttheta}/& \\
\end{tabular}

\vspace{20pt}\hline

\end{nopagebreak}
\filbreak



\vspace{15pt}
\begin{nopagebreak}
\noindent{\fontsize{20pt}{10pt}\textbf{ysanin} } \textit{41 of autumn/October 11} (N)\\
\noindent {\tovian \fontsize{20pt}{10pt} \textbf{ysanin} }\\
\noindent /ys{\textprimstress}anin/\\
\noindent lit. moon+berry\\
\noindent \textit{Moonmelon}\\


\noindent History:

\vspace{-0pt}
\hspace{40pt}
\begin{tabular}{ccc}
\textit{6000} & /aliys{\textyogh}lanina/&$\rightarrow$ & \textit{6300} & /alys{\textyogh}lanina/&$\rightarrow$ & \textit{7500} & /lys{\textyogh}lanina/&$\rightarrow$ & \textit{7501} & /ys{\textyogh}lanina/&$\rightarrow$ & \textit{9500} & /ys{\textyogh}lanin/&$\rightarrow$ & \textit{12001} & /ys{\textesh}lanin/&$\rightarrow$ & \textit{12500} & /ys{\textesh}anin/&$\rightarrow$ & \textit{13000} & /ysanin/& \\
\end{tabular}

\vspace{20pt}\hline

\end{nopagebreak}
\filbreak



\vspace{15pt}
\begin{nopagebreak}
\noindent{\fontsize{20pt}{10pt}\textbf{lanatan} } \textit{41 of spring/April 11} (N)\\
\noindent {\tovian \fontsize{20pt}{10pt} \textbf{lanatan} }\\
\noindent /lan{\textprimstress}a{\textsubbridge{t}}an/\\
\noindent lit. twilight+face\\
\noindent \textit{Duskwhísker}\\


\noindent History:

\vspace{-0pt}
\hspace{40pt}
\begin{tabular}{ccc}
\textit{6000} & /brlana{\textsubbridge{t}}arn/&$\rightarrow$ & \textit{7501} & /rlana{\textsubbridge{t}}arn/&$\rightarrow$ & \textit{7600} & /lana{\textsubbridge{t}}an/& \\
\end{tabular}

\vspace{20pt}\hline

\end{nopagebreak}
\filbreak



\vspace{15pt}
\begin{nopagebreak}
\noindent{\fontsize{20pt}{10pt}\textbf{melsyaleth} } \textit{41 of summer/July 11} (N)\\
\noindent {\tovian \fontsize{20pt}{10pt} \textbf{melsyaleth} }\\
\noindent /melsy{\textprimstress}ale{\texttheta}/\\
\noindent lit. life+food\\
\noindent \textit{Heartbread}\\


\noindent History:

\vspace{-0pt}
\hspace{40pt}
\begin{tabular}{ccc}
\textit{6000} & /melsyale{\dh}a/&$\rightarrow$ & \textit{9500} & /melsyale{\dh}/&$\rightarrow$ & \textit{12002} & /melsyale{\texttheta}/& \\
\end{tabular}

\vspace{20pt}\hline

\end{nopagebreak}
\filbreak



\vspace{15pt}
\begin{nopagebreak}
\noindent{\fontsize{20pt}{10pt}\textbf{nanaingtlun} } \textit{41 of winter/January 41} (N)\\
\noindent {\tovian \fontsize{20pt}{10pt} \textbf{nanaingtlun} }\\
\noindent /nana{\textprimstress}i{\ng}{\texttoptiebar{t\textbeltl}}un/\\
\noindent lit. star+tar\\
\noindent \textit{Starwax}\\


\noindent History:

\vspace{-0pt}
\hspace{40pt}
\begin{tabular}{ccc}
\textit{6000} & /nnaanai{\ng}g{\texttoptiebar{t\textbeltl}}una/&$\rightarrow$ & \textit{6500} & /naanai{\ng}g{\texttoptiebar{t\textbeltl}}una/&$\rightarrow$ & \textit{8500} & /naanai{\ng}{\texttoptiebar{t\textbeltl}}una/&$\rightarrow$ & \textit{9500} & /naanai{\ng}{\texttoptiebar{t\textbeltl}}un/&$\rightarrow$ & \textit{12005} & /nanai{\ng}{\texttoptiebar{t\textbeltl}}un/& \\
\end{tabular}

\vspace{20pt}\hline

\end{nopagebreak}
\filbreak



\vspace{15pt}
\begin{nopagebreak}
\noindent{\fontsize{20pt}{10pt}\textbf{kelumral} } \textit{42 of autumn/October 12} (N)\\
\noindent {\tovian \fontsize{20pt}{10pt} \textbf{kelumral} }\\
\noindent /kel{\textprimstress}umral/\\
\noindent lit. branch+darkness\\
\noindent \textit{Starfennel}\\


\noindent History:

\vspace{-0pt}
\hspace{40pt}
\begin{tabular}{ccc}
\textit{6000} & /kelubral/&$\rightarrow$ & \textit{11000} & /kelumral/& \\
\end{tabular}

\vspace{20pt}\hline

\end{nopagebreak}
\filbreak



\vspace{15pt}
\begin{nopagebreak}
\noindent{\fontsize{20pt}{10pt}\textbf{ysotel} } \textit{42 of spring/April 12} (N)\\
\noindent {\tovian \fontsize{20pt}{10pt} \textbf{ysotel} }\\
\noindent /ys{\textprimstress}o{\textsubbridge{t}}el/\\
\noindent lit. moon+lake\\
\noindent \textit{Moonpool}\\


\noindent History:

\vspace{-0pt}
\hspace{40pt}
\begin{tabular}{ccc}
\textit{6000} & /aliyso{\textsubbridge{t}}ele/&$\rightarrow$ & \textit{6300} & /alyso{\textsubbridge{t}}ele/&$\rightarrow$ & \textit{7500} & /lyso{\textsubbridge{t}}ele/&$\rightarrow$ & \textit{7501} & /yso{\textsubbridge{t}}ele/&$\rightarrow$ & \textit{9500} & /yso{\textsubbridge{t}}el/& \\
\end{tabular}

\vspace{20pt}\hline

\end{nopagebreak}
\filbreak



\vspace{15pt}
\begin{nopagebreak}
\noindent{\fontsize{20pt}{10pt}\textbf{nomen} } \textit{42 of summer/July 12} (N)\\
\noindent {\tovian \fontsize{20pt}{10pt} \textbf{nomen} }\\
\noindent /n{\textprimstress}omen/\\
\noindent lit. sea+stone\\
\noindent \textit{Seaspire}\\


\noindent History:

\vspace{-0pt}
\hspace{40pt}
\begin{tabular}{ccc}
\textit{6000} & /normern/&$\rightarrow$ & \textit{7600} & /nomen/& \\
\end{tabular}

\vspace{20pt}\hline

\end{nopagebreak}
\filbreak



\vspace{15pt}
\begin{nopagebreak}
\noindent{\fontsize{20pt}{10pt}\textbf{???} } \textit{42 of winter/January 42} (N)\\
\noindent {\tovian \fontsize{20pt}{10pt} \textbf{???} }\\
\noindent /???/\\
\noindent lit. ???\\
\noindent \textit{Namiloray}\\


\noindent History:

\vspace{-0pt}
\hspace{40pt}
\begin{tabular}{ccc}
\textit{6000} & /???/& \\
\end{tabular}

\vspace{20pt}\hline

\end{nopagebreak}
\filbreak



\vspace{15pt}
\begin{nopagebreak}
\noindent{\fontsize{20pt}{10pt}\textbf{shithol} } \textit{43 of autumn/October 13} (N)\\
\noindent {\tovian \fontsize{20pt}{10pt} \textbf{shithol} }\\
\noindent /{\textesh}{\textprimstress}i{\texttheta}ol/\\
\noindent lit. sharpness+sweet\\
\noindent \textit{Thornsweet}\\


\noindent History:

\vspace{-0pt}
\hspace{40pt}
\begin{tabular}{ccc}
\textit{6000} & /{\textyogh}i{\texttheta}hola/&$\rightarrow$ & \textit{9500} & /{\textyogh}i{\texttheta}hol/&$\rightarrow$ & \textit{12001} & /{\textesh}i{\texttheta}hol/&$\rightarrow$ & \textit{13000} & /{\textesh}i{\texttheta}ol/& \\
\end{tabular}

\vspace{20pt}\hline

\end{nopagebreak}
\filbreak



\vspace{15pt}
\begin{nopagebreak}
\noindent{\fontsize{20pt}{10pt}\textbf{mrilaglar} } \textit{43 of spring/April 13} (N)\\
\noindent {\tovian \fontsize{20pt}{10pt} \textbf{mrilaglar} }\\
\noindent /mril{\textprimstress}aglar/\\
\noindent lit. sight+tree\\
\noindent \textit{Mrilaple}\\


\noindent History:

\vspace{-0pt}
\hspace{40pt}
\begin{tabular}{ccc}
\textit{6000} & /emrilaaglare/&$\rightarrow$ & \textit{7500} & /mrilaaglare/&$\rightarrow$ & \textit{9500} & /mrilaaglar/&$\rightarrow$ & \textit{12005} & /mrilaglar/& \\
\end{tabular}

\vspace{20pt}\hline

\end{nopagebreak}
\filbreak



\vspace{15pt}
\begin{nopagebreak}
\noindent{\fontsize{20pt}{10pt}\textbf{mloanafor} } \textit{43 of summer/July 13} (N)\\
\noindent {\tovian \fontsize{20pt}{10pt} \textbf{mloanafor} }\\
\noindent /mloan{\textprimstress}afor/\\
\noindent lit. morning+writing\\
\noindent \textit{Dawnfelor}\\


\noindent History:

\vspace{-0pt}
\hspace{40pt}
\begin{tabular}{ccc}
\textit{6000} & /ploanaflore/&$\rightarrow$ & \textit{9500} & /ploanaflor/&$\rightarrow$ & \textit{11000} & /mloanaflor/&$\rightarrow$ & \textit{12500} & /mloanafor/& \\
\end{tabular}

\vspace{20pt}\hline

\end{nopagebreak}
\filbreak



\vspace{15pt}
\begin{nopagebreak}
\noindent{\fontsize{20pt}{10pt}\textbf{nanador} } \textit{43 of winter/January 43} (N)\\
\noindent {\tovian \fontsize{20pt}{10pt} \textbf{nanador} }\\
\noindent /nan{\textprimstress}ador/\\
\noindent lit. star+blue\\
\noindent \textit{Starwoad}\\


\noindent History:

\vspace{-0pt}
\hspace{40pt}
\begin{tabular}{ccc}
\textit{6000} & /nnaanadora/&$\rightarrow$ & \textit{6500} & /naanadora/&$\rightarrow$ & \textit{9500} & /naanador/&$\rightarrow$ & \textit{12005} & /nanador/& \\
\end{tabular}

\vspace{20pt}\hline

\end{nopagebreak}
\filbreak



\vspace{15pt}
\begin{nopagebreak}
\noindent{\fontsize{20pt}{10pt}\textbf{munorhanashol} } \textit{44 of autumn/October 14} (N)\\
\noindent {\tovian \fontsize{20pt}{10pt} \textbf{munorhanashol} }\\
\noindent /munorħan{\textprimstress}a{\textesh}ol/\\
\noindent lit. dusk+fruit\\
\noindent \textit{Dusknut}\\


\noindent History:

\vspace{-0pt}
\hspace{40pt}
\begin{tabular}{ccc}
\textit{6000} & /pnorhana{\textyogh}ola/&$\rightarrow$ & \textit{9500} & /pnorhana{\textyogh}ol/&$\rightarrow$ & \textit{11000} & /mnorhana{\textyogh}ol/&$\rightarrow$ & \textit{11996} & /munorhana{\textyogh}ol/&$\rightarrow$ & \textit{12001} & /munorhana{\textesh}ol/&$\rightarrow$ & \textit{14006} & /munorħana{\textesh}ol/& \\
\end{tabular}

\vspace{20pt}\hline

\end{nopagebreak}
\filbreak



\vspace{15pt}
\begin{nopagebreak}
\noindent{\fontsize{20pt}{10pt}\textbf{rithilefethr} } \textit{44 of spring/April 14} (N)\\
\noindent {\tovian \fontsize{20pt}{10pt} \textbf{rithilefethr} }\\
\noindent /ri{\texttheta}il{\textprimstress}efe{\texttheta}r/\\
\noindent lit. sky+speech\\
\noindent \textit{Skyfethr}\\


\noindent History:

\vspace{-0pt}
\hspace{40pt}
\begin{tabular}{ccc}
\textit{6000} & /ri{\texttheta}ilefe{\dh}ra/&$\rightarrow$ & \textit{9500} & /ri{\texttheta}ilefe{\dh}r/&$\rightarrow$ & \textit{12002} & /ri{\texttheta}ilefe{\texttheta}r/& \\
\end{tabular}

\vspace{20pt}\hline

\end{nopagebreak}
\filbreak



\vspace{15pt}
\begin{nopagebreak}
\noindent{\fontsize{20pt}{10pt}\textbf{thangekel} } \textit{44 of summer/July 14} (N)\\
\noindent {\tovian \fontsize{20pt}{10pt} \textbf{thangekel} }\\
\noindent /{\texttheta}a{\ng}{\textprimstress}ekel/\\
\noindent lit. leaf+branch\\
\noindent \textit{Starherb}\\


\noindent History:

\vspace{-0pt}
\hspace{40pt}
\begin{tabular}{ccc}
\textit{6000} & /a{\dh}a{\ng}gekelu/&$\rightarrow$ & \textit{7500} & /{\dh}a{\ng}gekelu/&$\rightarrow$ & \textit{8750} & /{\dh}a{\ng}ekelu/&$\rightarrow$ & \textit{9500} & /{\dh}a{\ng}ekel/&$\rightarrow$ & \textit{12002} & /{\texttheta}a{\ng}ekel/& \\
\end{tabular}

\vspace{20pt}\hline

\end{nopagebreak}
\filbreak



\vspace{15pt}
\begin{nopagebreak}
\noindent{\fontsize{20pt}{10pt}\textbf{shakel} } \textit{44 of winter/January 44} (N)\\
\noindent {\tovian \fontsize{20pt}{10pt} \textbf{shakel} }\\
\noindent /{\textesh}{\textprimstress}akel/\\
\noindent lit. ice+branch\\
\noindent \textit{Frostpine}\\


\noindent History:

\vspace{-0pt}
\hspace{40pt}
\begin{tabular}{ccc}
\textit{6000} & /{\textesh}u{\textesh}akelu/&$\rightarrow$ & \textit{9200} & /{\textesh}{\textschwa}{\textesh}akelu/&$\rightarrow$ & \textit{9500} & /{\textesh}{\textschwa}{\textesh}akel/&$\rightarrow$ & \textit{11990} & /{\textesh}{\textesh}akel/&$\rightarrow$ & \textit{12007} & /{\textesh}akel/& \\
\end{tabular}

\vspace{20pt}\hline

\end{nopagebreak}
\filbreak



\vspace{15pt}
\begin{nopagebreak}
\noindent{\fontsize{20pt}{10pt}\textbf{???} } \textit{45 of autumn/October 15} (N)\\
\noindent {\tovian \fontsize{20pt}{10pt} \textbf{???} }\\
\noindent /???/\\
\noindent lit. ???\\
\noindent \textit{Stamglith}\\


\noindent History:

\vspace{-0pt}
\hspace{40pt}
\begin{tabular}{ccc}
\textit{6000} & /???/& \\
\end{tabular}

\vspace{20pt}\hline

\end{nopagebreak}
\filbreak



\vspace{15pt}
\begin{nopagebreak}
\noindent{\fontsize{20pt}{10pt}\textbf{lrithitan} } \textit{45 of spring/April 15} (N)\\
\noindent {\tovian \fontsize{20pt}{10pt} \textbf{lrithitan} }\\
\noindent /lri{\texttheta}{\textprimstress}i{\textsubbridge{t}}an/\\
\noindent lit. storm+bird\\
\noindent \textit{Stormbird}\\


\noindent History:

\vspace{-0pt}
\hspace{40pt}
\begin{tabular}{ccc}
\textit{6000} & /kwlyri{\texttheta}lri{\texttheta}lamn/&$\rightarrow$ & \textit{6503} & /kwlyri{\texttheta}rli{\texttheta}lamn/&$\rightarrow$ & \textit{7501} & /wlyri{\texttheta}rli{\texttheta}lamn/&$\rightarrow$ & \textit{7600} & /lri{\texttheta}li{\texttheta}lan/&$\rightarrow$ & \textit{12500} & /lri{\texttheta}i{\texttheta}an/&$\rightarrow$ & \textit{12501} & /lri{\texttheta}i{\textsubbridge{t}}an/& \\
\end{tabular}

\vspace{20pt}\hline

\end{nopagebreak}
\filbreak



\vspace{15pt}
\begin{nopagebreak}
\noindent{\fontsize{20pt}{10pt}\textbf{lhenaysan} } \textit{45 of summer/July 15} (N)\\
\noindent {\tovian \fontsize{20pt}{10pt} \textbf{lhenaysan} }\\
\noindent /{\textbeltl}en{\textprimstress}aysan/\\
\noindent lit. dream+moon+animal\\
\noindent \textit{Moonewe}\\


\noindent History:

\vspace{-0pt}
\hspace{40pt}
\begin{tabular}{ccc}
\textit{6000} & /{\textbeltl}enaaliys{\textsubbridge{t}}lamn/&$\rightarrow$ & \textit{6300} & /{\textbeltl}enaalys{\textsubbridge{t}}lamn/&$\rightarrow$ & \textit{6401} & /{\textbeltl}enaalys{\texttoptiebar{t\textbeltl}}amn/&$\rightarrow$ & \textit{7600} & /{\textbeltl}enaays{\texttoptiebar{t\textbeltl}}an/&$\rightarrow$ & \textit{12005} & /{\textbeltl}enays{\texttoptiebar{t\textbeltl}}an/&$\rightarrow$ & \textit{14001} & /{\textbeltl}enaysan/& \\
\end{tabular}

\vspace{20pt}\hline

\end{nopagebreak}
\filbreak



\vspace{15pt}
\begin{nopagebreak}
\noindent{\fontsize{20pt}{10pt}\textbf{rishwen} } \textit{45 of winter/January 45} (N)\\
\noindent {\tovian \fontsize{20pt}{10pt} \textbf{rishwen} }\\
\noindent /r{\textprimstress}i{\textesh}wen/\\
\noindent lit. kindness+beauty\\
\noindent \textit{Rashelwen}\\


\noindent History:

\vspace{-0pt}
\hspace{40pt}
\begin{tabular}{ccc}
\textit{6000} & /ri{\textesh}lwen/&$\rightarrow$ & \textit{7600} & /ri{\textesh}wen/& \\
\end{tabular}

\vspace{20pt}\hline

\end{nopagebreak}
\filbreak



\vspace{15pt}
\begin{nopagebreak}
\noindent{\fontsize{20pt}{10pt}\textbf{for} } \textit{46 of autumn/October 16} (N)\\
\noindent {\tovian \fontsize{20pt}{10pt} \textbf{for} }\\
\noindent /f{\textprimstress}or/\\
\noindent lit. sun+action\\
\noindent \textit{Sunzest}\\


\noindent History:

\vspace{-0pt}
\hspace{40pt}
\begin{tabular}{ccc}
\textit{6000} & /efllore/&$\rightarrow$ & \textit{6500} & /eflore/&$\rightarrow$ & \textit{7500} & /flore/&$\rightarrow$ & \textit{9500} & /flor/&$\rightarrow$ & \textit{12500} & /for/& \\
\end{tabular}

\vspace{20pt}\hline

\end{nopagebreak}
\filbreak



\vspace{15pt}
\begin{nopagebreak}
\noindent{\fontsize{20pt}{10pt}\textbf{nanatom} } \textit{46 of spring/April 16} (N)\\
\noindent {\tovian \fontsize{20pt}{10pt} \textbf{nanatom} }\\
\noindent /nan{\textprimstress}a{\textsubbridge{t}}om/\\
\noindent lit. star+root\\
\noindent \textit{Starroot}\\


\noindent History:

\vspace{-0pt}
\hspace{40pt}
\begin{tabular}{ccc}
\textit{6000} & /nnaana{\textsubbridge{t}}omo/&$\rightarrow$ & \textit{6500} & /naana{\textsubbridge{t}}omo/&$\rightarrow$ & \textit{9500} & /naana{\textsubbridge{t}}om/&$\rightarrow$ & \textit{12005} & /nana{\textsubbridge{t}}om/& \\
\end{tabular}

\vspace{20pt}\hline

\end{nopagebreak}
\filbreak



\vspace{15pt}
\begin{nopagebreak}
\noindent{\fontsize{20pt}{10pt}\textbf{lhenalor} } \textit{46 of summer/July 16} (N)\\
\noindent {\tovian \fontsize{20pt}{10pt} \textbf{lhenalor} }\\
\noindent /{\textbeltl}en{\textprimstress}alor/\\
\noindent lit. dream+plant\\
\noindent \textit{Dreammallow}\\


\noindent History:

\vspace{-0pt}
\hspace{40pt}
\begin{tabular}{ccc}
\textit{6000} & /{\textbeltl}enamllore/&$\rightarrow$ & \textit{6500} & /{\textbeltl}enamlore/&$\rightarrow$ & \textit{7600} & /{\textbeltl}enalore/&$\rightarrow$ & \textit{9500} & /{\textbeltl}enalor/& \\
\end{tabular}

\vspace{20pt}\hline

\end{nopagebreak}
\filbreak



\vspace{15pt}
\begin{nopagebreak}
\noindent{\fontsize{20pt}{10pt}\textbf{thatlun} } \textit{46 of winter/January 46} (N)\\
\noindent {\tovian \fontsize{20pt}{10pt} \textbf{thatlun} }\\
\noindent /{\texttheta}{\textprimstress}a{\texttoptiebar{t\textbeltl}}un/\\
\noindent lit. tool+joining\\
\noindent \textit{Thornweave}\\


\noindent History:

\vspace{-0pt}
\hspace{40pt}
\begin{tabular}{ccc}
\textit{6000} & /{\dh}an{\texttoptiebar{t\textbeltl}}una/&$\rightarrow$ & \textit{6400} & /{\dh}a{\texttoptiebar{t\textbeltl}}una/&$\rightarrow$ & \textit{9500} & /{\dh}a{\texttoptiebar{t\textbeltl}}un/&$\rightarrow$ & \textit{12002} & /{\texttheta}a{\texttoptiebar{t\textbeltl}}un/& \\
\end{tabular}

\vspace{20pt}\hline

\end{nopagebreak}
\filbreak



\vspace{15pt}
\begin{nopagebreak}
\noindent{\fontsize{20pt}{10pt}\textbf{lenan} } \textit{47 of autumn/October 17} (N)\\
\noindent {\tovian \fontsize{20pt}{10pt} \textbf{lenan} }\\
\noindent /l{\textprimstress}enan/\\
\noindent lit. hoe+covering\\
\noindent \textit{Melorethangor}\\


\noindent History:

\vspace{-0pt}
\hspace{40pt}
\begin{tabular}{ccc}
\textit{6000} & /mllre{\dh}nbnani/&$\rightarrow$ & \textit{6500} & /mlre{\dh}nbnani/&$\rightarrow$ & \textit{6503} & /mrle{\dh}nbnani/&$\rightarrow$ & \textit{7501} & /rle{\dh}nbnani/&$\rightarrow$ & \textit{7600} & /le{\dh}nbnani/&$\rightarrow$ & \textit{8500} & /le{\dh}nnani/&$\rightarrow$ & \textit{9500} & /le{\dh}nnan/&$\rightarrow$ & \textit{12002} & /le{\texttheta}nnan/&$\rightarrow$ & \textit{12007} & /le{\texttheta}nan/&$\rightarrow$ & \textit{13002} & /lenan/& \\
\end{tabular}

\vspace{20pt}\hline

\end{nopagebreak}
\filbreak



\vspace{15pt}
\begin{nopagebreak}
\noindent{\fontsize{20pt}{10pt}\textbf{shasel} } \textit{47 of spring/April 17} (N)\\
\noindent {\tovian \fontsize{20pt}{10pt} \textbf{shasel} }\\
\noindent /{\textesh}{\textprimstress}asel/\\
\noindent lit. ice+water\\
\noindent \textit{Mooncress}\\


\noindent History:

\vspace{-0pt}
\hspace{40pt}
\begin{tabular}{ccc}
\textit{6000} & /{\textesh}u{\textesh}azel/&$\rightarrow$ & \textit{9200} & /{\textesh}{\textschwa}{\textesh}azel/&$\rightarrow$ & \textit{11990} & /{\textesh}{\textesh}azel/&$\rightarrow$ & \textit{12000} & /{\textesh}{\textesh}asel/&$\rightarrow$ & \textit{12007} & /{\textesh}asel/& \\
\end{tabular}

\vspace{20pt}\hline

\end{nopagebreak}
\filbreak



\vspace{15pt}
\begin{nopagebreak}
\noindent{\fontsize{20pt}{10pt}\textbf{shaninakel} } \textit{47 of summer/July 17} (N)\\
\noindent {\tovian \fontsize{20pt}{10pt} \textbf{shaninakel} }\\
\noindent /{\textesh}anin{\textprimstress}akel/\\
\noindent lit. berry+branch\\
\noindent \textit{Starflax}\\


\noindent History:

\vspace{-0pt}
\hspace{40pt}
\begin{tabular}{ccc}
\textit{6000} & /{\textyogh}laninakelu/&$\rightarrow$ & \textit{9500} & /{\textyogh}laninakel/&$\rightarrow$ & \textit{12001} & /{\textesh}laninakel/&$\rightarrow$ & \textit{12500} & /{\textesh}aninakel/& \\
\end{tabular}

\vspace{20pt}\hline

\end{nopagebreak}
\filbreak



\vspace{15pt}
\begin{nopagebreak}
\noindent{\fontsize{20pt}{10pt}\textbf{???} } \textit{47 of winter/January 47} (N)\\
\noindent {\tovian \fontsize{20pt}{10pt} \textbf{???} }\\
\noindent /???/\\
\noindent lit. ???\\
\noindent \textit{Nicypres}\\


\noindent History:

\vspace{-0pt}
\hspace{40pt}
\begin{tabular}{ccc}
\textit{6000} & /???/& \\
\end{tabular}

\vspace{20pt}\hline

\end{nopagebreak}
\filbreak



\vspace{15pt}
\begin{nopagebreak}
\noindent{\fontsize{20pt}{10pt}\textbf{mloanashith} } \textit{48 of autumn/October 18} (N)\\
\noindent {\tovian \fontsize{20pt}{10pt} \textbf{mloanashith} }\\
\noindent /mloan{\textprimstress}a{\textesh}i{\texttheta}/\\
\noindent lit. morning+sharpness\\
\noindent \textit{Dawnthor}\\


\noindent History:

\vspace{-0pt}
\hspace{40pt}
\begin{tabular}{ccc}
\textit{6000} & /ploana{\textyogh}i{\texttheta}/&$\rightarrow$ & \textit{11000} & /mloana{\textyogh}i{\texttheta}/&$\rightarrow$ & \textit{12001} & /mloana{\textesh}i{\texttheta}/& \\
\end{tabular}

\vspace{20pt}\hline

\end{nopagebreak}
\filbreak



\vspace{15pt}
\begin{nopagebreak}
\noindent{\fontsize{20pt}{10pt}\textbf{shagang} } \textit{48 of spring/April 18} (N)\\
\noindent {\tovian \fontsize{20pt}{10pt} \textbf{shagang} }\\
\noindent /{\textesh}{\textprimstress}aga{\ng}/\\
\noindent lit. ice+tooth\\
\noindent \textit{Frostdent}\\


\noindent History:

\vspace{-0pt}
\hspace{40pt}
\begin{tabular}{ccc}
\textit{6000} & /{\textesh}u{\textesh}aga{\ng}ga/&$\rightarrow$ & \textit{8750} & /{\textesh}u{\textesh}aga{\ng}a/&$\rightarrow$ & \textit{9200} & /{\textesh}{\textschwa}{\textesh}aga{\ng}a/&$\rightarrow$ & \textit{9500} & /{\textesh}{\textschwa}{\textesh}aga{\ng}/&$\rightarrow$ & \textit{11990} & /{\textesh}{\textesh}aga{\ng}/&$\rightarrow$ & \textit{12007} & /{\textesh}aga{\ng}/& \\
\end{tabular}

\vspace{20pt}\hline

\end{nopagebreak}
\filbreak



\vspace{15pt}
\begin{nopagebreak}
\noindent{\fontsize{20pt}{10pt}\textbf{holashol} } \textit{48 of summer/July 18} (N)\\
\noindent {\tovian \fontsize{20pt}{10pt} \textbf{holashol} }\\
\noindent /ħol{\textprimstress}a{\textesh}ol/\\
\noindent lit. sweet+fruit\\
\noindent \textit{Sweetnut}\\


\noindent History:

\vspace{-0pt}
\hspace{40pt}
\begin{tabular}{ccc}
\textit{6000} & /hola{\textyogh}ola/&$\rightarrow$ & \textit{9500} & /hola{\textyogh}ol/&$\rightarrow$ & \textit{12001} & /hola{\textesh}ol/&$\rightarrow$ & \textit{14006} & /ħola{\textesh}ol/& \\
\end{tabular}

\vspace{20pt}\hline

\end{nopagebreak}
\filbreak



\vspace{15pt}
\begin{nopagebreak}
\noindent{\fontsize{20pt}{10pt}\textbf{klumloretush} } \textit{48 of winter/January 48} (N)\\
\noindent {\tovian \fontsize{20pt}{10pt} \textbf{klumloretush} }\\
\noindent /klumlor{\textprimstress}e{\textsubbridge{t}}u{\textesh}/\\
\noindent lit. vine+cold\\
\noindent \textit{Winterivy}\\


\noindent History:

\vspace{-0pt}
\hspace{40pt}
\begin{tabular}{ccc}
\textit{6000} & /kluwmllore{\textsubbridge{t}}u{\textesh}u/&$\rightarrow$ & \textit{6500} & /kluwmlore{\textsubbridge{t}}u{\textesh}u/&$\rightarrow$ & \textit{7600} & /klumlore{\textsubbridge{t}}u{\textesh}u/&$\rightarrow$ & \textit{9500} & /klumlore{\textsubbridge{t}}u{\textesh}/& \\
\end{tabular}

\vspace{20pt}\hline

\end{nopagebreak}
\filbreak



\vspace{15pt}
\begin{nopagebreak}
\noindent{\fontsize{20pt}{10pt}\textbf{varethan} } \textit{49 of autumn/October 19} (N)\\
\noindent {\tovian \fontsize{20pt}{10pt} \textbf{varethan} }\\
\noindent /var{\textprimstress}e{\texttheta}an/\\
\noindent lit. value+tool\\
\noindent \textit{Golace}\\


\noindent History:

\vspace{-0pt}
\hspace{40pt}
\begin{tabular}{ccc}
\textit{6000} & /vlare{\dh}an/&$\rightarrow$ & \textit{12002} & /vlare{\texttheta}an/&$\rightarrow$ & \textit{12500} & /vare{\texttheta}an/& \\
\end{tabular}

\vspace{20pt}\hline

\end{nopagebreak}
\filbreak



\vspace{15pt}
\begin{nopagebreak}
\noindent{\fontsize{20pt}{10pt}\textbf{nanashol} } \textit{49 of spring/April 19} (N)\\
\noindent {\tovian \fontsize{20pt}{10pt} \textbf{nanashol} }\\
\noindent /nan{\textprimstress}a{\textesh}ol/\\
\noindent lit. star+fruit\\
\noindent \textit{Starfruit}\\


\noindent History:

\vspace{-0pt}
\hspace{40pt}
\begin{tabular}{ccc}
\textit{6000} & /nnaana{\textyogh}ola/&$\rightarrow$ & \textit{6500} & /naana{\textyogh}ola/&$\rightarrow$ & \textit{9500} & /naana{\textyogh}ol/&$\rightarrow$ & \textit{12001} & /naana{\textesh}ol/&$\rightarrow$ & \textit{12005} & /nana{\textesh}ol/& \\
\end{tabular}

\vspace{20pt}\hline

\end{nopagebreak}
\filbreak



\vspace{15pt}
\begin{nopagebreak}
\noindent{\fontsize{20pt}{10pt}\textbf{dolaglar} } \textit{49 of summer/July 19} (N)\\
\noindent {\tovian \fontsize{20pt}{10pt} \textbf{dolaglar} }\\
\noindent /dol{\textprimstress}aglar/\\
\noindent lit. death+tree\\
\noindent \textit{Twilightroot}\\


\noindent History:

\vspace{-0pt}
\hspace{40pt}
\begin{tabular}{ccc}
\textit{6000} & /dolaglare/&$\rightarrow$ & \textit{9500} & /dolaglar/& \\
\end{tabular}

\vspace{20pt}\hline

\end{nopagebreak}
\filbreak



\vspace{15pt}
\begin{nopagebreak}
\noindent{\fontsize{20pt}{10pt}\textbf{aymen} } \textit{49 of winter/January 49} (N)\\
\noindent {\tovian \fontsize{20pt}{10pt} \textbf{aymen} }\\
\noindent /{\textprimstress}aymen/\\
\noindent lit. moon+earth\\
\noindent \textit{Moonjuniper}\\


\noindent History:

\vspace{-0pt}
\hspace{40pt}
\begin{tabular}{ccc}
\textit{6000} & /aliysmern/&$\rightarrow$ & \textit{6300} & /alysmern/&$\rightarrow$ & \textit{7600} & /aysmen/&$\rightarrow$ & \textit{13002} & /aymen/& \\
\end{tabular}

\vspace{20pt}\hline

\end{nopagebreak}
\filbreak



\vspace{15pt}
\begin{nopagebreak}
\noindent{\fontsize{20pt}{10pt}\textbf{nanamen} } \textit{5 of autumn/September 5} (N)\\
\noindent {\tovian \fontsize{20pt}{10pt} \textbf{nanamen} }\\
\noindent /nan{\textprimstress}amen/\\
\noindent lit. star+earth\\
\noindent \textit{Starmul}\\


\noindent History:

\vspace{-0pt}
\hspace{40pt}
\begin{tabular}{ccc}
\textit{6000} & /nnaanamern/&$\rightarrow$ & \textit{6500} & /naanamern/&$\rightarrow$ & \textit{7600} & /naanamen/&$\rightarrow$ & \textit{12005} & /nanamen/& \\
\end{tabular}

\vspace{20pt}\hline

\end{nopagebreak}
\filbreak



\vspace{15pt}
\begin{nopagebreak}
\noindent{\fontsize{20pt}{10pt}\textbf{lrithan} } \textit{5 of spring/March 5} (N)\\
\noindent {\tovian \fontsize{20pt}{10pt} \textbf{lrithan} }\\
\noindent /lr{\textprimstress}i{\texttheta}an/\\
\noindent lit. wind+animal\\
\noindent \textit{Vindsteed}\\


\noindent History:

\vspace{-0pt}
\hspace{40pt}
\begin{tabular}{ccc}
\textit{6000} & /kwlyri{\texttheta}{\textsubbridge{t}}lamn/&$\rightarrow$ & \textit{6401} & /kwlyri{\texttheta}{\texttoptiebar{t\textbeltl}}amn/&$\rightarrow$ & \textit{7501} & /wlyri{\texttheta}{\texttoptiebar{t\textbeltl}}amn/&$\rightarrow$ & \textit{7600} & /lri{\texttheta}{\texttoptiebar{t\textbeltl}}an/&$\rightarrow$ & \textit{14001} & /lri{\texttheta}an/& \\
\end{tabular}

\vspace{20pt}\hline

\end{nopagebreak}
\filbreak



\vspace{15pt}
\begin{nopagebreak}
\noindent{\fontsize{20pt}{10pt}\textbf{mloanatlan} } \textit{5 of summer/June 5} (N)\\
\noindent {\tovian \fontsize{20pt}{10pt} \textbf{mloanatlan} }\\
\noindent /mloan{\textprimstress}a{\texttoptiebar{t\textbeltl}}an/\\
\noindent lit. morning+animal\\
\noindent \textit{Dawncluck}\\


\noindent History:

\vspace{-0pt}
\hspace{40pt}
\begin{tabular}{ccc}
\textit{6000} & /ploana{\textsubbridge{t}}lamn/&$\rightarrow$ & \textit{6401} & /ploana{\texttoptiebar{t\textbeltl}}amn/&$\rightarrow$ & \textit{7600} & /ploana{\texttoptiebar{t\textbeltl}}an/&$\rightarrow$ & \textit{11000} & /mloana{\texttoptiebar{t\textbeltl}}an/& \\
\end{tabular}

\vspace{20pt}\hline

\end{nopagebreak}
\filbreak



\vspace{15pt}
\begin{nopagebreak}
\noindent{\fontsize{20pt}{10pt}\textbf{lrithan} } \textit{5 of winter/december 5} (N)\\
\noindent {\tovian \fontsize{20pt}{10pt} \textbf{lrithan} }\\
\noindent /lr{\textprimstress}i{\texttheta}an/\\
\noindent lit. breath+tool\\
\noindent \textit{Windhorn}\\


\noindent History:

\vspace{-0pt}
\hspace{40pt}
\begin{tabular}{ccc}
\textit{6000} & /kwlyri{\texttheta}{\dh}an/&$\rightarrow$ & \textit{7501} & /wlyri{\texttheta}{\dh}an/&$\rightarrow$ & \textit{7600} & /lri{\texttheta}{\dh}an/&$\rightarrow$ & \textit{12002} & /lri{\texttheta}{\texttheta}an/&$\rightarrow$ & \textit{12007} & /lri{\texttheta}an/& \\
\end{tabular}

\vspace{20pt}\hline

\end{nopagebreak}
\filbreak



\vspace{15pt}
\begin{nopagebreak}
\noindent{\fontsize{20pt}{10pt}\textbf{harathan} } \textit{50 of autumn/October 20} (N)\\
\noindent {\tovian \fontsize{20pt}{10pt} \textbf{harathan} }\\
\noindent /ħar{\textprimstress}a{\texttheta}an/\\
\noindent lit. pain+tool\\
\noindent \textit{Harsack}\\


\noindent History:

\vspace{-0pt}
\hspace{40pt}
\begin{tabular}{ccc}
\textit{6000} & /hara{\dh}an/&$\rightarrow$ & \textit{12002} & /hara{\texttheta}an/&$\rightarrow$ & \textit{14006} & /ħara{\texttheta}an/& \\
\end{tabular}

\vspace{20pt}\hline

\end{nopagebreak}
\filbreak



\vspace{15pt}
\begin{nopagebreak}
\noindent{\fontsize{20pt}{10pt}\textbf{nanamnan} } \textit{50 of spring/April 20} (N)\\
\noindent {\tovian \fontsize{20pt}{10pt} \textbf{nanamnan} }\\
\noindent /nan{\textprimstress}amnan/\\
\noindent lit. star+covering\\
\noindent \textit{Starharrow}\\


\noindent History:

\vspace{-0pt}
\hspace{40pt}
\begin{tabular}{ccc}
\textit{6000} & /nnaanabnani/&$\rightarrow$ & \textit{6500} & /naanabnani/&$\rightarrow$ & \textit{9500} & /naanabnan/&$\rightarrow$ & \textit{11000} & /naanamnan/&$\rightarrow$ & \textit{12005} & /nanamnan/& \\
\end{tabular}

\vspace{20pt}\hline

\end{nopagebreak}
\filbreak



\vspace{15pt}
\begin{nopagebreak}
\noindent{\fontsize{20pt}{10pt}\textbf{selyethil} } \textit{50 of summer/July 20} (N)\\
\noindent {\tovian \fontsize{20pt}{10pt} \textbf{selyethil} }\\
\noindent /sely{\textprimstress}e{\texttheta}il/\\
\noindent lit. water+home\\
\noindent \textit{Wetgate}\\


\noindent History:

\vspace{-0pt}
\hspace{40pt}
\begin{tabular}{ccc}
\textit{6000} & /zeldryel{\texttheta}ile/&$\rightarrow$ & \textit{7600} & /zeldyel{\texttheta}ile/&$\rightarrow$ & \textit{8750} & /zelyel{\texttheta}ile/&$\rightarrow$ & \textit{9500} & /zelyel{\texttheta}il/&$\rightarrow$ & \textit{12000} & /selyel{\texttheta}il/&$\rightarrow$ & \textit{12502} & /selye{\texttheta}il/& \\
\end{tabular}

\vspace{20pt}\hline

\end{nopagebreak}
\filbreak



\vspace{15pt}
\begin{nopagebreak}
\noindent{\fontsize{20pt}{10pt}\textbf{lethiledrem} } \textit{51 of autumn/October 21} (N)\\
\noindent {\tovian \fontsize{20pt}{10pt} \textbf{lethiledrem} }\\
\noindent /le{\texttheta}il{\textprimstress}edrem/\\
\noindent lit. forest+flower\\
\noindent \textit{Wildrose}\\


\noindent History:

\vspace{-0pt}
\hspace{40pt}
\begin{tabular}{ccc}
\textit{6000} & /aglre{\texttheta}iledremp/&$\rightarrow$ & \textit{6503} & /agrle{\texttheta}iledremp/&$\rightarrow$ & \textit{7500} & /grle{\texttheta}iledremp/&$\rightarrow$ & \textit{7501} & /rle{\texttheta}iledremp/&$\rightarrow$ & \textit{7600} & /le{\texttheta}iledremp/&$\rightarrow$ & \textit{8750} & /le{\texttheta}iledrem/& \\
\end{tabular}

\vspace{20pt}\hline

\end{nopagebreak}
\filbreak



\vspace{15pt}
\begin{nopagebreak}
\noindent{\fontsize{20pt}{10pt}\textbf{shimisel} } \textit{51 of spring/April 21} (N)\\
\noindent {\tovian \fontsize{20pt}{10pt} \textbf{shimisel} }\\
\noindent /{\textesh}im{\textprimstress}isel/\\
\noindent lit. mystery+water\\
\noindent \textit{Mystbach}\\


\noindent History:

\vspace{-0pt}
\hspace{40pt}
\begin{tabular}{ccc}
\textit{6000} & /{\textesh}imizel/&$\rightarrow$ & \textit{12000} & /{\textesh}imisel/& \\
\end{tabular}

\vspace{20pt}\hline

\end{nopagebreak}
\filbreak



\vspace{15pt}
\begin{nopagebreak}
\noindent{\fontsize{20pt}{10pt}\textbf{felakorim} } \textit{51 of summer/July 21} (N)\\
\noindent {\tovian \fontsize{20pt}{10pt} \textbf{felakorim} }\\
\noindent /felak{\textprimstress}orim/\\
\noindent lit. sun+cycle\\
\noindent \textit{Sunwheel}\\


\noindent History:

\vspace{-0pt}
\hspace{40pt}
\begin{tabular}{ccc}
\textit{6000} & /eflelakorim/&$\rightarrow$ & \textit{7500} & /flelakorim/&$\rightarrow$ & \textit{12500} & /felakorim/& \\
\end{tabular}

\vspace{20pt}\hline

\end{nopagebreak}
\filbreak



\vspace{15pt}
\begin{nopagebreak}
\noindent{\fontsize{20pt}{10pt}\textbf{nanaglar} } \textit{51 of winter/January 51} (N)\\
\noindent {\tovian \fontsize{20pt}{10pt} \textbf{nanaglar} }\\
\noindent /nan{\textprimstress}aglar/\\
\noindent lit. star+tree\\
\noindent \textit{Starmaple}\\


\noindent History:

\vspace{-0pt}
\hspace{40pt}
\begin{tabular}{ccc}
\textit{6000} & /nnaanaaglare/&$\rightarrow$ & \textit{6500} & /naanaaglare/&$\rightarrow$ & \textit{9500} & /naanaaglar/&$\rightarrow$ & \textit{12005} & /nanaglar/& \\
\end{tabular}

\vspace{20pt}\hline

\end{nopagebreak}
\filbreak



\vspace{15pt}
\begin{nopagebreak}
\noindent{\fontsize{20pt}{10pt}\textbf{lhenaysol} } \textit{52 of autumn/October 22} (N)\\
\noindent {\tovian \fontsize{20pt}{10pt} \textbf{lhenaysol} }\\
\noindent /{\textbeltl}en{\textprimstress}aysol/\\
\noindent lit. dream+moon+fruit\\
\noindent \textit{Monut}\\


\noindent History:

\vspace{-0pt}
\hspace{40pt}
\begin{tabular}{ccc}
\textit{6000} & /{\textbeltl}enaaliys{\textyogh}ola/&$\rightarrow$ & \textit{6300} & /{\textbeltl}enaalys{\textyogh}ola/&$\rightarrow$ & \textit{7600} & /{\textbeltl}enaays{\textyogh}ola/&$\rightarrow$ & \textit{9500} & /{\textbeltl}enaays{\textyogh}ol/&$\rightarrow$ & \textit{12001} & /{\textbeltl}enaays{\textesh}ol/&$\rightarrow$ & \textit{12005} & /{\textbeltl}enays{\textesh}ol/&$\rightarrow$ & \textit{13000} & /{\textbeltl}enaysol/& \\
\end{tabular}

\vspace{20pt}\hline

\end{nopagebreak}
\filbreak



\vspace{15pt}
\begin{nopagebreak}
\noindent{\fontsize{20pt}{10pt}\textbf{mralshanin} } \textit{52 of spring/April 22} (N)\\
\noindent {\tovian \fontsize{20pt}{10pt} \textbf{mralshanin} }\\
\noindent /mral{\textesh}{\textprimstress}anin/\\
\noindent lit. darkness+berry\\
\noindent \textit{Nightberry}\\


\noindent History:

\vspace{-0pt}
\hspace{40pt}
\begin{tabular}{ccc}
\textit{6000} & /bral{\textyogh}lanina/&$\rightarrow$ & \textit{9500} & /bral{\textyogh}lanin/&$\rightarrow$ & \textit{11000} & /mral{\textyogh}lanin/&$\rightarrow$ & \textit{12001} & /mral{\textesh}lanin/&$\rightarrow$ & \textit{12500} & /mral{\textesh}anin/& \\
\end{tabular}

\vspace{20pt}\hline

\end{nopagebreak}
\filbreak



\vspace{15pt}
\begin{nopagebreak}
\noindent{\fontsize{20pt}{10pt}\textbf{mralshol} } \textit{52 of summer/July 22} (N)\\
\noindent {\tovian \fontsize{20pt}{10pt} \textbf{mralshol} }\\
\noindent /mr{\textprimstress}al{\textesh}ol/\\
\noindent lit. dark+fruit\\
\noindent \textit{Duskberry}\\


\noindent History:

\vspace{-0pt}
\hspace{40pt}
\begin{tabular}{ccc}
\textit{6000} & /bral{\textyogh}ola/&$\rightarrow$ & \textit{9500} & /bral{\textyogh}ol/&$\rightarrow$ & \textit{11000} & /mral{\textyogh}ol/&$\rightarrow$ & \textit{12001} & /mral{\textesh}ol/& \\
\end{tabular}

\vspace{20pt}\hline

\end{nopagebreak}
\filbreak



\vspace{15pt}
\begin{nopagebreak}
\noindent{\fontsize{20pt}{10pt}\textbf{lanalor} } \textit{52 of winter/January 52} (N)\\
\noindent {\tovian \fontsize{20pt}{10pt} \textbf{lanalor} }\\
\noindent /lan{\textprimstress}alor/\\
\noindent lit. twilight+plant\\
\noindent \textit{Duskheather}\\


\noindent History:

\vspace{-0pt}
\hspace{40pt}
\begin{tabular}{ccc}
\textit{6000} & /brlanamllore/&$\rightarrow$ & \textit{6500} & /brlanamlore/&$\rightarrow$ & \textit{7501} & /rlanamlore/&$\rightarrow$ & \textit{7600} & /lanalore/&$\rightarrow$ & \textit{9500} & /lanalor/& \\
\end{tabular}

\vspace{20pt}\hline

\end{nopagebreak}
\filbreak



\vspace{15pt}
\begin{nopagebreak}
\noindent{\fontsize{20pt}{10pt}\textbf{thamlor} } \textit{53 of autumn/October 23} (N)\\
\noindent {\tovian \fontsize{20pt}{10pt} \textbf{thamlor} }\\
\noindent /{\texttheta}{\textprimstress}amlor/\\
\noindent lit. tool+plant\\
\noindent \textit{Brewvine}\\


\noindent History:

\vspace{-0pt}
\hspace{40pt}
\begin{tabular}{ccc}
\textit{6000} & /{\dh}anmllore/&$\rightarrow$ & \textit{6500} & /{\dh}anmlore/&$\rightarrow$ & \textit{7600} & /{\dh}amlore/&$\rightarrow$ & \textit{9500} & /{\dh}amlor/&$\rightarrow$ & \textit{12002} & /{\texttheta}amlor/& \\
\end{tabular}

\vspace{20pt}\hline

\end{nopagebreak}
\filbreak



\vspace{15pt}
\begin{nopagebreak}
\noindent{\fontsize{20pt}{10pt}\textbf{femlor} } \textit{53 of spring/April 23} (N)\\
\noindent {\tovian \fontsize{20pt}{10pt} \textbf{femlor} }\\
\noindent /f{\textprimstress}emlor/\\
\noindent lit. speaking+plant\\
\noindent \textit{Duskmadder}\\


\noindent History:

\vspace{-0pt}
\hspace{40pt}
\begin{tabular}{ccc}
\textit{6000} & /fe{\texttheta}rmllore/&$\rightarrow$ & \textit{6500} & /fe{\texttheta}rmlore/&$\rightarrow$ & \textit{7600} & /fe{\texttheta}mlore/&$\rightarrow$ & \textit{9500} & /fe{\texttheta}mlor/&$\rightarrow$ & \textit{13002} & /femlor/& \\
\end{tabular}

\vspace{20pt}\hline

\end{nopagebreak}
\filbreak



\vspace{15pt}
\begin{nopagebreak}
\noindent{\fontsize{20pt}{10pt}\textbf{ysayeol} } \textit{53 of summer/July 23} (N)\\
\noindent {\tovian \fontsize{20pt}{10pt} \textbf{ysayeol} }\\
\noindent /ysay{\textprimstress}eol/\\
\noindent lit. moon+wisdom\\
\noindent \textit{Moonlent}\\


\noindent History:

\vspace{-0pt}
\hspace{40pt}
\begin{tabular}{ccc}
\textit{6000} & /aliysaryeol/&$\rightarrow$ & \textit{6300} & /alysaryeol/&$\rightarrow$ & \textit{7500} & /lysaryeol/&$\rightarrow$ & \textit{7501} & /ysaryeol/&$\rightarrow$ & \textit{7600} & /ysayeol/& \\
\end{tabular}

\vspace{20pt}\hline

\end{nopagebreak}
\filbreak



\vspace{15pt}
\begin{nopagebreak}
\noindent{\fontsize{20pt}{10pt}\textbf{shalor} } \textit{53 of winter/January 53} (N)\\
\noindent {\tovian \fontsize{20pt}{10pt} \textbf{shalor} }\\
\noindent /{\textesh}{\textprimstress}alor/\\
\noindent lit. ice+plant\\
\noindent \textit{Frostreed}\\


\noindent History:

\vspace{-0pt}
\hspace{40pt}
\begin{tabular}{ccc}
\textit{6000} & /{\textesh}u{\textesh}amllore/&$\rightarrow$ & \textit{6500} & /{\textesh}u{\textesh}amlore/&$\rightarrow$ & \textit{7600} & /{\textesh}u{\textesh}alore/&$\rightarrow$ & \textit{9200} & /{\textesh}{\textschwa}{\textesh}alore/&$\rightarrow$ & \textit{9500} & /{\textesh}{\textschwa}{\textesh}alor/&$\rightarrow$ & \textit{11990} & /{\textesh}{\textesh}alor/&$\rightarrow$ & \textit{12007} & /{\textesh}alor/& \\
\end{tabular}

\vspace{20pt}\hline

\end{nopagebreak}
\filbreak



\vspace{15pt}
\begin{nopagebreak}
\noindent{\fontsize{20pt}{10pt}\textbf{meloresyalethavar} } \textit{54 of autumn/October 24} (N)\\
\noindent {\tovian \fontsize{20pt}{10pt} \textbf{meloresyalethavar} }\\
\noindent /meloresyale{\texttheta}{\textprimstress}avar/\\
\noindent lit. grain+value\\
\noindent \textit{Goldgrain}\\


\noindent History:

\vspace{-0pt}
\hspace{40pt}
\begin{tabular}{ccc}
\textit{6000} & /mlloresyale{\dh}avlare/&$\rightarrow$ & \textit{6500} & /mloresyale{\dh}avlare/&$\rightarrow$ & \textit{7502} & /meloresyale{\dh}avlare/&$\rightarrow$ & \textit{9500} & /meloresyale{\dh}avlar/&$\rightarrow$ & \textit{12002} & /meloresyale{\texttheta}avlar/&$\rightarrow$ & \textit{12500} & /meloresyale{\texttheta}avar/& \\
\end{tabular}

\vspace{20pt}\hline

\end{nopagebreak}
\filbreak



\vspace{15pt}
\begin{nopagebreak}
\noindent{\fontsize{20pt}{10pt}\textbf{lunum} } \textit{54 of spring/April 24} (N)\\
\noindent {\tovian \fontsize{20pt}{10pt} \textbf{lunum} }\\
\noindent /l{\textprimstress}unum/\\
\noindent lit. sun+TEMP+present\\
\noindent \textit{Sunorb}\\


\noindent History:

\vspace{-0pt}
\hspace{40pt}
\begin{tabular}{ccc}
\textit{6000} & /efl{\textsubbridge{t}}u-num{\textbeltl}le/&$\rightarrow$ & \textit{7500} & /fl{\textsubbridge{t}}u-num{\textbeltl}le/&$\rightarrow$ & \textit{7501} & /l{\textsubbridge{t}}u-num{\textbeltl}le/&$\rightarrow$ & \textit{8500} & /l{\textsubbridge{t}}unum{\textbeltl}le/&$\rightarrow$ & \textit{8750} & /lunum{\textbeltl}le/&$\rightarrow$ & \textit{9500} & /lunum{\textbeltl}l/&$\rightarrow$ & \textit{12500} & /lunum{\textbeltl}/&$\rightarrow$ & \textit{13001} & /lunum/& \\
\end{tabular}

\vspace{20pt}\hline

\end{nopagebreak}
\filbreak



\vspace{15pt}
\begin{nopagebreak}
\noindent{\fontsize{20pt}{10pt}\textbf{kelulor} } \textit{54 of summer/July 24} (N)\\
\noindent {\tovian \fontsize{20pt}{10pt} \textbf{kelulor} }\\
\noindent /kel{\textprimstress}ulor/\\
\noindent lit. branch+plant\\
\noindent \textit{Staranise}\\


\noindent History:

\vspace{-0pt}
\hspace{40pt}
\begin{tabular}{ccc}
\textit{6000} & /kelumllore/&$\rightarrow$ & \textit{6500} & /kelumlore/&$\rightarrow$ & \textit{7600} & /kelulore/&$\rightarrow$ & \textit{9500} & /kelulor/& \\
\end{tabular}

\vspace{20pt}\hline

\end{nopagebreak}
\filbreak



\vspace{15pt}
\begin{nopagebreak}
\noindent{\fontsize{20pt}{10pt}\textbf{kelusyaleth} } \textit{54 of winter/January 54} (N)\\
\noindent {\tovian \fontsize{20pt}{10pt} \textbf{kelusyaleth} }\\
\noindent /kelusy{\textprimstress}ale{\texttheta}/\\
\noindent lit. branch+food\\
\noindent \textit{Startherb}\\


\noindent History:

\vspace{-0pt}
\hspace{40pt}
\begin{tabular}{ccc}
\textit{6000} & /kelusyale{\dh}a/&$\rightarrow$ & \textit{9500} & /kelusyale{\dh}/&$\rightarrow$ & \textit{12002} & /kelusyale{\texttheta}/& \\
\end{tabular}

\vspace{20pt}\hline

\end{nopagebreak}
\filbreak



\vspace{15pt}
\begin{nopagebreak}
\noindent{\fontsize{20pt}{10pt}\textbf{shotlesh} } \textit{55 of autumn/October 25} (N)\\
\noindent {\tovian \fontsize{20pt}{10pt} \textbf{shotlesh} }\\
\noindent /{\textesh}{\textprimstress}o{\texttoptiebar{t\textbeltl}}e{\textesh}/\\
\noindent lit. river+sound\\
\noindent \textit{Rivercaw}\\


\noindent History:

\vspace{-0pt}
\hspace{40pt}
\begin{tabular}{ccc}
\textit{6000} & /{\textesh}o{\texttoptiebar{t\textbeltl}}zl{\textsubbridge{t}}e{\textyogh}/&$\rightarrow$ & \textit{8750} & /{\textesh}o{\texttoptiebar{t\textbeltl}}zle{\textyogh}/&$\rightarrow$ & \textit{12000} & /{\textesh}o{\texttoptiebar{t\textbeltl}}sle{\textyogh}/&$\rightarrow$ & \textit{12001} & /{\textesh}o{\texttoptiebar{t\textbeltl}}sle{\textesh}/&$\rightarrow$ & \textit{12500} & /{\textesh}o{\texttoptiebar{t\textbeltl}}se{\textesh}/&$\rightarrow$ & \textit{14000} & /{\textesh}o{\texttoptiebar{t\textbeltl}}e{\textesh}/& \\
\end{tabular}

\vspace{20pt}\hline

\end{nopagebreak}
\filbreak



\vspace{15pt}
\begin{nopagebreak}
\noindent{\fontsize{20pt}{10pt}\textbf{rithiletlan} } \textit{55 of spring/April 25} (N)\\
\noindent {\tovian \fontsize{20pt}{10pt} \textbf{rithiletlan} }\\
\noindent /ri{\texttheta}il{\textprimstress}e{\texttoptiebar{t\textbeltl}}an/\\
\noindent lit. sky+animal\\
\noindent \textit{Skyfowl}\\


\noindent History:

\vspace{-0pt}
\hspace{40pt}
\begin{tabular}{ccc}
\textit{6000} & /ri{\texttheta}ile{\textsubbridge{t}}lamn/&$\rightarrow$ & \textit{6401} & /ri{\texttheta}ile{\texttoptiebar{t\textbeltl}}amn/&$\rightarrow$ & \textit{7600} & /ri{\texttheta}ile{\texttoptiebar{t\textbeltl}}an/& \\
\end{tabular}

\vspace{20pt}\hline

\end{nopagebreak}
\filbreak



\vspace{15pt}
\begin{nopagebreak}
\noindent{\fontsize{20pt}{10pt}\textbf{sethan} } \textit{55 of summer/July 25} (N)\\
\noindent {\tovian \fontsize{20pt}{10pt} \textbf{sethan} }\\
\noindent /s{\textprimstress}e{\texttheta}an/\\
\noindent lit. water+tool\\
\noindent \textit{Waterwhisker}\\


\noindent History:

\vspace{-0pt}
\hspace{40pt}
\begin{tabular}{ccc}
\textit{6000} & /zel{\dh}an/&$\rightarrow$ & \textit{12000} & /sel{\dh}an/&$\rightarrow$ & \textit{12002} & /sel{\texttheta}an/&$\rightarrow$ & \textit{12502} & /se{\texttheta}an/& \\
\end{tabular}

\vspace{20pt}\hline

\end{nopagebreak}
\filbreak



\vspace{15pt}
\begin{nopagebreak}
\noindent{\fontsize{20pt}{10pt}\textbf{mralesh} } \textit{55 of winter/January 55} (N)\\
\noindent {\tovian \fontsize{20pt}{10pt} \textbf{mralesh} }\\
\noindent /mr{\textprimstress}ale{\textesh}/\\
\noindent lit. darkness+sound\\
\noindent \textit{Nightchirp}\\


\noindent History:

\vspace{-0pt}
\hspace{40pt}
\begin{tabular}{ccc}
\textit{6000} & /bral{\textsubbridge{t}}e{\textyogh}/&$\rightarrow$ & \textit{8750} & /brale{\textyogh}/&$\rightarrow$ & \textit{11000} & /mrale{\textyogh}/&$\rightarrow$ & \textit{12001} & /mrale{\textesh}/& \\
\end{tabular}

\vspace{20pt}\hline

\end{nopagebreak}
\filbreak



\vspace{15pt}
\begin{nopagebreak}
\noindent{\fontsize{20pt}{10pt}\textbf{haraon} } \textit{56 of autumn/October 26} (N)\\
\noindent {\tovian \fontsize{20pt}{10pt} \textbf{haraon} }\\
\noindent /ħar{\textprimstress}aon/\\
\noindent lit. pain+mouth\\
\noindent \textit{Bitterpearl}\\


\noindent History:

\vspace{-0pt}
\hspace{40pt}
\begin{tabular}{ccc}
\textit{6000} & /haraeoln/&$\rightarrow$ & \textit{7600} & /haraeon/&$\rightarrow$ & \textit{10000} & /haraon/&$\rightarrow$ & \textit{14006} & /ħaraon/& \\
\end{tabular}

\vspace{20pt}\hline

\end{nopagebreak}
\filbreak



\vspace{15pt}
\begin{nopagebreak}
\noindent{\fontsize{20pt}{10pt}\textbf{mralsyalethashol} } \textit{56 of spring/April 26} (N)\\
\noindent {\tovian \fontsize{20pt}{10pt} \textbf{mralsyalethashol} }\\
\noindent /mralsyale{\texttheta}{\textprimstress}a{\textesh}ol/\\
\noindent lit. darkness+food+fruit\\
\noindent \textit{Nightnut}\\


\noindent History:

\vspace{-0pt}
\hspace{40pt}
\begin{tabular}{ccc}
\textit{6000} & /bralsyale{\dh}a{\textyogh}ola/&$\rightarrow$ & \textit{9500} & /bralsyale{\dh}a{\textyogh}ol/&$\rightarrow$ & \textit{11000} & /mralsyale{\dh}a{\textyogh}ol/&$\rightarrow$ & \textit{12001} & /mralsyale{\dh}a{\textesh}ol/&$\rightarrow$ & \textit{12002} & /mralsyale{\texttheta}a{\textesh}ol/& \\
\end{tabular}

\vspace{20pt}\hline

\end{nopagebreak}
\filbreak



\vspace{15pt}
\begin{nopagebreak}
\noindent{\fontsize{20pt}{10pt}\textbf{mramlor} } \textit{56 of summer/July 26} (N)\\
\noindent {\tovian \fontsize{20pt}{10pt} \textbf{mramlor} }\\
\noindent /mr{\textprimstress}amlor/\\
\noindent lit. darkness+plant\\
\noindent \textit{Nightmyrtle}\\


\noindent History:

\vspace{-0pt}
\hspace{40pt}
\begin{tabular}{ccc}
\textit{6000} & /bralmllore/&$\rightarrow$ & \textit{6500} & /bralmlore/&$\rightarrow$ & \textit{7600} & /bramlore/&$\rightarrow$ & \textit{9500} & /bramlor/&$\rightarrow$ & \textit{11000} & /mramlor/& \\
\end{tabular}

\vspace{20pt}\hline

\end{nopagebreak}
\filbreak



\vspace{15pt}
\begin{nopagebreak}
\noindent{\fontsize{20pt}{10pt}\textbf{nanad} } \textit{56 of winter/January 56} (N)\\
\noindent {\tovian \fontsize{20pt}{10pt} \textbf{nanad} }\\
\noindent /n{\textprimstress}anad/\\
\noindent lit. star+showing\\
\noindent \textit{Stardush}\\


\noindent History:

\vspace{-0pt}
\hspace{40pt}
\begin{tabular}{ccc}
\textit{6000} & /nnaanafede/&$\rightarrow$ & \textit{6500} & /naanafede/&$\rightarrow$ & \textit{9500} & /naanafed/&$\rightarrow$ & \textit{12003} & /naanad/&$\rightarrow$ & \textit{12005} & /nanad/& \\
\end{tabular}

\vspace{20pt}\hline

\end{nopagebreak}
\filbreak



\vspace{15pt}
\begin{nopagebreak}
\noindent{\fontsize{20pt}{10pt}\textbf{lethilekel} } \textit{57 of autumn/October 27} (N)\\
\noindent {\tovian \fontsize{20pt}{10pt} \textbf{lethilekel} }\\
\noindent /le{\texttheta}il{\textprimstress}ekel/\\
\noindent lit. forest+branch\\
\noindent \textit{Starblum}\\


\noindent History:

\vspace{-0pt}
\hspace{40pt}
\begin{tabular}{ccc}
\textit{6000} & /aglre{\texttheta}ilekelu/&$\rightarrow$ & \textit{6503} & /agrle{\texttheta}ilekelu/&$\rightarrow$ & \textit{7500} & /grle{\texttheta}ilekelu/&$\rightarrow$ & \textit{7501} & /rle{\texttheta}ilekelu/&$\rightarrow$ & \textit{7600} & /le{\texttheta}ilekelu/&$\rightarrow$ & \textit{9500} & /le{\texttheta}ilekel/& \\
\end{tabular}

\vspace{20pt}\hline

\end{nopagebreak}
\filbreak



\vspace{15pt}
\begin{nopagebreak}
\noindent{\fontsize{20pt}{10pt}\textbf{yelathan} } \textit{57 of spring/April 27} (N)\\
\noindent {\tovian \fontsize{20pt}{10pt} \textbf{yelathan} }\\
\noindent /yel{\textprimstress}a{\texttheta}an/\\
\noindent lit. person+tool\\
\noindent \textit{Drielthorn}\\


\noindent History:

\vspace{-0pt}
\hspace{40pt}
\begin{tabular}{ccc}
\textit{6000} & /dryela{\dh}an/&$\rightarrow$ & \textit{7501} & /ryela{\dh}an/&$\rightarrow$ & \textit{7600} & /yela{\dh}an/&$\rightarrow$ & \textit{12002} & /yela{\texttheta}an/& \\
\end{tabular}

\vspace{20pt}\hline

\end{nopagebreak}
\filbreak



\vspace{15pt}
\begin{nopagebreak}
\noindent{\fontsize{20pt}{10pt}\textbf{varedrem} } \textit{57 of summer/July 27} (N)\\
\noindent {\tovian \fontsize{20pt}{10pt} \textbf{varedrem} }\\
\noindent /var{\textprimstress}edrem/\\
\noindent lit. value+flower\\
\noindent \textit{Goldbloom}\\


\noindent History:

\vspace{-0pt}
\hspace{40pt}
\begin{tabular}{ccc}
\textit{6000} & /vlaredremp/&$\rightarrow$ & \textit{8750} & /vlaredrem/&$\rightarrow$ & \textit{12500} & /varedrem/& \\
\end{tabular}

\vspace{20pt}\hline

\end{nopagebreak}
\filbreak



\vspace{15pt}
\begin{nopagebreak}
\noindent{\fontsize{20pt}{10pt}\textbf{ysaglar} } \textit{57 of winter/January 57} (N)\\
\noindent {\tovian \fontsize{20pt}{10pt} \textbf{ysaglar} }\\
\noindent /ys{\textprimstress}aglar/\\
\noindent lit. moon+tree\\
\noindent \textit{Moonbark}\\


\noindent History:

\vspace{-0pt}
\hspace{40pt}
\begin{tabular}{ccc}
\textit{6000} & /aliysaglare/&$\rightarrow$ & \textit{6300} & /alysaglare/&$\rightarrow$ & \textit{7500} & /lysaglare/&$\rightarrow$ & \textit{7501} & /ysaglare/&$\rightarrow$ & \textit{9500} & /ysaglar/& \\
\end{tabular}

\vspace{20pt}\hline

\end{nopagebreak}
\filbreak



\vspace{15pt}
\begin{nopagebreak}
\noindent{\fontsize{20pt}{10pt}\textbf{mlor} } \textit{58 of autumn/October 28} (N)\\
\noindent {\tovian \fontsize{20pt}{10pt} \textbf{mlor} }\\
\noindent /ml{\textprimstress}or/\\
\noindent lit. sun+plant\\
\noindent \textit{Sunstalk}\\


\noindent History:

\vspace{-0pt}
\hspace{40pt}
\begin{tabular}{ccc}
\textit{6000} & /eflmllore/&$\rightarrow$ & \textit{6500} & /eflmlore/&$\rightarrow$ & \textit{7500} & /flmlore/&$\rightarrow$ & \textit{7501} & /lmlore/&$\rightarrow$ & \textit{7600} & /mlore/&$\rightarrow$ & \textit{9500} & /mlor/& \\
\end{tabular}

\vspace{20pt}\hline

\end{nopagebreak}
\filbreak



\vspace{15pt}
\begin{nopagebreak}
\noindent{\fontsize{20pt}{10pt}\textbf{ysol} } \textit{58 of spring/April 28} (N)\\
\noindent {\tovian \fontsize{20pt}{10pt} \textbf{ysol} }\\
\noindent /ys{\textprimstress}ol/\\
\noindent lit. moon+fruit\\
\noindent \textit{Moonquince}\\


\noindent History:

\vspace{-0pt}
\hspace{40pt}
\begin{tabular}{ccc}
\textit{6000} & /aliys{\textyogh}ola/&$\rightarrow$ & \textit{6300} & /alys{\textyogh}ola/&$\rightarrow$ & \textit{7500} & /lys{\textyogh}ola/&$\rightarrow$ & \textit{7501} & /ys{\textyogh}ola/&$\rightarrow$ & \textit{9500} & /ys{\textyogh}ol/&$\rightarrow$ & \textit{12001} & /ys{\textesh}ol/&$\rightarrow$ & \textit{13000} & /ysol/& \\
\end{tabular}

\vspace{20pt}\hline

\end{nopagebreak}
\filbreak



\vspace{15pt}
\begin{nopagebreak}
\noindent{\fontsize{20pt}{10pt}\textbf{ysdmfed} } \textit{58 of summer/July 28} (N)\\
\noindent {\tovian \fontsize{20pt}{10pt} \textbf{ysdmfed} }\\
\noindent /ysdmf{\textprimstress}ed/\\
\noindent lit. moon+bloom\\
\noindent \textit{Moonbloom}\\


\noindent History:

\vspace{-0pt}
\hspace{40pt}
\begin{tabular}{ccc}
\textit{6000} & /aliysdrmpfede/&$\rightarrow$ & \textit{6300} & /alysdrmpfede/&$\rightarrow$ & \textit{7500} & /lysdrmpfede/&$\rightarrow$ & \textit{7501} & /ysdrmpfede/&$\rightarrow$ & \textit{7600} & /ysdmpfede/&$\rightarrow$ & \textit{8500} & /ysdmfede/&$\rightarrow$ & \textit{9500} & /ysdmfed/& \\
\end{tabular}

\vspace{20pt}\hline

\end{nopagebreak}
\filbreak



\vspace{15pt}
\begin{nopagebreak}
\noindent{\fontsize{20pt}{10pt}\textbf{mer} } \textit{58 of winter/January 58} (N)\\
\noindent {\tovian \fontsize{20pt}{10pt} \textbf{mer} }\\
\noindent /m{\textprimstress}er/\\
\noindent lit. earth+showing\\
\noindent \textit{Earthem}\\


\noindent History:

\vspace{-0pt}
\hspace{40pt}
\begin{tabular}{ccc}
\textit{6000} & /mernfede/&$\rightarrow$ & \textit{6400} & /merfede/&$\rightarrow$ & \textit{9500} & /merfed/&$\rightarrow$ & \textit{12003} & /merd/&$\rightarrow$ & \textit{13001} & /mer/& \\
\end{tabular}

\vspace{20pt}\hline

\end{nopagebreak}
\filbreak



\vspace{15pt}
\begin{nopagebreak}
\noindent{\fontsize{20pt}{10pt}\textbf{glareshol} } \textit{59 of autumn/October 29} (N)\\
\noindent {\tovian \fontsize{20pt}{10pt} \textbf{glareshol} }\\
\noindent /glar{\textprimstress}e{\textesh}ol/\\
\noindent lit. tree+fruit\\
\noindent \textit{Hovenut}\\


\noindent History:

\vspace{-0pt}
\hspace{40pt}
\begin{tabular}{ccc}
\textit{6000} & /aglare{\textyogh}ola/&$\rightarrow$ & \textit{7500} & /glare{\textyogh}ola/&$\rightarrow$ & \textit{9500} & /glare{\textyogh}ol/&$\rightarrow$ & \textit{12001} & /glare{\textesh}ol/& \\
\end{tabular}

\vspace{20pt}\hline

\end{nopagebreak}
\filbreak



\vspace{15pt}
\begin{nopagebreak}
\noindent{\fontsize{20pt}{10pt}\textbf{keluaglar} } \textit{59 of spring/April 29} (N)\\
\noindent {\tovian \fontsize{20pt}{10pt} \textbf{keluaglar} }\\
\noindent /kelu{\textprimstress}aglar/\\
\noindent lit. branch+tree\\
\noindent \textit{Starwood}\\


\noindent History:

\vspace{-0pt}
\hspace{40pt}
\begin{tabular}{ccc}
\textit{6000} & /keluaglare/&$\rightarrow$ & \textit{9500} & /keluaglar/& \\
\end{tabular}

\vspace{20pt}\hline

\end{nopagebreak}
\filbreak



\vspace{15pt}
\begin{nopagebreak}
\noindent{\fontsize{20pt}{10pt}\textbf{rithiledyel} } \textit{59 of summer/July 29} (N)\\
\noindent {\tovian \fontsize{20pt}{10pt} \textbf{rithiledyel} }\\
\noindent /ri{\texttheta}il{\textprimstress}edyel/\\
\noindent lit. sky+person\\
\noindent \textit{Cloudmun}\\


\noindent History:

\vspace{-0pt}
\hspace{40pt}
\begin{tabular}{ccc}
\textit{6000} & /ri{\texttheta}iledryela/&$\rightarrow$ & \textit{7600} & /ri{\texttheta}iledyela/&$\rightarrow$ & \textit{9500} & /ri{\texttheta}iledyel/& \\
\end{tabular}

\vspace{20pt}\hline

\end{nopagebreak}
\filbreak



\vspace{15pt}
\begin{nopagebreak}
\noindent{\fontsize{20pt}{10pt}\textbf{varefethr} } \textit{59 of winter/January 59} (N)\\
\noindent {\tovian \fontsize{20pt}{10pt} \textbf{varefethr} }\\
\noindent /var{\textprimstress}efe{\texttheta}r/\\
\noindent lit. value+speaking\\
\noindent \textit{Velarum}\\


\noindent History:

\vspace{-0pt}
\hspace{40pt}
\begin{tabular}{ccc}
\textit{6000} & /vlarefe{\texttheta}r/&$\rightarrow$ & \textit{12500} & /varefe{\texttheta}r/& \\
\end{tabular}

\vspace{20pt}\hline

\end{nopagebreak}
\filbreak



\vspace{15pt}
\begin{nopagebreak}
\noindent{\fontsize{20pt}{10pt}\textbf{mar} } \textit{6 of autumn/September 6} (N)\\
\noindent {\tovian \fontsize{20pt}{10pt} \textbf{mar} }\\
\noindent /m{\textprimstress}ar/\\
\noindent lit. sun+mother\\
\noindent \textit{Sunmary}\\


\noindent History:

\vspace{-0pt}
\hspace{40pt}
\begin{tabular}{ccc}
\textit{6000} & /eflmari/&$\rightarrow$ & \textit{7500} & /flmari/&$\rightarrow$ & \textit{7501} & /lmari/&$\rightarrow$ & \textit{7600} & /mari/&$\rightarrow$ & \textit{9500} & /mar/& \\
\end{tabular}

\vspace{20pt}\hline

\end{nopagebreak}
\filbreak



\vspace{15pt}
\begin{nopagebreak}
\noindent{\fontsize{20pt}{10pt}\textbf{shimilor} } \textit{6 of spring/March 6} (N)\\
\noindent {\tovian \fontsize{20pt}{10pt} \textbf{shimilor} }\\
\noindent /{\textesh}im{\textprimstress}ilor/\\
\noindent lit. mystery+plant\\
\noindent \textit{Mystsap}\\


\noindent History:

\vspace{-0pt}
\hspace{40pt}
\begin{tabular}{ccc}
\textit{6000} & /{\textesh}imimllore/&$\rightarrow$ & \textit{6500} & /{\textesh}imimlore/&$\rightarrow$ & \textit{7600} & /{\textesh}imilore/&$\rightarrow$ & \textit{9500} & /{\textesh}imilor/& \\
\end{tabular}

\vspace{20pt}\hline

\end{nopagebreak}
\filbreak



\vspace{15pt}
\begin{nopagebreak}
\noindent{\fontsize{20pt}{10pt}\textbf{meloreathang} } \textit{6 of summer/June 6} (N)\\
\noindent {\tovian \fontsize{20pt}{10pt} \textbf{meloreathang} }\\
\noindent /melore{\textprimstress}a{\texttheta}a{\ng}/\\
\noindent lit. plant+leaf\\
\noindent \textit{Melorleaf}\\


\noindent History:

\vspace{-0pt}
\hspace{40pt}
\begin{tabular}{ccc}
\textit{6000} & /mllorea{\dh}a{\ng}ge/&$\rightarrow$ & \textit{6500} & /mlorea{\dh}a{\ng}ge/&$\rightarrow$ & \textit{7502} & /melorea{\dh}a{\ng}ge/&$\rightarrow$ & \textit{8750} & /melorea{\dh}a{\ng}e/&$\rightarrow$ & \textit{9500} & /melorea{\dh}a{\ng}/&$\rightarrow$ & \textit{12002} & /melorea{\texttheta}a{\ng}/& \\
\end{tabular}

\vspace{20pt}\hline

\end{nopagebreak}
\filbreak



\vspace{15pt}
\begin{nopagebreak}
\noindent{\fontsize{20pt}{10pt}\textbf{shotltom} } \textit{6 of winter/december 6} (N)\\
\noindent {\tovian \fontsize{20pt}{10pt} \textbf{shotltom} }\\
\noindent /{\textesh}{\textprimstress}o{\texttoptiebar{t\textbeltl}}{\textsubbridge{t}}om/\\
\noindent lit. flowing+root\\
\noindent \textit{Mistroot}\\


\noindent History:

\vspace{-0pt}
\hspace{40pt}
\begin{tabular}{ccc}
\textit{6000} & /{\textesh}o{\texttoptiebar{t\textbeltl}}{\textsubbridge{t}}omo/&$\rightarrow$ & \textit{9500} & /{\textesh}o{\texttoptiebar{t\textbeltl}}{\textsubbridge{t}}om/& \\
\end{tabular}

\vspace{20pt}\hline

\end{nopagebreak}
\filbreak



\vspace{15pt}
\begin{nopagebreak}
\noindent{\fontsize{20pt}{10pt}\textbf{leshor} } \textit{60 of autumn/October 30} (N)\\
\noindent {\tovian \fontsize{20pt}{10pt} \textbf{leshor} }\\
\noindent /l{\textprimstress}e{\textesh}or/\\
\noindent lit. harvest+action\\
\noindent \textit{Harvesh}\\


\noindent History:

\vspace{-0pt}
\hspace{40pt}
\begin{tabular}{ccc}
\textit{6000} & /mllre{\textyogh}lore/&$\rightarrow$ & \textit{6500} & /mlre{\textyogh}lore/&$\rightarrow$ & \textit{6503} & /mrle{\textyogh}lore/&$\rightarrow$ & \textit{7501} & /rle{\textyogh}lore/&$\rightarrow$ & \textit{7600} & /le{\textyogh}lore/&$\rightarrow$ & \textit{9500} & /le{\textyogh}lor/&$\rightarrow$ & \textit{12001} & /le{\textesh}lor/&$\rightarrow$ & \textit{12500} & /le{\textesh}or/& \\
\end{tabular}

\vspace{20pt}\hline

\end{nopagebreak}
\filbreak



\vspace{15pt}
\begin{nopagebreak}
\noindent{\fontsize{20pt}{10pt}\textbf{rithilekwil} } \textit{60 of spring/April 30} (N)\\
\noindent {\tovian \fontsize{20pt}{10pt} \textbf{rithilekwil} }\\
\noindent /ri{\texttheta}il{\textprimstress}ekwil/\\
\noindent lit. sky+movement\\
\noindent \textit{Cloudroll}\\


\noindent History:

\vspace{-0pt}
\hspace{40pt}
\begin{tabular}{ccc}
\textit{6000} & /ri{\texttheta}ilekwil/& \\
\end{tabular}

\vspace{20pt}\hline

\end{nopagebreak}
\filbreak



\vspace{15pt}
\begin{nopagebreak}
\noindent{\fontsize{20pt}{10pt}\textbf{lhenumlhen} } \textit{60 of summer/July 30} (N)\\
\noindent {\tovian \fontsize{20pt}{10pt} \textbf{lhenumlhen} }\\
\noindent /{\textbeltl}en{\textprimstress}um{\textbeltl}en/\\
\noindent lit. center+path+dream\\
\noindent \textit{Lhelnumly}\\


\noindent History:

\vspace{-0pt}
\hspace{40pt}
\begin{tabular}{ccc}
\textit{6000} & /{\textbeltl}elnum{\textbeltl}ena/&$\rightarrow$ & \textit{7600} & /{\textbeltl}enum{\textbeltl}ena/&$\rightarrow$ & \textit{9500} & /{\textbeltl}enum{\textbeltl}en/& \\
\end{tabular}

\vspace{20pt}\hline

\end{nopagebreak}
\filbreak



\vspace{15pt}
\begin{nopagebreak}
\noindent{\fontsize{20pt}{10pt}\textbf{nanadyethil} } \textit{60 of winter/January 60} (N)\\
\noindent {\tovian \fontsize{20pt}{10pt} \textbf{nanadyethil} }\\
\noindent /nanady{\textprimstress}e{\texttheta}il/\\
\noindent lit. star+home\\
\noindent \textit{Starhovel}\\


\noindent History:

\vspace{-0pt}
\hspace{40pt}
\begin{tabular}{ccc}
\textit{6000} & /nnaanadryel{\texttheta}ile/&$\rightarrow$ & \textit{6500} & /naanadryel{\texttheta}ile/&$\rightarrow$ & \textit{7600} & /naanadyel{\texttheta}ile/&$\rightarrow$ & \textit{9500} & /naanadyel{\texttheta}il/&$\rightarrow$ & \textit{12005} & /nanadyel{\texttheta}il/&$\rightarrow$ & \textit{12502} & /nanadye{\texttheta}il/& \\
\end{tabular}

\vspace{20pt}\hline

\end{nopagebreak}
\filbreak



\vspace{15pt}
\begin{nopagebreak}
\noindent{\fontsize{20pt}{10pt}\textbf{nanasyaleth} } \textit{61 of autumn/November 1} (N)\\
\noindent {\tovian \fontsize{20pt}{10pt} \textbf{nanasyaleth} }\\
\noindent /nanasy{\textprimstress}ale{\texttheta}/\\
\noindent lit. star+food\\
\noindent \textit{Starfodder}\\


\noindent History:

\vspace{-0pt}
\hspace{40pt}
\begin{tabular}{ccc}
\textit{6000} & /nnaanasyale{\dh}a/&$\rightarrow$ & \textit{6500} & /naanasyale{\dh}a/&$\rightarrow$ & \textit{9500} & /naanasyale{\dh}/&$\rightarrow$ & \textit{12002} & /naanasyale{\texttheta}/&$\rightarrow$ & \textit{12005} & /nanasyale{\texttheta}/& \\
\end{tabular}

\vspace{20pt}\hline

\end{nopagebreak}
\filbreak



\vspace{15pt}
\begin{nopagebreak}
\noindent{\fontsize{20pt}{10pt}\textbf{selaglar} } \textit{61 of spring/May 1} (N)\\
\noindent {\tovian \fontsize{20pt}{10pt} \textbf{selaglar} }\\
\noindent /sel{\textprimstress}aglar/\\
\noindent lit. water+tree\\
\noindent \textit{Rainlaurel}\\


\noindent History:

\vspace{-0pt}
\hspace{40pt}
\begin{tabular}{ccc}
\textit{6000} & /zelaglare/&$\rightarrow$ & \textit{9500} & /zelaglar/&$\rightarrow$ & \textit{12000} & /selaglar/& \\
\end{tabular}

\vspace{20pt}\hline

\end{nopagebreak}
\filbreak



\vspace{15pt}
\begin{nopagebreak}
\noindent{\fontsize{20pt}{10pt}\textbf{mloanadrem} } \textit{61 of summer/August 1} (N)\\
\noindent {\tovian \fontsize{20pt}{10pt} \textbf{mloanadrem} }\\
\noindent /mloan{\textprimstress}adrem/\\
\noindent lit. morning+flower\\
\noindent \textit{Dawnrose}\\


\noindent History:

\vspace{-0pt}
\hspace{40pt}
\begin{tabular}{ccc}
\textit{6000} & /ploanadremp/&$\rightarrow$ & \textit{8750} & /ploanadrem/&$\rightarrow$ & \textit{11000} & /mloanadrem/& \\
\end{tabular}

\vspace{20pt}\hline

\end{nopagebreak}
\filbreak



\vspace{15pt}
\begin{nopagebreak}
\noindent{\fontsize{20pt}{10pt}\textbf{tushulor} } \textit{61 of winter/February 1} (N)\\
\noindent {\tovian \fontsize{20pt}{10pt} \textbf{tushulor} }\\
\noindent /{\textsubbridge{t}}u{\textesh}{\textprimstress}ulor/\\
\noindent lit. cold+plant\\
\noindent \textit{Froshpea}\\


\noindent History:

\vspace{-0pt}
\hspace{40pt}
\begin{tabular}{ccc}
\textit{6000} & /{\textsubbridge{t}}u{\textesh}umllore/&$\rightarrow$ & \textit{6500} & /{\textsubbridge{t}}u{\textesh}umlore/&$\rightarrow$ & \textit{7600} & /{\textsubbridge{t}}u{\textesh}ulore/&$\rightarrow$ & \textit{9500} & /{\textsubbridge{t}}u{\textesh}ulor/& \\
\end{tabular}

\vspace{20pt}\hline

\end{nopagebreak}
\filbreak



\vspace{15pt}
\begin{nopagebreak}
\noindent{\fontsize{20pt}{10pt}\textbf{lhoman} } \textit{62 of autumn/November 2} (N)\\
\noindent {\tovian \fontsize{20pt}{10pt} \textbf{lhoman} }\\
\noindent /{\textbeltl}{\textprimstress}oman/\\
\noindent lit. day+light\\
\noindent \textit{Dayglem}\\


\noindent History:

\vspace{-0pt}
\hspace{40pt}
\begin{tabular}{ccc}
\textit{6000} & /{\textbeltl}onmana/&$\rightarrow$ & \textit{7600} & /{\textbeltl}omana/&$\rightarrow$ & \textit{9500} & /{\textbeltl}oman/& \\
\end{tabular}

\vspace{20pt}\hline

\end{nopagebreak}
\filbreak



\vspace{15pt}
\begin{nopagebreak}
\noindent{\fontsize{20pt}{10pt}\textbf{meloresel} } \textit{62 of spring/May 2} (N)\\
\noindent {\tovian \fontsize{20pt}{10pt} \textbf{meloresel} }\\
\noindent /melor{\textprimstress}esel/\\
\noindent lit. plant+water\\
\noindent \textit{Dewmoss}\\


\noindent History:

\vspace{-0pt}
\hspace{40pt}
\begin{tabular}{ccc}
\textit{6000} & /mllorezel/&$\rightarrow$ & \textit{6500} & /mlorezel/&$\rightarrow$ & \textit{7502} & /melorezel/&$\rightarrow$ & \textit{12000} & /meloresel/& \\
\end{tabular}

\vspace{20pt}\hline

\end{nopagebreak}
\filbreak



\vspace{15pt}
\begin{nopagebreak}
\noindent{\fontsize{20pt}{10pt}\textbf{teshaglar} } \textit{62 of summer/August 2} (N)\\
\noindent {\tovian \fontsize{20pt}{10pt} \textbf{teshaglar} }\\
\noindent /{\textsubbridge{t}}e{\textesh}{\textprimstress}aglar/\\
\noindent lit. sound+tree\\
\noindent \textit{Thunderoak}\\


\noindent History:

\vspace{-0pt}
\hspace{40pt}
\begin{tabular}{ccc}
\textit{6000} & /{\textsubbridge{t}}e{\textyogh}aglare/&$\rightarrow$ & \textit{9500} & /{\textsubbridge{t}}e{\textyogh}aglar/&$\rightarrow$ & \textit{12001} & /{\textsubbridge{t}}e{\textesh}aglar/& \\
\end{tabular}

\vspace{20pt}\hline

\end{nopagebreak}
\filbreak



\vspace{15pt}
\begin{nopagebreak}
\noindent{\fontsize{20pt}{10pt}\textbf{sulramen} } \textit{62 of winter/February 2} (N)\\
\noindent {\tovian \fontsize{20pt}{10pt} \textbf{sulramen} }\\
\noindent /sulr{\textprimstress}amen/\\
\noindent lit. bottom+darkness+stone\\
\noindent \textit{Nightcoal}\\


\noindent History:

\vspace{-0pt}
\hspace{40pt}
\begin{tabular}{ccc}
\textit{6000} & /zulbralmern/&$\rightarrow$ & \textit{7600} & /zulbramen/&$\rightarrow$ & \textit{8750} & /zulramen/&$\rightarrow$ & \textit{12000} & /sulramen/& \\
\end{tabular}

\vspace{20pt}\hline

\end{nopagebreak}
\filbreak



\vspace{15pt}
\begin{nopagebreak}
\noindent{\fontsize{20pt}{10pt}\textbf{haralor} } \textit{63 of autumn/November 3} (N)\\
\noindent {\tovian \fontsize{20pt}{10pt} \textbf{haralor} }\\
\noindent /ħar{\textprimstress}alor/\\
\noindent lit. pain+plant\\
\noindent \textit{Harclover}\\


\noindent History:

\vspace{-0pt}
\hspace{40pt}
\begin{tabular}{ccc}
\textit{6000} & /haramllore/&$\rightarrow$ & \textit{6500} & /haramlore/&$\rightarrow$ & \textit{7600} & /haralore/&$\rightarrow$ & \textit{9500} & /haralor/&$\rightarrow$ & \textit{14006} & /ħaralor/& \\
\end{tabular}

\vspace{20pt}\hline

\end{nopagebreak}
\filbreak



\vspace{15pt}
\begin{nopagebreak}
\noindent{\fontsize{20pt}{10pt}\textbf{shotlen} } \textit{63 of spring/May 3} (N)\\
\noindent {\tovian \fontsize{20pt}{10pt} \textbf{shotlen} }\\
\noindent /{\textesh}{\textprimstress}o{\texttoptiebar{t\textbeltl}}en/\\
\noindent lit. flowing+dream\\
\noindent \textit{Mistthorn}\\


\noindent History:

\vspace{-0pt}
\hspace{40pt}
\begin{tabular}{ccc}
\textit{6000} & /{\textesh}o{\texttoptiebar{t\textbeltl}}{\textbeltl}ena/&$\rightarrow$ & \textit{9500} & /{\textesh}o{\texttoptiebar{t\textbeltl}}{\textbeltl}en/&$\rightarrow$ & \textit{14000} & /{\textesh}o{\texttoptiebar{t\textbeltl}}en/& \\
\end{tabular}

\vspace{20pt}\hline

\end{nopagebreak}
\filbreak



\vspace{15pt}
\begin{nopagebreak}
\noindent{\fontsize{20pt}{10pt}\textbf{munanilor} } \textit{63 of summer/August 3} (N)\\
\noindent {\tovian \fontsize{20pt}{10pt} \textbf{munanilor} }\\
\noindent /munan{\textprimstress}ilor/\\
\noindent lit. covering+plant\\
\noindent \textit{Veilfern}\\


\noindent History:

\vspace{-0pt}
\hspace{40pt}
\begin{tabular}{ccc}
\textit{6000} & /bnanimllore/&$\rightarrow$ & \textit{6500} & /bnanimlore/&$\rightarrow$ & \textit{7600} & /bnanilore/&$\rightarrow$ & \textit{9500} & /bnanilor/&$\rightarrow$ & \textit{11000} & /mnanilor/&$\rightarrow$ & \textit{11996} & /munanilor/& \\
\end{tabular}

\vspace{20pt}\hline

\end{nopagebreak}
\filbreak



\vspace{15pt}
\begin{nopagebreak}
\noindent{\fontsize{20pt}{10pt}\textbf{mralingtlun} } \textit{63 of winter/February 3} (N)\\
\noindent {\tovian \fontsize{20pt}{10pt} \textbf{mralingtlun} }\\
\noindent /mral{\textprimstress}i{\ng}{\texttoptiebar{t\textbeltl}}un/\\
\noindent lit. darkness+tar\\
\noindent \textit{Shadowtar}\\


\noindent History:

\vspace{-0pt}
\hspace{40pt}
\begin{tabular}{ccc}
\textit{6000} & /brali{\ng}g{\texttoptiebar{t\textbeltl}}una/&$\rightarrow$ & \textit{8500} & /brali{\ng}{\texttoptiebar{t\textbeltl}}una/&$\rightarrow$ & \textit{9500} & /brali{\ng}{\texttoptiebar{t\textbeltl}}un/&$\rightarrow$ & \textit{11000} & /mrali{\ng}{\texttoptiebar{t\textbeltl}}un/& \\
\end{tabular}

\vspace{20pt}\hline

\end{nopagebreak}
\filbreak



\vspace{15pt}
\begin{nopagebreak}
\noindent{\fontsize{20pt}{10pt}\textbf{ay} } \textit{64 of autumn/November 4} (N)\\
\noindent {\tovian \fontsize{20pt}{10pt} \textbf{ay} }\\
\noindent /{\textprimstress}ay/\\
\noindent lit. moon+???\\
\noindent \textit{Moonerb}\\


\noindent History:

\vspace{-0pt}
\hspace{40pt}
\begin{tabular}{ccc}
\textit{6000} & /aliys???/&$\rightarrow$ & \textit{6300} & /alys???/&$\rightarrow$ & \textit{7600} & /ays???/&$\rightarrow$ & \textit{8500} & /ays/&$\rightarrow$ & \textit{13001} & /ay/& \\
\end{tabular}

\vspace{20pt}\hline

\end{nopagebreak}
\filbreak



\vspace{15pt}
\begin{nopagebreak}
\noindent{\fontsize{20pt}{10pt}\textbf{shashithan} } \textit{64 of spring/May 4} (N)\\
\noindent {\tovian \fontsize{20pt}{10pt} \textbf{shashithan} }\\
\noindent /{\textesh}a{\textesh}{\textprimstress}i{\texttheta}an/\\
\noindent lit. ice+sharpness+tool\\
\noindent \textit{Snowpierce}\\


\noindent History:

\vspace{-0pt}
\hspace{40pt}
\begin{tabular}{ccc}
\textit{6000} & /{\textesh}u{\textesh}a{\textyogh}i{\texttheta}{\dh}an/&$\rightarrow$ & \textit{9200} & /{\textesh}{\textschwa}{\textesh}a{\textyogh}i{\texttheta}{\dh}an/&$\rightarrow$ & \textit{11990} & /{\textesh}{\textesh}a{\textyogh}i{\texttheta}{\dh}an/&$\rightarrow$ & \textit{12001} & /{\textesh}{\textesh}a{\textesh}i{\texttheta}{\dh}an/&$\rightarrow$ & \textit{12002} & /{\textesh}{\textesh}a{\textesh}i{\texttheta}{\texttheta}an/&$\rightarrow$ & \textit{12004} & /{\textesh}a{\textesh}i{\texttheta}{\texttheta}an/&$\rightarrow$ & \textit{12007} & /{\textesh}a{\textesh}i{\texttheta}an/& \\
\end{tabular}

\vspace{20pt}\hline

\end{nopagebreak}
\filbreak



\vspace{15pt}
\begin{nopagebreak}
\noindent{\fontsize{20pt}{10pt}\textbf{kelushith} } \textit{64 of summer/August 4} (N)\\
\noindent {\tovian \fontsize{20pt}{10pt} \textbf{kelushith} }\\
\noindent /kel{\textprimstress}u{\textesh}i{\texttheta}/\\
\noindent lit. branch+sharpness\\
\noindent \textit{Starthorn}\\


\noindent History:

\vspace{-0pt}
\hspace{40pt}
\begin{tabular}{ccc}
\textit{6000} & /kelu{\textyogh}i{\texttheta}/&$\rightarrow$ & \textit{12001} & /kelu{\textesh}i{\texttheta}/& \\
\end{tabular}

\vspace{20pt}\hline

\end{nopagebreak}
\filbreak



\vspace{15pt}
\begin{nopagebreak}
\noindent{\fontsize{20pt}{10pt}\textbf{nanamen} } \textit{64 of winter/February 4} (N)\\
\noindent {\tovian \fontsize{20pt}{10pt} \textbf{nanamen} }\\
\noindent /nan{\textprimstress}amen/\\
\noindent lit. star+stone\\
\noindent \textit{Starbrimstone}\\


\noindent History:

\vspace{-0pt}
\hspace{40pt}
\begin{tabular}{ccc}
\textit{6000} & /nnaanamern/&$\rightarrow$ & \textit{6500} & /naanamern/&$\rightarrow$ & \textit{7600} & /naanamen/&$\rightarrow$ & \textit{12005} & /nanamen/& \\
\end{tabular}

\vspace{20pt}\hline

\end{nopagebreak}
\filbreak



\vspace{15pt}
\begin{nopagebreak}
\noindent{\fontsize{20pt}{10pt}\textbf{teleay} } \textit{65 of autumn/November 5} (N)\\
\noindent {\tovian \fontsize{20pt}{10pt} \textbf{teleay} }\\
\noindent /{\textsubbridge{t}}el{\textprimstress}eay/\\
\noindent lit. lake+feather\\
\noindent \textit{Lakequill}\\


\noindent History:

\vspace{-0pt}
\hspace{40pt}
\begin{tabular}{ccc}
\textit{6000} & /o{\textsubbridge{t}}elearya/&$\rightarrow$ & \textit{7500} & /{\textsubbridge{t}}elearya/&$\rightarrow$ & \textit{7600} & /{\textsubbridge{t}}eleaya/&$\rightarrow$ & \textit{9500} & /{\textsubbridge{t}}eleay/& \\
\end{tabular}

\vspace{20pt}\hline

\end{nopagebreak}
\filbreak



\vspace{15pt}
\begin{nopagebreak}
\noindent{\fontsize{20pt}{10pt}\textbf{doltlan} } \textit{65 of spring/May 5} (N)\\
\noindent {\tovian \fontsize{20pt}{10pt} \textbf{doltlan} }\\
\noindent /d{\textprimstress}ol{\texttoptiebar{t\textbeltl}}an/\\
\noindent lit. death+animal\\
\noindent \textit{Stambul}\\


\noindent History:

\vspace{-0pt}
\hspace{40pt}
\begin{tabular}{ccc}
\textit{6000} & /dol{\textsubbridge{t}}lamn/&$\rightarrow$ & \textit{6401} & /dol{\texttoptiebar{t\textbeltl}}amn/&$\rightarrow$ & \textit{7600} & /dol{\texttoptiebar{t\textbeltl}}an/& \\
\end{tabular}

\vspace{20pt}\hline

\end{nopagebreak}
\filbreak



\vspace{15pt}
\begin{nopagebreak}
\noindent{\fontsize{20pt}{10pt}\textbf{mral} } \textit{65 of summer/August 5} (N)\\
\noindent {\tovian \fontsize{20pt}{10pt} \textbf{mral} }\\
\noindent /mr{\textprimstress}al/\\
\noindent lit. darkness+???\\
\noindent \textit{Nightahn}\\


\noindent History:

\vspace{-0pt}
\hspace{40pt}
\begin{tabular}{ccc}
\textit{6000} & /bral???/&$\rightarrow$ & \textit{8500} & /bral/&$\rightarrow$ & \textit{11000} & /mral/& \\
\end{tabular}

\vspace{20pt}\hline

\end{nopagebreak}
\filbreak



\vspace{15pt}
\begin{nopagebreak}
\noindent{\fontsize{20pt}{10pt}\textbf{ysanon} } \textit{65 of winter/February 5} (N)\\
\noindent {\tovian \fontsize{20pt}{10pt} \textbf{ysanon} }\\
\noindent /ys{\textprimstress}anon/\\
\noindent lit. moon+dog\\
\noindent \textit{Moonhound}\\


\noindent History:

\vspace{-0pt}
\hspace{40pt}
\begin{tabular}{ccc}
\textit{6000} & /aliysandorn/&$\rightarrow$ & \textit{6300} & /alysandorn/&$\rightarrow$ & \textit{7500} & /lysandorn/&$\rightarrow$ & \textit{7501} & /ysandorn/&$\rightarrow$ & \textit{7600} & /ysandon/&$\rightarrow$ & \textit{8750} & /ysanon/& \\
\end{tabular}

\vspace{20pt}\hline

\end{nopagebreak}
\filbreak



\vspace{15pt}
\begin{nopagebreak}
\noindent{\fontsize{20pt}{10pt}\textbf{holalor} } \textit{66 of autumn/November 6} (N)\\
\noindent {\tovian \fontsize{20pt}{10pt} \textbf{holalor} }\\
\noindent /ħol{\textprimstress}alor/\\
\noindent lit. sweet+plant\\
\noindent \textit{Sweetbalm}\\


\noindent History:

\vspace{-0pt}
\hspace{40pt}
\begin{tabular}{ccc}
\textit{6000} & /holamllore/&$\rightarrow$ & \textit{6500} & /holamlore/&$\rightarrow$ & \textit{7600} & /holalore/&$\rightarrow$ & \textit{9500} & /holalor/&$\rightarrow$ & \textit{14006} & /ħolalor/& \\
\end{tabular}

\vspace{20pt}\hline

\end{nopagebreak}
\filbreak



\vspace{15pt}
\begin{nopagebreak}
\noindent{\fontsize{20pt}{10pt}\textbf{rithor} } \textit{66 of spring/May 6} (N)\\
\noindent {\tovian \fontsize{20pt}{10pt} \textbf{rithor} }\\
\noindent /r{\textprimstress}i{\texttheta}or/\\
\noindent lit. air+plant\\
\noindent \textit{Rinthyme}\\


\noindent History:

\vspace{-0pt}
\hspace{40pt}
\begin{tabular}{ccc}
\textit{6000} & /ri{\texttheta}mllore/&$\rightarrow$ & \textit{6500} & /ri{\texttheta}mlore/&$\rightarrow$ & \textit{7600} & /ri{\texttheta}lore/&$\rightarrow$ & \textit{9500} & /ri{\texttheta}lor/&$\rightarrow$ & \textit{12500} & /ri{\texttheta}or/& \\
\end{tabular}

\vspace{20pt}\hline

\end{nopagebreak}
\filbreak



\vspace{15pt}
\begin{nopagebreak}
\noindent{\fontsize{20pt}{10pt}\textbf{lhenathan} } \textit{66 of summer/August 6} (N)\\
\noindent {\tovian \fontsize{20pt}{10pt} \textbf{lhenathan} }\\
\noindent /{\textbeltl}en{\textprimstress}a{\texttheta}an/\\
\noindent lit. dream+tool\\
\noindent \textit{Drembell}\\


\noindent History:

\vspace{-0pt}
\hspace{40pt}
\begin{tabular}{ccc}
\textit{6000} & /{\textbeltl}ena{\dh}an/&$\rightarrow$ & \textit{12002} & /{\textbeltl}ena{\texttheta}an/& \\
\end{tabular}

\vspace{20pt}\hline

\end{nopagebreak}
\filbreak



\vspace{15pt}
\begin{nopagebreak}
\noindent{\fontsize{20pt}{10pt}\textbf{lhenashotl} } \textit{66 of winter/February 6} (N)\\
\noindent {\tovian \fontsize{20pt}{10pt} \textbf{lhenashotl} }\\
\noindent /{\textbeltl}en{\textprimstress}a{\textesh}o{\texttoptiebar{t\textbeltl}}/\\
\noindent lit. dream+flow\\
\noindent \textit{Dreamlava}\\


\noindent History:

\vspace{-0pt}
\hspace{40pt}
\begin{tabular}{ccc}
\textit{6000} & /{\textbeltl}ena{\textesh}o{\texttoptiebar{t\textbeltl}}/& \\
\end{tabular}

\vspace{20pt}\hline

\end{nopagebreak}
\filbreak



\vspace{15pt}
\begin{nopagebreak}
\noindent{\fontsize{20pt}{10pt}\textbf{lhuloresyaleth} } \textit{67 of autumn/November 7} (N)\\
\noindent {\tovian \fontsize{20pt}{10pt} \textbf{lhuloresyaleth} }\\
\noindent /{\textbeltl}uloresy{\textprimstress}ale{\texttheta}/\\
\noindent lit. goodness+grain\\
\noindent \textit{Sugargrain}\\


\noindent History:

\vspace{-0pt}
\hspace{40pt}
\begin{tabular}{ccc}
\textit{6000} & /{\textbeltl}ummlloresyale{\dh}a/&$\rightarrow$ & \textit{6500} & /{\textbeltl}umloresyale{\dh}a/&$\rightarrow$ & \textit{7600} & /{\textbeltl}uloresyale{\dh}a/&$\rightarrow$ & \textit{9500} & /{\textbeltl}uloresyale{\dh}/&$\rightarrow$ & \textit{12002} & /{\textbeltl}uloresyale{\texttheta}/& \\
\end{tabular}

\vspace{20pt}\hline

\end{nopagebreak}
\filbreak



\vspace{15pt}
\begin{nopagebreak}
\noindent{\fontsize{20pt}{10pt}\textbf{munaniay} } \textit{67 of spring/May 7} (N)\\
\noindent {\tovian \fontsize{20pt}{10pt} \textbf{munaniay} }\\
\noindent /munan{\textprimstress}iay/\\
\noindent lit. protection+moon\\
\noindent \textit{Moonsponge}\\


\noindent History:

\vspace{-0pt}
\hspace{40pt}
\begin{tabular}{ccc}
\textit{6000} & /bnanialiys/&$\rightarrow$ & \textit{6300} & /bnanialys/&$\rightarrow$ & \textit{7600} & /bnaniays/&$\rightarrow$ & \textit{11000} & /mnaniays/&$\rightarrow$ & \textit{11996} & /munaniays/&$\rightarrow$ & \textit{13001} & /munaniay/& \\
\end{tabular}

\vspace{20pt}\hline

\end{nopagebreak}
\filbreak



\vspace{15pt}
\begin{nopagebreak}
\noindent{\fontsize{20pt}{10pt}\textbf{aysdrem} } \textit{67 of summer/August 7} (N)\\
\noindent {\tovian \fontsize{20pt}{10pt} \textbf{aysdrem} }\\
\noindent /{\textprimstress}aysdrem/\\
\noindent lit. moon+flower\\
\noindent \textit{Moonlily}\\


\noindent History:

\vspace{-0pt}
\hspace{40pt}
\begin{tabular}{ccc}
\textit{6000} & /aliysdremp/&$\rightarrow$ & \textit{6300} & /alysdremp/&$\rightarrow$ & \textit{7600} & /aysdremp/&$\rightarrow$ & \textit{8750} & /aysdrem/& \\
\end{tabular}

\vspace{20pt}\hline

\end{nopagebreak}
\filbreak



\vspace{15pt}
\begin{nopagebreak}
\noindent{\fontsize{20pt}{10pt}\textbf{menlor} } \textit{67 of winter/February 7} (N)\\
\noindent {\tovian \fontsize{20pt}{10pt} \textbf{menlor} }\\
\noindent /m{\textprimstress}enlor/\\
\noindent lit. earth+plant\\
\noindent \textit{Earthmoss}\\


\noindent History:

\vspace{-0pt}
\hspace{40pt}
\begin{tabular}{ccc}
\textit{6000} & /mernmllore/&$\rightarrow$ & \textit{6500} & /mernmlore/&$\rightarrow$ & \textit{7600} & /menlore/&$\rightarrow$ & \textit{9500} & /menlor/& \\
\end{tabular}

\vspace{20pt}\hline

\end{nopagebreak}
\filbreak



\vspace{15pt}
\begin{nopagebreak}
\noindent{\fontsize{20pt}{10pt}\textbf{dolag} } \textit{68 of autumn/November 8} (N)\\
\noindent {\tovian \fontsize{20pt}{10pt} \textbf{dolag} }\\
\noindent /d{\textprimstress}olag/\\
\noindent lit. death+hotness\\
\noindent \textit{Twilily}\\


\noindent History:

\vspace{-0pt}
\hspace{40pt}
\begin{tabular}{ccc}
\textit{6000} & /dolaga/&$\rightarrow$ & \textit{9500} & /dolag/& \\
\end{tabular}

\vspace{20pt}\hline

\end{nopagebreak}
\filbreak



\vspace{15pt}
\begin{nopagebreak}
\noindent{\fontsize{20pt}{10pt}\textbf{munorhanalor} } \textit{68 of spring/May 8} (N)\\
\noindent {\tovian \fontsize{20pt}{10pt} \textbf{munorhanalor} }\\
\noindent /munorħan{\textprimstress}alor/\\
\noindent lit. dusk+plant\\
\noindent \textit{Duskdaphne}\\


\noindent History:

\vspace{-0pt}
\hspace{40pt}
\begin{tabular}{ccc}
\textit{6000} & /pnorhanamllore/&$\rightarrow$ & \textit{6500} & /pnorhanamlore/&$\rightarrow$ & \textit{7600} & /pnorhanalore/&$\rightarrow$ & \textit{9500} & /pnorhanalor/&$\rightarrow$ & \textit{11000} & /mnorhanalor/&$\rightarrow$ & \textit{11996} & /munorhanalor/&$\rightarrow$ & \textit{14006} & /munorħanalor/& \\
\end{tabular}

\vspace{20pt}\hline

\end{nopagebreak}
\filbreak



\vspace{15pt}
\begin{nopagebreak}
\noindent{\fontsize{20pt}{10pt}\textbf{menrunan} } \textit{68 of summer/August 8} (N)\\
\noindent {\tovian \fontsize{20pt}{10pt} \textbf{menrunan} }\\
\noindent /menr{\textprimstress}unan/\\
\noindent lit. earth+hat\\
\noindent \textit{Earthcap}\\


\noindent History:

\vspace{-0pt}
\hspace{40pt}
\begin{tabular}{ccc}
\textit{6000} & /merndrunbnani/&$\rightarrow$ & \textit{7600} & /mendrunbnani/&$\rightarrow$ & \textit{8500} & /menrunnani/&$\rightarrow$ & \textit{9500} & /menrunnan/&$\rightarrow$ & \textit{12007} & /menrunan/& \\
\end{tabular}

\vspace{20pt}\hline

\end{nopagebreak}
\filbreak



\vspace{15pt}
\begin{nopagebreak}
\noindent{\fontsize{20pt}{10pt}\textbf{mersyaleth} } \textit{68 of winter/February 8} (N)\\
\noindent {\tovian \fontsize{20pt}{10pt} \textbf{mersyaleth} }\\
\noindent /mersy{\textprimstress}ale{\texttheta}/\\
\noindent lit. earth+food\\
\noindent \textit{Starfern}\\


\noindent History:

\vspace{-0pt}
\hspace{40pt}
\begin{tabular}{ccc}
\textit{6000} & /mernsyale{\dh}a/&$\rightarrow$ & \textit{6400} & /mersyale{\dh}a/&$\rightarrow$ & \textit{9500} & /mersyale{\dh}/&$\rightarrow$ & \textit{12002} & /mersyale{\texttheta}/& \\
\end{tabular}

\vspace{20pt}\hline

\end{nopagebreak}
\filbreak



\vspace{15pt}
\begin{nopagebreak}
\noindent{\fontsize{20pt}{10pt}\textbf{rithtom} } \textit{69 of autumn/November 9} (N)\\
\noindent {\tovian \fontsize{20pt}{10pt} \textbf{rithtom} }\\
\noindent /r{\textprimstress}i{\texttheta}{\textsubbridge{t}}om/\\
\noindent lit. air+root\\
\noindent \textit{Windthyme}\\


\noindent History:

\vspace{-0pt}
\hspace{40pt}
\begin{tabular}{ccc}
\textit{6000} & /ri{\texttheta}{\textsubbridge{t}}omo/&$\rightarrow$ & \textit{9500} & /ri{\texttheta}{\textsubbridge{t}}om/& \\
\end{tabular}

\vspace{20pt}\hline

\end{nopagebreak}
\filbreak



\vspace{15pt}
\begin{nopagebreak}
\noindent{\fontsize{20pt}{10pt}\textbf{teleaglar} } \textit{69 of spring/May 9} (N)\\
\noindent {\tovian \fontsize{20pt}{10pt} \textbf{teleaglar} }\\
\noindent /{\textsubbridge{t}}ele{\textprimstress}aglar/\\
\noindent lit. lake+tree\\
\noindent \textit{Telpoplar}\\


\noindent History:

\vspace{-0pt}
\hspace{40pt}
\begin{tabular}{ccc}
\textit{6000} & /o{\textsubbridge{t}}eleaglare/&$\rightarrow$ & \textit{7500} & /{\textsubbridge{t}}eleaglare/&$\rightarrow$ & \textit{9500} & /{\textsubbridge{t}}eleaglar/& \\
\end{tabular}

\vspace{20pt}\hline

\end{nopagebreak}
\filbreak



\vspace{15pt}
\begin{nopagebreak}
\noindent{\fontsize{20pt}{10pt}\textbf{selrem} } \textit{69 of summer/August 9} (N)\\
\noindent {\tovian \fontsize{20pt}{10pt} \textbf{selrem} }\\
\noindent /s{\textprimstress}elrem/\\
\noindent lit. water+flower\\
\noindent \textit{Rainblossom}\\


\noindent History:

\vspace{-0pt}
\hspace{40pt}
\begin{tabular}{ccc}
\textit{6000} & /zeldremp/&$\rightarrow$ & \textit{8750} & /zelrem/&$\rightarrow$ & \textit{12000} & /selrem/& \\
\end{tabular}

\vspace{20pt}\hline

\end{nopagebreak}
\filbreak



\vspace{15pt}
\begin{nopagebreak}
\noindent{\fontsize{20pt}{10pt}\textbf{shalhen} } \textit{69 of winter/February 9} (N)\\
\noindent {\tovian \fontsize{20pt}{10pt} \textbf{shalhen} }\\
\noindent /{\textesh}{\textprimstress}a{\textbeltl}en/\\
\noindent lit. ice+dream\\
\noindent \textit{Frostniter}\\


\noindent History:

\vspace{-0pt}
\hspace{40pt}
\begin{tabular}{ccc}
\textit{6000} & /{\textesh}u{\textesh}a{\textbeltl}ena/&$\rightarrow$ & \textit{9200} & /{\textesh}{\textschwa}{\textesh}a{\textbeltl}ena/&$\rightarrow$ & \textit{9500} & /{\textesh}{\textschwa}{\textesh}a{\textbeltl}en/&$\rightarrow$ & \textit{11990} & /{\textesh}{\textesh}a{\textbeltl}en/&$\rightarrow$ & \textit{12007} & /{\textesh}a{\textbeltl}en/& \\
\end{tabular}

\vspace{20pt}\hline

\end{nopagebreak}
\filbreak



\vspace{15pt}
\begin{nopagebreak}
\noindent{\fontsize{20pt}{10pt}\textbf{ay} } \textit{7 of autumn/September 7} (N)\\
\noindent {\tovian \fontsize{20pt}{10pt} \textbf{ay} }\\
\noindent /{\textprimstress}ay/\\
\noindent lit. moon+???\\
\noindent \textit{Mooncoke}\\


\noindent History:

\vspace{-0pt}
\hspace{40pt}
\begin{tabular}{ccc}
\textit{6000} & /aliys???/&$\rightarrow$ & \textit{6300} & /alys???/&$\rightarrow$ & \textit{7600} & /ays???/&$\rightarrow$ & \textit{8500} & /ays/&$\rightarrow$ & \textit{13001} & /ay/& \\
\end{tabular}

\vspace{20pt}\hline

\end{nopagebreak}
\filbreak



\vspace{15pt}
\begin{nopagebreak}
\noindent{\fontsize{20pt}{10pt}\textbf{mrilamen} } \textit{7 of spring/March 7} (N)\\
\noindent {\tovian \fontsize{20pt}{10pt} \textbf{mrilamen} }\\
\noindent /mril{\textprimstress}amen/\\
\noindent lit. sight+stone\\
\noindent \textit{Mrilgem}\\


\noindent History:

\vspace{-0pt}
\hspace{40pt}
\begin{tabular}{ccc}
\textit{6000} & /emrilamern/&$\rightarrow$ & \textit{7500} & /mrilamern/&$\rightarrow$ & \textit{7600} & /mrilamen/& \\
\end{tabular}

\vspace{20pt}\hline

\end{nopagebreak}
\filbreak



\vspace{15pt}
\begin{nopagebreak}
\noindent{\fontsize{20pt}{10pt}\textbf{glarenan} } \textit{7 of summer/June 7} (N)\\
\noindent {\tovian \fontsize{20pt}{10pt} \textbf{glarenan} }\\
\noindent /glar{\textprimstress}enan/\\
\noindent lit. tree+star\\
\noindent \textit{Starbirch}\\


\noindent History:

\vspace{-0pt}
\hspace{40pt}
\begin{tabular}{ccc}
\textit{6000} & /aglarennaana/&$\rightarrow$ & \textit{6500} & /aglarenaana/&$\rightarrow$ & \textit{7500} & /glarenaana/&$\rightarrow$ & \textit{9500} & /glarenaan/&$\rightarrow$ & \textit{12005} & /glarenan/& \\
\end{tabular}

\vspace{20pt}\hline

\end{nopagebreak}
\filbreak



\vspace{15pt}
\begin{nopagebreak}
\noindent{\fontsize{20pt}{10pt}\textbf{mloanathan} } \textit{7 of winter/december 7} (N)\\
\noindent {\tovian \fontsize{20pt}{10pt} \textbf{mloanathan} }\\
\noindent /mloan{\textprimstress}a{\texttheta}an/\\
\noindent lit. morning+tool\\
\noindent \textit{Dawnthorn}\\


\noindent History:

\vspace{-0pt}
\hspace{40pt}
\begin{tabular}{ccc}
\textit{6000} & /ploana{\dh}an/&$\rightarrow$ & \textit{11000} & /mloana{\dh}an/&$\rightarrow$ & \textit{12002} & /mloana{\texttheta}an/& \\
\end{tabular}

\vspace{20pt}\hline

\end{nopagebreak}
\filbreak



\vspace{15pt}
\begin{nopagebreak}
\noindent{\fontsize{20pt}{10pt}\textbf{mrallhenathan} } \textit{70 of autumn/November 10} (N)\\
\noindent {\tovian \fontsize{20pt}{10pt} \textbf{mrallhenathan} }\\
\noindent /mral{\textbeltl}en{\textprimstress}a{\texttheta}an/\\
\noindent lit. darkness+dream+tool\\
\noindent \textit{Shedscythe}\\


\noindent History:

\vspace{-0pt}
\hspace{40pt}
\begin{tabular}{ccc}
\textit{6000} & /bral{\textbeltl}ena{\dh}an/&$\rightarrow$ & \textit{11000} & /mral{\textbeltl}ena{\dh}an/&$\rightarrow$ & \textit{12002} & /mral{\textbeltl}ena{\texttheta}an/& \\
\end{tabular}

\vspace{20pt}\hline

\end{nopagebreak}
\filbreak



\vspace{15pt}
\begin{nopagebreak}
\noindent{\fontsize{20pt}{10pt}\textbf{shotlol} } \textit{70 of spring/May 10} (N)\\
\noindent {\tovian \fontsize{20pt}{10pt} \textbf{shotlol} }\\
\noindent /{\textesh}{\textprimstress}o{\texttoptiebar{t\textbeltl}}ol/\\
\noindent lit. flowing+fruit\\
\noindent \textit{Mistquince}\\


\noindent History:

\vspace{-0pt}
\hspace{40pt}
\begin{tabular}{ccc}
\textit{6000} & /{\textesh}o{\texttoptiebar{t\textbeltl}}{\textyogh}ola/&$\rightarrow$ & \textit{9500} & /{\textesh}o{\texttoptiebar{t\textbeltl}}{\textyogh}ol/&$\rightarrow$ & \textit{12001} & /{\textesh}o{\texttoptiebar{t\textbeltl}}{\textesh}ol/&$\rightarrow$ & \textit{14000} & /{\textesh}o{\texttoptiebar{t\textbeltl}}ol/& \\
\end{tabular}

\vspace{20pt}\hline

\end{nopagebreak}
\filbreak



\vspace{15pt}
\begin{nopagebreak}
\noindent{\fontsize{20pt}{10pt}\textbf{nanalor} } \textit{70 of summer/August 10} (N)\\
\noindent {\tovian \fontsize{20pt}{10pt} \textbf{nanalor} }\\
\noindent /nan{\textprimstress}alor/\\
\noindent lit. star+action\\
\noindent \textit{Starharvest}\\


\noindent History:

\vspace{-0pt}
\hspace{40pt}
\begin{tabular}{ccc}
\textit{6000} & /nnaanalore/&$\rightarrow$ & \textit{6500} & /naanalore/&$\rightarrow$ & \textit{9500} & /naanalor/&$\rightarrow$ & \textit{12005} & /nanalor/& \\
\end{tabular}

\vspace{20pt}\hline

\end{nopagebreak}
\filbreak



\vspace{15pt}
\begin{nopagebreak}
\noindent{\fontsize{20pt}{10pt}\textbf{tushuthan} } \textit{70 of winter/February 10} (N)\\
\noindent {\tovian \fontsize{20pt}{10pt} \textbf{tushuthan} }\\
\noindent /{\textsubbridge{t}}u{\textesh}{\textprimstress}u{\texttheta}an/\\
\noindent lit. cold+tool\\
\noindent \textit{Winterflail}\\


\noindent History:

\vspace{-0pt}
\hspace{40pt}
\begin{tabular}{ccc}
\textit{6000} & /{\textsubbridge{t}}u{\textesh}u{\dh}an/&$\rightarrow$ & \textit{12002} & /{\textsubbridge{t}}u{\textesh}u{\texttheta}an/& \\
\end{tabular}

\vspace{20pt}\hline

\end{nopagebreak}
\filbreak



\vspace{15pt}
\begin{nopagebreak}
\noindent{\fontsize{20pt}{10pt}\textbf{shithanin} } \textit{71 of autumn/November 11} (N)\\
\noindent {\tovian \fontsize{20pt}{10pt} \textbf{shithanin} }\\
\noindent /{\textesh}i{\texttheta}{\textprimstress}anin/\\
\noindent lit. sharpness+berry\\
\noindent \textit{Stefbery}\\


\noindent History:

\vspace{-0pt}
\hspace{40pt}
\begin{tabular}{ccc}
\textit{6000} & /{\textyogh}i{\texttheta}{\textyogh}lanina/&$\rightarrow$ & \textit{9500} & /{\textyogh}i{\texttheta}{\textyogh}lanin/&$\rightarrow$ & \textit{12001} & /{\textesh}i{\texttheta}{\textesh}lanin/&$\rightarrow$ & \textit{12500} & /{\textesh}i{\texttheta}{\textesh}anin/&$\rightarrow$ & \textit{13000} & /{\textesh}i{\texttheta}anin/& \\
\end{tabular}

\vspace{20pt}\hline

\end{nopagebreak}
\filbreak



\vspace{15pt}
\begin{nopagebreak}
\noindent{\fontsize{20pt}{10pt}\textbf{mramlorean} } \textit{71 of spring/May 11} (N)\\
\noindent {\tovian \fontsize{20pt}{10pt} \textbf{mramlorean} }\\
\noindent /mramlor{\textprimstress}ean/\\
\noindent lit. darkness+plant+light\\
\noindent \textit{Nightshade}\\


\noindent History:

\vspace{-0pt}
\hspace{40pt}
\begin{tabular}{ccc}
\textit{6000} & /bralmlloreana/&$\rightarrow$ & \textit{6500} & /bralmloreana/&$\rightarrow$ & \textit{7600} & /bramloreana/&$\rightarrow$ & \textit{9500} & /bramlorean/&$\rightarrow$ & \textit{11000} & /mramlorean/& \\
\end{tabular}

\vspace{20pt}\hline

\end{nopagebreak}
\filbreak



\vspace{15pt}
\begin{nopagebreak}
\noindent{\fontsize{20pt}{10pt}\textbf{lanathan} } \textit{71 of summer/August 11} (N)\\
\noindent {\tovian \fontsize{20pt}{10pt} \textbf{lanathan} }\\
\noindent /lan{\textprimstress}a{\texttheta}an/\\
\noindent lit. twilight+tool\\
\noindent \textit{Dusktang}\\


\noindent History:

\vspace{-0pt}
\hspace{40pt}
\begin{tabular}{ccc}
\textit{6000} & /brlana{\dh}an/&$\rightarrow$ & \textit{7501} & /rlana{\dh}an/&$\rightarrow$ & \textit{7600} & /lana{\dh}an/&$\rightarrow$ & \textit{12002} & /lana{\texttheta}an/& \\
\end{tabular}

\vspace{20pt}\hline

\end{nopagebreak}
\filbreak



\vspace{15pt}
\begin{nopagebreak}
\noindent{\fontsize{20pt}{10pt}\textbf{thanramen} } \textit{71 of winter/February 11} (N)\\
\noindent {\tovian \fontsize{20pt}{10pt} \textbf{thanramen} }\\
\noindent /{\texttheta}anr{\textprimstress}amen/\\
\noindent lit. tool+darkness+stone\\
\noindent \textit{Nightstone}\\


\noindent History:

\vspace{-0pt}
\hspace{40pt}
\begin{tabular}{ccc}
\textit{6000} & /{\dh}anbralmern/&$\rightarrow$ & \textit{7600} & /{\dh}anbramen/&$\rightarrow$ & \textit{8500} & /{\dh}anramen/&$\rightarrow$ & \textit{12002} & /{\texttheta}anramen/& \\
\end{tabular}

\vspace{20pt}\hline

\end{nopagebreak}
\filbreak



\vspace{15pt}
\begin{nopagebreak}
\noindent{\fontsize{20pt}{10pt}\textbf{lethileaglar} } \textit{72 of autumn/November 12} (N)\\
\noindent {\tovian \fontsize{20pt}{10pt} \textbf{lethileaglar} }\\
\noindent /le{\texttheta}ile{\textprimstress}aglar/\\
\noindent lit. forest+tree\\
\noindent \textit{Wildroot}\\


\noindent History:

\vspace{-0pt}
\hspace{40pt}
\begin{tabular}{ccc}
\textit{6000} & /aglre{\texttheta}ileaglare/&$\rightarrow$ & \textit{6503} & /agrle{\texttheta}ileaglare/&$\rightarrow$ & \textit{7500} & /grle{\texttheta}ileaglare/&$\rightarrow$ & \textit{7501} & /rle{\texttheta}ileaglare/&$\rightarrow$ & \textit{7600} & /le{\texttheta}ileaglare/&$\rightarrow$ & \textit{9500} & /le{\texttheta}ileaglar/& \\
\end{tabular}

\vspace{20pt}\hline

\end{nopagebreak}
\filbreak



\vspace{15pt}
\begin{nopagebreak}
\noindent{\fontsize{20pt}{10pt}\textbf{dleleon} } \textit{72 of spring/May 12} (N)\\
\noindent {\tovian \fontsize{20pt}{10pt} \textbf{dleleon} }\\
\noindent /dlel{\textprimstress}eon/\\
\noindent lit. trunk+mouth\\
\noindent \textit{Starstalk}\\


\noindent History:

\vspace{-0pt}
\hspace{40pt}
\begin{tabular}{ccc}
\textit{6000} & /adllenleoln/&$\rightarrow$ & \textit{6500} & /adlenleoln/&$\rightarrow$ & \textit{7500} & /dlenleoln/&$\rightarrow$ & \textit{7600} & /dleleon/& \\
\end{tabular}

\vspace{20pt}\hline

\end{nopagebreak}
\filbreak



\vspace{15pt}
\begin{nopagebreak}
\noindent{\fontsize{20pt}{10pt}\textbf{mloanatesh} } \textit{72 of summer/August 12} (N)\\
\noindent {\tovian \fontsize{20pt}{10pt} \textbf{mloanatesh} }\\
\noindent /mloan{\textprimstress}a{\textsubbridge{t}}e{\textesh}/\\
\noindent lit. morning+sound\\
\noindent \textit{Dawnvox}\\


\noindent History:

\vspace{-0pt}
\hspace{40pt}
\begin{tabular}{ccc}
\textit{6000} & /ploana{\textsubbridge{t}}e{\textyogh}/&$\rightarrow$ & \textit{11000} & /mloana{\textsubbridge{t}}e{\textyogh}/&$\rightarrow$ & \textit{12001} & /mloana{\textsubbridge{t}}e{\textesh}/& \\
\end{tabular}

\vspace{20pt}\hline

\end{nopagebreak}
\filbreak



\vspace{15pt}
\begin{nopagebreak}
\noindent{\fontsize{20pt}{10pt}\textbf{numvar} } \textit{72 of winter/February 12} (N)\\
\noindent {\tovian \fontsize{20pt}{10pt} \textbf{numvar} }\\
\noindent /n{\textprimstress}umvar/\\
\noindent lit. path+value\\
\noindent \textit{Numkar}\\


\noindent History:

\vspace{-0pt}
\hspace{40pt}
\begin{tabular}{ccc}
\textit{6000} & /numvlare/&$\rightarrow$ & \textit{9500} & /numvlar/&$\rightarrow$ & \textit{12500} & /numvar/& \\
\end{tabular}

\vspace{20pt}\hline

\end{nopagebreak}
\filbreak



\vspace{15pt}
\begin{nopagebreak}
\noindent{\fontsize{20pt}{10pt}\textbf{kanaithr} } \textit{73 of autumn/November 13} (N)\\
\noindent {\tovian \fontsize{20pt}{10pt} \textbf{kanaithr} }\\
\noindent /kan{\textprimstress}ai{\texttheta}r/\\
\noindent lit. change+thought\\
\noindent \textit{Kanalh}\\


\noindent History:

\vspace{-0pt}
\hspace{40pt}
\begin{tabular}{ccc}
\textit{6000} & /kanai{\texttheta}r/& \\
\end{tabular}

\vspace{20pt}\hline

\end{nopagebreak}
\filbreak



\vspace{15pt}
\begin{nopagebreak}
\noindent{\fontsize{20pt}{10pt}\textbf{???} } \textit{73 of spring/May 13} (N)\\
\noindent {\tovian \fontsize{20pt}{10pt} \textbf{???} }\\
\noindent /???/\\
\noindent lit. ???\\
\noindent \textit{Bayra}\\


\noindent History:

\vspace{-0pt}
\hspace{40pt}
\begin{tabular}{ccc}
\textit{6000} & /???/& \\
\end{tabular}

\vspace{20pt}\hline

\end{nopagebreak}
\filbreak



\vspace{15pt}
\begin{nopagebreak}
\noindent{\fontsize{20pt}{10pt}\textbf{efan} } \textit{73 of summer/August 13} (N)\\
\noindent {\tovian \fontsize{20pt}{10pt} \textbf{efan} }\\
\noindent /{\textprimstress}efan/\\
\noindent lit. sun+tool\\
\noindent \textit{Sunstaff}\\


\noindent History:

\vspace{-0pt}
\hspace{40pt}
\begin{tabular}{ccc}
\textit{6000} & /efl{\dh}an/&$\rightarrow$ & \textit{12002} & /efl{\texttheta}an/&$\rightarrow$ & \textit{12500} & /ef{\texttheta}an/&$\rightarrow$ & \textit{13000} & /efan/& \\
\end{tabular}

\vspace{20pt}\hline

\end{nopagebreak}
\filbreak



\vspace{15pt}
\begin{nopagebreak}
\noindent{\fontsize{20pt}{10pt}\textbf{mramen} } \textit{73 of winter/February 13} (N)\\
\noindent {\tovian \fontsize{20pt}{10pt} \textbf{mramen} }\\
\noindent /mr{\textprimstress}amen/\\
\noindent lit. darkness+stone\\
\noindent \textit{Shedowslate}\\


\noindent History:

\vspace{-0pt}
\hspace{40pt}
\begin{tabular}{ccc}
\textit{6000} & /bralmern/&$\rightarrow$ & \textit{7600} & /bramen/&$\rightarrow$ & \textit{11000} & /mramen/& \\
\end{tabular}

\vspace{20pt}\hline

\end{nopagebreak}
\filbreak



\vspace{15pt}
\begin{nopagebreak}
\noindent{\fontsize{20pt}{10pt}\textbf{sillhum} } \textit{74 of autumn/November 14} (N)\\
\noindent {\tovian \fontsize{20pt}{10pt} \textbf{sillhum} }\\
\noindent /s{\textprimstress}il{\textbeltl}um/\\
\noindent lit. existence+goodness\\
\noindent \textit{Silshede}\\


\noindent History:

\vspace{-0pt}
\hspace{40pt}
\begin{tabular}{ccc}
\textit{6000} & /sil{\textbeltl}um/& \\
\end{tabular}

\vspace{20pt}\hline

\end{nopagebreak}
\filbreak



\vspace{15pt}
\begin{nopagebreak}
\noindent{\fontsize{20pt}{10pt}\textbf{shotl} } \textit{74 of spring/May 14} (N)\\
\noindent {\tovian \fontsize{20pt}{10pt} \textbf{shotl} }\\
\noindent /{\textesh}{\textprimstress}o{\texttoptiebar{t\textbeltl}}/\\
\noindent lit. flowing+???\\
\noindent \textit{Misazel}\\


\noindent History:

\vspace{-0pt}
\hspace{40pt}
\begin{tabular}{ccc}
\textit{6000} & /{\textesh}o{\texttoptiebar{t\textbeltl}}???/&$\rightarrow$ & \textit{8500} & /{\textesh}o{\texttoptiebar{t\textbeltl}}/& \\
\end{tabular}

\vspace{20pt}\hline

\end{nopagebreak}
\filbreak



\vspace{15pt}
\begin{nopagebreak}
\noindent{\fontsize{20pt}{10pt}\textbf{glaretesh} } \textit{74 of summer/August 14} (N)\\
\noindent {\tovian \fontsize{20pt}{10pt} \textbf{glaretesh} }\\
\noindent /glar{\textprimstress}e{\textsubbridge{t}}e{\textesh}/\\
\noindent lit. tree+sound\\
\noindent \textit{Palmwhisper}\\


\noindent History:

\vspace{-0pt}
\hspace{40pt}
\begin{tabular}{ccc}
\textit{6000} & /aglare{\textsubbridge{t}}e{\textyogh}/&$\rightarrow$ & \textit{7500} & /glare{\textsubbridge{t}}e{\textyogh}/&$\rightarrow$ & \textit{12001} & /glare{\textsubbridge{t}}e{\textesh}/& \\
\end{tabular}

\vspace{20pt}\hline

\end{nopagebreak}
\filbreak



\vspace{15pt}
\begin{nopagebreak}
\noindent{\fontsize{20pt}{10pt}\textbf{menmen} } \textit{74 of winter/February 14} (N)\\
\noindent {\tovian \fontsize{20pt}{10pt} \textbf{menmen} }\\
\noindent /m{\textprimstress}enmen/\\
\noindent lit. earth+stone\\
\noindent \textit{Peatstone}\\


\noindent History:

\vspace{-0pt}
\hspace{40pt}
\begin{tabular}{ccc}
\textit{6000} & /mernmern/&$\rightarrow$ & \textit{7600} & /menmen/& \\
\end{tabular}

\vspace{20pt}\hline

\end{nopagebreak}
\filbreak



\vspace{15pt}
\begin{nopagebreak}
\noindent{\fontsize{20pt}{10pt}\textbf{mloanarithan} } \textit{75 of autumn/November 15} (N)\\
\noindent {\tovian \fontsize{20pt}{10pt} \textbf{mloanarithan} }\\
\noindent /mloanar{\textprimstress}i{\texttheta}an/\\
\noindent lit. morning+bird\\
\noindent \textit{Dawnquail}\\


\noindent History:

\vspace{-0pt}
\hspace{40pt}
\begin{tabular}{ccc}
\textit{6000} & /ploanari{\texttheta}lamn/&$\rightarrow$ & \textit{7600} & /ploanari{\texttheta}lan/&$\rightarrow$ & \textit{11000} & /mloanari{\texttheta}lan/&$\rightarrow$ & \textit{12500} & /mloanari{\texttheta}an/& \\
\end{tabular}

\vspace{20pt}\hline

\end{nopagebreak}
\filbreak



\vspace{15pt}
\begin{nopagebreak}
\noindent{\fontsize{20pt}{10pt}\textbf{aysan} } \textit{75 of spring/May 15} (N)\\
\noindent {\tovian \fontsize{20pt}{10pt} \textbf{aysan} }\\
\noindent /{\textprimstress}aysan/\\
\noindent lit. moon+animal\\
\noindent \textit{Mooncalf}\\


\noindent History:

\vspace{-0pt}
\hspace{40pt}
\begin{tabular}{ccc}
\textit{6000} & /aliys{\textsubbridge{t}}lamn/&$\rightarrow$ & \textit{6300} & /alys{\textsubbridge{t}}lamn/&$\rightarrow$ & \textit{6401} & /alys{\texttoptiebar{t\textbeltl}}amn/&$\rightarrow$ & \textit{7600} & /ays{\texttoptiebar{t\textbeltl}}an/&$\rightarrow$ & \textit{14001} & /aysan/& \\
\end{tabular}

\vspace{20pt}\hline

\end{nopagebreak}
\filbreak



\vspace{15pt}
\begin{nopagebreak}
\noindent{\fontsize{20pt}{10pt}\textbf{holalor} } \textit{75 of summer/August 15} (N)\\
\noindent {\tovian \fontsize{20pt}{10pt} \textbf{holalor} }\\
\noindent /ħol{\textprimstress}alor/\\
\noindent lit. sweet+action\\
\noindent \textit{Silkspinner}\\


\noindent History:

\vspace{-0pt}
\hspace{40pt}
\begin{tabular}{ccc}
\textit{6000} & /holalore/&$\rightarrow$ & \textit{9500} & /holalor/&$\rightarrow$ & \textit{14006} & /ħolalor/& \\
\end{tabular}

\vspace{20pt}\hline

\end{nopagebreak}
\filbreak



\vspace{15pt}
\begin{nopagebreak}
\noindent{\fontsize{20pt}{10pt}\textbf{shashol} } \textit{75 of winter/February 15} (N)\\
\noindent {\tovian \fontsize{20pt}{10pt} \textbf{shashol} }\\
\noindent /{\textesh}{\textprimstress}a{\textesh}ol/\\
\noindent lit. ice+fruit\\
\noindent \textit{Frostbun}\\


\noindent History:

\vspace{-0pt}
\hspace{40pt}
\begin{tabular}{ccc}
\textit{6000} & /{\textesh}u{\textesh}a{\textyogh}ola/&$\rightarrow$ & \textit{9200} & /{\textesh}{\textschwa}{\textesh}a{\textyogh}ola/&$\rightarrow$ & \textit{9500} & /{\textesh}{\textschwa}{\textesh}a{\textyogh}ol/&$\rightarrow$ & \textit{11990} & /{\textesh}{\textesh}a{\textyogh}ol/&$\rightarrow$ & \textit{12001} & /{\textesh}{\textesh}a{\textesh}ol/&$\rightarrow$ & \textit{12007} & /{\textesh}a{\textesh}ol/& \\
\end{tabular}

\vspace{20pt}\hline

\end{nopagebreak}
\filbreak



\vspace{15pt}
\begin{nopagebreak}
\noindent{\fontsize{20pt}{10pt}\textbf{fedemralrem} } \textit{76 of autumn/November 16} (N)\\
\noindent {\tovian \fontsize{20pt}{10pt} \textbf{fedemralrem} }\\
\noindent /fedemr{\textprimstress}alrem/\\
\noindent lit. showing+darkness+flower\\
\noindent \textit{Duflor}\\


\noindent History:

\vspace{-0pt}
\hspace{40pt}
\begin{tabular}{ccc}
\textit{6000} & /fedebraldremp/&$\rightarrow$ & \textit{8750} & /fedebralrem/&$\rightarrow$ & \textit{11000} & /fedemralrem/& \\
\end{tabular}

\vspace{20pt}\hline

\end{nopagebreak}
\filbreak



\vspace{15pt}
\begin{nopagebreak}
\noindent{\fontsize{20pt}{10pt}\textbf{rithilesyaleth} } \textit{76 of spring/May 16} (N)\\
\noindent {\tovian \fontsize{20pt}{10pt} \textbf{rithilesyaleth} }\\
\noindent /ri{\texttheta}ilesy{\textprimstress}ale{\texttheta}/\\
\noindent lit. sky+food\\
\noindent \textit{Sky's Feed}\\


\noindent History:

\vspace{-0pt}
\hspace{40pt}
\begin{tabular}{ccc}
\textit{6000} & /ri{\texttheta}ilesyale{\dh}a/&$\rightarrow$ & \textit{9500} & /ri{\texttheta}ilesyale{\dh}/&$\rightarrow$ & \textit{12002} & /ri{\texttheta}ilesyale{\texttheta}/& \\
\end{tabular}

\vspace{20pt}\hline

\end{nopagebreak}
\filbreak



\vspace{15pt}
\begin{nopagebreak}
\noindent{\fontsize{20pt}{10pt}\textbf{melrem} } \textit{76 of summer/August 16} (N)\\
\noindent {\tovian \fontsize{20pt}{10pt} \textbf{melrem} }\\
\noindent /m{\textprimstress}elrem/\\
\noindent lit. life+flower\\
\noindent \textit{Healbloom}\\


\noindent History:

\vspace{-0pt}
\hspace{40pt}
\begin{tabular}{ccc}
\textit{6000} & /meldremp/&$\rightarrow$ & \textit{8750} & /melrem/& \\
\end{tabular}

\vspace{20pt}\hline

\end{nopagebreak}
\filbreak



\vspace{15pt}
\begin{nopagebreak}
\noindent{\fontsize{20pt}{10pt}\textbf{namen} } \textit{76 of winter/February 16} (N)\\
\noindent {\tovian \fontsize{20pt}{10pt} \textbf{namen} }\\
\noindent /n{\textprimstress}amen/\\
\noindent lit. light+stone\\
\noindent \textit{Sparkstone}\\


\noindent History:

\vspace{-0pt}
\hspace{40pt}
\begin{tabular}{ccc}
\textit{6000} & /anamern/&$\rightarrow$ & \textit{7500} & /namern/&$\rightarrow$ & \textit{7600} & /namen/& \\
\end{tabular}

\vspace{20pt}\hline

\end{nopagebreak}
\filbreak



\vspace{15pt}
\begin{nopagebreak}
\noindent{\fontsize{20pt}{10pt}\textbf{lhiladrem} } \textit{77 of autumn/November 17} (N)\\
\noindent {\tovian \fontsize{20pt}{10pt} \textbf{lhiladrem} }\\
\noindent /{\textbeltl}il{\textprimstress}adrem/\\
\noindent lit. night+flower\\
\noindent \textit{Niteblum}\\


\noindent History:

\vspace{-0pt}
\hspace{40pt}
\begin{tabular}{ccc}
\textit{6000} & /a{\textbeltl}inladremp/&$\rightarrow$ & \textit{7500} & /{\textbeltl}inladremp/&$\rightarrow$ & \textit{7600} & /{\textbeltl}iladremp/&$\rightarrow$ & \textit{8750} & /{\textbeltl}iladrem/& \\
\end{tabular}

\vspace{20pt}\hline

\end{nopagebreak}
\filbreak



\vspace{15pt}
\begin{nopagebreak}
\noindent{\fontsize{20pt}{10pt}\textbf{kelueon} } \textit{77 of spring/May 17} (N)\\
\noindent {\tovian \fontsize{20pt}{10pt} \textbf{kelueon} }\\
\noindent /kelu{\textprimstress}eon/\\
\noindent lit. branch+mouth\\
\noindent \textit{Starlichen}\\


\noindent History:

\vspace{-0pt}
\hspace{40pt}
\begin{tabular}{ccc}
\textit{6000} & /kelueoln/&$\rightarrow$ & \textit{7600} & /kelueon/& \\
\end{tabular}

\vspace{20pt}\hline

\end{nopagebreak}
\filbreak



\vspace{15pt}
\begin{nopagebreak}
\noindent{\fontsize{20pt}{10pt}\textbf{keluthan} } \textit{77 of summer/August 17} (N)\\
\noindent {\tovian \fontsize{20pt}{10pt} \textbf{keluthan} }\\
\noindent /kel{\textprimstress}u{\texttheta}an/\\
\noindent lit. branch+tool\\
\noindent \textit{Starpim}\\


\noindent History:

\vspace{-0pt}
\hspace{40pt}
\begin{tabular}{ccc}
\textit{6000} & /kelu{\dh}an/&$\rightarrow$ & \textit{12002} & /kelu{\texttheta}an/& \\
\end{tabular}

\vspace{20pt}\hline

\end{nopagebreak}
\filbreak



\vspace{15pt}
\begin{nopagebreak}
\noindent{\fontsize{20pt}{10pt}\textbf{lhenamen} } \textit{77 of winter/February 17} (N)\\
\noindent {\tovian \fontsize{20pt}{10pt} \textbf{lhenamen} }\\
\noindent /{\textbeltl}en{\textprimstress}amen/\\
\noindent lit. dream+earth\\
\noindent \textit{Dreamclay}\\


\noindent History:

\vspace{-0pt}
\hspace{40pt}
\begin{tabular}{ccc}
\textit{6000} & /{\textbeltl}enamern/&$\rightarrow$ & \textit{7600} & /{\textbeltl}enamen/& \\
\end{tabular}

\vspace{20pt}\hline

\end{nopagebreak}
\filbreak



\vspace{15pt}
\begin{nopagebreak}
\noindent{\fontsize{20pt}{10pt}\textbf{lhenakel} } \textit{78 of autumn/November 18} (N)\\
\noindent {\tovian \fontsize{20pt}{10pt} \textbf{lhenakel} }\\
\noindent /{\textbeltl}en{\textprimstress}akel/\\
\noindent lit. dream+branch\\
\noindent \textit{Dremsweed}\\


\noindent History:

\vspace{-0pt}
\hspace{40pt}
\begin{tabular}{ccc}
\textit{6000} & /{\textbeltl}enakelu/&$\rightarrow$ & \textit{9500} & /{\textbeltl}enakel/& \\
\end{tabular}

\vspace{20pt}\hline

\end{nopagebreak}
\filbreak



\vspace{15pt}
\begin{nopagebreak}
\noindent{\fontsize{20pt}{10pt}\textbf{mralel} } \textit{78 of spring/May 18} (N)\\
\noindent {\tovian \fontsize{20pt}{10pt} \textbf{mralel} }\\
\noindent /mr{\textprimstress}alel/\\
\noindent lit. darkness+branch\\
\noindent \textit{Nightyew}\\


\noindent History:

\vspace{-0pt}
\hspace{40pt}
\begin{tabular}{ccc}
\textit{6000} & /bralkelu/&$\rightarrow$ & \textit{8750} & /bralelu/&$\rightarrow$ & \textit{9500} & /bralel/&$\rightarrow$ & \textit{11000} & /mralel/& \\
\end{tabular}

\vspace{20pt}\hline

\end{nopagebreak}
\filbreak



\vspace{15pt}
\begin{nopagebreak}
\noindent{\fontsize{20pt}{10pt}\textbf{varetlun} } \textit{78 of summer/August 18} (N)\\
\noindent {\tovian \fontsize{20pt}{10pt} \textbf{varetlun} }\\
\noindent /var{\textprimstress}e{\texttoptiebar{t\textbeltl}}un/\\
\noindent lit. value+joining\\
\noindent \textit{Goldweave}\\


\noindent History:

\vspace{-0pt}
\hspace{40pt}
\begin{tabular}{ccc}
\textit{6000} & /vlare{\texttoptiebar{t\textbeltl}}una/&$\rightarrow$ & \textit{9500} & /vlare{\texttoptiebar{t\textbeltl}}un/&$\rightarrow$ & \textit{12500} & /vare{\texttoptiebar{t\textbeltl}}un/& \\
\end{tabular}

\vspace{20pt}\hline

\end{nopagebreak}
\filbreak



\vspace{15pt}
\begin{nopagebreak}
\noindent{\fontsize{20pt}{10pt}\textbf{ysdeshol} } \textit{78 of winter/February 18} (N)\\
\noindent {\tovian \fontsize{20pt}{10pt} \textbf{ysdeshol} }\\
\noindent /ysd{\textprimstress}e{\textesh}ol/\\
\noindent lit. moon+showing+fruit\\
\noindent \textit{Moonlime}\\


\noindent History:

\vspace{-0pt}
\hspace{40pt}
\begin{tabular}{ccc}
\textit{6000} & /aliysfede{\textyogh}ola/&$\rightarrow$ & \textit{6300} & /alysfede{\textyogh}ola/&$\rightarrow$ & \textit{7500} & /lysfede{\textyogh}ola/&$\rightarrow$ & \textit{7501} & /ysfede{\textyogh}ola/&$\rightarrow$ & \textit{9500} & /ysfede{\textyogh}ol/&$\rightarrow$ & \textit{12001} & /ysfede{\textesh}ol/&$\rightarrow$ & \textit{12003} & /ysde{\textesh}ol/& \\
\end{tabular}

\vspace{20pt}\hline

\end{nopagebreak}
\filbreak



\vspace{15pt}
\begin{nopagebreak}
\noindent{\fontsize{20pt}{10pt}\textbf{nanathang} } \textit{79 of autumn/November 19} (N)\\
\noindent {\tovian \fontsize{20pt}{10pt} \textbf{nanathang} }\\
\noindent /nan{\textprimstress}a{\texttheta}a{\ng}/\\
\noindent lit. star+leaf\\
\noindent \textit{Starlef}\\


\noindent History:

\vspace{-0pt}
\hspace{40pt}
\begin{tabular}{ccc}
\textit{6000} & /nnaanaa{\dh}a{\ng}ge/&$\rightarrow$ & \textit{6500} & /naanaa{\dh}a{\ng}ge/&$\rightarrow$ & \textit{8750} & /naanaa{\dh}a{\ng}e/&$\rightarrow$ & \textit{9500} & /naanaa{\dh}a{\ng}/&$\rightarrow$ & \textit{12002} & /naanaa{\texttheta}a{\ng}/&$\rightarrow$ & \textit{12005} & /nana{\texttheta}a{\ng}/& \\
\end{tabular}

\vspace{20pt}\hline

\end{nopagebreak}
\filbreak



\vspace{15pt}
\begin{nopagebreak}
\noindent{\fontsize{20pt}{10pt}\textbf{ysor} } \textit{79 of spring/May 19} (N)\\
\noindent {\tovian \fontsize{20pt}{10pt} \textbf{ysor} }\\
\noindent /ys{\textprimstress}or/\\
\noindent lit. moon+plant\\
\noindent \textit{Moongrass}\\


\noindent History:

\vspace{-0pt}
\hspace{40pt}
\begin{tabular}{ccc}
\textit{6000} & /aliysmllore/&$\rightarrow$ & \textit{6300} & /alysmllore/&$\rightarrow$ & \textit{6500} & /alysmlore/&$\rightarrow$ & \textit{7500} & /lysmlore/&$\rightarrow$ & \textit{7501} & /ysmlore/&$\rightarrow$ & \textit{7600} & /yslore/&$\rightarrow$ & \textit{9500} & /yslor/&$\rightarrow$ & \textit{12500} & /ysor/& \\
\end{tabular}

\vspace{20pt}\hline

\end{nopagebreak}
\filbreak



\vspace{15pt}
\begin{nopagebreak}
\noindent{\fontsize{20pt}{10pt}\textbf{mralaysor} } \textit{79 of summer/August 19} (N)\\
\noindent {\tovian \fontsize{20pt}{10pt} \textbf{mralaysor} }\\
\noindent /mral{\textprimstress}aysor/\\
\noindent lit. darkness+moon+plant\\
\noindent \textit{Shadegreen}\\


\noindent History:

\vspace{-0pt}
\hspace{40pt}
\begin{tabular}{ccc}
\textit{6000} & /bralaliysmllore/&$\rightarrow$ & \textit{6300} & /bralalysmllore/&$\rightarrow$ & \textit{6500} & /bralalysmlore/&$\rightarrow$ & \textit{7600} & /bralayslore/&$\rightarrow$ & \textit{9500} & /bralayslor/&$\rightarrow$ & \textit{11000} & /mralayslor/&$\rightarrow$ & \textit{12500} & /mralaysor/& \\
\end{tabular}

\vspace{20pt}\hline

\end{nopagebreak}
\filbreak



\vspace{15pt}
\begin{nopagebreak}
\noindent{\fontsize{20pt}{10pt}\textbf{nanalhen} } \textit{79 of winter/February 19} (N)\\
\noindent {\tovian \fontsize{20pt}{10pt} \textbf{nanalhen} }\\
\noindent /nan{\textprimstress}a{\textbeltl}en/\\
\noindent lit. star+dream\\
\noindent \textit{Starble}\\


\noindent History:

\vspace{-0pt}
\hspace{40pt}
\begin{tabular}{ccc}
\textit{6000} & /nnaana{\textbeltl}ena/&$\rightarrow$ & \textit{6500} & /naana{\textbeltl}ena/&$\rightarrow$ & \textit{9500} & /naana{\textbeltl}en/&$\rightarrow$ & \textit{12005} & /nana{\textbeltl}en/& \\
\end{tabular}

\vspace{20pt}\hline

\end{nopagebreak}
\filbreak



\vspace{15pt}
\begin{nopagebreak}
\noindent{\fontsize{20pt}{10pt}\textbf{munorhanapal} } \textit{8 of autumn/September 8} (N)\\
\noindent {\tovian \fontsize{20pt}{10pt} \textbf{munorhanapal} }\\
\noindent /munorħan{\textprimstress}apal/\\
\noindent lit. dusk+stopping\\
\noindent \textit{Duskshalt}\\


\noindent History:

\vspace{-0pt}
\hspace{40pt}
\begin{tabular}{ccc}
\textit{6000} & /pnorhanapalr/&$\rightarrow$ & \textit{6503} & /pnorhanaparl/&$\rightarrow$ & \textit{7600} & /pnorhanapal/&$\rightarrow$ & \textit{11000} & /mnorhanapal/&$\rightarrow$ & \textit{11996} & /munorhanapal/&$\rightarrow$ & \textit{14006} & /munorħanapal/& \\
\end{tabular}

\vspace{20pt}\hline

\end{nopagebreak}
\filbreak



\vspace{15pt}
\begin{nopagebreak}
\noindent{\fontsize{20pt}{10pt}\textbf{todakesh} } \textit{8 of spring/March 8} (N)\\
\noindent {\tovian \fontsize{20pt}{10pt} \textbf{todakesh} }\\
\noindent /{\textsubbridge{t}}od{\textprimstress}ake{\textesh}/\\
\noindent lit. eternal+fire\\
\noindent \textit{Everflame}\\


\noindent History:

\vspace{-0pt}
\hspace{40pt}
\begin{tabular}{ccc}
\textit{6000} & /{\textsubbridge{t}}odake{\textesh}/& \\
\end{tabular}

\vspace{20pt}\hline

\end{nopagebreak}
\filbreak



\vspace{15pt}
\begin{nopagebreak}
\noindent{\fontsize{20pt}{10pt}\textbf{nalhom} } \textit{8 of summer/June 8} (N)\\
\noindent {\tovian \fontsize{20pt}{10pt} \textbf{nalhom} }\\
\noindent /n{\textprimstress}a{\textbeltl}om/\\
\noindent lit. light+container\\
\noindent \textit{Lightcup}\\


\noindent History:

\vspace{-0pt}
\hspace{40pt}
\begin{tabular}{ccc}
\textit{6000} & /ana{\textbeltl}oma/&$\rightarrow$ & \textit{7500} & /na{\textbeltl}oma/&$\rightarrow$ & \textit{9500} & /na{\textbeltl}om/& \\
\end{tabular}

\vspace{20pt}\hline

\end{nopagebreak}
\filbreak



\vspace{15pt}
\begin{nopagebreak}
\noindent{\fontsize{20pt}{10pt}\textbf{dolrem} } \textit{8 of winter/december 8} (N)\\
\noindent {\tovian \fontsize{20pt}{10pt} \textbf{dolrem} }\\
\noindent /d{\textprimstress}olrem/\\
\noindent lit. death+flower\\
\noindent \textit{Dappledviolet}\\


\noindent History:

\vspace{-0pt}
\hspace{40pt}
\begin{tabular}{ccc}
\textit{6000} & /doldremp/&$\rightarrow$ & \textit{8750} & /dolrem/& \\
\end{tabular}

\vspace{20pt}\hline

\end{nopagebreak}
\filbreak



\vspace{15pt}
\begin{nopagebreak}
\noindent{\fontsize{20pt}{10pt}\textbf{nathan} } \textit{80 of autumn/November 20} (N)\\
\noindent {\tovian \fontsize{20pt}{10pt} \textbf{nathan} }\\
\noindent /n{\textprimstress}a{\texttheta}an/\\
\noindent lit. light+tool\\
\noindent \textit{Glemfork}\\


\noindent History:

\vspace{-0pt}
\hspace{40pt}
\begin{tabular}{ccc}
\textit{6000} & /ana{\dh}an/&$\rightarrow$ & \textit{7500} & /na{\dh}an/&$\rightarrow$ & \textit{12002} & /na{\texttheta}an/& \\
\end{tabular}

\vspace{20pt}\hline

\end{nopagebreak}
\filbreak



\vspace{15pt}
\begin{nopagebreak}
\noindent{\fontsize{20pt}{10pt}\textbf{shotlrleth} } \textit{80 of spring/May 20} (N)\\
\noindent {\tovian \fontsize{20pt}{10pt} \textbf{shotlrleth} }\\
\noindent /{\textesh}{\textprimstress}o{\texttoptiebar{t\textbeltl}}rle{\texttheta}/\\
\noindent lit. flowing+hoe\\
\noindent \textit{Mistsickle}\\


\noindent History:

\vspace{-0pt}
\hspace{40pt}
\begin{tabular}{ccc}
\textit{6000} & /{\textesh}o{\texttoptiebar{t\textbeltl}}mllre{\dh}n/&$\rightarrow$ & \textit{6500} & /{\textesh}o{\texttoptiebar{t\textbeltl}}mlre{\dh}n/&$\rightarrow$ & \textit{6503} & /{\textesh}o{\texttoptiebar{t\textbeltl}}mrle{\dh}n/&$\rightarrow$ & \textit{7600} & /{\textesh}o{\texttoptiebar{t\textbeltl}}rle{\dh}n/&$\rightarrow$ & \textit{12002} & /{\textesh}o{\texttoptiebar{t\textbeltl}}rle{\texttheta}n/&$\rightarrow$ & \textit{13001} & /{\textesh}o{\texttoptiebar{t\textbeltl}}rle{\texttheta}/& \\
\end{tabular}

\vspace{20pt}\hline

\end{nopagebreak}
\filbreak



\vspace{15pt}
\begin{nopagebreak}
\noindent{\fontsize{20pt}{10pt}\textbf{thangethan} } \textit{80 of summer/August 20} (N)\\
\noindent {\tovian \fontsize{20pt}{10pt} \textbf{thangethan} }\\
\noindent /{\texttheta}a{\ng}{\textprimstress}e{\texttheta}an/\\
\noindent lit. leaf+tool\\
\noindent \textit{Leafcut}\\


\noindent History:

\vspace{-0pt}
\hspace{40pt}
\begin{tabular}{ccc}
\textit{6000} & /a{\dh}a{\ng}ge{\dh}an/&$\rightarrow$ & \textit{7500} & /{\dh}a{\ng}ge{\dh}an/&$\rightarrow$ & \textit{8750} & /{\dh}a{\ng}e{\dh}an/&$\rightarrow$ & \textit{12002} & /{\texttheta}a{\ng}e{\texttheta}an/& \\
\end{tabular}

\vspace{20pt}\hline

\end{nopagebreak}
\filbreak



\vspace{15pt}
\begin{nopagebreak}
\noindent{\fontsize{20pt}{10pt}\textbf{mrathan} } \textit{80 of winter/February 20} (N)\\
\noindent {\tovian \fontsize{20pt}{10pt} \textbf{mrathan} }\\
\noindent /mr{\textprimstress}a{\texttheta}an/\\
\noindent lit. darkness+tool\\
\noindent \textit{Nighbasket}\\


\noindent History:

\vspace{-0pt}
\hspace{40pt}
\begin{tabular}{ccc}
\textit{6000} & /bral{\dh}an/&$\rightarrow$ & \textit{11000} & /mral{\dh}an/&$\rightarrow$ & \textit{12002} & /mral{\texttheta}an/&$\rightarrow$ & \textit{12502} & /mra{\texttheta}an/& \\
\end{tabular}

\vspace{20pt}\hline

\end{nopagebreak}
\filbreak



\vspace{15pt}
\begin{nopagebreak}
\noindent{\fontsize{20pt}{10pt}\textbf{???} } \textit{81 of autumn/November 21} (N)\\
\noindent {\tovian \fontsize{20pt}{10pt} \textbf{???} }\\
\noindent /???/\\
\noindent lit. ???\\
\noindent \textit{Skoble}\\


\noindent History:

\vspace{-0pt}
\hspace{40pt}
\begin{tabular}{ccc}
\textit{6000} & /???/& \\
\end{tabular}

\vspace{20pt}\hline

\end{nopagebreak}
\filbreak



\vspace{15pt}
\begin{nopagebreak}
\noindent{\fontsize{20pt}{10pt}\textbf{sethok} } \textit{81 of spring/May 21} (N)\\
\noindent {\tovian \fontsize{20pt}{10pt} \textbf{sethok} }\\
\noindent /s{\textprimstress}e{\texttheta}ok/\\
\noindent lit. water+down\\
\noindent \textit{Raincress}\\


\noindent History:

\vspace{-0pt}
\hspace{40pt}
\begin{tabular}{ccc}
\textit{6000} & /zel{\texttheta}ok/&$\rightarrow$ & \textit{12000} & /sel{\texttheta}ok/&$\rightarrow$ & \textit{12502} & /se{\texttheta}ok/& \\
\end{tabular}

\vspace{20pt}\hline

\end{nopagebreak}
\filbreak



\vspace{15pt}
\begin{nopagebreak}
\noindent{\fontsize{20pt}{10pt}\textbf{lhenashotlodrem} } \textit{81 of summer/August 21} (N)\\
\noindent {\tovian \fontsize{20pt}{10pt} \textbf{lhenashotlodrem} }\\
\noindent /{\textbeltl}ena{\textesh}o{\texttoptiebar{t\textbeltl}}{\textprimstress}odrem/\\
\noindent lit. dream+flowing+flower\\
\noindent \textit{Tideflower}\\


\noindent History:

\vspace{-0pt}
\hspace{40pt}
\begin{tabular}{ccc}
\textit{6000} & /{\textbeltl}ena{\textesh}o{\texttoptiebar{t\textbeltl}}dremp/&$\rightarrow$ & \textit{8750} & /{\textbeltl}ena{\textesh}o{\texttoptiebar{t\textbeltl}}drem/&$\rightarrow$ & \textit{9300} & /{\textbeltl}ena{\textesh}o{\texttoptiebar{t\textbeltl}}odrem/& \\
\end{tabular}

\vspace{20pt}\hline

\end{nopagebreak}
\filbreak



\vspace{15pt}
\begin{nopagebreak}
\noindent{\fontsize{20pt}{10pt}\textbf{dremen} } \textit{81 of winter/February 1} (N)\\
\noindent {\tovian \fontsize{20pt}{10pt} \textbf{dremen} }\\
\noindent /dr{\textprimstress}emen/\\
\noindent lit. flower+stone\\
\noindent \textit{Dremgypsum}\\


\noindent History:

\vspace{-0pt}
\hspace{40pt}
\begin{tabular}{ccc}
\textit{6000} & /drempmern/&$\rightarrow$ & \textit{7600} & /drempmen/&$\rightarrow$ & \textit{8500} & /dremmen/&$\rightarrow$ & \textit{12007} & /dremen/& \\
\end{tabular}

\vspace{20pt}\hline

\end{nopagebreak}
\filbreak



\vspace{15pt}
\begin{nopagebreak}
\noindent{\fontsize{20pt}{10pt}\textbf{???} } \textit{82 of autumn/November 22} (N)\\
\noindent {\tovian \fontsize{20pt}{10pt} \textbf{???} }\\
\noindent /???/\\
\noindent lit. ???\\
\noindent \textit{Muncal}\\


\noindent History:

\vspace{-0pt}
\hspace{40pt}
\begin{tabular}{ccc}
\textit{6000} & /???/& \\
\end{tabular}

\vspace{20pt}\hline

\end{nopagebreak}
\filbreak



\vspace{15pt}
\begin{nopagebreak}
\noindent{\fontsize{20pt}{10pt}\textbf{shotlom} } \textit{82 of spring/May 22} (N)\\
\noindent {\tovian \fontsize{20pt}{10pt} \textbf{shotlom} }\\
\noindent /{\textesh}{\textprimstress}o{\texttoptiebar{t\textbeltl}}om/\\
\noindent lit. flowing+bloom\\
\noindent \textit{Mistbloom}\\


\noindent History:

\vspace{-0pt}
\hspace{40pt}
\begin{tabular}{ccc}
\textit{6000} & /{\textesh}o{\texttoptiebar{t\textbeltl}}drmpfede/&$\rightarrow$ & \textit{7600} & /{\textesh}o{\texttoptiebar{t\textbeltl}}dmpfede/&$\rightarrow$ & \textit{8500} & /{\textesh}o{\texttoptiebar{t\textbeltl}}dmfede/&$\rightarrow$ & \textit{9300} & /{\textesh}o{\texttoptiebar{t\textbeltl}}odmfede/&$\rightarrow$ & \textit{9500} & /{\textesh}o{\texttoptiebar{t\textbeltl}}odmfed/&$\rightarrow$ & \textit{12003} & /{\textesh}o{\texttoptiebar{t\textbeltl}}odmd/&$\rightarrow$ & \textit{13001} & /{\textesh}o{\texttoptiebar{t\textbeltl}}odm/&$\rightarrow$ & \textit{13004} & /{\textesh}o{\texttoptiebar{t\textbeltl}}om/& \\
\end{tabular}

\vspace{20pt}\hline

\end{nopagebreak}
\filbreak



\vspace{15pt}
\begin{nopagebreak}
\noindent{\fontsize{20pt}{10pt}\textbf{munorhanatesh} } \textit{82 of summer/August 22} (N)\\
\noindent {\tovian \fontsize{20pt}{10pt} \textbf{munorhanatesh} }\\
\noindent /munorħan{\textprimstress}a{\textsubbridge{t}}e{\textesh}/\\
\noindent lit. dusk+sound\\
\noindent \textit{Duskbell}\\


\noindent History:

\vspace{-0pt}
\hspace{40pt}
\begin{tabular}{ccc}
\textit{6000} & /pnorhana{\textsubbridge{t}}e{\textyogh}/&$\rightarrow$ & \textit{11000} & /mnorhana{\textsubbridge{t}}e{\textyogh}/&$\rightarrow$ & \textit{11996} & /munorhana{\textsubbridge{t}}e{\textyogh}/&$\rightarrow$ & \textit{12001} & /munorhana{\textsubbridge{t}}e{\textesh}/&$\rightarrow$ & \textit{14006} & /munorħana{\textsubbridge{t}}e{\textesh}/& \\
\end{tabular}

\vspace{20pt}\hline

\end{nopagebreak}
\filbreak



\vspace{15pt}
\begin{nopagebreak}
\noindent{\fontsize{20pt}{10pt}\textbf{syalethanan} } \textit{82 of winter/February 22} (N)\\
\noindent {\tovian \fontsize{20pt}{10pt} \textbf{syalethanan} }\\
\noindent /syale{\texttheta}{\textprimstress}anan/\\
\noindent lit. food+star\\
\noindent \textit{Starssalt}\\


\noindent History:

\vspace{-0pt}
\hspace{40pt}
\begin{tabular}{ccc}
\textit{6000} & /syale{\dh}annaana/&$\rightarrow$ & \textit{6500} & /syale{\dh}anaana/&$\rightarrow$ & \textit{9500} & /syale{\dh}anaan/&$\rightarrow$ & \textit{12002} & /syale{\texttheta}anaan/&$\rightarrow$ & \textit{12005} & /syale{\texttheta}anan/& \\
\end{tabular}

\vspace{20pt}\hline

\end{nopagebreak}
\filbreak



\vspace{15pt}
\begin{nopagebreak}
\noindent{\fontsize{20pt}{10pt}\textbf{holamel} } \textit{83 of autumn/November 23} (N)\\
\noindent {\tovian \fontsize{20pt}{10pt} \textbf{holamel} }\\
\noindent /ħol{\textprimstress}amel/\\
\noindent lit. sweet+life\\
\noindent \textit{Sweevin}\\


\noindent History:

\vspace{-0pt}
\hspace{40pt}
\begin{tabular}{ccc}
\textit{6000} & /holamel/&$\rightarrow$ & \textit{14006} & /ħolamel/& \\
\end{tabular}

\vspace{20pt}\hline

\end{nopagebreak}
\filbreak



\vspace{15pt}
\begin{nopagebreak}
\noindent{\fontsize{20pt}{10pt}\textbf{mloanalor} } \textit{83 of spring/May 23} (N)\\
\noindent {\tovian \fontsize{20pt}{10pt} \textbf{mloanalor} }\\
\noindent /mloan{\textprimstress}alor/\\
\noindent lit. morning+plant\\
\noindent \textit{Dawngrass}\\


\noindent History:

\vspace{-0pt}
\hspace{40pt}
\begin{tabular}{ccc}
\textit{6000} & /ploanamllore/&$\rightarrow$ & \textit{6500} & /ploanamlore/&$\rightarrow$ & \textit{7600} & /ploanalore/&$\rightarrow$ & \textit{9500} & /ploanalor/&$\rightarrow$ & \textit{11000} & /mloanalor/& \\
\end{tabular}

\vspace{20pt}\hline

\end{nopagebreak}
\filbreak



\vspace{15pt}
\begin{nopagebreak}
\noindent{\fontsize{20pt}{10pt}\textbf{yelan} } \textit{83 of summer/August 23} (N)\\
\noindent {\tovian \fontsize{20pt}{10pt} \textbf{yelan} }\\
\noindent /y{\textprimstress}elan/\\
\noindent lit. person+light\\
\noindent \textit{Selitha}\\


\noindent History:

\vspace{-0pt}
\hspace{40pt}
\begin{tabular}{ccc}
\textit{6000} & /dryelaana/&$\rightarrow$ & \textit{7501} & /ryelaana/&$\rightarrow$ & \textit{7600} & /yelaana/&$\rightarrow$ & \textit{9500} & /yelaan/&$\rightarrow$ & \textit{12005} & /yelan/& \\
\end{tabular}

\vspace{20pt}\hline

\end{nopagebreak}
\filbreak



\vspace{15pt}
\begin{nopagebreak}
\noindent{\fontsize{20pt}{10pt}\textbf{aymen} } \textit{83 of winter/February 23} (N)\\
\noindent {\tovian \fontsize{20pt}{10pt} \textbf{aymen} }\\
\noindent /{\textprimstress}aymen/\\
\noindent lit. moon+stone\\
\noindent \textit{Mooniron}\\


\noindent History:

\vspace{-0pt}
\hspace{40pt}
\begin{tabular}{ccc}
\textit{6000} & /aliysmern/&$\rightarrow$ & \textit{6300} & /alysmern/&$\rightarrow$ & \textit{7600} & /aysmen/&$\rightarrow$ & \textit{13002} & /aymen/& \\
\end{tabular}

\vspace{20pt}\hline

\end{nopagebreak}
\filbreak



\vspace{15pt}
\begin{nopagebreak}
\noindent{\fontsize{20pt}{10pt}\textbf{rithdrem} } \textit{84 of autumn/November 24} (N)\\
\noindent {\tovian \fontsize{20pt}{10pt} \textbf{rithdrem} }\\
\noindent /r{\textprimstress}i{\texttheta}drem/\\
\noindent lit. air+flower\\
\noindent \textit{Frothfloer}\\


\noindent History:

\vspace{-0pt}
\hspace{40pt}
\begin{tabular}{ccc}
\textit{6000} & /ri{\texttheta}dremp/&$\rightarrow$ & \textit{8750} & /ri{\texttheta}drem/& \\
\end{tabular}

\vspace{20pt}\hline

\end{nopagebreak}
\filbreak



\vspace{15pt}
\begin{nopagebreak}
\noindent{\fontsize{20pt}{10pt}\textbf{selumlor} } \textit{84 of spring/May 24} (N)\\
\noindent {\tovian \fontsize{20pt}{10pt} \textbf{selumlor} }\\
\noindent /sel{\textprimstress}umlor/\\
\noindent lit. water+vine\\
\noindent \textit{Rainvin}\\


\noindent History:

\vspace{-0pt}
\hspace{40pt}
\begin{tabular}{ccc}
\textit{6000} & /zelkluwmllore/&$\rightarrow$ & \textit{6500} & /zelkluwmlore/&$\rightarrow$ & \textit{7600} & /zelklumlore/&$\rightarrow$ & \textit{8750} & /zellumlore/&$\rightarrow$ & \textit{9500} & /zellumlor/&$\rightarrow$ & \textit{12000} & /sellumlor/&$\rightarrow$ & \textit{12007} & /selumlor/& \\
\end{tabular}

\vspace{20pt}\hline

\end{nopagebreak}
\filbreak



\vspace{15pt}
\begin{nopagebreak}
\noindent{\fontsize{20pt}{10pt}\textbf{lhenaglar} } \textit{84 of summer/August 24} (N)\\
\noindent {\tovian \fontsize{20pt}{10pt} \textbf{lhenaglar} }\\
\noindent /{\textbeltl}en{\textprimstress}aglar/\\
\noindent lit. dream+tree\\
\noindent \textit{Dremroot}\\


\noindent History:

\vspace{-0pt}
\hspace{40pt}
\begin{tabular}{ccc}
\textit{6000} & /{\textbeltl}enaaglare/&$\rightarrow$ & \textit{9500} & /{\textbeltl}enaaglar/&$\rightarrow$ & \textit{12005} & /{\textbeltl}enaglar/& \\
\end{tabular}

\vspace{20pt}\hline

\end{nopagebreak}
\filbreak



\vspace{15pt}
\begin{nopagebreak}
\noindent{\fontsize{20pt}{10pt}\textbf{mloanamen} } \textit{84 of winter/February 24} (N)\\
\noindent {\tovian \fontsize{20pt}{10pt} \textbf{mloanamen} }\\
\noindent /mloan{\textprimstress}amen/\\
\noindent lit. morning+metal\\
\noindent \textit{Dawnbrass}\\


\noindent History:

\vspace{-0pt}
\hspace{40pt}
\begin{tabular}{ccc}
\textit{6000} & /ploana{\dh}nhmern/&$\rightarrow$ & \textit{6400} & /ploana{\dh}hmern/&$\rightarrow$ & \textit{7600} & /ploana{\dh}hmen/&$\rightarrow$ & \textit{11000} & /mloana{\dh}hmen/&$\rightarrow$ & \textit{12002} & /mloana{\texttheta}hmen/&$\rightarrow$ & \textit{13000} & /mloana{\texttheta}men/&$\rightarrow$ & \textit{13002} & /mloanamen/& \\
\end{tabular}

\vspace{20pt}\hline

\end{nopagebreak}
\filbreak



\vspace{15pt}
\begin{nopagebreak}
\noindent{\fontsize{20pt}{10pt}\textbf{lhenarith} } \textit{85 of autumn/November 25} (N)\\
\noindent {\tovian \fontsize{20pt}{10pt} \textbf{lhenarith} }\\
\noindent /{\textbeltl}en{\textprimstress}ari{\texttheta}/\\
\noindent lit. dream+air\\
\noindent \textit{Depswe}\\


\noindent History:

\vspace{-0pt}
\hspace{40pt}
\begin{tabular}{ccc}
\textit{6000} & /{\textbeltl}enari{\texttheta}/& \\
\end{tabular}

\vspace{20pt}\hline

\end{nopagebreak}
\filbreak



\vspace{15pt}
\begin{nopagebreak}
\noindent{\fontsize{20pt}{10pt}\textbf{nanatlan} } \textit{85 of spring/May 25} (N)\\
\noindent {\tovian \fontsize{20pt}{10pt} \textbf{nanatlan} }\\
\noindent /nan{\textprimstress}a{\texttoptiebar{t\textbeltl}}an/\\
\noindent lit. star+animal\\
\noindent \textit{Starhare}\\


\noindent History:

\vspace{-0pt}
\hspace{40pt}
\begin{tabular}{ccc}
\textit{6000} & /nnaana{\textsubbridge{t}}lamn/&$\rightarrow$ & \textit{6401} & /nnaana{\texttoptiebar{t\textbeltl}}amn/&$\rightarrow$ & \textit{6500} & /naana{\texttoptiebar{t\textbeltl}}amn/&$\rightarrow$ & \textit{7600} & /naana{\texttoptiebar{t\textbeltl}}an/&$\rightarrow$ & \textit{12005} & /nana{\texttoptiebar{t\textbeltl}}an/& \\
\end{tabular}

\vspace{20pt}\hline

\end{nopagebreak}
\filbreak



\vspace{15pt}
\begin{nopagebreak}
\noindent{\fontsize{20pt}{10pt}\textbf{yselom} } \textit{85 of summer/August 25} (N)\\
\noindent {\tovian \fontsize{20pt}{10pt} \textbf{yselom} }\\
\noindent /ys{\textprimstress}elom/\\
\noindent lit. moon+memory\\
\noindent \textit{Tlunalorum}\\


\noindent History:

\vspace{-0pt}
\hspace{40pt}
\begin{tabular}{ccc}
\textit{6000} & /aliyselorm/&$\rightarrow$ & \textit{6300} & /alyselorm/&$\rightarrow$ & \textit{7500} & /lyselorm/&$\rightarrow$ & \textit{7501} & /yselorm/&$\rightarrow$ & \textit{7600} & /yselom/& \\
\end{tabular}

\vspace{20pt}\hline

\end{nopagebreak}
\filbreak



\vspace{15pt}
\begin{nopagebreak}
\noindent{\fontsize{20pt}{10pt}\textbf{nanatan} } \textit{85 of winter/February 25} (N)\\
\noindent {\tovian \fontsize{20pt}{10pt} \textbf{nanatan} }\\
\noindent /nan{\textprimstress}a{\textsubbridge{t}}an/\\
\noindent lit. star+face\\
\noindent \textit{Starwhisker}\\


\noindent History:

\vspace{-0pt}
\hspace{40pt}
\begin{tabular}{ccc}
\textit{6000} & /nnaana{\textsubbridge{t}}arn/&$\rightarrow$ & \textit{6500} & /naana{\textsubbridge{t}}arn/&$\rightarrow$ & \textit{7600} & /naana{\textsubbridge{t}}an/&$\rightarrow$ & \textit{12005} & /nana{\textsubbridge{t}}an/& \\
\end{tabular}

\vspace{20pt}\hline

\end{nopagebreak}
\filbreak



\vspace{15pt}
\begin{nopagebreak}
\noindent{\fontsize{20pt}{10pt}\textbf{nunasil} } \textit{86 of autumn/November 26} (N)\\
\noindent {\tovian \fontsize{20pt}{10pt} \textbf{nunasil} }\\
\noindent /nun{\textprimstress}asil/\\
\noindent lit. small+existence\\
\noindent \textit{Niwhis}\\


\noindent History:

\vspace{-0pt}
\hspace{40pt}
\begin{tabular}{ccc}
\textit{6000} & /ninasil/&$\rightarrow$ & \textit{9200} & /n{\textschwa}nasil/&$\rightarrow$ & \textit{11990} & /nnasil/&$\rightarrow$ & \textit{11996} & /nunasil/& \\
\end{tabular}

\vspace{20pt}\hline

\end{nopagebreak}
\filbreak



\vspace{15pt}
\begin{nopagebreak}
\noindent{\fontsize{20pt}{10pt}\textbf{dolor} } \textit{86 of spring/May 26} (N)\\
\noindent {\tovian \fontsize{20pt}{10pt} \textbf{dolor} }\\
\noindent /d{\textprimstress}olor/\\
\noindent lit. death+blue\\
\noindent \textit{Twiliwoad}\\


\noindent History:

\vspace{-0pt}
\hspace{40pt}
\begin{tabular}{ccc}
\textit{6000} & /doldora/&$\rightarrow$ & \textit{8750} & /dolora/&$\rightarrow$ & \textit{9500} & /dolor/& \\
\end{tabular}

\vspace{20pt}\hline

\end{nopagebreak}
\filbreak



\vspace{15pt}
\begin{nopagebreak}
\noindent{\fontsize{20pt}{10pt}\textbf{keshaglar} } \textit{86 of summer/August 26} (N)\\
\noindent {\tovian \fontsize{20pt}{10pt} \textbf{keshaglar} }\\
\noindent /ke{\textesh}{\textprimstress}aglar/\\
\noindent lit. fire+tree\\
\noindent \textit{Charwood}\\


\noindent History:

\vspace{-0pt}
\hspace{40pt}
\begin{tabular}{ccc}
\textit{6000} & /ke{\textesh}aglare/&$\rightarrow$ & \textit{9500} & /ke{\textesh}aglar/& \\
\end{tabular}

\vspace{20pt}\hline

\end{nopagebreak}
\filbreak



\vspace{15pt}
\begin{nopagebreak}
\noindent{\fontsize{20pt}{10pt}\textbf{lanamen} } \textit{86 of winter/February 26} (N)\\
\noindent {\tovian \fontsize{20pt}{10pt} \textbf{lanamen} }\\
\noindent /lan{\textprimstress}amen/\\
\noindent lit. twilight+metal\\
\noindent \textit{Twilighttin}\\


\noindent History:

\vspace{-0pt}
\hspace{40pt}
\begin{tabular}{ccc}
\textit{6000} & /brlana{\dh}nhmern/&$\rightarrow$ & \textit{6400} & /brlana{\dh}hmern/&$\rightarrow$ & \textit{7501} & /rlana{\dh}hmern/&$\rightarrow$ & \textit{7600} & /lana{\dh}hmen/&$\rightarrow$ & \textit{12002} & /lana{\texttheta}hmen/&$\rightarrow$ & \textit{13000} & /lana{\texttheta}men/&$\rightarrow$ & \textit{13002} & /lanamen/& \\
\end{tabular}

\vspace{20pt}\hline

\end{nopagebreak}
\filbreak



\vspace{15pt}
\begin{nopagebreak}
\noindent{\fontsize{20pt}{10pt}\textbf{ysen} } \textit{87 of autumn/November 27} (N)\\
\noindent {\tovian \fontsize{20pt}{10pt} \textbf{ysen} }\\
\noindent /ys{\textprimstress}en/\\
\noindent lit. moon+dream\\
\noindent \textit{Munverbena}\\


\noindent History:

\vspace{-0pt}
\hspace{40pt}
\begin{tabular}{ccc}
\textit{6000} & /aliys{\textbeltl}ena/&$\rightarrow$ & \textit{6300} & /alys{\textbeltl}ena/&$\rightarrow$ & \textit{7500} & /lys{\textbeltl}ena/&$\rightarrow$ & \textit{7501} & /ys{\textbeltl}ena/&$\rightarrow$ & \textit{9500} & /ys{\textbeltl}en/&$\rightarrow$ & \textit{13000} & /ysen/& \\
\end{tabular}

\vspace{20pt}\hline

\end{nopagebreak}
\filbreak



\vspace{15pt}
\begin{nopagebreak}
\noindent{\fontsize{20pt}{10pt}\textbf{yskel} } \textit{87 of spring/May 27} (N)\\
\noindent {\tovian \fontsize{20pt}{10pt} \textbf{yskel} }\\
\noindent /ysk{\textprimstress}el/\\
\noindent lit. moon+branch\\
\noindent \textit{Moonhazel}\\


\noindent History:

\vspace{-0pt}
\hspace{40pt}
\begin{tabular}{ccc}
\textit{6000} & /aliyskelu/&$\rightarrow$ & \textit{6300} & /alyskelu/&$\rightarrow$ & \textit{7500} & /lyskelu/&$\rightarrow$ & \textit{7501} & /yskelu/&$\rightarrow$ & \textit{9500} & /yskel/& \\
\end{tabular}

\vspace{20pt}\hline

\end{nopagebreak}
\filbreak



\vspace{15pt}
\begin{nopagebreak}
\noindent{\fontsize{20pt}{10pt}\textbf{rithenalor} } \textit{87 of summer/August 27} (N)\\
\noindent {\tovian \fontsize{20pt}{10pt} \textbf{rithenalor} }\\
\noindent /ri{\texttheta}en{\textprimstress}alor/\\
\noindent lit. air+dream+plant\\
\noindent \textit{Zephyrherb}\\


\noindent History:

\vspace{-0pt}
\hspace{40pt}
\begin{tabular}{ccc}
\textit{6000} & /ri{\texttheta}{\textbeltl}enamllore/&$\rightarrow$ & \textit{6500} & /ri{\texttheta}{\textbeltl}enamlore/&$\rightarrow$ & \textit{7600} & /ri{\texttheta}{\textbeltl}enalore/&$\rightarrow$ & \textit{9500} & /ri{\texttheta}{\textbeltl}enalor/&$\rightarrow$ & \textit{13000} & /ri{\texttheta}enalor/& \\
\end{tabular}

\vspace{20pt}\hline

\end{nopagebreak}
\filbreak



\vspace{15pt}
\begin{nopagebreak}
\noindent{\fontsize{20pt}{10pt}\textbf{mramen} } \textit{87 of winter/February 27} (N)\\
\noindent {\tovian \fontsize{20pt}{10pt} \textbf{mramen} }\\
\noindent /mr{\textprimstress}amen/\\
\noindent lit. darkness+metal\\
\noindent \textit{Shadowlead}\\


\noindent History:

\vspace{-0pt}
\hspace{40pt}
\begin{tabular}{ccc}
\textit{6000} & /bral{\dh}nhmern/&$\rightarrow$ & \textit{6400} & /bral{\dh}hmern/&$\rightarrow$ & \textit{7600} & /bral{\dh}hmen/&$\rightarrow$ & \textit{11000} & /mral{\dh}hmen/&$\rightarrow$ & \textit{12002} & /mral{\texttheta}hmen/&$\rightarrow$ & \textit{12502} & /mra{\texttheta}hmen/&$\rightarrow$ & \textit{13000} & /mra{\texttheta}men/&$\rightarrow$ & \textit{13002} & /mramen/& \\
\end{tabular}

\vspace{20pt}\hline

\end{nopagebreak}
\filbreak



\vspace{15pt}
\begin{nopagebreak}
\noindent{\fontsize{20pt}{10pt}\textbf{engothr} } \textit{88 of autumn/November 28} (N)\\
\noindent {\tovian \fontsize{20pt}{10pt} \textbf{engothr} }\\
\noindent /{\textprimstress}e{\ng}o{\texttheta}r/\\
\noindent lit. sun+time\\
\noindent \textit{Sunethime}\\


\noindent History:

\vspace{-0pt}
\hspace{40pt}
\begin{tabular}{ccc}
\textit{6000} & /efl{\ng}go{\texttheta}r/&$\rightarrow$ & \textit{7600} & /ef{\ng}go{\texttheta}r/&$\rightarrow$ & \textit{8750} & /ef{\ng}o{\texttheta}r/&$\rightarrow$ & \textit{13002} & /e{\ng}o{\texttheta}r/& \\
\end{tabular}

\vspace{20pt}\hline

\end{nopagebreak}
\filbreak



\vspace{15pt}
\begin{nopagebreak}
\noindent{\fontsize{20pt}{10pt}\textbf{mloanakel} } \textit{88 of spring/May 28} (N)\\
\noindent {\tovian \fontsize{20pt}{10pt} \textbf{mloanakel} }\\
\noindent /mloan{\textprimstress}akel/\\
\noindent lit. dawn+branch\\
\noindent \textit{Duslower}\\


\noindent History:

\vspace{-0pt}
\hspace{40pt}
\begin{tabular}{ccc}
\textit{6000} & /ploanakelu/&$\rightarrow$ & \textit{9500} & /ploanakel/&$\rightarrow$ & \textit{11000} & /mloanakel/& \\
\end{tabular}

\vspace{20pt}\hline

\end{nopagebreak}
\filbreak



\vspace{15pt}
\begin{nopagebreak}
\noindent{\fontsize{20pt}{10pt}\textbf{mralor} } \textit{88 of summer/August 28} (N)\\
\noindent {\tovian \fontsize{20pt}{10pt} \textbf{mralor} }\\
\noindent /mr{\textprimstress}alor/\\
\noindent lit. darkness+blue\\
\noindent \textit{Nightblue}\\


\noindent History:

\vspace{-0pt}
\hspace{40pt}
\begin{tabular}{ccc}
\textit{6000} & /braldora/&$\rightarrow$ & \textit{8750} & /bralora/&$\rightarrow$ & \textit{9500} & /bralor/&$\rightarrow$ & \textit{11000} & /mralor/& \\
\end{tabular}

\vspace{20pt}\hline

\end{nopagebreak}
\filbreak



\vspace{15pt}
\begin{nopagebreak}
\noindent{\fontsize{20pt}{10pt}\textbf{lan} } \textit{88 of winter/February 28} (N)\\
\noindent {\tovian \fontsize{20pt}{10pt} \textbf{lan} }\\
\noindent /l{\textprimstress}an/\\
\noindent lit. twilight+???\\
\noindent \textit{Duskin}\\


\noindent History:

\vspace{-0pt}
\hspace{40pt}
\begin{tabular}{ccc}
\textit{6000} & /brlana???/&$\rightarrow$ & \textit{7501} & /rlana???/&$\rightarrow$ & \textit{7600} & /lana???/&$\rightarrow$ & \textit{8500} & /lana/&$\rightarrow$ & \textit{9500} & /lan/& \\
\end{tabular}

\vspace{20pt}\hline

\end{nopagebreak}
\filbreak



\vspace{15pt}
\begin{nopagebreak}
\noindent{\fontsize{20pt}{10pt}\textbf{haradrem} } \textit{89 of autumn/November 29} (N)\\
\noindent {\tovian \fontsize{20pt}{10pt} \textbf{haradrem} }\\
\noindent /ħar{\textprimstress}adrem/\\
\noindent lit. pain+flower\\
\noindent \textit{Harpetal}\\


\noindent History:

\vspace{-0pt}
\hspace{40pt}
\begin{tabular}{ccc}
\textit{6000} & /haradremp/&$\rightarrow$ & \textit{8750} & /haradrem/&$\rightarrow$ & \textit{14006} & /ħaradrem/& \\
\end{tabular}

\vspace{20pt}\hline

\end{nopagebreak}
\filbreak



\vspace{15pt}
\begin{nopagebreak}
\noindent{\fontsize{20pt}{10pt}\textbf{shotleon} } \textit{89 of spring/May 29} (N)\\
\noindent {\tovian \fontsize{20pt}{10pt} \textbf{shotleon} }\\
\noindent /{\textesh}o{\texttoptiebar{t\textbeltl}}{\textprimstress}eon/\\
\noindent lit. flowing+mouth\\
\noindent \textit{Starwort}\\


\noindent History:

\vspace{-0pt}
\hspace{40pt}
\begin{tabular}{ccc}
\textit{6000} & /{\textesh}o{\texttoptiebar{t\textbeltl}}eoln/&$\rightarrow$ & \textit{7600} & /{\textesh}o{\texttoptiebar{t\textbeltl}}eon/& \\
\end{tabular}

\vspace{20pt}\hline

\end{nopagebreak}
\filbreak



\vspace{15pt}
\begin{nopagebreak}
\noindent{\fontsize{20pt}{10pt}\textbf{fidol} } \textit{89 of summer/August 29} (N)\\
\noindent {\tovian \fontsize{20pt}{10pt} \textbf{fidol} }\\
\noindent /f{\textprimstress}idol/\\
\noindent lit. sun+GEN+death\\
\noindent \textit{Sunseed}\\


\noindent History:

\vspace{-0pt}
\hspace{40pt}
\begin{tabular}{ccc}
\textit{6000} & /eflhi-dol/&$\rightarrow$ & \textit{6402} & /ef{\textbeltl}i-dol/&$\rightarrow$ & \textit{7500} & /f{\textbeltl}i-dol/&$\rightarrow$ & \textit{8500} & /f{\textbeltl}idol/&$\rightarrow$ & \textit{13000} & /fidol/& \\
\end{tabular}

\vspace{20pt}\hline

\end{nopagebreak}
\filbreak



\vspace{15pt}
\begin{nopagebreak}
\noindent{\fontsize{20pt}{10pt}\textbf{tleshmen} } \textit{89 of winter/February 29} (N)\\
\noindent {\tovian \fontsize{20pt}{10pt} \textbf{tleshmen} }\\
\noindent /{\texttoptiebar{t\textbeltl}}{\textprimstress}e{\textesh}men/\\
\noindent lit. speed+metal\\
\noindent \textit{Quicksilver}\\


\noindent History:

\vspace{-0pt}
\hspace{40pt}
\begin{tabular}{ccc}
\textit{6000} & /{\texttoptiebar{t\textbeltl}}e{\textyogh}{\dh}nhmern/&$\rightarrow$ & \textit{6400} & /{\texttoptiebar{t\textbeltl}}e{\textyogh}{\dh}hmern/&$\rightarrow$ & \textit{7600} & /{\texttoptiebar{t\textbeltl}}e{\textyogh}{\dh}hmen/&$\rightarrow$ & \textit{12001} & /{\texttoptiebar{t\textbeltl}}e{\textesh}{\dh}hmen/&$\rightarrow$ & \textit{12002} & /{\texttoptiebar{t\textbeltl}}e{\textesh}{\texttheta}hmen/&$\rightarrow$ & \textit{13000} & /{\texttoptiebar{t\textbeltl}}e{\textesh}hmen/&$\rightarrow$ & \textit{13002} & /{\texttoptiebar{t\textbeltl}}e{\textesh}men/& \\
\end{tabular}

\vspace{20pt}\hline

\end{nopagebreak}
\filbreak



\vspace{15pt}
\begin{nopagebreak}
\noindent{\fontsize{20pt}{10pt}\textbf{lhenashith} } \textit{9 of autumn/September 9} (N)\\
\noindent {\tovian \fontsize{20pt}{10pt} \textbf{lhenashith} }\\
\noindent /{\textbeltl}en{\textprimstress}a{\textesh}i{\texttheta}/\\
\noindent lit. dream+sharpness\\
\noindent \textit{Dremthorn}\\


\noindent History:

\vspace{-0pt}
\hspace{40pt}
\begin{tabular}{ccc}
\textit{6000} & /{\textbeltl}ena{\textyogh}i{\texttheta}/&$\rightarrow$ & \textit{12001} & /{\textbeltl}ena{\textesh}i{\texttheta}/& \\
\end{tabular}

\vspace{20pt}\hline

\end{nopagebreak}
\filbreak



\vspace{15pt}
\begin{nopagebreak}
\noindent{\fontsize{20pt}{10pt}\textbf{shatom} } \textit{9 of spring/March 9} (N)\\
\noindent {\tovian \fontsize{20pt}{10pt} \textbf{shatom} }\\
\noindent /{\textesh}{\textprimstress}a{\textsubbridge{t}}om/\\
\noindent lit. ice+root\\
\noindent \textit{Frostroot}\\


\noindent History:

\vspace{-0pt}
\hspace{40pt}
\begin{tabular}{ccc}
\textit{6000} & /{\textesh}u{\textesh}a{\textsubbridge{t}}omo/&$\rightarrow$ & \textit{9200} & /{\textesh}{\textschwa}{\textesh}a{\textsubbridge{t}}omo/&$\rightarrow$ & \textit{9500} & /{\textesh}{\textschwa}{\textesh}a{\textsubbridge{t}}om/&$\rightarrow$ & \textit{11990} & /{\textesh}{\textesh}a{\textsubbridge{t}}om/&$\rightarrow$ & \textit{12007} & /{\textesh}a{\textsubbridge{t}}om/& \\
\end{tabular}

\vspace{20pt}\hline

\end{nopagebreak}
\filbreak



\vspace{15pt}
\begin{nopagebreak}
\noindent{\fontsize{20pt}{10pt}\textbf{temen} } \textit{9 of summer/June 9} (N)\\
\noindent {\tovian \fontsize{20pt}{10pt} \textbf{temen} }\\
\noindent /{\textsubbridge{t}}{\textprimstress}emen/\\
\noindent lit. sound+stone\\
\noindent \textit{Whisperald}\\


\noindent History:

\vspace{-0pt}
\hspace{40pt}
\begin{tabular}{ccc}
\textit{6000} & /{\textsubbridge{t}}e{\textyogh}mern/&$\rightarrow$ & \textit{7600} & /{\textsubbridge{t}}e{\textyogh}men/&$\rightarrow$ & \textit{12001} & /{\textsubbridge{t}}e{\textesh}men/&$\rightarrow$ & \textit{13002} & /{\textsubbridge{t}}emen/& \\
\end{tabular}

\vspace{20pt}\hline

\end{nopagebreak}
\filbreak



\vspace{15pt}
\begin{nopagebreak}
\noindent{\fontsize{20pt}{10pt}\textbf{lrithaglar} } \textit{9 of winter/december 9} (N)\\
\noindent {\tovian \fontsize{20pt}{10pt} \textbf{lrithaglar} }\\
\noindent /lri{\texttheta}{\textprimstress}aglar/\\
\noindent lit. breath+tree\\
\noindent \textit{Windtwig}\\


\noindent History:

\vspace{-0pt}
\hspace{40pt}
\begin{tabular}{ccc}
\textit{6000} & /kwlyri{\texttheta}aglare/&$\rightarrow$ & \textit{7501} & /wlyri{\texttheta}aglare/&$\rightarrow$ & \textit{7600} & /lri{\texttheta}aglare/&$\rightarrow$ & \textit{9500} & /lri{\texttheta}aglar/& \\
\end{tabular}

\vspace{20pt}\hline

\end{nopagebreak}
\filbreak



\vspace{15pt}
\begin{nopagebreak}
\noindent{\fontsize{20pt}{10pt}\textbf{meloresyalethalor} } \textit{90 of autumn/November 30} (N)\\
\noindent {\tovian \fontsize{20pt}{10pt} \textbf{meloresyalethalor} }\\
\noindent /meloresyale{\texttheta}{\textprimstress}alor/\\
\noindent lit. grain+plant\\
\noindent \textit{Sorgho}\\


\noindent History:

\vspace{-0pt}
\hspace{40pt}
\begin{tabular}{ccc}
\textit{6000} & /mlloresyale{\dh}amllore/&$\rightarrow$ & \textit{6500} & /mloresyale{\dh}amlore/&$\rightarrow$ & \textit{7502} & /meloresyale{\dh}amlore/&$\rightarrow$ & \textit{7600} & /meloresyale{\dh}alore/&$\rightarrow$ & \textit{9500} & /meloresyale{\dh}alor/&$\rightarrow$ & \textit{12002} & /meloresyale{\texttheta}alor/& \\
\end{tabular}

\vspace{20pt}\hline

\end{nopagebreak}
\filbreak



\vspace{15pt}
\begin{nopagebreak}
\noindent{\fontsize{20pt}{10pt}\textbf{shakwil} } \textit{90 of spring/May 30} (N)\\
\noindent {\tovian \fontsize{20pt}{10pt} \textbf{shakwil} }\\
\noindent /{\textesh}{\textprimstress}akwil/\\
\noindent lit. ice+movement\\
\noindent \textit{Frostslide}\\


\noindent History:

\vspace{-0pt}
\hspace{40pt}
\begin{tabular}{ccc}
\textit{6000} & /{\textesh}u{\textesh}akwil/&$\rightarrow$ & \textit{9200} & /{\textesh}{\textschwa}{\textesh}akwil/&$\rightarrow$ & \textit{11990} & /{\textesh}{\textesh}akwil/&$\rightarrow$ & \textit{12007} & /{\textesh}akwil/& \\
\end{tabular}

\vspace{20pt}\hline

\end{nopagebreak}
\filbreak



\vspace{15pt}
\begin{nopagebreak}
\noindent{\fontsize{20pt}{10pt}\textbf{falnan} } \textit{90 of summer/August 30} (N)\\
\noindent {\tovian \fontsize{20pt}{10pt} \textbf{falnan} }\\
\noindent /f{\textprimstress}alnan/\\
\noindent lit. leader+star\\
\noindent \textit{Shepherd's Star}\\


\noindent History:

\vspace{-0pt}
\hspace{40pt}
\begin{tabular}{ccc}
\textit{6000} & /falrnnaana/&$\rightarrow$ & \textit{6500} & /falrnaana/&$\rightarrow$ & \textit{6503} & /farlnaana/&$\rightarrow$ & \textit{7600} & /falnaana/&$\rightarrow$ & \textit{9500} & /falnaan/&$\rightarrow$ & \textit{12005} & /falnan/& \\
\end{tabular}

\vspace{20pt}\hline

\end{nopagebreak}
\filbreak



\vspace{15pt}
\begin{nopagebreak}
\noindent{\fontsize{20pt}{10pt}\textbf{shathan} } \textit{90 of winter/February 30} (N)\\
\noindent {\tovian \fontsize{20pt}{10pt} \textbf{shathan} }\\
\noindent /{\textesh}{\textprimstress}a{\texttheta}an/\\
\noindent lit. ice+tool\\
\noindent \textit{Frostscreen}\\


\noindent History:

\vspace{-0pt}
\hspace{40pt}
\begin{tabular}{ccc}
\textit{6000} & /{\textesh}u{\textesh}a{\dh}an/&$\rightarrow$ & \textit{9200} & /{\textesh}{\textschwa}{\textesh}a{\dh}an/&$\rightarrow$ & \textit{11990} & /{\textesh}{\textesh}a{\dh}an/&$\rightarrow$ & \textit{12002} & /{\textesh}{\textesh}a{\texttheta}an/&$\rightarrow$ & \textit{12007} & /{\textesh}a{\texttheta}an/& \\
\end{tabular}

\vspace{20pt}\hline

\end{nopagebreak}
\filbreak



\vspace{15pt}
\begin{nopagebreak}
\noindent{\fontsize{20pt}{10pt}\textbf{???} } \textit{???} (PLACEHOLDER)\\
\noindent {\tovian \fontsize{20pt}{10pt} \textbf{???} }\\
\noindent /???/\\


\noindent History:

\vspace{-0pt}
\hspace{40pt}
\begin{tabular}{ccc}
\textit{0} & /???/& \\
\end{tabular}

\vspace{20pt}\hline

\end{nopagebreak}
\filbreak



\vspace{15pt}
\begin{nopagebreak}
\noindent{\fontsize{20pt}{10pt}\textbf{di-} } \textit{ABL/ablative case marker/away from} (CASE)\\
\noindent {\tovian \fontsize{20pt}{10pt} \textbf{di-} }\\
\noindent /d{\textprimstress}i/\\


\noindent History:

\vspace{-0pt}
\hspace{40pt}
\begin{tabular}{ccc}
\textit{0} & /di-/& \\
\end{tabular}

\vspace{20pt}\hline

\end{nopagebreak}
\filbreak



\vspace{15pt}
\begin{nopagebreak}
\noindent{\fontsize{20pt}{10pt}\textbf{le-} } \textit{abstract noun class} (CLASS)\\
\noindent {\tovian \fontsize{20pt}{10pt} \textbf{le-} }\\
\noindent /l{\textprimstress}e/\\


\noindent History:

\vspace{-0pt}
\hspace{40pt}
\begin{tabular}{ccc}
\textit{0} & /le-/& \\
\end{tabular}

\vspace{20pt}\hline

\end{nopagebreak}
\filbreak



\vspace{15pt}
\begin{nopagebreak}
\noindent{\fontsize{20pt}{10pt}\textbf{yu-} } \textit{ACC/accusative case marker} (CASE)\\
\noindent {\tovian \fontsize{20pt}{10pt} \textbf{yu-} }\\
\noindent /y{\textprimstress}u/\\


\noindent History:

\vspace{-0pt}
\hspace{40pt}
\begin{tabular}{ccc}
\textit{0} & /yu-/& \\
\end{tabular}

\vspace{20pt}\hline

\end{nopagebreak}
\filbreak



\vspace{15pt}
\begin{nopagebreak}
\noindent{\fontsize{20pt}{10pt}\textbf{lor} } \textit{action/doing} (N)\\
\noindent {\tovian \fontsize{20pt}{10pt} \textbf{lor} }\\
\noindent /l{\textprimstress}or/\\


\noindent History:

\vspace{-0pt}
\hspace{40pt}
\begin{tabular}{ccc}
\textit{0} & /lore/&$\rightarrow$ & \textit{9500} & /lor/& \\
\end{tabular}

\vspace{20pt}\hline

\end{nopagebreak}
\filbreak



\vspace{15pt}
\begin{nopagebreak}
\noindent{\fontsize{20pt}{10pt}\textbf{thim} } \textit{adventure} (N)\\
\noindent {\tovian \fontsize{20pt}{10pt} \textbf{thim} }\\
\noindent /{\texttheta}{\textprimstress}im/\\
\noindent lit. going+unknown\\


\noindent History:

\vspace{-0pt}
\hspace{40pt}
\begin{tabular}{ccc}
\textit{1000} & /ki{\texttheta}sa{\textesh}imi/&$\rightarrow$ & \textit{1000} & /k{\texttheta}s{\textesh}imi/&$\rightarrow$ & \textit{4502} & /k{\texttheta}{\textesh}imi/&$\rightarrow$ & \textit{7501} & /{\texttheta}{\textesh}imi/&$\rightarrow$ & \textit{9500} & /{\texttheta}{\textesh}im/&$\rightarrow$ & \textit{13000} & /{\texttheta}im/& \\
\end{tabular}

\vspace{20pt}\hline

\end{nopagebreak}
\filbreak



\vspace{15pt}
\begin{nopagebreak}
\noindent{\fontsize{20pt}{10pt}\textbf{rith} } \textit{air} (N)\\
\noindent {\tovian \fontsize{20pt}{10pt} \textbf{rith} }\\
\noindent /r{\textprimstress}i{\texttheta}/\\


\noindent History:

\vspace{-0pt}
\hspace{40pt}
\begin{tabular}{ccc}
\textit{0} & /rai{\texttheta}/&$\rightarrow$ & \textit{5500} & /ri{\texttheta}/& \\
\end{tabular}

\vspace{20pt}\hline

\end{nopagebreak}
\filbreak



\vspace{15pt}
\begin{nopagebreak}
\noindent{\fontsize{20pt}{10pt}\textbf{su-} } \textit{ALL/allative case marker} (CASE)\\
\noindent {\tovian \fontsize{20pt}{10pt} \textbf{su-} }\\
\noindent /s{\textprimstress}u/\\


\noindent History:

\vspace{-0pt}
\hspace{40pt}
\begin{tabular}{ccc}
\textit{0} & /su-/& \\
\end{tabular}

\vspace{20pt}\hline

\end{nopagebreak}
\filbreak



\vspace{15pt}
\begin{nopagebreak}
\noindent{\fontsize{20pt}{10pt}\textbf{nan} } \textit{ancient/old} (N)\\
\noindent {\tovian \fontsize{20pt}{10pt} \textbf{nan} }\\
\noindent /n{\textprimstress}an/\\


\noindent History:

\vspace{-0pt}
\hspace{40pt}
\begin{tabular}{ccc}
\textit{0} & /nan/& \\
\end{tabular}

\vspace{20pt}\hline

\end{nopagebreak}
\filbreak



\vspace{15pt}
\begin{nopagebreak}
\noindent{\fontsize{20pt}{10pt}\textbf{tlan} } \textit{animal/creature/beast} (N)\\
\noindent {\tovian \fontsize{20pt}{10pt} \textbf{tlan} }\\
\noindent /{\texttoptiebar{t\textbeltl}}{\textprimstress}an/\\


\noindent History:

\vspace{-0pt}
\hspace{40pt}
\begin{tabular}{ccc}
\textit{0} & /{\textsubbridge{t}}elaman/&$\rightarrow$ & \textit{3200} & /{\textsubbridge{t}}lamn/&$\rightarrow$ & \textit{6401} & /{\texttoptiebar{t\textbeltl}}amn/&$\rightarrow$ & \textit{7600} & /{\texttoptiebar{t\textbeltl}}an/& \\
\end{tabular}

\vspace{20pt}\hline

\end{nopagebreak}
\filbreak



\vspace{15pt}
\begin{nopagebreak}
\noindent{\fontsize{20pt}{10pt}\textbf{la-} } \textit{animate noun class} (CLASS)\\
\noindent {\tovian \fontsize{20pt}{10pt} \textbf{la-} }\\
\noindent /l{\textprimstress}a/\\


\noindent History:

\vspace{-0pt}
\hspace{40pt}
\begin{tabular}{ccc}
\textit{0} & /la-/& \\
\end{tabular}

\vspace{20pt}\hline

\end{nopagebreak}
\filbreak



\vspace{15pt}
\begin{nopagebreak}
\noindent{\fontsize{20pt}{10pt}\textbf{loelor} } \textit{April} (N)\\
\noindent {\tovian \fontsize{20pt}{10pt} \textbf{loelor} }\\
\noindent /lo{\textprimstress}elor/\\
\noindent lit. two+plant\\
\noindent \textit{Foegvel}\\


\noindent History:

\vspace{-0pt}
\hspace{40pt}
\begin{tabular}{ccc}
\textit{6000} & /loemllore/&$\rightarrow$ & \textit{6500} & /loemlore/&$\rightarrow$ & \textit{7600} & /loelore/&$\rightarrow$ & \textit{9500} & /loelor/& \\
\end{tabular}

\vspace{20pt}\hline

\end{nopagebreak}
\filbreak



\vspace{15pt}
\begin{nopagebreak}
\noindent{\fontsize{20pt}{10pt}\textbf{femthulef} } \textit{August} (N)\\
\noindent {\tovian \fontsize{20pt}{10pt} \textbf{femthulef} }\\
\noindent /fem{\texttheta}{\textprimstress}ulef/\\
\noindent lit. three+height+sun\\
\noindent \textit{Floreart}\\


\noindent History:

\vspace{-0pt}
\hspace{40pt}
\begin{tabular}{ccc}
\textit{6000} & /fem{\dh}ulefl/&$\rightarrow$ & \textit{12002} & /fem{\texttheta}ulefl/&$\rightarrow$ & \textit{12500} & /fem{\texttheta}ulef/& \\
\end{tabular}

\vspace{20pt}\hline

\end{nopagebreak}
\filbreak



\vspace{15pt}
\begin{nopagebreak}
\noindent{\fontsize{20pt}{10pt}\textbf{mlherlesh} } \textit{Autumn} (N)\\
\noindent {\tovian \fontsize{20pt}{10pt} \textbf{mlherlesh} }\\
\noindent /m{\textbeltl}{\textprimstress}erle{\textesh}/\\
\noindent lit. present+harvest\\


\noindent History:

\vspace{-0pt}
\hspace{40pt}
\begin{tabular}{ccc}
\textit{0} & /num{\textbeltl}ele{\textschwa}mellore{\textyogh}i{\texttheta}/&$\rightarrow$ & \textit{2300} & /num{\textbeltl}elemellore{\textyogh}i{\texttheta}/&$\rightarrow$ & \textit{3200} & /nm{\textbeltl}lemllre{\textyogh}{\texttheta}/&$\rightarrow$ & \textit{4502} & /nm{\textbeltl}lemllre{\textyogh}/&$\rightarrow$ & \textit{6500} & /nm{\textbeltl}lemlre{\textyogh}/&$\rightarrow$ & \textit{6503} & /nm{\textbeltl}lemrle{\textyogh}/&$\rightarrow$ & \textit{7501} & /m{\textbeltl}lemrle{\textyogh}/&$\rightarrow$ & \textit{7600} & /m{\textbeltl}lerle{\textyogh}/&$\rightarrow$ & \textit{12001} & /m{\textbeltl}lerle{\textesh}/&$\rightarrow$ & \textit{12500} & /m{\textbeltl}erle{\textesh}/& \\
\end{tabular}

\vspace{20pt}\hline

\end{nopagebreak}
\filbreak



\vspace{15pt}
\begin{nopagebreak}
\noindent{\fontsize{20pt}{10pt}\textbf{kar} } \textit{balance} (N)\\
\noindent {\tovian \fontsize{20pt}{10pt} \textbf{kar} }\\
\noindent /k{\textprimstress}ar/\\


\noindent History:

\vspace{-0pt}
\hspace{40pt}
\begin{tabular}{ccc}
\textit{0} & /kara/&$\rightarrow$ & \textit{9500} & /kar/& \\
\end{tabular}

\vspace{20pt}\hline

\end{nopagebreak}
\filbreak



\vspace{15pt}
\begin{nopagebreak}
\noindent{\fontsize{20pt}{10pt}\textbf{wen} } \textit{beauty} (N)\\
\noindent {\tovian \fontsize{20pt}{10pt} \textbf{wen} }\\
\noindent /w{\textprimstress}en/\\


\noindent History:

\vspace{-0pt}
\hspace{40pt}
\begin{tabular}{ccc}
\textit{0} & /wen/& \\
\end{tabular}

\vspace{20pt}\hline

\end{nopagebreak}
\filbreak



\vspace{15pt}
\begin{nopagebreak}
\noindent{\fontsize{20pt}{10pt}\textbf{kan} } \textit{becoming} (N)\\
\noindent {\tovian \fontsize{20pt}{10pt} \textbf{kan} }\\
\noindent /k{\textprimstress}an/\\
\noindent lit. arriving+change\\


\noindent History:

\vspace{-0pt}
\hspace{40pt}
\begin{tabular}{ccc}
\textit{0} & /ki{\texttheta}sakana/&$\rightarrow$ & \textit{1000} & /k{\texttheta}skana/&$\rightarrow$ & \textit{3000} & /k{\texttheta}sana/&$\rightarrow$ & \textit{4502} & /k{\texttheta}ana/&$\rightarrow$ & \textit{9500} & /k{\texttheta}an/&$\rightarrow$ & \textit{10700} & /k{\textsubbridge{t}}an/&$\rightarrow$ & \textit{11500} & /kan/& \\
\end{tabular}

\vspace{20pt}\hline

\end{nopagebreak}
\filbreak



\vspace{15pt}
\begin{nopagebreak}
\noindent{\fontsize{20pt}{10pt}\textbf{lon} } \textit{beginning/starting} (N)\\
\noindent {\tovian \fontsize{20pt}{10pt} \textbf{lon} }\\
\noindent /l{\textprimstress}on/\\


\noindent History:

\vspace{-0pt}
\hspace{40pt}
\begin{tabular}{ccc}
\textit{0} & /yaloran/&$\rightarrow$ & \textit{3200} & /ylorn/&$\rightarrow$ & \textit{7600} & /lon/& \\
\end{tabular}

\vspace{20pt}\hline

\end{nopagebreak}
\filbreak



\vspace{15pt}
\begin{nopagebreak}
\noindent{\fontsize{20pt}{10pt}\textbf{hul} } \textit{behind} (N)\\
\noindent {\tovian \fontsize{20pt}{10pt} \textbf{hul} }\\
\noindent /ħ{\textprimstress}ul/\\


\noindent History:

\vspace{-0pt}
\hspace{40pt}
\begin{tabular}{ccc}
\textit{0} & /hul/&$\rightarrow$ & \textit{14006} & /ħul/& \\
\end{tabular}

\vspace{20pt}\hline

\end{nopagebreak}
\filbreak



\vspace{15pt}
\begin{nopagebreak}
\noindent{\fontsize{20pt}{10pt}\textbf{sil} } \textit{being/existence/created thing} (N)\\
\noindent {\tovian \fontsize{20pt}{10pt} \textbf{sil} }\\
\noindent /s{\textprimstress}il/\\


\noindent History:

\vspace{-0pt}
\hspace{40pt}
\begin{tabular}{ccc}
\textit{0} & /sil/& \\
\end{tabular}

\vspace{20pt}\hline

\end{nopagebreak}
\filbreak



\vspace{15pt}
\begin{nopagebreak}
\noindent{\fontsize{20pt}{10pt}\textbf{shanin} } \textit{berry} (N)\\
\noindent {\tovian \fontsize{20pt}{10pt} \textbf{shanin} }\\
\noindent /{\textesh}{\textprimstress}anin/\\
\noindent lit. fruit+small\\


\noindent History:

\vspace{-0pt}
\hspace{40pt}
\begin{tabular}{ccc}
\textit{1000} & /{\textyogh}olanina/&$\rightarrow$ & \textit{3200} & /{\textyogh}lanina/&$\rightarrow$ & \textit{9500} & /{\textyogh}lanin/&$\rightarrow$ & \textit{12001} & /{\textesh}lanin/&$\rightarrow$ & \textit{12500} & /{\textesh}anin/& \\
\end{tabular}

\vspace{20pt}\hline

\end{nopagebreak}
\filbreak



\vspace{15pt}
\begin{nopagebreak}
\noindent{\fontsize{20pt}{10pt}\textbf{mot} } \textit{big/large} (N)\\
\noindent {\tovian \fontsize{20pt}{10pt} \textbf{mot} }\\
\noindent /m{\textprimstress}o{\textsubbridge{t}}/\\


\noindent History:

\vspace{-0pt}
\hspace{40pt}
\begin{tabular}{ccc}
\textit{0} & /mo{\textsubbridge{t}}/& \\
\end{tabular}

\vspace{20pt}\hline

\end{nopagebreak}
\filbreak



\vspace{15pt}
\begin{nopagebreak}
\noindent{\fontsize{20pt}{10pt}\textbf{rithan} } \textit{bird} (N)\\
\noindent {\tovian \fontsize{20pt}{10pt} \textbf{rithan} }\\
\noindent /r{\textprimstress}i{\texttheta}an/\\
\noindent lit. air+animal\\


\noindent History:

\vspace{-0pt}
\hspace{40pt}
\begin{tabular}{ccc}
\textit{100} & /rai{\texttheta}{\textsubbridge{t}}elaman/&$\rightarrow$ & \textit{3000} & /rai{\texttheta}elaman/&$\rightarrow$ & \textit{3200} & /rai{\texttheta}lamn/&$\rightarrow$ & \textit{5500} & /ri{\texttheta}lamn/&$\rightarrow$ & \textit{7600} & /ri{\texttheta}lan/&$\rightarrow$ & \textit{12500} & /ri{\texttheta}an/& \\
\end{tabular}

\vspace{20pt}\hline

\end{nopagebreak}
\filbreak



\vspace{15pt}
\begin{nopagebreak}
\noindent{\fontsize{20pt}{10pt}\textbf{lhan} } \textit{birth} (N)\\
\noindent {\tovian \fontsize{20pt}{10pt} \textbf{lhan} }\\
\noindent /{\textbeltl}{\textprimstress}an/\\
\noindent lit. bringing+light\\


\noindent History:

\vspace{-0pt}
\hspace{40pt}
\begin{tabular}{ccc}
\textit{50} & /{\textbeltl}ogaahana/&$\rightarrow$ & \textit{3200} & /{\textbeltl}gaahana/&$\rightarrow$ & \textit{4500} & /{\textbeltl}aahana/&$\rightarrow$ & \textit{5000} & /{\textbeltl}aaana/&$\rightarrow$ & \textit{5500} & /{\textbeltl}aana/&$\rightarrow$ & \textit{9500} & /{\textbeltl}aan/&$\rightarrow$ & \textit{12005} & /{\textbeltl}an/& \\
\end{tabular}

\vspace{20pt}\hline

\end{nopagebreak}
\filbreak



\vspace{15pt}
\begin{nopagebreak}
\noindent{\fontsize{20pt}{10pt}\textbf{drik} } \textit{blood} (N)\\
\noindent {\tovian \fontsize{20pt}{10pt} \textbf{drik} }\\
\noindent /dr{\textprimstress}ik/\\


\noindent History:

\vspace{-0pt}
\hspace{40pt}
\begin{tabular}{ccc}
\textit{0} & /driks/&$\rightarrow$ & \textit{13001} & /drik/& \\
\end{tabular}

\vspace{20pt}\hline

\end{nopagebreak}
\filbreak



\vspace{15pt}
\begin{nopagebreak}
\noindent{\fontsize{20pt}{10pt}\textbf{med} } \textit{bloom} (N)\\
\noindent {\tovian \fontsize{20pt}{10pt} \textbf{med} }\\
\noindent /m{\textprimstress}ed/\\
\noindent lit. flower+showing\\


\noindent History:

\vspace{-0pt}
\hspace{40pt}
\begin{tabular}{ccc}
\textit{1000} & /drenpkfe{\textsubbridge{t}}e/&$\rightarrow$ & \textit{2000} & /drenpkfede/&$\rightarrow$ & \textit{3000} & /drenpfede/&$\rightarrow$ & \textit{3200} & /drnpfede/&$\rightarrow$ & \textit{6000} & /drmpfede/&$\rightarrow$ & \textit{7501} & /rmpfede/&$\rightarrow$ & \textit{7600} & /mpfede/&$\rightarrow$ & \textit{8750} & /mfede/&$\rightarrow$ & \textit{9500} & /mfed/&$\rightarrow$ & \textit{11995} & /med/& \\
\end{tabular}

\vspace{20pt}\hline

\end{nopagebreak}
\filbreak



\vspace{15pt}
\begin{nopagebreak}
\noindent{\fontsize{20pt}{10pt}\textbf{dor} } \textit{blueness/blue} (N)\\
\noindent {\tovian \fontsize{20pt}{10pt} \textbf{dor} }\\
\noindent /d{\textprimstress}or/\\


\noindent History:

\vspace{-0pt}
\hspace{40pt}
\begin{tabular}{ccc}
\textit{0} & /dora/&$\rightarrow$ & \textit{9500} & /dor/& \\
\end{tabular}

\vspace{20pt}\hline

\end{nopagebreak}
\filbreak



\vspace{15pt}
\begin{nopagebreak}
\noindent{\fontsize{20pt}{10pt}\textbf{sith} } \textit{boat/ship/canoe} (N)\\
\noindent {\tovian \fontsize{20pt}{10pt} \textbf{sith} }\\
\noindent /s{\textprimstress}i{\texttheta}/\\
\noindent lit. water+traveling\\


\noindent History:

\vspace{-0pt}
\hspace{40pt}
\begin{tabular}{ccc}
\textit{3000} & /zdelki{\texttheta}sa/&$\rightarrow$ & \textit{3000} & /zdeli{\texttheta}sa/&$\rightarrow$ & \textit{3200} & /zdli{\texttheta}sa/&$\rightarrow$ & \textit{4500} & /zli{\texttheta}sa/&$\rightarrow$ & \textit{4502} & /zli{\texttheta}a/&$\rightarrow$ & \textit{9500} & /zli{\texttheta}/&$\rightarrow$ & \textit{12000} & /sli{\texttheta}/&$\rightarrow$ & \textit{12500} & /si{\texttheta}/& \\
\end{tabular}

\vspace{20pt}\hline

\end{nopagebreak}
\filbreak



\vspace{15pt}
\begin{nopagebreak}
\noindent{\fontsize{20pt}{10pt}\textbf{adlel} } \textit{body/trunk} (N)\\
\noindent {\tovian \fontsize{20pt}{10pt} \textbf{adlel} }\\
\noindent /{\textprimstress}adlel/\\


\noindent History:

\vspace{-0pt}
\hspace{40pt}
\begin{tabular}{ccc}
\textit{0} & /a{\textsubbridge{t}}ellenal/&$\rightarrow$ & \textit{2000} & /adellenal/&$\rightarrow$ & \textit{3200} & /adllenl/&$\rightarrow$ & \textit{6500} & /adlenl/&$\rightarrow$ & \textit{7600} & /adlel/& \\
\end{tabular}

\vspace{20pt}\hline

\end{nopagebreak}
\filbreak



\vspace{15pt}
\begin{nopagebreak}
\noindent{\fontsize{20pt}{10pt}\textbf{sul} } \textit{bottom/shallows/shortness} (N)\\
\noindent {\tovian \fontsize{20pt}{10pt} \textbf{sul} }\\
\noindent /s{\textprimstress}ul/\\


\noindent History:

\vspace{-0pt}
\hspace{40pt}
\begin{tabular}{ccc}
\textit{0} & /zul/&$\rightarrow$ & \textit{12000} & /sul/& \\
\end{tabular}

\vspace{20pt}\hline

\end{nopagebreak}
\filbreak



\vspace{15pt}
\begin{nopagebreak}
\noindent{\fontsize{20pt}{10pt}\textbf{lhog} } \textit{bringing/carrying} (N)\\
\noindent {\tovian \fontsize{20pt}{10pt} \textbf{lhog} }\\
\noindent /{\textbeltl}{\textprimstress}og/\\


\noindent History:

\vspace{-0pt}
\hspace{40pt}
\begin{tabular}{ccc}
\textit{0} & /{\textbeltl}oga/&$\rightarrow$ & \textit{9500} & /{\textbeltl}og/& \\
\end{tabular}

\vspace{20pt}\hline

\end{nopagebreak}
\filbreak



\vspace{15pt}
\begin{nopagebreak}
\noindent{\fontsize{20pt}{10pt}\textbf{ke-} } \textit{CAUS/causative voice marker} (VOICE)\\
\noindent {\tovian \fontsize{20pt}{10pt} \textbf{ke-} }\\
\noindent /k{\textprimstress}e/\\


\noindent History:

\vspace{-0pt}
\hspace{40pt}
\begin{tabular}{ccc}
\textit{0} & /ke-/& \\
\end{tabular}

\vspace{20pt}\hline

\end{nopagebreak}
\filbreak



\vspace{15pt}
\begin{nopagebreak}
\noindent{\fontsize{20pt}{10pt}\textbf{lhel} } \textit{center} (N)\\
\noindent {\tovian \fontsize{20pt}{10pt} \textbf{lhel} }\\
\noindent /{\textbeltl}{\textprimstress}el/\\


\noindent History:

\vspace{-0pt}
\hspace{40pt}
\begin{tabular}{ccc}
\textit{0} & /{\textbeltl}el/& \\
\end{tabular}

\vspace{20pt}\hline

\end{nopagebreak}
\filbreak



\vspace{15pt}
\begin{nopagebreak}
\noindent{\fontsize{20pt}{10pt}\textbf{kan} } \textit{change} (N)\\
\noindent {\tovian \fontsize{20pt}{10pt} \textbf{kan} }\\
\noindent /k{\textprimstress}an/\\


\noindent History:

\vspace{-0pt}
\hspace{40pt}
\begin{tabular}{ccc}
\textit{0} & /kana/&$\rightarrow$ & \textit{9500} & /kan/& \\
\end{tabular}

\vspace{20pt}\hline

\end{nopagebreak}
\filbreak



\vspace{15pt}
\begin{nopagebreak}
\noindent{\fontsize{20pt}{10pt}\textbf{melhen} } \textit{child} (N)\\
\noindent {\tovian \fontsize{20pt}{10pt} \textbf{melhen} }\\
\noindent /m{\textprimstress}e{\textbeltl}en/\\
\noindent lit. newness+dream\\


\noindent History:

\vspace{-0pt}
\hspace{40pt}
\begin{tabular}{ccc}
\textit{0} & /{\textyogh}emeno{\textbeltl}ena/&$\rightarrow$ & \textit{3200} & /{\textyogh}men{\textbeltl}ena/&$\rightarrow$ & \textit{6400} & /{\textyogh}me{\textbeltl}ena/&$\rightarrow$ & \textit{9500} & /{\textyogh}me{\textbeltl}en/&$\rightarrow$ & \textit{12001} & /{\textesh}me{\textbeltl}en/&$\rightarrow$ & \textit{13002} & /me{\textbeltl}en/& \\
\end{tabular}

\vspace{20pt}\hline

\end{nopagebreak}
\filbreak



\vspace{15pt}
\begin{nopagebreak}
\noindent{\fontsize{20pt}{10pt}\textbf{shithil} } \textit{city} (N)\\
\noindent {\tovian \fontsize{20pt}{10pt} \textbf{shithil} }\\
\noindent /{\textesh}{\textprimstress}i{\texttheta}il/\\
\noindent lit. wakefulness+place\\


\noindent History:

\vspace{-0pt}
\hspace{40pt}
\begin{tabular}{ccc}
\textit{6000} & /{\textyogh}i{\texttheta}{\texttheta}ile/&$\rightarrow$ & \textit{6500} & /{\textyogh}i{\texttheta}ile/&$\rightarrow$ & \textit{9500} & /{\textyogh}i{\texttheta}il/&$\rightarrow$ & \textit{12001} & /{\textesh}i{\texttheta}il/& \\
\end{tabular}

\vspace{20pt}\hline

\end{nopagebreak}
\filbreak



\vspace{15pt}
\begin{nopagebreak}
\noindent{\fontsize{20pt}{10pt}\textbf{tam} } \textit{closure/shutting} (N)\\
\noindent {\tovian \fontsize{20pt}{10pt} \textbf{tam} }\\
\noindent /{\textsubbridge{t}}{\textprimstress}am/\\


\noindent History:

\vspace{-0pt}
\hspace{40pt}
\begin{tabular}{ccc}
\textit{0} & /{\textsubbridge{t}}anum/&$\rightarrow$ & \textit{3200} & /{\textsubbridge{t}}anm/&$\rightarrow$ & \textit{7600} & /{\textsubbridge{t}}am/& \\
\end{tabular}

\vspace{20pt}\hline

\end{nopagebreak}
\filbreak



\vspace{15pt}
\begin{nopagebreak}
\noindent{\fontsize{20pt}{10pt}\textbf{tush} } \textit{coldness/cold} (N)\\
\noindent {\tovian \fontsize{20pt}{10pt} \textbf{tush} }\\
\noindent /{\textsubbridge{t}}{\textprimstress}u{\textesh}/\\


\noindent History:

\vspace{-0pt}
\hspace{40pt}
\begin{tabular}{ccc}
\textit{0} & /{\textsubbridge{t}}ku{\textesh}u/&$\rightarrow$ & \textit{3000} & /{\textsubbridge{t}}u{\textesh}u/&$\rightarrow$ & \textit{9500} & /{\textsubbridge{t}}u{\textesh}/& \\
\end{tabular}

\vspace{20pt}\hline

\end{nopagebreak}
\filbreak



\vspace{15pt}
\begin{nopagebreak}
\noindent{\fontsize{20pt}{10pt}\textbf{yi-} } \textit{COM/comitative case marker/with/together with} (CASE)\\
\noindent {\tovian \fontsize{20pt}{10pt} \textbf{yi-} }\\
\noindent /y{\textprimstress}i/\\


\noindent History:

\vspace{-0pt}
\hspace{40pt}
\begin{tabular}{ccc}
\textit{0} & /yi-/& \\
\end{tabular}

\vspace{20pt}\hline

\end{nopagebreak}
\filbreak



\vspace{15pt}
\begin{nopagebreak}
\noindent{\fontsize{20pt}{10pt}\textbf{yalrith} } \textit{common bird} (N)\\
\noindent {\tovian \fontsize{20pt}{10pt} \textbf{yalrith} }\\
\noindent /y{\textprimstress}alri{\texttheta}/\\
\noindent lit. normality+air\\


\noindent History:

\vspace{-0pt}
\hspace{40pt}
\begin{tabular}{ccc}
\textit{12000} & /yalri{\texttheta}/& \\
\end{tabular}

\vspace{20pt}\hline

\end{nopagebreak}
\filbreak



\vspace{15pt}
\begin{nopagebreak}
\noindent{\fontsize{20pt}{10pt}\textbf{yal} } \textit{commonness/mundanity/normality} (N)\\
\noindent {\tovian \fontsize{20pt}{10pt} \textbf{yal} }\\
\noindent /y{\textprimstress}al/\\


\noindent History:

\vspace{-0pt}
\hspace{40pt}
\begin{tabular}{ccc}
\textit{0} & /ileriala/&$\rightarrow$ & \textit{2201} & /ileriyala/&$\rightarrow$ & \textit{3200} & /ilriyala/&$\rightarrow$ & \textit{6300} & /ilryala/&$\rightarrow$ & \textit{6503} & /irlyala/&$\rightarrow$ & \textit{7500} & /rlyala/&$\rightarrow$ & \textit{7501} & /lyala/&$\rightarrow$ & \textit{7600} & /yala/&$\rightarrow$ & \textit{9500} & /yal/& \\
\end{tabular}

\vspace{20pt}\hline

\end{nopagebreak}
\filbreak



\vspace{15pt}
\begin{nopagebreak}
\noindent{\fontsize{20pt}{10pt}\textbf{lad} } \textit{comparison} (N)\\
\noindent {\tovian \fontsize{20pt}{10pt} \textbf{lad} }\\
\noindent /l{\textprimstress}ad/\\


\noindent History:

\vspace{-0pt}
\hspace{40pt}
\begin{tabular}{ccc}
\textit{0} & /roladil/&$\rightarrow$ & \textit{3200} & /rladl/&$\rightarrow$ & \textit{7600} & /ladl/&$\rightarrow$ & \textit{13001} & /lad/& \\
\end{tabular}

\vspace{20pt}\hline

\end{nopagebreak}
\filbreak



\vspace{15pt}
\begin{nopagebreak}
\noindent{\fontsize{20pt}{10pt}\textbf{om} } \textit{completeness/wholeness} (N)\\
\noindent {\tovian \fontsize{20pt}{10pt} \textbf{om} }\\
\noindent /{\textprimstress}om/\\


\noindent History:

\vspace{-0pt}
\hspace{40pt}
\begin{tabular}{ccc}
\textit{0} & /oma/&$\rightarrow$ & \textit{9500} & /om/& \\
\end{tabular}

\vspace{20pt}\hline

\end{nopagebreak}
\filbreak



\vspace{15pt}
\begin{nopagebreak}
\noindent{\fontsize{20pt}{10pt}\textbf{ko-} } \textit{COND/conditional mood marker} (MOOD)\\
\noindent {\tovian \fontsize{20pt}{10pt} \textbf{ko-} }\\
\noindent /k{\textprimstress}o/\\


\noindent History:

\vspace{-0pt}
\hspace{40pt}
\begin{tabular}{ccc}
\textit{0} & /ko-/& \\
\end{tabular}

\vspace{20pt}\hline

\end{nopagebreak}
\filbreak



\vspace{15pt}
\begin{nopagebreak}
\noindent{\fontsize{20pt}{10pt}\textbf{glor} } \textit{conflict/fighting} (N)\\
\noindent {\tovian \fontsize{20pt}{10pt} \textbf{glor} }\\
\noindent /gl{\textprimstress}or/\\


\noindent History:

\vspace{-0pt}
\hspace{40pt}
\begin{tabular}{ccc}
\textit{0} & /gilore/&$\rightarrow$ & \textit{3200} & /glore/&$\rightarrow$ & \textit{9500} & /glor/& \\
\end{tabular}

\vspace{20pt}\hline

\end{nopagebreak}
\filbreak



\vspace{15pt}
\begin{nopagebreak}
\noindent{\fontsize{20pt}{10pt}\textbf{leth} } \textit{consumption/tasting/eating} (N)\\
\noindent {\tovian \fontsize{20pt}{10pt} \textbf{leth} }\\
\noindent /l{\textprimstress}e{\texttheta}/\\


\noindent History:

\vspace{-0pt}
\hspace{40pt}
\begin{tabular}{ccc}
\textit{0} & /ale{\dh}a/&$\rightarrow$ & \textit{7500} & /le{\dh}a/&$\rightarrow$ & \textit{9500} & /le{\dh}/&$\rightarrow$ & \textit{12002} & /le{\texttheta}/& \\
\end{tabular}

\vspace{20pt}\hline

\end{nopagebreak}
\filbreak



\vspace{15pt}
\begin{nopagebreak}
\noindent{\fontsize{20pt}{10pt}\textbf{lhom} } \textit{container/vessel} (N)\\
\noindent {\tovian \fontsize{20pt}{10pt} \textbf{lhom} }\\
\noindent /{\textbeltl}{\textprimstress}om/\\


\noindent History:

\vspace{-0pt}
\hspace{40pt}
\begin{tabular}{ccc}
\textit{0} & /{\textbeltl}oma/&$\rightarrow$ & \textit{9500} & /{\textbeltl}om/& \\
\end{tabular}

\vspace{20pt}\hline

\end{nopagebreak}
\filbreak



\vspace{15pt}
\begin{nopagebreak}
\noindent{\fontsize{20pt}{10pt}\textbf{lhun} } \textit{council} (N)\\
\noindent {\tovian \fontsize{20pt}{10pt} \textbf{lhun} }\\
\noindent /{\textbeltl}{\textprimstress}un/\\
\noindent lit. thinking+joining\\


\noindent History:

\vspace{-0pt}
\hspace{40pt}
\begin{tabular}{ccc}
\textit{0} & /na{\textbeltl}{\texttoptiebar{t\textbeltl}}una/&$\rightarrow$ & \textit{3200} & /n{\textbeltl}{\texttoptiebar{t\textbeltl}}una/&$\rightarrow$ & \textit{6400} & /{\textbeltl}{\texttoptiebar{t\textbeltl}}una/&$\rightarrow$ & \textit{9500} & /{\textbeltl}{\texttoptiebar{t\textbeltl}}un/&$\rightarrow$ & \textit{14001} & /{\textbeltl}un/& \\
\end{tabular}

\vspace{20pt}\hline

\end{nopagebreak}
\filbreak



\vspace{15pt}
\begin{nopagebreak}
\noindent{\fontsize{20pt}{10pt}\textbf{ks-o} } \textit{COUNTER/counterfactual mood marker} (MOOD)\\
\noindent {\tovian \fontsize{20pt}{10pt} \textbf{ks-o} }\\
\noindent /{\textprimstress}o/\\


\noindent History:

\vspace{-0pt}
\hspace{40pt}
\begin{tabular}{ccc}
\textit{0} & /ks-o/& \\
\end{tabular}

\vspace{20pt}\hline

\end{nopagebreak}
\filbreak



\vspace{15pt}
\begin{nopagebreak}
\noindent{\fontsize{20pt}{10pt}\textbf{lor} } \textit{dance} (N)\\
\noindent {\tovian \fontsize{20pt}{10pt} \textbf{lor} }\\
\noindent /l{\textprimstress}or/\\
\noindent lit. movement+action\\


\noindent History:

\vspace{-0pt}
\hspace{40pt}
\begin{tabular}{ccc}
\textit{0} & /kwillore/&$\rightarrow$ & \textit{3200} & /kwllore/&$\rightarrow$ & \textit{6500} & /kwlore/&$\rightarrow$ & \textit{7501} & /wlore/&$\rightarrow$ & \textit{7600} & /lore/&$\rightarrow$ & \textit{9500} & /lor/& \\
\end{tabular}

\vspace{20pt}\hline

\end{nopagebreak}
\filbreak



\vspace{15pt}
\begin{nopagebreak}
\noindent{\fontsize{20pt}{10pt}\textbf{mral} } \textit{darkness/dark} (N)\\
\noindent {\tovian \fontsize{20pt}{10pt} \textbf{mral} }\\
\noindent /mr{\textprimstress}al/\\


\noindent History:

\vspace{-0pt}
\hspace{40pt}
\begin{tabular}{ccc}
\textit{0} & /bral/&$\rightarrow$ & \textit{11000} & /mral/& \\
\end{tabular}

\vspace{20pt}\hline

\end{nopagebreak}
\filbreak



\vspace{15pt}
\begin{nopagebreak}
\noindent{\fontsize{20pt}{10pt}\textbf{mi-} } \textit{DAT/dative case marker} (CASE)\\
\noindent {\tovian \fontsize{20pt}{10pt} \textbf{mi-} }\\
\noindent /m{\textprimstress}i/\\


\noindent History:

\vspace{-0pt}
\hspace{40pt}
\begin{tabular}{ccc}
\textit{0} & /mi-/& \\
\end{tabular}

\vspace{20pt}\hline

\end{nopagebreak}
\filbreak



\vspace{15pt}
\begin{nopagebreak}
\noindent{\fontsize{20pt}{10pt}\textbf{mloan} } \textit{dawn/morning} (N)\\
\noindent {\tovian \fontsize{20pt}{10pt} \textbf{mloan} }\\
\noindent /ml{\textprimstress}oan/\\
\noindent lit. first+light\\


\noindent History:

\vspace{-0pt}
\hspace{40pt}
\begin{tabular}{ccc}
\textit{500} & /pu-loahana/&$\rightarrow$ & \textit{1000} & /puloahana/&$\rightarrow$ & \textit{3200} & /ploahana/&$\rightarrow$ & \textit{5000} & /ploaana/&$\rightarrow$ & \textit{5500} & /ploana/&$\rightarrow$ & \textit{9500} & /ploan/&$\rightarrow$ & \textit{11000} & /mloan/& \\
\end{tabular}

\vspace{20pt}\hline

\end{nopagebreak}
\filbreak



\vspace{15pt}
\begin{nopagebreak}
\noindent{\fontsize{20pt}{10pt}\textbf{lhom} } \textit{day} (N)\\
\noindent {\tovian \fontsize{20pt}{10pt} \textbf{lhom} }\\
\noindent /{\textbeltl}{\textprimstress}om/\\
\noindent \textit{Period of time}\\


\noindent History:

\vspace{-0pt}
\hspace{40pt}
\begin{tabular}{ccc}
\textit{0} & /{\textbeltl}onum/&$\rightarrow$ & \textit{3200} & /{\textbeltl}onm/&$\rightarrow$ & \textit{7600} & /{\textbeltl}om/& \\
\end{tabular}

\vspace{20pt}\hline

\end{nopagebreak}
\filbreak



\vspace{15pt}
\begin{nopagebreak}
\noindent{\fontsize{20pt}{10pt}\textbf{twimen} } \textit{day/daytime} (N)\\
\noindent {\tovian \fontsize{20pt}{10pt} \textbf{twimen} }\\
\noindent /{\textsubbridge{t}}w{\textprimstress}imen/\\
\noindent \textit{Time of day}\\


\noindent History:

\vspace{-0pt}
\hspace{40pt}
\begin{tabular}{ccc}
\textit{0} & /{\textsubbridge{t}}uimena/&$\rightarrow$ & \textit{2202} & /{\textsubbridge{t}}uwimena/&$\rightarrow$ & \textit{3200} & /{\textsubbridge{t}}wimena/&$\rightarrow$ & \textit{9500} & /{\textsubbridge{t}}wimen/& \\
\end{tabular}

\vspace{20pt}\hline

\end{nopagebreak}
\filbreak



\vspace{15pt}
\begin{nopagebreak}
\noindent{\fontsize{20pt}{10pt}\textbf{dol} } \textit{death/dying} (N)\\
\noindent {\tovian \fontsize{20pt}{10pt} \textbf{dol} }\\
\noindent /d{\textprimstress}ol/\\


\noindent History:

\vspace{-0pt}
\hspace{40pt}
\begin{tabular}{ccc}
\textit{0} & /{\textschwa}{\textsubbridge{t}}ol/&$\rightarrow$ & \textit{2000} & /{\textschwa}dol/&$\rightarrow$ & \textit{2300} & /dol/& \\
\end{tabular}

\vspace{20pt}\hline

\end{nopagebreak}
\filbreak



\vspace{15pt}
\begin{nopagebreak}
\noindent{\fontsize{20pt}{10pt}\textbf{lotush} } \textit{December} (N)\\
\noindent {\tovian \fontsize{20pt}{10pt} \textbf{lotush} }\\
\noindent /l{\textprimstress}o{\textsubbridge{t}}u{\textesh}/\\
\noindent lit. one+cold\\
\noindent \textit{Ventiosh}\\


\noindent History:

\vspace{-0pt}
\hspace{40pt}
\begin{tabular}{ccc}
\textit{6000} & /lo{\textsubbridge{t}}u{\textesh}u/&$\rightarrow$ & \textit{9500} & /lo{\textsubbridge{t}}u{\textesh}/& \\
\end{tabular}

\vspace{20pt}\hline

\end{nopagebreak}
\filbreak



\vspace{15pt}
\begin{nopagebreak}
\noindent{\fontsize{20pt}{10pt}\textbf{tlil} } \textit{desert} (N)\\
\noindent {\tovian \fontsize{20pt}{10pt} \textbf{tlil} }\\
\noindent /{\texttoptiebar{t\textbeltl}}{\textprimstress}il/\\
\noindent lit. sand+place\\


\noindent History:

\vspace{-0pt}
\hspace{40pt}
\begin{tabular}{ccc}
\textit{1000} & /{\textesh}o{\texttoptiebar{t\textbeltl}}{\texttheta}k{\textesh}{\textsubbridge{t}}{\texttheta}ile/&$\rightarrow$ & \textit{1000} & /{\textesh}{\texttoptiebar{t\textbeltl}}{\texttheta}k{\textesh}{\textsubbridge{t}}{\texttheta}ile/&$\rightarrow$ & \textit{3000} & /{\textesh}{\texttoptiebar{t\textbeltl}}{\texttheta}k{\textesh}{\texttheta}ile/&$\rightarrow$ & \textit{4500} & /{\textesh}{\texttoptiebar{t\textbeltl}}{\texttheta}{\textesh}{\texttheta}ile/&$\rightarrow$ & \textit{4502} & /{\textesh}{\texttoptiebar{t\textbeltl}}{\texttheta}{\texttheta}ile/&$\rightarrow$ & \textit{6500} & /{\textesh}{\texttoptiebar{t\textbeltl}}{\texttheta}ile/&$\rightarrow$ & \textit{7501} & /{\texttoptiebar{t\textbeltl}}{\texttheta}ile/&$\rightarrow$ & \textit{9500} & /{\texttoptiebar{t\textbeltl}}{\texttheta}il/&$\rightarrow$ & \textit{14000} & /{\texttoptiebar{t\textbeltl}}il/& \\
\end{tabular}

\vspace{20pt}\hline

\end{nopagebreak}
\filbreak



\vspace{15pt}
\begin{nopagebreak}
\noindent{\fontsize{20pt}{10pt}\textbf{sahyeol} } \textit{discovery} (N)\\
\noindent {\tovian \fontsize{20pt}{10pt} \textbf{sahyeol} }\\
\noindent /saħy{\textprimstress}eol/\\
\noindent lit. make+wisdom\\


\noindent History:

\vspace{-0pt}
\hspace{40pt}
\begin{tabular}{ccc}
\textit{50} & /silahe{\textsubbridge{t}}arayeol/&$\rightarrow$ & \textit{1000} & /silah{\textsubbridge{t}}arayeol/&$\rightarrow$ & \textit{2100} & /silah{\textsubbridge{t}}aryeol/&$\rightarrow$ & \textit{3000} & /silaharyeol/&$\rightarrow$ & \textit{3200} & /slahryeol/&$\rightarrow$ & \textit{7600} & /slahyeol/&$\rightarrow$ & \textit{12500} & /sahyeol/&$\rightarrow$ & \textit{14006} & /saħyeol/& \\
\end{tabular}

\vspace{20pt}\hline

\end{nopagebreak}
\filbreak



\vspace{15pt}
\begin{nopagebreak}
\noindent{\fontsize{20pt}{10pt}\textbf{anon} } \textit{dog} (N)\\
\noindent {\tovian \fontsize{20pt}{10pt} \textbf{anon} }\\
\noindent /{\textprimstress}anon/\\


\noindent History:

\vspace{-0pt}
\hspace{40pt}
\begin{tabular}{ccc}
\textit{0} & /am{\textsubbridge{t}}orin/&$\rightarrow$ & \textit{2000} & /amdorin/&$\rightarrow$ & \textit{3200} & /amdorn/&$\rightarrow$ & \textit{6000} & /andorn/&$\rightarrow$ & \textit{7600} & /andon/&$\rightarrow$ & \textit{8750} & /anon/& \\
\end{tabular}

\vspace{20pt}\hline

\end{nopagebreak}
\filbreak



\vspace{15pt}
\begin{nopagebreak}
\noindent{\fontsize{20pt}{10pt}\textbf{thok} } \textit{down/underside/bottom} (N)\\
\noindent {\tovian \fontsize{20pt}{10pt} \textbf{thok} }\\
\noindent /{\texttheta}{\textprimstress}ok/\\


\noindent History:

\vspace{-0pt}
\hspace{40pt}
\begin{tabular}{ccc}
\textit{0} & /{\texttheta}ok/& \\
\end{tabular}

\vspace{20pt}\hline

\end{nopagebreak}
\filbreak



\vspace{15pt}
\begin{nopagebreak}
\noindent{\fontsize{20pt}{10pt}\textbf{lhen} } \textit{dream/spirit/soul} (N)\\
\noindent {\tovian \fontsize{20pt}{10pt} \textbf{lhen} }\\
\noindent /{\textbeltl}{\textprimstress}en/\\


\noindent History:

\vspace{-0pt}
\hspace{40pt}
\begin{tabular}{ccc}
\textit{0} & /{\textbeltl}ena/&$\rightarrow$ & \textit{9500} & /{\textbeltl}en/& \\
\end{tabular}

\vspace{20pt}\hline

\end{nopagebreak}
\filbreak



\vspace{15pt}
\begin{nopagebreak}
\noindent{\fontsize{20pt}{10pt}\textbf{anolhyel} } \textit{dreambringer} (N)\\
\noindent {\tovian \fontsize{20pt}{10pt} \textbf{anolhyel} }\\
\noindent /an{\textprimstress}o{\textbeltl}yel/\\
\noindent lit. dream+COM+walking+person\\


\noindent History:

\vspace{-0pt}
\hspace{40pt}
\begin{tabular}{ccc}
\textit{10} & /{\textbeltl}enayi-no{\textbeltl}driela/&$\rightarrow$ & \textit{1000} & /{\textbeltl}enayino{\textbeltl}driela/&$\rightarrow$ & \textit{2100} & /{\textbeltl}enyino{\textbeltl}driela/&$\rightarrow$ & \textit{2201} & /{\textbeltl}enyino{\textbeltl}driyela/&$\rightarrow$ & \textit{3200} & /{\textbeltl}nyno{\textbeltl}dryela/&$\rightarrow$ & \textit{4500} & /{\textbeltl}nyno{\textbeltl}ryela/&$\rightarrow$ & \textit{6000} & /{\textbeltl}nd{\textyogh}no{\textbeltl}ryela/&$\rightarrow$ & \textit{7501} & /nd{\textyogh}no{\textbeltl}ryela/&$\rightarrow$ & \textit{7600} & /nd{\textyogh}no{\textbeltl}yela/&$\rightarrow$ & \textit{8750} & /n{\textyogh}no{\textbeltl}yela/&$\rightarrow$ & \textit{8760} & /an{\textyogh}no{\textbeltl}yela/&$\rightarrow$ & \textit{9500} & /an{\textyogh}no{\textbeltl}yel/&$\rightarrow$ & \textit{12001} & /an{\textesh}no{\textbeltl}yel/&$\rightarrow$ & \textit{13002} & /anno{\textbeltl}yel/&$\rightarrow$ & \textit{15000} & /ano{\textbeltl}yel/& \\
\end{tabular}

\vspace{20pt}\hline

\end{nopagebreak}
\filbreak



\vspace{15pt}
\begin{nopagebreak}
\noindent{\fontsize{20pt}{10pt}\textbf{munorhan} } \textit{dusk/evening/sunset} (N)\\
\noindent {\tovian \fontsize{20pt}{10pt} \textbf{munorhan} }\\
\noindent /mun{\textprimstress}orħan/\\
\noindent lit. last+light\\


\noindent History:

\vspace{-0pt}
\hspace{40pt}
\begin{tabular}{ccc}
\textit{500} & /panorahana/&$\rightarrow$ & \textit{3200} & /pnorhana/&$\rightarrow$ & \textit{9500} & /pnorhan/&$\rightarrow$ & \textit{11000} & /mnorhan/&$\rightarrow$ & \textit{11996} & /munorhan/&$\rightarrow$ & \textit{14006} & /munorħan/& \\
\end{tabular}

\vspace{20pt}\hline

\end{nopagebreak}
\filbreak



\vspace{15pt}
\begin{nopagebreak}
\noindent{\fontsize{20pt}{10pt}\textbf{thek} } \textit{ear} (N)\\
\noindent {\tovian \fontsize{20pt}{10pt} \textbf{thek} }\\
\noindent /{\texttheta}{\textprimstress}ek/\\


\noindent History:

\vspace{-0pt}
\hspace{40pt}
\begin{tabular}{ccc}
\textit{0} & /{\texttheta}ek/& \\
\end{tabular}

\vspace{20pt}\hline

\end{nopagebreak}
\filbreak



\vspace{15pt}
\begin{nopagebreak}
\noindent{\fontsize{20pt}{10pt}\textbf{tal} } \textit{east} (N)\\
\noindent {\tovian \fontsize{20pt}{10pt} \textbf{tal} }\\
\noindent /{\textsubbridge{t}}{\textprimstress}al/\\


\noindent History:

\vspace{-0pt}
\hspace{40pt}
\begin{tabular}{ccc}
\textit{0} & /{\textsubbridge{t}}al/& \\
\end{tabular}

\vspace{20pt}\hline

\end{nopagebreak}
\filbreak



\vspace{15pt}
\begin{nopagebreak}
\noindent{\fontsize{20pt}{10pt}\textbf{ath} } \textit{eight/8} (N)\\
\noindent {\tovian \fontsize{20pt}{10pt} \textbf{ath} }\\
\noindent /{\textprimstress}a{\texttheta}/\\


\noindent History:

\vspace{-0pt}
\hspace{40pt}
\begin{tabular}{ccc}
\textit{0} & /a{\texttheta}/& \\
\end{tabular}

\vspace{20pt}\hline

\end{nopagebreak}
\filbreak



\vspace{15pt}
\begin{nopagebreak}
\noindent{\fontsize{20pt}{10pt}\textbf{myar} } \textit{eleven/11} (N)\\
\noindent {\tovian \fontsize{20pt}{10pt} \textbf{myar} }\\
\noindent /my{\textprimstress}ar/\\


\noindent History:

\vspace{-0pt}
\hspace{40pt}
\begin{tabular}{ccc}
\textit{0} & /beyar/&$\rightarrow$ & \textit{2100} & /byar/&$\rightarrow$ & \textit{11000} & /myar/& \\
\end{tabular}

\vspace{20pt}\hline

\end{nopagebreak}
\filbreak



\vspace{15pt}
\begin{nopagebreak}
\noindent{\fontsize{20pt}{10pt}\textbf{remi} } \textit{empty/hollow} (N)\\
\noindent {\tovian \fontsize{20pt}{10pt} \textbf{remi} }\\
\noindent /r{\textprimstress}emi/\\


\noindent History:

\vspace{-0pt}
\hspace{40pt}
\begin{tabular}{ccc}
\textit{0} & /oneremiha/&$\rightarrow$ & \textit{3200} & /onremiha/&$\rightarrow$ & \textit{5000} & /onremia/&$\rightarrow$ & \textit{7500} & /nremia/&$\rightarrow$ & \textit{7600} & /remia/&$\rightarrow$ & \textit{9500} & /remi/& \\
\end{tabular}

\vspace{20pt}\hline

\end{nopagebreak}
\filbreak



\vspace{15pt}
\begin{nopagebreak}
\noindent{\fontsize{20pt}{10pt}\textbf{par} } \textit{end/final/last} (N)\\
\noindent {\tovian \fontsize{20pt}{10pt} \textbf{par} }\\
\noindent /p{\textprimstress}ar/\\


\noindent History:

\vspace{-0pt}
\hspace{40pt}
\begin{tabular}{ccc}
\textit{0} & /panor/&$\rightarrow$ & \textit{3200} & /panr/&$\rightarrow$ & \textit{7600} & /par/& \\
\end{tabular}

\vspace{20pt}\hline

\end{nopagebreak}
\filbreak



\vspace{15pt}
\begin{nopagebreak}
\noindent{\fontsize{20pt}{10pt}\textbf{yey} } \textit{enemy} (N)\\
\noindent {\tovian \fontsize{20pt}{10pt} \textbf{yey} }\\
\noindent /y{\textprimstress}ey/\\
\noindent lit. GEN+person+evil\\


\noindent History:

\vspace{-0pt}
\hspace{40pt}
\begin{tabular}{ccc}
\textit{0} & /hi-drielaoyal/&$\rightarrow$ & \textit{1000} & /hidrielaoyal/&$\rightarrow$ & \textit{2100} & /hidrielayal/&$\rightarrow$ & \textit{2201} & /hidriyelayal/&$\rightarrow$ & \textit{3200} & /hdryelayl/&$\rightarrow$ & \textit{4500} & /hryelayl/&$\rightarrow$ & \textit{6300} & /hryelyl/&$\rightarrow$ & \textit{7501} & /ryelyl/&$\rightarrow$ & \textit{7600} & /yeyl/&$\rightarrow$ & \textit{13001} & /yey/& \\
\end{tabular}

\vspace{20pt}\hline

\end{nopagebreak}
\filbreak



\vspace{15pt}
\begin{nopagebreak}
\noindent{\fontsize{20pt}{10pt}\textbf{pu-} } \textit{ESS/essive case marker} (CASE)\\
\noindent {\tovian \fontsize{20pt}{10pt} \textbf{pu-} }\\
\noindent /p{\textprimstress}u/\\


\noindent History:

\vspace{-0pt}
\hspace{40pt}
\begin{tabular}{ccc}
\textit{0} & /pu-/& \\
\end{tabular}

\vspace{20pt}\hline

\end{nopagebreak}
\filbreak



\vspace{15pt}
\begin{nopagebreak}
\noindent{\fontsize{20pt}{10pt}\textbf{yal} } \textit{evil/wickedness/badness} (N)\\
\noindent {\tovian \fontsize{20pt}{10pt} \textbf{yal} }\\
\noindent /y{\textprimstress}al/\\


\noindent History:

\vspace{-0pt}
\hspace{40pt}
\begin{tabular}{ccc}
\textit{0} & /oyal/&$\rightarrow$ & \textit{2100} & /yal/& \\
\end{tabular}

\vspace{20pt}\hline

\end{nopagebreak}
\filbreak



\vspace{15pt}
\begin{nopagebreak}
\noindent{\fontsize{20pt}{10pt}\textbf{mril} } \textit{eye/sight/seeing/watching} (N)\\
\noindent {\tovian \fontsize{20pt}{10pt} \textbf{mril} }\\
\noindent /mr{\textprimstress}il/\\


\noindent History:

\vspace{-0pt}
\hspace{40pt}
\begin{tabular}{ccc}
\textit{0} & /em{\textschwa}rila/&$\rightarrow$ & \textit{2300} & /emrila/&$\rightarrow$ & \textit{7500} & /mrila/&$\rightarrow$ & \textit{9500} & /mril/& \\
\end{tabular}

\vspace{20pt}\hline

\end{nopagebreak}
\filbreak



\vspace{15pt}
\begin{nopagebreak}
\noindent{\fontsize{20pt}{10pt}\textbf{tan} } \textit{face} (N)\\
\noindent {\tovian \fontsize{20pt}{10pt} \textbf{tan} }\\
\noindent /{\textsubbridge{t}}{\textprimstress}an/\\


\noindent History:

\vspace{-0pt}
\hspace{40pt}
\begin{tabular}{ccc}
\textit{0} & /{\textsubbridge{t}}aran/&$\rightarrow$ & \textit{3200} & /{\textsubbridge{t}}arn/&$\rightarrow$ & \textit{7600} & /{\textsubbridge{t}}an/& \\
\end{tabular}

\vspace{20pt}\hline

\end{nopagebreak}
\filbreak



\vspace{15pt}
\begin{nopagebreak}
\noindent{\fontsize{20pt}{10pt}\textbf{pan} } \textit{father} (N)\\
\noindent {\tovian \fontsize{20pt}{10pt} \textbf{pan} }\\
\noindent /p{\textprimstress}an/\\


\noindent History:

\vspace{-0pt}
\hspace{40pt}
\begin{tabular}{ccc}
\textit{0} & /pander/&$\rightarrow$ & \textit{3200} & /pandr/&$\rightarrow$ & \textit{8500} & /panr/&$\rightarrow$ & \textit{13001} & /pan/& \\
\end{tabular}

\vspace{20pt}\hline

\end{nopagebreak}
\filbreak



\vspace{15pt}
\begin{nopagebreak}
\noindent{\fontsize{20pt}{10pt}\textbf{ay} } \textit{feather} (N)\\
\noindent {\tovian \fontsize{20pt}{10pt} \textbf{ay} }\\
\noindent /{\textprimstress}ay/\\


\noindent History:

\vspace{-0pt}
\hspace{40pt}
\begin{tabular}{ccc}
\textit{0} & /aruya/&$\rightarrow$ & \textit{2100} & /arya/&$\rightarrow$ & \textit{7600} & /aya/&$\rightarrow$ & \textit{9500} & /ay/& \\
\end{tabular}

\vspace{20pt}\hline

\end{nopagebreak}
\filbreak



\vspace{15pt}
\begin{nopagebreak}
\noindent{\fontsize{20pt}{10pt}\textbf{femush} } \textit{February} (N)\\
\noindent {\tovian \fontsize{20pt}{10pt} \textbf{femush} }\\
\noindent /f{\textprimstress}emu{\textesh}/\\
\noindent lit. three+cold\\
\noindent \textit{Nivost}\\


\noindent History:

\vspace{-0pt}
\hspace{40pt}
\begin{tabular}{ccc}
\textit{6000} & /fem{\textsubbridge{t}}u{\textesh}u/&$\rightarrow$ & \textit{8750} & /femu{\textesh}u/&$\rightarrow$ & \textit{9500} & /femu{\textesh}/& \\
\end{tabular}

\vspace{20pt}\hline

\end{nopagebreak}
\filbreak



\vspace{15pt}
\begin{nopagebreak}
\noindent{\fontsize{20pt}{10pt}\textbf{lhyam} } \textit{fifteen/15} (N)\\
\noindent {\tovian \fontsize{20pt}{10pt} \textbf{lhyam} }\\
\noindent /{\textbeltl}y{\textprimstress}am/\\
\noindent lit. twelve+three\\


\noindent History:

\vspace{-0pt}
\hspace{40pt}
\begin{tabular}{ccc}
\textit{4000} & /{\textbeltl}yafem/&$\rightarrow$ & \textit{12003} & /{\textbeltl}yam/& \\
\end{tabular}

\vspace{20pt}\hline

\end{nopagebreak}
\filbreak



\vspace{15pt}
\begin{nopagebreak}
\noindent{\fontsize{20pt}{10pt}\textbf{tom} } \textit{finger/root} (N)\\
\noindent {\tovian \fontsize{20pt}{10pt} \textbf{tom} }\\
\noindent /{\textsubbridge{t}}{\textprimstress}om/\\


\noindent History:

\vspace{-0pt}
\hspace{40pt}
\begin{tabular}{ccc}
\textit{0} & /{\textsubbridge{t}}omo/&$\rightarrow$ & \textit{9500} & /{\textsubbridge{t}}om/& \\
\end{tabular}

\vspace{20pt}\hline

\end{nopagebreak}
\filbreak



\vspace{15pt}
\begin{nopagebreak}
\noindent{\fontsize{20pt}{10pt}\textbf{kesh} } \textit{fire/burning} (N)\\
\noindent {\tovian \fontsize{20pt}{10pt} \textbf{kesh} }\\
\noindent /k{\textprimstress}e{\textesh}/\\


\noindent History:

\vspace{-0pt}
\hspace{40pt}
\begin{tabular}{ccc}
\textit{0} & /ke{\textesh}/& \\
\end{tabular}

\vspace{20pt}\hline

\end{nopagebreak}
\filbreak



\vspace{15pt}
\begin{nopagebreak}
\noindent{\fontsize{20pt}{10pt}\textbf{pul} } \textit{first} (N)\\
\noindent {\tovian \fontsize{20pt}{10pt} \textbf{pul} }\\
\noindent /p{\textprimstress}ul/\\
\noindent lit. ESS+one\\


\noindent History:

\vspace{-0pt}
\hspace{40pt}
\begin{tabular}{ccc}
\textit{0} & /pu-lo/&$\rightarrow$ & \textit{1000} & /pulo/&$\rightarrow$ & \textit{9500} & /pul/& \\
\end{tabular}

\vspace{20pt}\hline

\end{nopagebreak}
\filbreak



\vspace{15pt}
\begin{nopagebreak}
\noindent{\fontsize{20pt}{10pt}\textbf{rem} } \textit{five/5} (N)\\
\noindent {\tovian \fontsize{20pt}{10pt} \textbf{rem} }\\
\noindent /r{\textprimstress}em/\\


\noindent History:

\vspace{-0pt}
\hspace{40pt}
\begin{tabular}{ccc}
\textit{0} & /rep/&$\rightarrow$ & \textit{11000} & /rem/& \\
\end{tabular}

\vspace{20pt}\hline

\end{nopagebreak}
\filbreak



\vspace{15pt}
\begin{nopagebreak}
\noindent{\fontsize{20pt}{10pt}\textbf{tum} } \textit{flatness} (N)\\
\noindent {\tovian \fontsize{20pt}{10pt} \textbf{tum} }\\
\noindent /{\textsubbridge{t}}{\textprimstress}um/\\


\noindent History:

\vspace{-0pt}
\hspace{40pt}
\begin{tabular}{ccc}
\textit{0} & /{\textsubbridge{t}}um/& \\
\end{tabular}

\vspace{20pt}\hline

\end{nopagebreak}
\filbreak



\vspace{15pt}
\begin{nopagebreak}
\noindent{\fontsize{20pt}{10pt}\textbf{thok} } \textit{floor/ground} (N)\\
\noindent {\tovian \fontsize{20pt}{10pt} \textbf{thok} }\\
\noindent /{\texttheta}{\textprimstress}ok/\\
\noindent lit. down+foot\\


\noindent History:

\vspace{-0pt}
\hspace{40pt}
\begin{tabular}{ccc}
\textit{10} & /{\texttheta}ok{\textesh}o{\textsubbridge{t}}/&$\rightarrow$ & \textit{1000} & /{\texttheta}ok{\textesh}{\textsubbridge{t}}/&$\rightarrow$ & \textit{3000} & /{\texttheta}ok{\textesh}/&$\rightarrow$ & \textit{13001} & /{\texttheta}ok/& \\
\end{tabular}

\vspace{20pt}\hline

\end{nopagebreak}
\filbreak



\vspace{15pt}
\begin{nopagebreak}
\noindent{\fontsize{20pt}{10pt}\textbf{shotl} } \textit{flow} (N)\\
\noindent {\tovian \fontsize{20pt}{10pt} \textbf{shotl} }\\
\noindent /{\textesh}{\textprimstress}o{\texttoptiebar{t\textbeltl}}/\\


\noindent History:

\vspace{-0pt}
\hspace{40pt}
\begin{tabular}{ccc}
\textit{0} & /{\textesh}o{\texttoptiebar{t\textbeltl}}/& \\
\end{tabular}

\vspace{20pt}\hline

\end{nopagebreak}
\filbreak



\vspace{15pt}
\begin{nopagebreak}
\noindent{\fontsize{20pt}{10pt}\textbf{drem} } \textit{flower} (N)\\
\noindent {\tovian \fontsize{20pt}{10pt} \textbf{drem} }\\
\noindent /dr{\textprimstress}em/\\


\noindent History:

\vspace{-0pt}
\hspace{40pt}
\begin{tabular}{ccc}
\textit{0} & /drenpok/&$\rightarrow$ & \textit{1000} & /drenpk/&$\rightarrow$ & \textit{3000} & /drenp/&$\rightarrow$ & \textit{6000} & /dremp/&$\rightarrow$ & \textit{8750} & /drem/& \\
\end{tabular}

\vspace{20pt}\hline

\end{nopagebreak}
\filbreak



\vspace{15pt}
\begin{nopagebreak}
\noindent{\fontsize{20pt}{10pt}\textbf{shotl} } \textit{flowing} (N)\\
\noindent {\tovian \fontsize{20pt}{10pt} \textbf{shotl} }\\
\noindent /{\textesh}{\textprimstress}o{\texttoptiebar{t\textbeltl}}/\\


\noindent History:

\vspace{-0pt}
\hspace{40pt}
\begin{tabular}{ccc}
\textit{0} & /{\textesh}o{\texttoptiebar{t\textbeltl}}/& \\
\end{tabular}

\vspace{20pt}\hline

\end{nopagebreak}
\filbreak



\vspace{15pt}
\begin{nopagebreak}
\noindent{\fontsize{20pt}{10pt}\textbf{tlurith} } \textit{fog} (N)\\
\noindent {\tovian \fontsize{20pt}{10pt} \textbf{tlurith} }\\
\noindent /{\texttoptiebar{t\textbeltl}}{\textprimstress}uri{\texttheta}/\\
\noindent lit. tar+air\\


\noindent History:

\vspace{-0pt}
\hspace{40pt}
\begin{tabular}{ccc}
\textit{1000} & /inga{\texttoptiebar{t\textbeltl}}unarai{\texttheta}/&$\rightarrow$ & \textit{3200} & /ing{\texttoptiebar{t\textbeltl}}unrai{\texttheta}/&$\rightarrow$ & \textit{5500} & /ing{\texttoptiebar{t\textbeltl}}unri{\texttheta}/&$\rightarrow$ & \textit{6000} & /i{\ng}g{\texttoptiebar{t\textbeltl}}unri{\texttheta}/&$\rightarrow$ & \textit{7500} & /{\ng}g{\texttoptiebar{t\textbeltl}}unri{\texttheta}/&$\rightarrow$ & \textit{7501} & /g{\texttoptiebar{t\textbeltl}}unri{\texttheta}/&$\rightarrow$ & \textit{7600} & /g{\texttoptiebar{t\textbeltl}}uri{\texttheta}/&$\rightarrow$ & \textit{11501} & /{\texttoptiebar{t\textbeltl}}uri{\texttheta}/& \\
\end{tabular}

\vspace{20pt}\hline

\end{nopagebreak}
\filbreak



\vspace{15pt}
\begin{nopagebreak}
\noindent{\fontsize{20pt}{10pt}\textbf{syaleth} } \textit{food} (N)\\
\noindent {\tovian \fontsize{20pt}{10pt} \textbf{syaleth} }\\
\noindent /sy{\textprimstress}ale{\texttheta}/\\
\noindent lit. INS+eating\\


\noindent History:

\vspace{-0pt}
\hspace{40pt}
\begin{tabular}{ccc}
\textit{10} & /si-ale{\dh}a/&$\rightarrow$ & \textit{1000} & /siale{\dh}a/&$\rightarrow$ & \textit{2201} & /siyale{\dh}a/&$\rightarrow$ & \textit{3200} & /syale{\dh}a/&$\rightarrow$ & \textit{9500} & /syale{\dh}/&$\rightarrow$ & \textit{12002} & /syale{\texttheta}/& \\
\end{tabular}

\vspace{20pt}\hline

\end{nopagebreak}
\filbreak



\vspace{15pt}
\begin{nopagebreak}
\noindent{\fontsize{20pt}{10pt}\textbf{shot} } \textit{foot} (N)\\
\noindent {\tovian \fontsize{20pt}{10pt} \textbf{shot} }\\
\noindent /{\textesh}{\textprimstress}o{\textsubbridge{t}}/\\
\noindent \textit{_   }\\


\noindent History:

\vspace{-0pt}
\hspace{40pt}
\begin{tabular}{ccc}
\textit{0} & /{\textesh}o{\textsubbridge{t}}/& \\
\end{tabular}

\vspace{20pt}\hline

\end{nopagebreak}
\filbreak



\vspace{15pt}
\begin{nopagebreak}
\noindent{\fontsize{20pt}{10pt}\textbf{lethil} } \textit{forest} (N)\\
\noindent {\tovian \fontsize{20pt}{10pt} \textbf{lethil} }\\
\noindent /l{\textprimstress}e{\texttheta}il/\\
\noindent lit. tree+place\\


\noindent History:

\vspace{-0pt}
\hspace{40pt}
\begin{tabular}{ccc}
\textit{1000} & /aklare{\texttheta}ile/&$\rightarrow$ & \textit{2000} & /aglare{\texttheta}ile/&$\rightarrow$ & \textit{3200} & /aglre{\texttheta}ile/&$\rightarrow$ & \textit{6503} & /agrle{\texttheta}ile/&$\rightarrow$ & \textit{7500} & /grle{\texttheta}ile/&$\rightarrow$ & \textit{7501} & /rle{\texttheta}ile/&$\rightarrow$ & \textit{7600} & /le{\texttheta}ile/&$\rightarrow$ & \textit{9500} & /le{\texttheta}il/& \\
\end{tabular}

\vspace{20pt}\hline

\end{nopagebreak}
\filbreak



\vspace{15pt}
\begin{nopagebreak}
\noindent{\fontsize{20pt}{10pt}\textbf{tod} } \textit{forever/eternal/permanent} (N)\\
\noindent {\tovian \fontsize{20pt}{10pt} \textbf{tod} }\\
\noindent /{\textsubbridge{t}}{\textprimstress}od/\\


\noindent History:

\vspace{-0pt}
\hspace{40pt}
\begin{tabular}{ccc}
\textit{0} & /{\textsubbridge{t}}oda/&$\rightarrow$ & \textit{9500} & /{\textsubbridge{t}}od/& \\
\end{tabular}

\vspace{20pt}\hline

\end{nopagebreak}
\filbreak



\vspace{15pt}
\begin{nopagebreak}
\noindent{\fontsize{20pt}{10pt}\textbf{tev} } \textit{forty eight/48} (N)\\
\noindent {\tovian \fontsize{20pt}{10pt} \textbf{tev} }\\
\noindent /{\textsubbridge{t}}{\textprimstress}ev/\\
\noindent lit. four+twelve\\


\noindent History:

\vspace{-0pt}
\hspace{40pt}
\begin{tabular}{ccc}
\textit{4000} & /{\textsubbridge{t}}ev{\textbeltl}ya/&$\rightarrow$ & \textit{4502} & /{\textsubbridge{t}}evya/&$\rightarrow$ & \textit{9500} & /{\textsubbridge{t}}evy/&$\rightarrow$ & \textit{13001} & /{\textsubbridge{t}}ev/& \\
\end{tabular}

\vspace{20pt}\hline

\end{nopagebreak}
\filbreak



\vspace{15pt}
\begin{nopagebreak}
\noindent{\fontsize{20pt}{10pt}\textbf{tev} } \textit{four/4} (N)\\
\noindent {\tovian \fontsize{20pt}{10pt} \textbf{tev} }\\
\noindent /{\textsubbridge{t}}{\textprimstress}ev/\\


\noindent History:

\vspace{-0pt}
\hspace{40pt}
\begin{tabular}{ccc}
\textit{0} & /{\textsubbridge{t}}ev/& \\
\end{tabular}

\vspace{20pt}\hline

\end{nopagebreak}
\filbreak



\vspace{15pt}
\begin{nopagebreak}
\noindent{\fontsize{20pt}{10pt}\textbf{lhyalo} } \textit{fourteen/14} (N)\\
\noindent {\tovian \fontsize{20pt}{10pt} \textbf{lhyalo} }\\
\noindent /{\textbeltl}y{\textprimstress}alo/\\
\noindent lit. twelve+two\\


\noindent History:

\vspace{-0pt}
\hspace{40pt}
\begin{tabular}{ccc}
\textit{4000} & /{\textbeltl}yaloe/&$\rightarrow$ & \textit{9500} & /{\textbeltl}yalo/& \\
\end{tabular}

\vspace{20pt}\hline

\end{nopagebreak}
\filbreak



\vspace{15pt}
\begin{nopagebreak}
\noindent{\fontsize{20pt}{10pt}\textbf{yelalh} } \textit{friend/ally/friendship} (N)\\
\noindent {\tovian \fontsize{20pt}{10pt} \textbf{yelalh} }\\
\noindent /y{\textprimstress}ela{\textbeltl}/\\
\noindent lit. GEN+person+goodness\\


\noindent History:

\vspace{-0pt}
\hspace{40pt}
\begin{tabular}{ccc}
\textit{0} & /hi-driela{\textbeltl}um/&$\rightarrow$ & \textit{1000} & /hidriela{\textbeltl}um/&$\rightarrow$ & \textit{2201} & /hidriyela{\textbeltl}um/&$\rightarrow$ & \textit{3200} & /hdryela{\textbeltl}m/&$\rightarrow$ & \textit{4500} & /hryela{\textbeltl}m/&$\rightarrow$ & \textit{7501} & /ryela{\textbeltl}m/&$\rightarrow$ & \textit{7600} & /yela{\textbeltl}m/&$\rightarrow$ & \textit{13001} & /yela{\textbeltl}/& \\
\end{tabular}

\vspace{20pt}\hline

\end{nopagebreak}
\filbreak



\vspace{15pt}
\begin{nopagebreak}
\noindent{\fontsize{20pt}{10pt}\textbf{heth} } \textit{front} (N)\\
\noindent {\tovian \fontsize{20pt}{10pt} \textbf{heth} }\\
\noindent /ħ{\textprimstress}e{\texttheta}/\\


\noindent History:

\vspace{-0pt}
\hspace{40pt}
\begin{tabular}{ccc}
\textit{0} & /he{\texttheta}/&$\rightarrow$ & \textit{14006} & /ħe{\texttheta}/& \\
\end{tabular}

\vspace{20pt}\hline

\end{nopagebreak}
\filbreak



\vspace{15pt}
\begin{nopagebreak}
\noindent{\fontsize{20pt}{10pt}\textbf{laon} } \textit{frown} (N)\\
\noindent {\tovian \fontsize{20pt}{10pt} \textbf{laon} }\\
\noindent /l{\textprimstress}aon/\\
\noindent lit. sad+mouth\\


\noindent History:

\vspace{-0pt}
\hspace{40pt}
\begin{tabular}{ccc}
\textit{400} & /luraeolan/&$\rightarrow$ & \textit{3200} & /lraeoln/&$\rightarrow$ & \textit{6503} & /rlaeoln/&$\rightarrow$ & \textit{7600} & /laeon/&$\rightarrow$ & \textit{10000} & /laon/& \\
\end{tabular}

\vspace{20pt}\hline

\end{nopagebreak}
\filbreak



\vspace{15pt}
\begin{nopagebreak}
\noindent{\fontsize{20pt}{10pt}\textbf{shol} } \textit{fruit} (N)\\
\noindent {\tovian \fontsize{20pt}{10pt} \textbf{shol} }\\
\noindent /{\textesh}{\textprimstress}ol/\\


\noindent History:

\vspace{-0pt}
\hspace{40pt}
\begin{tabular}{ccc}
\textit{0} & /{\textyogh}ola/&$\rightarrow$ & \textit{9500} & /{\textyogh}ol/&$\rightarrow$ & \textit{12001} & /{\textesh}ol/& \\
\end{tabular}

\vspace{20pt}\hline

\end{nopagebreak}
\filbreak



\vspace{15pt}
\begin{nopagebreak}
\noindent{\fontsize{20pt}{10pt}\textbf{remhad} } \textit{full} (N)\\
\noindent {\tovian \fontsize{20pt}{10pt} \textbf{remhad} }\\
\noindent /r{\textprimstress}emħad/\\
\noindent lit. empty+not\\


\noindent History:

\vspace{-0pt}
\hspace{40pt}
\begin{tabular}{ccc}
\textit{0} & /oneremihado/&$\rightarrow$ & \textit{3200} & /onremhado/&$\rightarrow$ & \textit{7500} & /nremhado/&$\rightarrow$ & \textit{7600} & /remhado/&$\rightarrow$ & \textit{9500} & /remhad/&$\rightarrow$ & \textit{14006} & /remħad/& \\
\end{tabular}

\vspace{20pt}\hline

\end{nopagebreak}
\filbreak



\vspace{15pt}
\begin{nopagebreak}
\noindent{\fontsize{20pt}{10pt}\textbf{heth} } \textit{future} (Ns)\\
\noindent {\tovian \fontsize{20pt}{10pt} \textbf{heth} }\\
\noindent /ħ{\textprimstress}e{\texttheta}/\\
\noindent lit. front+path\\


\noindent History:

\vspace{-0pt}
\hspace{40pt}
\begin{tabular}{ccc}
\textit{10} & /he{\texttheta}num/&$\rightarrow$ & \textit{3200} & /he{\texttheta}nm/&$\rightarrow$ & \textit{7600} & /he{\texttheta}m/&$\rightarrow$ & \textit{13001} & /he{\texttheta}/&$\rightarrow$ & \textit{14006} & /ħe{\texttheta}/& \\
\end{tabular}

\vspace{20pt}\hline

\end{nopagebreak}
\filbreak



\vspace{15pt}
\begin{nopagebreak}
\noindent{\fontsize{20pt}{10pt}\textbf{hi-} } \textit{GEN/genitive case marker} (CASE)\\
\noindent {\tovian \fontsize{20pt}{10pt} \textbf{hi-} }\\
\noindent /h{\textprimstress}i/\\


\noindent History:

\vspace{-0pt}
\hspace{40pt}
\begin{tabular}{ccc}
\textit{0} & /hi-/& \\
\end{tabular}

\vspace{20pt}\hline

\end{nopagebreak}
\filbreak



\vspace{15pt}
\begin{nopagebreak}
\noindent{\fontsize{20pt}{10pt}\textbf{elakorim} } \textit{god of memory/cycle} (N)\\
\noindent {\tovian \fontsize{20pt}{10pt} \textbf{elakorim} }\\
\noindent /elak{\textprimstress}orim/\\


\noindent History:

\vspace{-0pt}
\hspace{40pt}
\begin{tabular}{ccc}
\textit{0} & /elakorim/& \\
\end{tabular}

\vspace{20pt}\hline

\end{nopagebreak}
\filbreak



\vspace{15pt}
\begin{nopagebreak}
\noindent{\fontsize{20pt}{10pt}\textbf{lhum} } \textit{goodness/virtue} (N)\\
\noindent {\tovian \fontsize{20pt}{10pt} \textbf{lhum} }\\
\noindent /{\textbeltl}{\textprimstress}um/\\


\noindent History:

\vspace{-0pt}
\hspace{40pt}
\begin{tabular}{ccc}
\textit{0} & /{\textbeltl}um/& \\
\end{tabular}

\vspace{20pt}\hline

\end{nopagebreak}
\filbreak



\vspace{15pt}
\begin{nopagebreak}
\noindent{\fontsize{20pt}{10pt}\textbf{meloresyaleth} } \textit{grain} (N)\\
\noindent {\tovian \fontsize{20pt}{10pt} \textbf{meloresyaleth} }\\
\noindent /meloresy{\textprimstress}ale{\texttheta}/\\
\noindent lit. plant+food\\


\noindent History:

\vspace{-0pt}
\hspace{40pt}
\begin{tabular}{ccc}
\textit{6000} & /mlloresyale{\dh}a/&$\rightarrow$ & \textit{6500} & /mloresyale{\dh}a/&$\rightarrow$ & \textit{7502} & /meloresyale{\dh}a/&$\rightarrow$ & \textit{9500} & /meloresyale{\dh}/&$\rightarrow$ & \textit{12002} & /meloresyale{\texttheta}/& \\
\end{tabular}

\vspace{20pt}\hline

\end{nopagebreak}
\filbreak



\vspace{15pt}
\begin{nopagebreak}
\noindent{\fontsize{20pt}{10pt}\textbf{nan} } \textit{grandfather} (N)\\
\noindent {\tovian \fontsize{20pt}{10pt} \textbf{nan} }\\
\noindent /n{\textprimstress}an/\\
\noindent lit. old+father\\


\noindent History:

\vspace{-0pt}
\hspace{40pt}
\begin{tabular}{ccc}
\textit{642} & /nanpander/&$\rightarrow$ & \textit{2000} & /nanbander/&$\rightarrow$ & \textit{3200} & /nnbandr/&$\rightarrow$ & \textit{6500} & /nbandr/&$\rightarrow$ & \textit{8500} & /nbanr/&$\rightarrow$ & \textit{8750} & /nanr/&$\rightarrow$ & \textit{13001} & /nan/& \\
\end{tabular}

\vspace{20pt}\hline

\end{nopagebreak}
\filbreak



\vspace{15pt}
\begin{nopagebreak}
\noindent{\fontsize{20pt}{10pt}\textbf{lothr} } \textit{grasp/grabbing} (N)\\
\noindent {\tovian \fontsize{20pt}{10pt} \textbf{lothr} }\\
\noindent /l{\textprimstress}o{\texttheta}r/\\


\noindent History:

\vspace{-0pt}
\hspace{40pt}
\begin{tabular}{ccc}
\textit{0} & /lo{\texttheta}ir/&$\rightarrow$ & \textit{3200} & /lo{\texttheta}r/& \\
\end{tabular}

\vspace{20pt}\hline

\end{nopagebreak}
\filbreak



\vspace{15pt}
\begin{nopagebreak}
\noindent{\fontsize{20pt}{10pt}\textbf{klat} } \textit{greenness/green} (N)\\
\noindent {\tovian \fontsize{20pt}{10pt} \textbf{klat} }\\
\noindent /kl{\textprimstress}a{\textsubbridge{t}}/\\


\noindent History:

\vspace{-0pt}
\hspace{40pt}
\begin{tabular}{ccc}
\textit{0} & /kla{\textsubbridge{t}}/& \\
\end{tabular}

\vspace{20pt}\hline

\end{nopagebreak}
\filbreak



\vspace{15pt}
\begin{nopagebreak}
\noindent{\fontsize{20pt}{10pt}\textbf{amal} } \textit{hair/person hair} (N)\\
\noindent {\tovian \fontsize{20pt}{10pt} \textbf{amal} }\\
\noindent /{\textprimstress}amal/\\


\noindent History:

\vspace{-0pt}
\hspace{40pt}
\begin{tabular}{ccc}
\textit{0} & /amalur/&$\rightarrow$ & \textit{3200} & /amalr/&$\rightarrow$ & \textit{6503} & /amarl/&$\rightarrow$ & \textit{7600} & /amal/& \\
\end{tabular}

\vspace{20pt}\hline

\end{nopagebreak}
\filbreak



\vspace{15pt}
\begin{nopagebreak}
\noindent{\fontsize{20pt}{10pt}\textbf{thang} } \textit{hand/leaf} (N)\\
\noindent {\tovian \fontsize{20pt}{10pt} \textbf{thang} }\\
\noindent /{\texttheta}{\textprimstress}a{\ng}/\\


\noindent History:

\vspace{-0pt}
\hspace{40pt}
\begin{tabular}{ccc}
\textit{0} & /a{\dh}ange/&$\rightarrow$ & \textit{6000} & /a{\dh}a{\ng}ge/&$\rightarrow$ & \textit{7500} & /{\dh}a{\ng}ge/&$\rightarrow$ & \textit{8750} & /{\dh}a{\ng}e/&$\rightarrow$ & \textit{9500} & /{\dh}a{\ng}/&$\rightarrow$ & \textit{12002} & /{\texttheta}a{\ng}/& \\
\end{tabular}

\vspace{20pt}\hline

\end{nopagebreak}
\filbreak



\vspace{15pt}
\begin{nopagebreak}
\noindent{\fontsize{20pt}{10pt}\textbf{nither} } \textit{hanging} (N)\\
\noindent {\tovian \fontsize{20pt}{10pt} \textbf{nither} }\\
\noindent /n{\textprimstress}i{\texttheta}er/\\
\noindent lit. putting+up\\


\noindent History:

\vspace{-0pt}
\hspace{40pt}
\begin{tabular}{ccc}
\textit{0} & /hemmani{\texttheta}era/&$\rightarrow$ & \textit{3200} & /hmmni{\texttheta}era/&$\rightarrow$ & \textit{6500} & /hmni{\texttheta}era/&$\rightarrow$ & \textit{7501} & /mni{\texttheta}era/&$\rightarrow$ & \textit{7600} & /ni{\texttheta}era/&$\rightarrow$ & \textit{9500} & /ni{\texttheta}er/& \\
\end{tabular}

\vspace{20pt}\hline

\end{nopagebreak}
\filbreak



\vspace{15pt}
\begin{nopagebreak}
\noindent{\fontsize{20pt}{10pt}\textbf{mat} } \textit{happy} (N)\\
\noindent {\tovian \fontsize{20pt}{10pt} \textbf{mat} }\\
\noindent /m{\textprimstress}a{\textsubbridge{t}}/\\


\noindent History:

\vspace{-0pt}
\hspace{40pt}
\begin{tabular}{ccc}
\textit{0} & /ma{\textsubbridge{t}}{\textsubbridge{t}}a/&$\rightarrow$ & \textit{3000} & /ma{\textsubbridge{t}}a/&$\rightarrow$ & \textit{9500} & /ma{\textsubbridge{t}}/& \\
\end{tabular}

\vspace{20pt}\hline

\end{nopagebreak}
\filbreak



\vspace{15pt}
\begin{nopagebreak}
\noindent{\fontsize{20pt}{10pt}\textbf{lha} } \textit{hardness/strength} (N)\\
\noindent {\tovian \fontsize{20pt}{10pt} \textbf{lha} }\\
\noindent /{\textbeltl}{\textprimstress}a/\\


\noindent History:

\vspace{-0pt}
\hspace{40pt}
\begin{tabular}{ccc}
\textit{0} & /{\textbeltl}aa/&$\rightarrow$ & \textit{5500} & /{\textbeltl}a/& \\
\end{tabular}

\vspace{20pt}\hline

\end{nopagebreak}
\filbreak



\vspace{15pt}
\begin{nopagebreak}
\noindent{\fontsize{20pt}{10pt}\textbf{lesh} } \textit{harvest} (N)\\
\noindent {\tovian \fontsize{20pt}{10pt} \textbf{lesh} }\\
\noindent /l{\textprimstress}e{\textesh}/\\
\noindent lit. plant+cutting\\


\noindent History:

\vspace{-0pt}
\hspace{40pt}
\begin{tabular}{ccc}
\textit{1000} & /{\textschwa}mellore{\textyogh}i{\texttheta}/&$\rightarrow$ & \textit{2300} & /mellore{\textyogh}i{\texttheta}/&$\rightarrow$ & \textit{3200} & /mllre{\textyogh}{\texttheta}/&$\rightarrow$ & \textit{4502} & /mllre{\textyogh}/&$\rightarrow$ & \textit{6500} & /mlre{\textyogh}/&$\rightarrow$ & \textit{6503} & /mrle{\textyogh}/&$\rightarrow$ & \textit{7501} & /rle{\textyogh}/&$\rightarrow$ & \textit{7600} & /le{\textyogh}/&$\rightarrow$ & \textit{12001} & /le{\textesh}/& \\
\end{tabular}

\vspace{20pt}\hline

\end{nopagebreak}
\filbreak



\vspace{15pt}
\begin{nopagebreak}
\noindent{\fontsize{20pt}{10pt}\textbf{drunan} } \textit{hat} (N)\\
\noindent {\tovian \fontsize{20pt}{10pt} \textbf{drunan} }\\
\noindent /dr{\textprimstress}unan/\\
\noindent lit. head+covering\\


\noindent History:

\vspace{-0pt}
\hspace{40pt}
\begin{tabular}{ccc}
\textit{3024} & /darunbenani/&$\rightarrow$ & \textit{3200} & /drunbnani/&$\rightarrow$ & \textit{8500} & /drunnani/&$\rightarrow$ & \textit{9500} & /drunnan/&$\rightarrow$ & \textit{12007} & /drunan/& \\
\end{tabular}

\vspace{20pt}\hline

\end{nopagebreak}
\filbreak



\vspace{15pt}
\begin{nopagebreak}
\noindent{\fontsize{20pt}{10pt}\textbf{ta} } \textit{he/she/them} (P)\\
\noindent {\tovian \fontsize{20pt}{10pt} \textbf{ta} }\\
\noindent /{\textsubbridge{t}}{\textprimstress}a/\\


\noindent History:

\vspace{-0pt}
\hspace{40pt}
\begin{tabular}{ccc}
\textit{0} & /{\textsubbridge{t}}a/& \\
\end{tabular}

\vspace{20pt}\hline

\end{nopagebreak}
\filbreak



\vspace{15pt}
\begin{nopagebreak}
\noindent{\fontsize{20pt}{10pt}\textbf{dan} } \textit{head} (N)\\
\noindent {\tovian \fontsize{20pt}{10pt} \textbf{dan} }\\
\noindent /d{\textprimstress}an/\\


\noindent History:

\vspace{-0pt}
\hspace{40pt}
\begin{tabular}{ccc}
\textit{0} & /darun/&$\rightarrow$ & \textit{3200} & /darn/&$\rightarrow$ & \textit{7600} & /dan/& \\
\end{tabular}

\vspace{20pt}\hline

\end{nopagebreak}
\filbreak



\vspace{15pt}
\begin{nopagebreak}
\noindent{\fontsize{20pt}{10pt}\textbf{dum} } \textit{hiding} ()\\
\noindent {\tovian \fontsize{20pt}{10pt} \textbf{dum} }\\
\noindent /d{\textprimstress}um/\\
\noindent lit. bbu\\


\noindent History:

\vspace{-0pt}
\hspace{40pt}
\begin{tabular}{ccc}
\textit{220} & /bubu/&$\rightarrow$ & \textit{9500} & /bub/&$\rightarrow$ & \textit{11000} & /bum/&$\rightarrow$ & \textit{11001} & /dum/& \\
\end{tabular}

\vspace{20pt}\hline

\end{nopagebreak}
\filbreak



\vspace{15pt}
\begin{nopagebreak}
\noindent{\fontsize{20pt}{10pt}\textbf{leth} } \textit{hoe} (N)\\
\noindent {\tovian \fontsize{20pt}{10pt} \textbf{leth} }\\
\noindent /l{\textprimstress}e{\texttheta}/\\
\noindent lit. plant+tool\\


\noindent History:

\vspace{-0pt}
\hspace{40pt}
\begin{tabular}{ccc}
\textit{2000} & /{\textschwa}mellore{\dh}an/&$\rightarrow$ & \textit{2300} & /mellore{\dh}an/&$\rightarrow$ & \textit{3200} & /mllre{\dh}n/&$\rightarrow$ & \textit{6500} & /mlre{\dh}n/&$\rightarrow$ & \textit{6503} & /mrle{\dh}n/&$\rightarrow$ & \textit{7501} & /rle{\dh}n/&$\rightarrow$ & \textit{7600} & /le{\dh}n/&$\rightarrow$ & \textit{12002} & /le{\texttheta}n/&$\rightarrow$ & \textit{13001} & /le{\texttheta}/& \\
\end{tabular}

\vspace{20pt}\hline

\end{nopagebreak}
\filbreak



\vspace{15pt}
\begin{nopagebreak}
\noindent{\fontsize{20pt}{10pt}\textbf{ag} } \textit{hotness/heat/warmth} (N)\\
\noindent {\tovian \fontsize{20pt}{10pt} \textbf{ag} }\\
\noindent /{\textprimstress}ag/\\


\noindent History:

\vspace{-0pt}
\hspace{40pt}
\begin{tabular}{ccc}
\textit{0} & /asaka/&$\rightarrow$ & \textit{2000} & /asaga/&$\rightarrow$ & \textit{5100} & /aaga/&$\rightarrow$ & \textit{5500} & /aga/&$\rightarrow$ & \textit{9500} & /ag/& \\
\end{tabular}

\vspace{20pt}\hline

\end{nopagebreak}
\filbreak



\vspace{15pt}
\begin{nopagebreak}
\noindent{\fontsize{20pt}{10pt}\textbf{yethil} } \textit{house/home} (N)\\
\noindent {\tovian \fontsize{20pt}{10pt} \textbf{yethil} }\\
\noindent /y{\textprimstress}e{\texttheta}il/\\
\noindent lit. person+place\\


\noindent History:

\vspace{-0pt}
\hspace{40pt}
\begin{tabular}{ccc}
\textit{10} & /driela{\texttheta}ile/&$\rightarrow$ & \textit{2201} & /driyela{\texttheta}ile/&$\rightarrow$ & \textit{3200} & /dryel{\texttheta}ile/&$\rightarrow$ & \textit{7501} & /ryel{\texttheta}ile/&$\rightarrow$ & \textit{7600} & /yel{\texttheta}ile/&$\rightarrow$ & \textit{9500} & /yel{\texttheta}il/&$\rightarrow$ & \textit{12502} & /ye{\texttheta}il/& \\
\end{tabular}

\vspace{20pt}\hline

\end{nopagebreak}
\filbreak



\vspace{15pt}
\begin{nopagebreak}
\noindent{\fontsize{20pt}{10pt}\textbf{lhan} } \textit{how} (INT)\\
\noindent {\tovian \fontsize{20pt}{10pt} \textbf{lhan} }\\
\noindent /{\textbeltl}{\textprimstress}an/\\


\noindent History:

\vspace{-0pt}
\hspace{40pt}
\begin{tabular}{ccc}
\textit{0} & /{\textbeltl}an/& \\
\end{tabular}

\vspace{20pt}\hline

\end{nopagebreak}
\filbreak



\vspace{15pt}
\begin{nopagebreak}
\noindent{\fontsize{20pt}{10pt}\textbf{shush} } \textit{ice/frost} (N)\\
\noindent {\tovian \fontsize{20pt}{10pt} \textbf{shush} }\\
\noindent /{\textesh}{\textprimstress}u{\textesh}/\\


\noindent History:

\vspace{-0pt}
\hspace{40pt}
\begin{tabular}{ccc}
\textit{0} & /{\textesh}u{\textesh}a/&$\rightarrow$ & \textit{9500} & /{\textesh}u{\textesh}/& \\
\end{tabular}

\vspace{20pt}\hline

\end{nopagebreak}
\filbreak



\vspace{15pt}
\begin{nopagebreak}
\noindent{\fontsize{20pt}{10pt}\textbf{se-} } \textit{IMP/imperative mood marker} (MOOD)\\
\noindent {\tovian \fontsize{20pt}{10pt} \textbf{se-} }\\
\noindent /s{\textprimstress}e/\\


\noindent History:

\vspace{-0pt}
\hspace{40pt}
\begin{tabular}{ccc}
\textit{0} & /se-/& \\
\end{tabular}

\vspace{20pt}\hline

\end{nopagebreak}
\filbreak



\vspace{15pt}
\begin{nopagebreak}
\noindent{\fontsize{20pt}{10pt}\textbf{gad} } \textit{importance/grandness/greatness} (N)\\
\noindent {\tovian \fontsize{20pt}{10pt} \textbf{gad} }\\
\noindent /g{\textprimstress}ad/\\


\noindent History:

\vspace{-0pt}
\hspace{40pt}
\begin{tabular}{ccc}
\textit{0} & /aga{\textsubbridge{t}}a/&$\rightarrow$ & \textit{2000} & /agada/&$\rightarrow$ & \textit{7500} & /gada/&$\rightarrow$ & \textit{9500} & /gad/& \\
\end{tabular}

\vspace{20pt}\hline

\end{nopagebreak}
\filbreak



\vspace{15pt}
\begin{nopagebreak}
\noindent{\fontsize{20pt}{10pt}\textbf{tid} } \textit{in (obsolete)} (became locative "ti-")\\
\noindent {\tovian \fontsize{20pt}{10pt} \textbf{tid} }\\
\noindent /{\textsubbridge{t}}{\textprimstress}id/\\


\noindent History:

\vspace{-0pt}
\hspace{40pt}
\begin{tabular}{ccc}
\textit{0} & /{\textsubbridge{t}}id/& \\
\end{tabular}

\vspace{20pt}\hline

\end{nopagebreak}
\filbreak



\vspace{15pt}
\begin{nopagebreak}
\noindent{\fontsize{20pt}{10pt}\textbf{lo-} } \textit{inanimate noun class} (CLASS)\\
\noindent {\tovian \fontsize{20pt}{10pt} \textbf{lo-} }\\
\noindent /l{\textprimstress}o/\\


\noindent History:

\vspace{-0pt}
\hspace{40pt}
\begin{tabular}{ccc}
\textit{0} & /lo-/& \\
\end{tabular}

\vspace{20pt}\hline

\end{nopagebreak}
\filbreak



\vspace{15pt}
\begin{nopagebreak}
\noindent{\fontsize{20pt}{10pt}\textbf{si-} } \textit{INS/instrumental case marker} (CASE)\\
\noindent {\tovian \fontsize{20pt}{10pt} \textbf{si-} }\\
\noindent /s{\textprimstress}i/\\


\noindent History:

\vspace{-0pt}
\hspace{40pt}
\begin{tabular}{ccc}
\textit{0} & /si-/& \\
\end{tabular}

\vspace{20pt}\hline

\end{nopagebreak}
\filbreak



\vspace{15pt}
\begin{nopagebreak}
\noindent{\fontsize{20pt}{10pt}\textbf{teshan} } \textit{instrument} (N)\\
\noindent {\tovian \fontsize{20pt}{10pt} \textbf{teshan} }\\
\noindent /{\textsubbridge{t}}{\textprimstress}e{\textesh}an/\\
\noindent lit. sound+tool\\


\noindent History:

\vspace{-0pt}
\hspace{40pt}
\begin{tabular}{ccc}
\textit{1500} & /{\textsubbridge{t}}{\textsubbridge{t}}e{\textyogh}ah{\dh}an/&$\rightarrow$ & \textit{3000} & /{\textsubbridge{t}}e{\textyogh}ah{\dh}an/&$\rightarrow$ & \textit{3200} & /{\textsubbridge{t}}{\textyogh}ah{\dh}n/&$\rightarrow$ & \textit{4502} & /{\textsubbridge{t}}{\textyogh}ahn/&$\rightarrow$ & \textit{6240} & /{\textsubbridge{t}}{\textyogh}an/&$\rightarrow$ & \textit{12001} & /{\textsubbridge{t}}{\textesh}an/&$\rightarrow$ & \textit{12400} & /{\textsubbridge{t}}e{\textesh}an/& \\
\end{tabular}

\vspace{20pt}\hline

\end{nopagebreak}
\filbreak



\vspace{15pt}
\begin{nopagebreak}
\noindent{\fontsize{20pt}{10pt}\textbf{lod} } \textit{it/it (inanimate)} (P)\\
\noindent {\tovian \fontsize{20pt}{10pt} \textbf{lod} }\\
\noindent /l{\textprimstress}od/\\


\noindent History:

\vspace{-0pt}
\hspace{40pt}
\begin{tabular}{ccc}
\textit{0} & /lo{\textsubbridge{t}}a/&$\rightarrow$ & \textit{2000} & /loda/&$\rightarrow$ & \textit{9500} & /lod/& \\
\end{tabular}

\vspace{20pt}\hline

\end{nopagebreak}
\filbreak



\vspace{15pt}
\begin{nopagebreak}
\noindent{\fontsize{20pt}{10pt}\textbf{loetush} } \textit{January} (N)\\
\noindent {\tovian \fontsize{20pt}{10pt} \textbf{loetush} }\\
\noindent /lo{\textprimstress}e{\textsubbridge{t}}u{\textesh}/\\
\noindent lit. two+cold\\
\noindent \textit{Frohsten}\\


\noindent History:

\vspace{-0pt}
\hspace{40pt}
\begin{tabular}{ccc}
\textit{6000} & /loe{\textsubbridge{t}}u{\textesh}u/&$\rightarrow$ & \textit{9500} & /loe{\textsubbridge{t}}u{\textesh}/& \\
\end{tabular}

\vspace{20pt}\hline

\end{nopagebreak}
\filbreak



\vspace{15pt}
\begin{nopagebreak}
\noindent{\fontsize{20pt}{10pt}\textbf{lad} } \textit{jealousy} (N)\\
\noindent {\tovian \fontsize{20pt}{10pt} \textbf{lad} }\\
\noindent /l{\textprimstress}ad/\\
\noindent lit. wickedness+comparison\\


\noindent History:

\vspace{-0pt}
\hspace{40pt}
\begin{tabular}{ccc}
\textit{402} & /oyalroladil/&$\rightarrow$ & \textit{2100} & /yalroladil/&$\rightarrow$ & \textit{3200} & /ylrladl/&$\rightarrow$ & \textit{6503} & /yrlladl/&$\rightarrow$ & \textit{7501} & /rlladl/&$\rightarrow$ & \textit{7600} & /lladl/&$\rightarrow$ & \textit{12007} & /ladl/&$\rightarrow$ & \textit{13001} & /lad/& \\
\end{tabular}

\vspace{20pt}\hline

\end{nopagebreak}
\filbreak



\vspace{15pt}
\begin{nopagebreak}
\noindent{\fontsize{20pt}{10pt}\textbf{tlun} } \textit{joining} (N)\\
\noindent {\tovian \fontsize{20pt}{10pt} \textbf{tlun} }\\
\noindent /{\texttoptiebar{t\textbeltl}}{\textprimstress}un/\\


\noindent History:

\vspace{-0pt}
\hspace{40pt}
\begin{tabular}{ccc}
\textit{0} & /{\texttoptiebar{t\textbeltl}}una/&$\rightarrow$ & \textit{9500} & /{\texttoptiebar{t\textbeltl}}un/& \\
\end{tabular}

\vspace{20pt}\hline

\end{nopagebreak}
\filbreak



\vspace{15pt}
\begin{nopagebreak}
\noindent{\fontsize{20pt}{10pt}\textbf{loethulef} } \textit{July} (N)\\
\noindent {\tovian \fontsize{20pt}{10pt} \textbf{loethulef} }\\
\noindent /loe{\texttheta}{\textprimstress}ulef/\\
\noindent lit. two+height+sun\\
\noindent \textit{Thermiblase}\\


\noindent History:

\vspace{-0pt}
\hspace{40pt}
\begin{tabular}{ccc}
\textit{6000} & /loe{\dh}ulefl/&$\rightarrow$ & \textit{12002} & /loe{\texttheta}ulefl/&$\rightarrow$ & \textit{12500} & /loe{\texttheta}ulef/& \\
\end{tabular}

\vspace{20pt}\hline

\end{nopagebreak}
\filbreak



\vspace{15pt}
\begin{nopagebreak}
\noindent{\fontsize{20pt}{10pt}\textbf{lothulef} } \textit{June} (N)\\
\noindent {\tovian \fontsize{20pt}{10pt} \textbf{lothulef} }\\
\noindent /lo{\texttheta}{\textprimstress}ulef/\\
\noindent lit. one+height+sun\\
\noindent \textit{Germspar}\\


\noindent History:

\vspace{-0pt}
\hspace{40pt}
\begin{tabular}{ccc}
\textit{6000} & /lo{\dh}ulefl/&$\rightarrow$ & \textit{12002} & /lo{\texttheta}ulefl/&$\rightarrow$ & \textit{12500} & /lo{\texttheta}ulef/& \\
\end{tabular}

\vspace{20pt}\hline

\end{nopagebreak}
\filbreak



\vspace{15pt}
\begin{nopagebreak}
\noindent{\fontsize{20pt}{10pt}\textbf{vek} } \textit{killing} (N)\\
\noindent {\tovian \fontsize{20pt}{10pt} \textbf{vek} }\\
\noindent /v{\textprimstress}ek/\\


\noindent History:

\vspace{-0pt}
\hspace{40pt}
\begin{tabular}{ccc}
\textit{0} & /vekh/&$\rightarrow$ & \textit{5000} & /vek/& \\
\end{tabular}

\vspace{20pt}\hline

\end{nopagebreak}
\filbreak



\vspace{15pt}
\begin{nopagebreak}
\noindent{\fontsize{20pt}{10pt}\textbf{rish} } \textit{kindness} (N)\\
\noindent {\tovian \fontsize{20pt}{10pt} \textbf{rish} }\\
\noindent /r{\textprimstress}i{\textesh}/\\


\noindent History:

\vspace{-0pt}
\hspace{40pt}
\begin{tabular}{ccc}
\textit{0} & /rai{\textesh}el/&$\rightarrow$ & \textit{3200} & /rai{\textesh}l/&$\rightarrow$ & \textit{5500} & /ri{\textesh}l/&$\rightarrow$ & \textit{12500} & /ri{\textesh}/& \\
\end{tabular}

\vspace{20pt}\hline

\end{nopagebreak}
\filbreak



\vspace{15pt}
\begin{nopagebreak}
\noindent{\fontsize{20pt}{10pt}\textbf{nashith} } \textit{knife} (N)\\
\noindent {\tovian \fontsize{20pt}{10pt} \textbf{nashith} }\\
\noindent /n{\textprimstress}a{\textesh}i{\texttheta}/\\
\noindent lit. small+sharpness+tool\\


\noindent History:

\vspace{-0pt}
\hspace{40pt}
\begin{tabular}{ccc}
\textit{50} & /nina{\textyogh}i{\texttheta}{\dh}an/&$\rightarrow$ & \textit{3200} & /nna{\textyogh}i{\texttheta}{\dh}n/&$\rightarrow$ & \textit{4502} & /nna{\textyogh}i{\texttheta}n/&$\rightarrow$ & \textit{6500} & /na{\textyogh}i{\texttheta}n/&$\rightarrow$ & \textit{12001} & /na{\textesh}i{\texttheta}n/&$\rightarrow$ & \textit{13001} & /na{\textesh}i{\texttheta}/& \\
\end{tabular}

\vspace{20pt}\hline

\end{nopagebreak}
\filbreak



\vspace{15pt}
\begin{nopagebreak}
\noindent{\fontsize{20pt}{10pt}\textbf{lhar} } \textit{knowledge} (N)\\
\noindent {\tovian \fontsize{20pt}{10pt} \textbf{lhar} }\\
\noindent /{\textbeltl}{\textprimstress}ar/\\


\noindent History:

\vspace{-0pt}
\hspace{40pt}
\begin{tabular}{ccc}
\textit{0} & /{\textbeltl}ara/&$\rightarrow$ & \textit{9500} & /{\textbeltl}ar/& \\
\end{tabular}

\vspace{20pt}\hline

\end{nopagebreak}
\filbreak



\vspace{15pt}
\begin{nopagebreak}
\noindent{\fontsize{20pt}{10pt}\textbf{tel} } \textit{lake} (N)\\
\noindent {\tovian \fontsize{20pt}{10pt} \textbf{tel} }\\
\noindent /{\textsubbridge{t}}{\textprimstress}el/\\


\noindent History:

\vspace{-0pt}
\hspace{40pt}
\begin{tabular}{ccc}
\textit{0} & /o{\textsubbridge{t}}kele/&$\rightarrow$ & \textit{3000} & /o{\textsubbridge{t}}ele/&$\rightarrow$ & \textit{7500} & /{\textsubbridge{t}}ele/&$\rightarrow$ & \textit{9500} & /{\textsubbridge{t}}el/& \\
\end{tabular}

\vspace{20pt}\hline

\end{nopagebreak}
\filbreak



\vspace{15pt}
\begin{nopagebreak}
\noindent{\fontsize{20pt}{10pt}\textbf{fal} } \textit{leader} (N)\\
\noindent {\tovian \fontsize{20pt}{10pt} \textbf{fal} }\\
\noindent /f{\textprimstress}al/\\


\noindent History:

\vspace{-0pt}
\hspace{40pt}
\begin{tabular}{ccc}
\textit{0} & /falor/&$\rightarrow$ & \textit{3200} & /falr/&$\rightarrow$ & \textit{6503} & /farl/&$\rightarrow$ & \textit{7600} & /fal/& \\
\end{tabular}

\vspace{20pt}\hline

\end{nopagebreak}
\filbreak



\vspace{15pt}
\begin{nopagebreak}
\noindent{\fontsize{20pt}{10pt}\textbf{hed} } \textit{left} (N)\\
\noindent {\tovian \fontsize{20pt}{10pt} \textbf{hed} }\\
\noindent /ħ{\textprimstress}ed/\\


\noindent History:

\vspace{-0pt}
\hspace{40pt}
\begin{tabular}{ccc}
\textit{0} & /he{\textsubbridge{t}}a/&$\rightarrow$ & \textit{2000} & /heda/&$\rightarrow$ & \textit{9500} & /hed/&$\rightarrow$ & \textit{14006} & /ħed/& \\
\end{tabular}

\vspace{20pt}\hline

\end{nopagebreak}
\filbreak



\vspace{15pt}
\begin{nopagebreak}
\noindent{\fontsize{20pt}{10pt}\textbf{kel} } \textit{leg/limb/branch/arm} (N)\\
\noindent {\tovian \fontsize{20pt}{10pt} \textbf{kel} }\\
\noindent /k{\textprimstress}el/\\


\noindent History:

\vspace{-0pt}
\hspace{40pt}
\begin{tabular}{ccc}
\textit{0} & /kelu/&$\rightarrow$ & \textit{9500} & /kel/& \\
\end{tabular}

\vspace{20pt}\hline

\end{nopagebreak}
\filbreak



\vspace{15pt}
\begin{nopagebreak}
\noindent{\fontsize{20pt}{10pt}\textbf{mel} } \textit{life} (N)\\
\noindent {\tovian \fontsize{20pt}{10pt} \textbf{mel} }\\
\noindent /m{\textprimstress}el/\\


\noindent History:

\vspace{-0pt}
\hspace{40pt}
\begin{tabular}{ccc}
\textit{0} & /{\textschwa}mel/&$\rightarrow$ & \textit{2300} & /mel/& \\
\end{tabular}

\vspace{20pt}\hline

\end{nopagebreak}
\filbreak



\vspace{15pt}
\begin{nopagebreak}
\noindent{\fontsize{20pt}{10pt}\textbf{an} } \textit{light} (N)\\
\noindent {\tovian \fontsize{20pt}{10pt} \textbf{an} }\\
\noindent /{\textprimstress}an/\\


\noindent History:

\vspace{-0pt}
\hspace{40pt}
\begin{tabular}{ccc}
\textit{0} & /ahana/&$\rightarrow$ & \textit{5000} & /aana/&$\rightarrow$ & \textit{5500} & /ana/&$\rightarrow$ & \textit{9500} & /an/& \\
\end{tabular}

\vspace{20pt}\hline

\end{nopagebreak}
\filbreak



\vspace{15pt}
\begin{nopagebreak}
\noindent{\fontsize{20pt}{10pt}\textbf{ti-} } \textit{LOC/locative case marker} (CASE)\\
\noindent {\tovian \fontsize{20pt}{10pt} \textbf{ti-} }\\
\noindent /{\textsubbridge{t}}{\textprimstress}i/\\


\noindent History:

\vspace{-0pt}
\hspace{40pt}
\begin{tabular}{ccc}
\textit{0} & /{\textsubbridge{t}}i-/& \\
\end{tabular}

\vspace{20pt}\hline

\end{nopagebreak}
\filbreak



\vspace{15pt}
\begin{nopagebreak}
\noindent{\fontsize{20pt}{10pt}\textbf{sil} } \textit{make/create} (N)\\
\noindent {\tovian \fontsize{20pt}{10pt} \textbf{sil} }\\
\noindent /s{\textprimstress}il/\\


\noindent History:

\vspace{-0pt}
\hspace{40pt}
\begin{tabular}{ccc}
\textit{0} & /silahe{\textsubbridge{t}}/&$\rightarrow$ & \textit{1000} & /silah{\textsubbridge{t}}/&$\rightarrow$ & \textit{3000} & /silah/&$\rightarrow$ & \textit{3200} & /silh/&$\rightarrow$ & \textit{5000} & /sil/& \\
\end{tabular}

\vspace{20pt}\hline

\end{nopagebreak}
\filbreak



\vspace{15pt}
\begin{nopagebreak}
\noindent{\fontsize{20pt}{10pt}\textbf{lor} } \textit{March} (N)\\
\noindent {\tovian \fontsize{20pt}{10pt} \textbf{lor} }\\
\noindent /l{\textprimstress}or/\\
\noindent lit. one+plant\\
\noindent \textit{Vindrem}\\


\noindent History:

\vspace{-0pt}
\hspace{40pt}
\begin{tabular}{ccc}
\textit{6000} & /lomllore/&$\rightarrow$ & \textit{6500} & /lomlore/&$\rightarrow$ & \textit{7600} & /lolore/&$\rightarrow$ & \textit{9200} & /l{\textschwa}lore/&$\rightarrow$ & \textit{9500} & /l{\textschwa}lor/&$\rightarrow$ & \textit{11990} & /llor/&$\rightarrow$ & \textit{12007} & /lor/& \\
\end{tabular}

\vspace{20pt}\hline

\end{nopagebreak}
\filbreak



\vspace{15pt}
\begin{nopagebreak}
\noindent{\fontsize{20pt}{10pt}\textbf{felor} } \textit{May} (N)\\
\noindent {\tovian \fontsize{20pt}{10pt} \textbf{felor} }\\
\noindent /f{\textprimstress}elor/\\
\noindent lit. three+plant\\
\noindent \textit{Luviosh}\\


\noindent History:

\vspace{-0pt}
\hspace{40pt}
\begin{tabular}{ccc}
\textit{6000} & /femmllore/&$\rightarrow$ & \textit{6500} & /femlore/&$\rightarrow$ & \textit{7600} & /felore/&$\rightarrow$ & \textit{9500} & /felor/& \\
\end{tabular}

\vspace{20pt}\hline

\end{nopagebreak}
\filbreak



\vspace{15pt}
\begin{nopagebreak}
\noindent{\fontsize{20pt}{10pt}\textbf{na} } \textit{me/I} (P)\\
\noindent {\tovian \fontsize{20pt}{10pt} \textbf{na} }\\
\noindent /n{\textprimstress}a/\\


\noindent History:

\vspace{-0pt}
\hspace{40pt}
\begin{tabular}{ccc}
\textit{0} & /na/& \\
\end{tabular}

\vspace{20pt}\hline

\end{nopagebreak}
\filbreak



\vspace{15pt}
\begin{nopagebreak}
\noindent{\fontsize{20pt}{10pt}\textbf{athir} } \textit{meat} (N)\\
\noindent {\tovian \fontsize{20pt}{10pt} \textbf{athir} }\\
\noindent /{\textprimstress}a{\texttheta}ir/\\
\noindent lit. tasting+GEN+blood\\


\noindent History:

\vspace{-0pt}
\hspace{40pt}
\begin{tabular}{ccc}
\textit{1000} & /ale{\dh}ahi-driks/&$\rightarrow$ & \textit{1000} & /ale{\dh}ahidriks/&$\rightarrow$ & \textit{3200} & /al{\dh}ahidrks/&$\rightarrow$ & \textit{4501} & /al{\dh}ahidrs/&$\rightarrow$ & \textit{5000} & /al{\dh}aidrs/&$\rightarrow$ & \textit{5500} & /al{\dh}idrs/&$\rightarrow$ & \textit{12002} & /al{\texttheta}idrs/&$\rightarrow$ & \textit{12502} & /a{\texttheta}idrs/&$\rightarrow$ & \textit{13001} & /a{\texttheta}idr/&$\rightarrow$ & \textit{13004} & /a{\texttheta}ir/& \\
\end{tabular}

\vspace{20pt}\hline

\end{nopagebreak}
\filbreak



\vspace{15pt}
\begin{nopagebreak}
\noindent{\fontsize{20pt}{10pt}\textbf{elom} } \textit{memory} (N)\\
\noindent {\tovian \fontsize{20pt}{10pt} \textbf{elom} }\\
\noindent /{\textprimstress}elom/\\


\noindent History:

\vspace{-0pt}
\hspace{40pt}
\begin{tabular}{ccc}
\textit{0} & /elakorim/&$\rightarrow$ & \textit{2000} & /elagorim/&$\rightarrow$ & \textit{3200} & /elgorm/&$\rightarrow$ & \textit{4501} & /elorm/&$\rightarrow$ & \textit{7600} & /elom/& \\
\end{tabular}

\vspace{20pt}\hline

\end{nopagebreak}
\filbreak



\vspace{15pt}
\begin{nopagebreak}
\noindent{\fontsize{20pt}{10pt}\textbf{korim} } \textit{memory-soul} (N)\\
\noindent {\tovian \fontsize{20pt}{10pt} \textbf{korim} }\\
\noindent /k{\textprimstress}orim/\\


\noindent History:

\vspace{-0pt}
\hspace{40pt}
\begin{tabular}{ccc}
\textit{0} & /korim/& \\
\end{tabular}

\vspace{20pt}\hline

\end{nopagebreak}
\filbreak



\vspace{15pt}
\begin{nopagebreak}
\noindent{\fontsize{20pt}{10pt}\textbf{men} } \textit{metal} (N)\\
\noindent {\tovian \fontsize{20pt}{10pt} \textbf{men} }\\
\noindent /m{\textprimstress}en/\\
\noindent lit. tool+GEN+rock\\


\noindent History:

\vspace{-0pt}
\hspace{40pt}
\begin{tabular}{ccc}
\textit{1000} & /{\dh}anhi-meron/&$\rightarrow$ & \textit{1000} & /{\dh}anhimeron/&$\rightarrow$ & \textit{3200} & /{\dh}nhmern/&$\rightarrow$ & \textit{6400} & /{\dh}hmern/&$\rightarrow$ & \textit{7501} & /hmern/&$\rightarrow$ & \textit{7600} & /hmen/&$\rightarrow$ & \textit{13002} & /men/& \\
\end{tabular}

\vspace{20pt}\hline

\end{nopagebreak}
\filbreak



\vspace{15pt}
\begin{nopagebreak}
\noindent{\fontsize{20pt}{10pt}\textbf{mo-} } \textit{MID/middle voice marker} (VOICE)\\
\noindent {\tovian \fontsize{20pt}{10pt} \textbf{mo-} }\\
\noindent /m{\textprimstress}o/\\


\noindent History:

\vspace{-0pt}
\hspace{40pt}
\begin{tabular}{ccc}
\textit{0} & /mo-/& \\
\end{tabular}

\vspace{20pt}\hline

\end{nopagebreak}
\filbreak



\vspace{15pt}
\begin{nopagebreak}
\noindent{\fontsize{20pt}{10pt}\textbf{tlurinin} } \textit{mist} (N)\\
\noindent {\tovian \fontsize{20pt}{10pt} \textbf{tlurinin} }\\
\noindent /{\texttoptiebar{t\textbeltl}}ur{\textprimstress}inin/\\
\noindent lit. fog+small\\


\noindent History:

\vspace{-0pt}
\hspace{40pt}
\begin{tabular}{ccc}
\textit{100} & /inga{\texttoptiebar{t\textbeltl}}unarai{\texttheta}nina/&$\rightarrow$ & \textit{3200} & /ing{\texttoptiebar{t\textbeltl}}unrai{\texttheta}nina/&$\rightarrow$ & \textit{5500} & /ing{\texttoptiebar{t\textbeltl}}unri{\texttheta}nina/&$\rightarrow$ & \textit{6000} & /i{\ng}g{\texttoptiebar{t\textbeltl}}unri{\texttheta}nina/&$\rightarrow$ & \textit{7500} & /{\ng}g{\texttoptiebar{t\textbeltl}}unri{\texttheta}nina/&$\rightarrow$ & \textit{7501} & /g{\texttoptiebar{t\textbeltl}}unri{\texttheta}nina/&$\rightarrow$ & \textit{7600} & /g{\texttoptiebar{t\textbeltl}}uri{\texttheta}nina/&$\rightarrow$ & \textit{9500} & /g{\texttoptiebar{t\textbeltl}}uri{\texttheta}nin/&$\rightarrow$ & \textit{11501} & /{\texttoptiebar{t\textbeltl}}uri{\texttheta}nin/&$\rightarrow$ & \textit{13002} & /{\texttoptiebar{t\textbeltl}}urinin/& \\
\end{tabular}

\vspace{20pt}\hline

\end{nopagebreak}
\filbreak



\vspace{15pt}
\begin{nopagebreak}
\noindent{\fontsize{20pt}{10pt}\textbf{nunalom} } \textit{month} (N)\\
\noindent {\tovian \fontsize{20pt}{10pt} \textbf{nunalom} }\\
\noindent /nun{\textprimstress}alom/\\
\noindent lit. small+small+cycle\\


\noindent History:

\vspace{-0pt}
\hspace{40pt}
\begin{tabular}{ccc}
\textit{0} & /ninaninaelakorim/&$\rightarrow$ & \textit{2000} & /ninaninaelagorim/&$\rightarrow$ & \textit{3200} & /nnannaelgorm/&$\rightarrow$ & \textit{4501} & /nnannaelorm/&$\rightarrow$ & \textit{6500} & /nanaelorm/&$\rightarrow$ & \textit{7600} & /nanaelom/&$\rightarrow$ & \textit{9200} & /n{\textschwa}naelom/&$\rightarrow$ & \textit{10000} & /n{\textschwa}nalom/&$\rightarrow$ & \textit{11990} & /nnalom/&$\rightarrow$ & \textit{11996} & /nunalom/& \\
\end{tabular}

\vspace{20pt}\hline

\end{nopagebreak}
\filbreak



\vspace{15pt}
\begin{nopagebreak}
\noindent{\fontsize{20pt}{10pt}\textbf{ay} } \textit{moon} (N)\\
\noindent {\tovian \fontsize{20pt}{10pt} \textbf{ay} }\\
\noindent /{\textprimstress}ay/\\


\noindent History:

\vspace{-0pt}
\hspace{40pt}
\begin{tabular}{ccc}
\textit{0} & /alius/&$\rightarrow$ & \textit{2201} & /aliyus/&$\rightarrow$ & \textit{3200} & /aliys/&$\rightarrow$ & \textit{6300} & /alys/&$\rightarrow$ & \textit{7600} & /ays/&$\rightarrow$ & \textit{13001} & /ay/& \\
\end{tabular}

\vspace{20pt}\hline

\end{nopagebreak}
\filbreak



\vspace{15pt}
\begin{nopagebreak}
\noindent{\fontsize{20pt}{10pt}\textbf{mar} } \textit{mother} (N)\\
\noindent {\tovian \fontsize{20pt}{10pt} \textbf{mar} }\\
\noindent /m{\textprimstress}ar/\\


\noindent History:

\vspace{-0pt}
\hspace{40pt}
\begin{tabular}{ccc}
\textit{0} & /mari/&$\rightarrow$ & \textit{9500} & /mar/& \\
\end{tabular}

\vspace{20pt}\hline

\end{nopagebreak}
\filbreak



\vspace{15pt}
\begin{nopagebreak}
\noindent{\fontsize{20pt}{10pt}\textbf{eon} } \textit{mouth} (N)\\
\noindent {\tovian \fontsize{20pt}{10pt} \textbf{eon} }\\
\noindent /{\textprimstress}eon/\\


\noindent History:

\vspace{-0pt}
\hspace{40pt}
\begin{tabular}{ccc}
\textit{0} & /eolan/&$\rightarrow$ & \textit{3200} & /eoln/&$\rightarrow$ & \textit{7600} & /eon/& \\
\end{tabular}

\vspace{20pt}\hline

\end{nopagebreak}
\filbreak



\vspace{15pt}
\begin{nopagebreak}
\noindent{\fontsize{20pt}{10pt}\textbf{kwil} } \textit{movement} (N)\\
\noindent {\tovian \fontsize{20pt}{10pt} \textbf{kwil} }\\
\noindent /kw{\textprimstress}il/\\


\noindent History:

\vspace{-0pt}
\hspace{40pt}
\begin{tabular}{ccc}
\textit{0} & /kwil/& \\
\end{tabular}

\vspace{20pt}\hline

\end{nopagebreak}
\filbreak



\vspace{15pt}
\begin{nopagebreak}
\noindent{\fontsize{20pt}{10pt}\textbf{lhemlor} } \textit{narcotic} (N)\\
\noindent {\tovian \fontsize{20pt}{10pt} \textbf{lhemlor} }\\
\noindent /{\textbeltl}{\textprimstress}emlor/\\
\noindent lit. sleep+plant\\


\noindent History:

\vspace{-0pt}
\hspace{40pt}
\begin{tabular}{ccc}
\textit{500} & /{\textbeltl}e{\dh}el{\textschwa}mellore/&$\rightarrow$ & \textit{2300} & /{\textbeltl}e{\dh}elmellore/&$\rightarrow$ & \textit{3200} & /{\textbeltl}{\dh}elmllore/&$\rightarrow$ & \textit{4502} & /{\textbeltl}elmllore/&$\rightarrow$ & \textit{6500} & /{\textbeltl}elmlore/&$\rightarrow$ & \textit{7600} & /{\textbeltl}emlore/&$\rightarrow$ & \textit{9500} & /{\textbeltl}emlor/& \\
\end{tabular}

\vspace{20pt}\hline

\end{nopagebreak}
\filbreak



\vspace{15pt}
\begin{nopagebreak}
\noindent{\fontsize{20pt}{10pt}\textbf{men} } \textit{newness} (N)\\
\noindent {\tovian \fontsize{20pt}{10pt} \textbf{men} }\\
\noindent /m{\textprimstress}en/\\


\noindent History:

\vspace{-0pt}
\hspace{40pt}
\begin{tabular}{ccc}
\textit{0} & /{\textyogh}emeno/&$\rightarrow$ & \textit{3200} & /{\textyogh}meno/&$\rightarrow$ & \textit{9500} & /{\textyogh}men/&$\rightarrow$ & \textit{12001} & /{\textesh}men/&$\rightarrow$ & \textit{13002} & /men/& \\
\end{tabular}

\vspace{20pt}\hline

\end{nopagebreak}
\filbreak



\vspace{15pt}
\begin{nopagebreak}
\noindent{\fontsize{20pt}{10pt}\textbf{lhil} } \textit{night/nighttime} (N)\\
\noindent {\tovian \fontsize{20pt}{10pt} \textbf{lhil} }\\
\noindent /{\textbeltl}{\textprimstress}il/\\
\noindent \textit{Time of day}\\


\noindent History:

\vspace{-0pt}
\hspace{40pt}
\begin{tabular}{ccc}
\textit{0} & /a{\textbeltl}e{\textsubbridge{t}}inla/&$\rightarrow$ & \textit{1000} & /a{\textbeltl}{\textsubbridge{t}}inla/&$\rightarrow$ & \textit{3000} & /a{\textbeltl}inla/&$\rightarrow$ & \textit{7500} & /{\textbeltl}inla/&$\rightarrow$ & \textit{7600} & /{\textbeltl}ila/&$\rightarrow$ & \textit{9500} & /{\textbeltl}il/& \\
\end{tabular}

\vspace{20pt}\hline

\end{nopagebreak}
\filbreak



\vspace{15pt}
\begin{nopagebreak}
\noindent{\fontsize{20pt}{10pt}\textbf{ki} } \textit{nine/9} (N)\\
\noindent {\tovian \fontsize{20pt}{10pt} \textbf{ki} }\\
\noindent /k{\textprimstress}i/\\


\noindent History:

\vspace{-0pt}
\hspace{40pt}
\begin{tabular}{ccc}
\textit{0} & /ki/& \\
\end{tabular}

\vspace{20pt}\hline

\end{nopagebreak}
\filbreak



\vspace{15pt}
\begin{nopagebreak}
\noindent{\fontsize{20pt}{10pt}\textbf{i-} } \textit{NOM/nomative case marker} (CASE)\\
\noindent {\tovian \fontsize{20pt}{10pt} \textbf{i-} }\\
\noindent /{\textprimstress}i/\\


\noindent History:

\vspace{-0pt}
\hspace{40pt}
\begin{tabular}{ccc}
\textit{0} & /i-/& \\
\end{tabular}

\vspace{20pt}\hline

\end{nopagebreak}
\filbreak



\vspace{15pt}
\begin{nopagebreak}
\noindent{\fontsize{20pt}{10pt}\textbf{anad} } \textit{non-person hair/fur/hide/skin/scales} (N)\\
\noindent {\tovian \fontsize{20pt}{10pt} \textbf{anad} }\\
\noindent /{\textprimstress}anad/\\


\noindent History:

\vspace{-0pt}
\hspace{40pt}
\begin{tabular}{ccc}
\textit{0} & /ana{\textsubbridge{t}}ur/&$\rightarrow$ & \textit{2000} & /anadur/&$\rightarrow$ & \textit{3200} & /anadr/&$\rightarrow$ & \textit{13001} & /anad/& \\
\end{tabular}

\vspace{20pt}\hline

\end{nopagebreak}
\filbreak



\vspace{15pt}
\begin{nopagebreak}
\noindent{\fontsize{20pt}{10pt}\textbf{lar} } \textit{north} (N)\\
\noindent {\tovian \fontsize{20pt}{10pt} \textbf{lar} }\\
\noindent /l{\textprimstress}ar/\\


\noindent History:

\vspace{-0pt}
\hspace{40pt}
\begin{tabular}{ccc}
\textit{0} & /lare/&$\rightarrow$ & \textit{9500} & /lar/& \\
\end{tabular}

\vspace{20pt}\hline

\end{nopagebreak}
\filbreak



\vspace{15pt}
\begin{nopagebreak}
\noindent{\fontsize{20pt}{10pt}\textbf{ven} } \textit{nose} (N)\\
\noindent {\tovian \fontsize{20pt}{10pt} \textbf{ven} }\\
\noindent /v{\textprimstress}en/\\


\noindent History:

\vspace{-0pt}
\hspace{40pt}
\begin{tabular}{ccc}
\textit{0} & /verun/&$\rightarrow$ & \textit{3200} & /vern/&$\rightarrow$ & \textit{7600} & /ven/& \\
\end{tabular}

\vspace{20pt}\hline

\end{nopagebreak}
\filbreak



\vspace{15pt}
\begin{nopagebreak}
\noindent{\fontsize{20pt}{10pt}\textbf{ferlesh} } \textit{November} (N)\\
\noindent {\tovian \fontsize{20pt}{10pt} \textbf{ferlesh} }\\
\noindent /f{\textprimstress}erle{\textesh}/\\
\noindent lit. three+harvest\\
\noindent \textit{Prairilight}\\


\noindent History:

\vspace{-0pt}
\hspace{40pt}
\begin{tabular}{ccc}
\textit{6000} & /femmllre{\textyogh}/&$\rightarrow$ & \textit{6500} & /femlre{\textyogh}/&$\rightarrow$ & \textit{6503} & /femrle{\textyogh}/&$\rightarrow$ & \textit{7600} & /ferle{\textyogh}/&$\rightarrow$ & \textit{12001} & /ferle{\textesh}/& \\
\end{tabular}

\vspace{20pt}\hline

\end{nopagebreak}
\filbreak



\vspace{15pt}
\begin{nopagebreak}
\noindent{\fontsize{20pt}{10pt}\textbf{tl0-} } \textit{OBL/obligative mood marker} (MOOD)\\
\noindent {\tovian \fontsize{20pt}{10pt} \textbf{tl0-} }\\
\noindent /{\texttoptiebar{t\textbeltl}}0-/\\


\noindent History:

\vspace{-0pt}
\hspace{40pt}
\begin{tabular}{ccc}
\textit{0} & /{\texttoptiebar{t\textbeltl}}0-/& \\
\end{tabular}

\vspace{20pt}\hline

\end{nopagebreak}
\filbreak



\vspace{15pt}
\begin{nopagebreak}
\noindent{\fontsize{20pt}{10pt}\textbf{loerlesh} } \textit{October} (N)\\
\noindent {\tovian \fontsize{20pt}{10pt} \textbf{loerlesh} }\\
\noindent /lo{\textprimstress}erle{\textesh}/\\
\noindent lit. two+harvest\\
\noindent \textit{Fructiwane}\\


\noindent History:

\vspace{-0pt}
\hspace{40pt}
\begin{tabular}{ccc}
\textit{6000} & /loemllre{\textyogh}/&$\rightarrow$ & \textit{6500} & /loemlre{\textyogh}/&$\rightarrow$ & \textit{6503} & /loemrle{\textyogh}/&$\rightarrow$ & \textit{7600} & /loerle{\textyogh}/&$\rightarrow$ & \textit{12001} & /loerle{\textesh}/& \\
\end{tabular}

\vspace{20pt}\hline

\end{nopagebreak}
\filbreak



\vspace{15pt}
\begin{nopagebreak}
\noindent{\fontsize{20pt}{10pt}\textbf{lhyalh} } \textit{one hundred forty four/144} (N)\\
\noindent {\tovian \fontsize{20pt}{10pt} \textbf{lhyalh} }\\
\noindent /{\textbeltl}y{\textprimstress}a{\textbeltl}/\\
\noindent lit. twelve+twelve\\


\noindent History:

\vspace{-0pt}
\hspace{40pt}
\begin{tabular}{ccc}
\textit{4000} & /{\textbeltl}ya{\textbeltl}ya/&$\rightarrow$ & \textit{9500} & /{\textbeltl}ya{\textbeltl}y/&$\rightarrow$ & \textit{13001} & /{\textbeltl}ya{\textbeltl}/& \\
\end{tabular}

\vspace{20pt}\hline

\end{nopagebreak}
\filbreak



\vspace{15pt}
\begin{nopagebreak}
\noindent{\fontsize{20pt}{10pt}\textbf{lo} } \textit{one/1} (N)\\
\noindent {\tovian \fontsize{20pt}{10pt} \textbf{lo} }\\
\noindent /l{\textprimstress}o/\\


\noindent History:

\vspace{-0pt}
\hspace{40pt}
\begin{tabular}{ccc}
\textit{0} & /lo/& \\
\end{tabular}

\vspace{20pt}\hline

\end{nopagebreak}
\filbreak



\vspace{15pt}
\begin{nopagebreak}
\noindent{\fontsize{20pt}{10pt}\textbf{dam} } \textit{openness} (N)\\
\noindent {\tovian \fontsize{20pt}{10pt} \textbf{dam} }\\
\noindent /d{\textprimstress}am/\\
\noindent lit. not+closure\\


\noindent History:

\vspace{-0pt}
\hspace{40pt}
\begin{tabular}{ccc}
\textit{1300} & /do{\textsubbridge{t}}anum/&$\rightarrow$ & \textit{2000} & /dodanum/&$\rightarrow$ & \textit{3200} & /ddanm/&$\rightarrow$ & \textit{6500} & /danm/&$\rightarrow$ & \textit{7600} & /dam/& \\
\end{tabular}

\vspace{20pt}\hline

\end{nopagebreak}
\filbreak



\vspace{15pt}
\begin{nopagebreak}
\noindent{\fontsize{20pt}{10pt}\textbf{do} } \textit{opposite/not/un-} (N)\\
\noindent {\tovian \fontsize{20pt}{10pt} \textbf{do} }\\
\noindent /d{\textprimstress}o/\\


\noindent History:

\vspace{-0pt}
\hspace{40pt}
\begin{tabular}{ccc}
\textit{0} & /do/& \\
\end{tabular}

\vspace{20pt}\hline

\end{nopagebreak}
\filbreak



\vspace{15pt}
\begin{nopagebreak}
\noindent{\fontsize{20pt}{10pt}\textbf{lhe-} } \textit{OPT/optative mood marker} (MOOD)\\
\noindent {\tovian \fontsize{20pt}{10pt} \textbf{lhe-} }\\
\noindent /{\textbeltl}{\textprimstress}e/\\


\noindent History:

\vspace{-0pt}
\hspace{40pt}
\begin{tabular}{ccc}
\textit{0} & /{\textbeltl}e-/& \\
\end{tabular}

\vspace{20pt}\hline

\end{nopagebreak}
\filbreak



\vspace{15pt}
\begin{nopagebreak}
\noindent{\fontsize{20pt}{10pt}\textbf{har} } \textit{pain} (N)\\
\noindent {\tovian \fontsize{20pt}{10pt} \textbf{har} }\\
\noindent /ħ{\textprimstress}ar/\\


\noindent History:

\vspace{-0pt}
\hspace{40pt}
\begin{tabular}{ccc}
\textit{0} & /hara/&$\rightarrow$ & \textit{9500} & /har/&$\rightarrow$ & \textit{14006} & /ħar/& \\
\end{tabular}

\vspace{20pt}\hline

\end{nopagebreak}
\filbreak



\vspace{15pt}
\begin{nopagebreak}
\noindent{\fontsize{20pt}{10pt}\textbf{pa-} } \textit{PASS/passive voice marker} (VOICE)\\
\noindent {\tovian \fontsize{20pt}{10pt} \textbf{pa-} }\\
\noindent /p{\textprimstress}a/\\


\noindent History:

\vspace{-0pt}
\hspace{40pt}
\begin{tabular}{ccc}
\textit{0} & /pa-/& \\
\end{tabular}

\vspace{20pt}\hline

\end{nopagebreak}
\filbreak



\vspace{15pt}
\begin{nopagebreak}
\noindent{\fontsize{20pt}{10pt}\textbf{hun} } \textit{past} (Ns)\\
\noindent {\tovian \fontsize{20pt}{10pt} \textbf{hun} }\\
\noindent /ħ{\textprimstress}un/\\
\noindent lit. behind+path\\


\noindent History:

\vspace{-0pt}
\hspace{40pt}
\begin{tabular}{ccc}
\textit{10} & /hulnum/&$\rightarrow$ & \textit{3200} & /hulnm/&$\rightarrow$ & \textit{7600} & /hunm/&$\rightarrow$ & \textit{13001} & /hun/&$\rightarrow$ & \textit{14006} & /ħun/& \\
\end{tabular}

\vspace{20pt}\hline

\end{nopagebreak}
\filbreak



\vspace{15pt}
\begin{nopagebreak}
\noindent{\fontsize{20pt}{10pt}\textbf{num} } \textit{path} (N)\\
\noindent {\tovian \fontsize{20pt}{10pt} \textbf{num} }\\
\noindent /n{\textprimstress}um/\\


\noindent History:

\vspace{-0pt}
\hspace{40pt}
\begin{tabular}{ccc}
\textit{0} & /num/& \\
\end{tabular}

\vspace{20pt}\hline

\end{nopagebreak}
\filbreak



\vspace{15pt}
\begin{nopagebreak}
\noindent{\fontsize{20pt}{10pt}\textbf{tla} } \textit{peace/peacefulness/calmness/calm} (N)\\
\noindent {\tovian \fontsize{20pt}{10pt} \textbf{tla} }\\
\noindent /{\texttoptiebar{t\textbeltl}}{\textprimstress}a/\\


\noindent History:

\vspace{-0pt}
\hspace{40pt}
\begin{tabular}{ccc}
\textit{0} & /{\texttoptiebar{t\textbeltl}}af/&$\rightarrow$ & \textit{7501} & /{\texttoptiebar{t\textbeltl}}a{\textphi}/&$\rightarrow$ & \textit{8500} & /{\texttoptiebar{t\textbeltl}}a/& \\
\end{tabular}

\vspace{20pt}\hline

\end{nopagebreak}
\filbreak



\vspace{15pt}
\begin{nopagebreak}
\noindent{\fontsize{20pt}{10pt}\textbf{yel} } \textit{person} (N)\\
\noindent {\tovian \fontsize{20pt}{10pt} \textbf{yel} }\\
\noindent /y{\textprimstress}el/\\


\noindent History:

\vspace{-0pt}
\hspace{40pt}
\begin{tabular}{ccc}
\textit{0} & /driela/&$\rightarrow$ & \textit{2201} & /driyela/&$\rightarrow$ & \textit{3200} & /dryela/&$\rightarrow$ & \textit{7501} & /ryela/&$\rightarrow$ & \textit{7600} & /yela/&$\rightarrow$ & \textit{9500} & /yel/& \\
\end{tabular}

\vspace{20pt}\hline

\end{nopagebreak}
\filbreak



\vspace{15pt}
\begin{nopagebreak}
\noindent{\fontsize{20pt}{10pt}\textbf{trel} } \textit{person skin} (N)\\
\noindent {\tovian \fontsize{20pt}{10pt} \textbf{trel} }\\
\noindent /{\textsubbridge{t}}r{\textprimstress}el/\\


\noindent History:

\vspace{-0pt}
\hspace{40pt}
\begin{tabular}{ccc}
\textit{0} & /{\textsubbridge{t}}relahe{\textsubbridge{t}}/&$\rightarrow$ & \textit{1000} & /{\textsubbridge{t}}relah{\textsubbridge{t}}/&$\rightarrow$ & \textit{3000} & /{\textsubbridge{t}}relah/&$\rightarrow$ & \textit{3200} & /{\textsubbridge{t}}relh/&$\rightarrow$ & \textit{5000} & /{\textsubbridge{t}}rel/& \\
\end{tabular}

\vspace{20pt}\hline

\end{nopagebreak}
\filbreak



\vspace{15pt}
\begin{nopagebreak}
\noindent{\fontsize{20pt}{10pt}\textbf{thil} } \textit{place} (N)\\
\noindent {\tovian \fontsize{20pt}{10pt} \textbf{thil} }\\
\noindent /{\texttheta}{\textprimstress}il/\\


\noindent History:

\vspace{-0pt}
\hspace{40pt}
\begin{tabular}{ccc}
\textit{0} & /{\texttheta}ile/&$\rightarrow$ & \textit{9500} & /{\texttheta}il/& \\
\end{tabular}

\vspace{20pt}\hline

\end{nopagebreak}
\filbreak



\vspace{15pt}
\begin{nopagebreak}
\noindent{\fontsize{20pt}{10pt}\textbf{til} } \textit{plain} (N)\\
\noindent {\tovian \fontsize{20pt}{10pt} \textbf{til} }\\
\noindent /{\textsubbridge{t}}{\textprimstress}il/\\
\noindent lit. flatness+place\\


\noindent History:

\vspace{-0pt}
\hspace{40pt}
\begin{tabular}{ccc}
\textit{50} & /{\textsubbridge{t}}um{\texttheta}ile/&$\rightarrow$ & \textit{3200} & /{\textsubbridge{t}}m{\texttheta}ile/&$\rightarrow$ & \textit{3503} & /{\textsubbridge{t}}n{\texttheta}ile/&$\rightarrow$ & \textit{6400} & /{\textsubbridge{t}}{\texttheta}ile/&$\rightarrow$ & \textit{9500} & /{\textsubbridge{t}}{\texttheta}il/&$\rightarrow$ & \textit{10700} & /{\textsubbridge{t}}{\textsubbridge{t}}il/&$\rightarrow$ & \textit{12007} & /{\textsubbridge{t}}il/& \\
\end{tabular}

\vspace{20pt}\hline

\end{nopagebreak}
\filbreak



\vspace{15pt}
\begin{nopagebreak}
\noindent{\fontsize{20pt}{10pt}\textbf{melor} } \textit{plant} (N)\\
\noindent {\tovian \fontsize{20pt}{10pt} \textbf{melor} }\\
\noindent /m{\textprimstress}elor/\\
\noindent lit. life+action\\


\noindent History:

\vspace{-0pt}
\hspace{40pt}
\begin{tabular}{ccc}
\textit{50} & /{\textschwa}mellore/&$\rightarrow$ & \textit{2300} & /mellore/&$\rightarrow$ & \textit{3200} & /mllore/&$\rightarrow$ & \textit{6500} & /mlore/&$\rightarrow$ & \textit{7502} & /melore/&$\rightarrow$ & \textit{9500} & /melor/& \\
\end{tabular}

\vspace{20pt}\hline

\end{nopagebreak}
\filbreak



\vspace{15pt}
\begin{nopagebreak}
\noindent{\fontsize{20pt}{10pt}\textbf{-e} } \textit{plural marker} (PLURAL)\\
\noindent {\tovian \fontsize{20pt}{10pt} \textbf{-e} }\\
\noindent /{\textprimstress}e/\\


\noindent History:

\vspace{-0pt}
\hspace{40pt}
\begin{tabular}{ccc}
\textit{0} & /-e/& \\
\end{tabular}

\vspace{20pt}\hline

\end{nopagebreak}
\filbreak



\vspace{15pt}
\begin{nopagebreak}
\noindent{\fontsize{20pt}{10pt}\textbf{el} } \textit{possession} (N)\\
\noindent {\tovian \fontsize{20pt}{10pt} \textbf{el} }\\
\noindent /{\textprimstress}el/\\


\noindent History:

\vspace{-0pt}
\hspace{40pt}
\begin{tabular}{ccc}
\textit{0} & /elis/&$\rightarrow$ & \textit{3200} & /els/&$\rightarrow$ & \textit{13001} & /el/& \\
\end{tabular}

\vspace{20pt}\hline

\end{nopagebreak}
\filbreak



\vspace{15pt}
\begin{nopagebreak}
\noindent{\fontsize{20pt}{10pt}\textbf{sho} } \textit{possibility} (N)\\
\noindent {\tovian \fontsize{20pt}{10pt} \textbf{sho} }\\
\noindent /{\textesh}{\textprimstress}o/\\


\noindent History:

\vspace{-0pt}
\hspace{40pt}
\begin{tabular}{ccc}
\textit{0} & /{\textesh}ohe/&$\rightarrow$ & \textit{5000} & /{\textesh}oe/&$\rightarrow$ & \textit{9500} & /{\textesh}o/& \\
\end{tabular}

\vspace{20pt}\hline

\end{nopagebreak}
\filbreak



\vspace{15pt}
\begin{nopagebreak}
\noindent{\fontsize{20pt}{10pt}\textbf{sho-} } \textit{POT/potential mood marker} (MOOD)\\
\noindent {\tovian \fontsize{20pt}{10pt} \textbf{sho-} }\\
\noindent /{\textesh}{\textprimstress}o/\\


\noindent History:

\vspace{-0pt}
\hspace{40pt}
\begin{tabular}{ccc}
\textit{0} & /{\textesh}o-/& \\
\end{tabular}

\vspace{20pt}\hline

\end{nopagebreak}
\filbreak



\vspace{15pt}
\begin{nopagebreak}
\noindent{\fontsize{20pt}{10pt}\textbf{num} } \textit{present} (Ns)\\
\noindent {\tovian \fontsize{20pt}{10pt} \textbf{num} }\\
\noindent /n{\textprimstress}um/\\
\noindent lit. path+center\\


\noindent History:

\vspace{-0pt}
\hspace{40pt}
\begin{tabular}{ccc}
\textit{10} & /num{\textbeltl}el/&$\rightarrow$ & \textit{3200} & /num{\textbeltl}l/&$\rightarrow$ & \textit{12500} & /num{\textbeltl}/&$\rightarrow$ & \textit{13001} & /num/& \\
\end{tabular}

\vspace{20pt}\hline

\end{nopagebreak}
\filbreak



\vspace{15pt}
\begin{nopagebreak}
\noindent{\fontsize{20pt}{10pt}\textbf{led} } \textit{present/demonstrate} (N)\\
\noindent {\tovian \fontsize{20pt}{10pt} \textbf{led} }\\
\noindent /l{\textprimstress}ed/\\
\noindent lit. action+showing\\


\noindent History:

\vspace{-0pt}
\hspace{40pt}
\begin{tabular}{ccc}
\textit{0} & /lorefe{\textsubbridge{t}}e/&$\rightarrow$ & \textit{2000} & /lorefede/&$\rightarrow$ & \textit{3200} & /lrefede/&$\rightarrow$ & \textit{6503} & /rlefede/&$\rightarrow$ & \textit{7600} & /lefede/&$\rightarrow$ & \textit{9500} & /lefed/&$\rightarrow$ & \textit{12003} & /led/& \\
\end{tabular}

\vspace{20pt}\hline

\end{nopagebreak}
\filbreak



\vspace{15pt}
\begin{nopagebreak}
\noindent{\fontsize{20pt}{10pt}\textbf{sor} } \textit{proximity/closeness} (N)\\
\noindent {\tovian \fontsize{20pt}{10pt} \textbf{sor} }\\
\noindent /s{\textprimstress}or/\\


\noindent History:

\vspace{-0pt}
\hspace{40pt}
\begin{tabular}{ccc}
\textit{0} & /sor/& \\
\end{tabular}

\vspace{20pt}\hline

\end{nopagebreak}
\filbreak



\vspace{15pt}
\begin{nopagebreak}
\noindent{\fontsize{20pt}{10pt}\textbf{lhu-} } \textit{PUR/purpose case marker/for} (CASE)\\
\noindent {\tovian \fontsize{20pt}{10pt} \textbf{lhu-} }\\
\noindent /{\textbeltl}{\textprimstress}u/\\


\noindent History:

\vspace{-0pt}
\hspace{40pt}
\begin{tabular}{ccc}
\textit{0} & /{\textbeltl}u-/& \\
\end{tabular}

\vspace{20pt}\hline

\end{nopagebreak}
\filbreak



\vspace{15pt}
\begin{nopagebreak}
\noindent{\fontsize{20pt}{10pt}\textbf{man} } \textit{putting/placing} (N)\\
\noindent {\tovian \fontsize{20pt}{10pt} \textbf{man} }\\
\noindent /m{\textprimstress}an/\\


\noindent History:

\vspace{-0pt}
\hspace{40pt}
\begin{tabular}{ccc}
\textit{0} & /hemmani/&$\rightarrow$ & \textit{3200} & /hmmani/&$\rightarrow$ & \textit{6500} & /hmani/&$\rightarrow$ & \textit{9500} & /hman/&$\rightarrow$ & \textit{13002} & /man/& \\
\end{tabular}

\vspace{20pt}\hline

\end{nopagebreak}
\filbreak



\vspace{15pt}
\begin{nopagebreak}
\noindent{\fontsize{20pt}{10pt}\textbf{kith} } \textit{reaching/arriving/traveling/going} (N)\\
\noindent {\tovian \fontsize{20pt}{10pt} \textbf{kith} }\\
\noindent /k{\textprimstress}i{\texttheta}/\\


\noindent History:

\vspace{-0pt}
\hspace{40pt}
\begin{tabular}{ccc}
\textit{0} & /ki{\texttheta}sa/&$\rightarrow$ & \textit{4502} & /ki{\texttheta}a/&$\rightarrow$ & \textit{9500} & /ki{\texttheta}/& \\
\end{tabular}

\vspace{20pt}\hline

\end{nopagebreak}
\filbreak



\vspace{15pt}
\begin{nopagebreak}
\noindent{\fontsize{20pt}{10pt}\textbf{mloelhan} } \textit{rebirth} (N)\\
\noindent {\tovian \fontsize{20pt}{10pt} \textbf{mloelhan} }\\
\noindent /mlo{\textprimstress}e{\textbeltl}an/\\
\noindent lit. second+birth\\


\noindent History:

\vspace{-0pt}
\hspace{40pt}
\begin{tabular}{ccc}
\textit{1500} & /puloe{\textbeltl}ogaahana/&$\rightarrow$ & \textit{3200} & /ploe{\textbeltl}gaahana/&$\rightarrow$ & \textit{4500} & /ploe{\textbeltl}aahana/&$\rightarrow$ & \textit{5000} & /ploe{\textbeltl}aaana/&$\rightarrow$ & \textit{5500} & /ploe{\textbeltl}aana/&$\rightarrow$ & \textit{9500} & /ploe{\textbeltl}aan/&$\rightarrow$ & \textit{11000} & /mloe{\textbeltl}aan/&$\rightarrow$ & \textit{12005} & /mloe{\textbeltl}an/& \\
\end{tabular}

\vspace{20pt}\hline

\end{nopagebreak}
\filbreak



\vspace{15pt}
\begin{nopagebreak}
\noindent{\fontsize{20pt}{10pt}\textbf{ra-} } \textit{REC/reciprocal voice marker} (VOICE)\\
\noindent {\tovian \fontsize{20pt}{10pt} \textbf{ra-} }\\
\noindent /r{\textprimstress}a/\\


\noindent History:

\vspace{-0pt}
\hspace{40pt}
\begin{tabular}{ccc}
\textit{0} & /ra-/& \\
\end{tabular}

\vspace{20pt}\hline

\end{nopagebreak}
\filbreak



\vspace{15pt}
\begin{nopagebreak}
\noindent{\fontsize{20pt}{10pt}\textbf{lhod} } \textit{redness/red} (N)\\
\noindent {\tovian \fontsize{20pt}{10pt} \textbf{lhod} }\\
\noindent /{\textbeltl}{\textprimstress}od/\\


\noindent History:

\vspace{-0pt}
\hspace{40pt}
\begin{tabular}{ccc}
\textit{0} & /{\textbeltl}eho{\textsubbridge{t}}en/&$\rightarrow$ & \textit{1000} & /{\textbeltl}ho{\textsubbridge{t}}en/&$\rightarrow$ & \textit{2000} & /{\textbeltl}hoden/&$\rightarrow$ & \textit{3200} & /{\textbeltl}hodn/&$\rightarrow$ & \textit{4502} & /{\textbeltl}odn/&$\rightarrow$ & \textit{13001} & /{\textbeltl}od/& \\
\end{tabular}

\vspace{20pt}\hline

\end{nopagebreak}
\filbreak



\vspace{15pt}
\begin{nopagebreak}
\noindent{\fontsize{20pt}{10pt}\textbf{te-} } \textit{REFL/reflexive voice marker} (VOICE)\\
\noindent {\tovian \fontsize{20pt}{10pt} \textbf{te-} }\\
\noindent /{\textsubbridge{t}}{\textprimstress}e/\\


\noindent History:

\vspace{-0pt}
\hspace{40pt}
\begin{tabular}{ccc}
\textit{0} & /{\textsubbridge{t}}e-/& \\
\end{tabular}

\vspace{20pt}\hline

\end{nopagebreak}
\filbreak



\vspace{15pt}
\begin{nopagebreak}
\noindent{\fontsize{20pt}{10pt}\textbf{yeolfethr} } \textit{report/confession} (N)\\
\noindent {\tovian \fontsize{20pt}{10pt} \textbf{yeolfethr} }\\
\noindent /ye{\textprimstress}olfe{\texttheta}r/\\
\noindent lit. truth+speech\\


\noindent History:

\vspace{-0pt}
\hspace{40pt}
\begin{tabular}{ccc}
\textit{7200} & /aryeolfe{\dh}ra/&$\rightarrow$ & \textit{7500} & /ryeolfe{\dh}ra/&$\rightarrow$ & \textit{7600} & /yeolfe{\dh}ra/&$\rightarrow$ & \textit{9500} & /yeolfe{\dh}r/&$\rightarrow$ & \textit{12002} & /yeolfe{\texttheta}r/& \\
\end{tabular}

\vspace{20pt}\hline

\end{nopagebreak}
\filbreak



\vspace{15pt}
\begin{nopagebreak}
\noindent{\fontsize{20pt}{10pt}\textbf{gith} } \textit{returning} (N)\\
\noindent {\tovian \fontsize{20pt}{10pt} \textbf{gith} }\\
\noindent /g{\textprimstress}i{\texttheta}/\\
\noindent lit. not+arriving\\


\noindent History:

\vspace{-0pt}
\hspace{40pt}
\begin{tabular}{ccc}
\textit{0} & /doki{\texttheta}sa/&$\rightarrow$ & \textit{2000} & /dogi{\texttheta}sa/&$\rightarrow$ & \textit{3200} & /dgi{\texttheta}sa/&$\rightarrow$ & \textit{4502} & /dgi{\texttheta}a/&$\rightarrow$ & \textit{9500} & /dgi{\texttheta}/&$\rightarrow$ & \textit{12503} & /gi{\texttheta}/& \\
\end{tabular}

\vspace{20pt}\hline

\end{nopagebreak}
\filbreak



\vspace{15pt}
\begin{nopagebreak}
\noindent{\fontsize{20pt}{10pt}\textbf{lhya} } \textit{right} (N)\\
\noindent {\tovian \fontsize{20pt}{10pt} \textbf{lhya} }\\
\noindent /{\textbeltl}y{\textprimstress}a/\\


\noindent History:

\vspace{-0pt}
\hspace{40pt}
\begin{tabular}{ccc}
\textit{0} & /{\textbeltl}oya/&$\rightarrow$ & \textit{2100} & /{\textbeltl}ya/& \\
\end{tabular}

\vspace{20pt}\hline

\end{nopagebreak}
\filbreak



\vspace{15pt}
\begin{nopagebreak}
\noindent{\fontsize{20pt}{10pt}\textbf{shotl} } \textit{river} (N)\\
\noindent {\tovian \fontsize{20pt}{10pt} \textbf{shotl} }\\
\noindent /{\textesh}{\textprimstress}o{\texttoptiebar{t\textbeltl}}/\\
\noindent lit. flow+water\\


\noindent History:

\vspace{-0pt}
\hspace{40pt}
\begin{tabular}{ccc}
\textit{500} & /{\textesh}o{\texttoptiebar{t\textbeltl}}z{\textsubbridge{t}}el/&$\rightarrow$ & \textit{2000} & /{\textesh}o{\texttoptiebar{t\textbeltl}}zdel/&$\rightarrow$ & \textit{3200} & /{\textesh}o{\texttoptiebar{t\textbeltl}}zdl/&$\rightarrow$ & \textit{4500} & /{\textesh}o{\texttoptiebar{t\textbeltl}}zl/&$\rightarrow$ & \textit{12000} & /{\textesh}o{\texttoptiebar{t\textbeltl}}sl/&$\rightarrow$ & \textit{12500} & /{\textesh}o{\texttoptiebar{t\textbeltl}}s/&$\rightarrow$ & \textit{13001} & /{\textesh}o{\texttoptiebar{t\textbeltl}}/& \\
\end{tabular}

\vspace{20pt}\hline

\end{nopagebreak}
\filbreak



\vspace{15pt}
\begin{nopagebreak}
\noindent{\fontsize{20pt}{10pt}\textbf{sethr} } \textit{running/moving quickly/rushing} (N)\\
\noindent {\tovian \fontsize{20pt}{10pt} \textbf{sethr} }\\
\noindent /s{\textprimstress}e{\texttheta}r/\\


\noindent History:

\vspace{-0pt}
\hspace{40pt}
\begin{tabular}{ccc}
\textit{0} & /se{\texttheta}ir/&$\rightarrow$ & \textit{3200} & /se{\texttheta}r/& \\
\end{tabular}

\vspace{20pt}\hline

\end{nopagebreak}
\filbreak



\vspace{15pt}
\begin{nopagebreak}
\noindent{\fontsize{20pt}{10pt}\textbf{lur} } \textit{sad} (N)\\
\noindent {\tovian \fontsize{20pt}{10pt} \textbf{lur} }\\
\noindent /l{\textprimstress}ur/\\


\noindent History:

\vspace{-0pt}
\hspace{40pt}
\begin{tabular}{ccc}
\textit{0} & /lura/&$\rightarrow$ & \textit{9500} & /lur/& \\
\end{tabular}

\vspace{20pt}\hline

\end{nopagebreak}
\filbreak



\vspace{15pt}
\begin{nopagebreak}
\noindent{\fontsize{20pt}{10pt}\textbf{shotl} } \textit{sand} (N)\\
\noindent {\tovian \fontsize{20pt}{10pt} \textbf{shotl} }\\
\noindent /{\textesh}{\textprimstress}o{\texttoptiebar{t\textbeltl}}/\\
\noindent lit. flow+ground\\
\noindent \textit{ Same as "flowing"}\\


\noindent History:

\vspace{-0pt}
\hspace{40pt}
\begin{tabular}{ccc}
\textit{1000} & /{\textesh}o{\texttoptiebar{t\textbeltl}}{\texttheta}ok{\textesh}{\textsubbridge{t}}/&$\rightarrow$ & \textit{1000} & /{\textesh}o{\texttoptiebar{t\textbeltl}}{\texttheta}k{\textesh}{\textsubbridge{t}}/&$\rightarrow$ & \textit{3000} & /{\textesh}o{\texttoptiebar{t\textbeltl}}{\texttheta}k{\textesh}/&$\rightarrow$ & \textit{4500} & /{\textesh}o{\texttoptiebar{t\textbeltl}}{\texttheta}{\textesh}/&$\rightarrow$ & \textit{4502} & /{\textesh}o{\texttoptiebar{t\textbeltl}}{\texttheta}/&$\rightarrow$ & \textit{13001} & /{\textesh}o{\texttoptiebar{t\textbeltl}}/& \\
\end{tabular}

\vspace{20pt}\hline

\end{nopagebreak}
\filbreak



\vspace{15pt}
\begin{nopagebreak}
\noindent{\fontsize{20pt}{10pt}\textbf{nor} } \textit{sea} (N)\\
\noindent {\tovian \fontsize{20pt}{10pt} \textbf{nor} }\\
\noindent /n{\textprimstress}or/\\


\noindent History:

\vspace{-0pt}
\hspace{40pt}
\begin{tabular}{ccc}
\textit{0} & /nor/& \\
\end{tabular}

\vspace{20pt}\hline

\end{nopagebreak}
\filbreak



\vspace{15pt}
\begin{nopagebreak}
\noindent{\fontsize{20pt}{10pt}\textbf{rithagom} } \textit{season} (N)\\
\noindent {\tovian \fontsize{20pt}{10pt} \textbf{rithagom} }\\
\noindent /ri{\texttheta}{\textprimstress}agom/\\
\noindent lit. air+cycle\\


\noindent History:

\vspace{-0pt}
\hspace{40pt}
\begin{tabular}{ccc}
\textit{1000} & /rai{\texttheta}elakorim/&$\rightarrow$ & \textit{2000} & /rai{\texttheta}elagorim/&$\rightarrow$ & \textit{3200} & /rai{\texttheta}lagorm/&$\rightarrow$ & \textit{5500} & /ri{\texttheta}lagorm/&$\rightarrow$ & \textit{7600} & /ri{\texttheta}lagom/&$\rightarrow$ & \textit{12500} & /ri{\texttheta}agom/& \\
\end{tabular}

\vspace{20pt}\hline

\end{nopagebreak}
\filbreak



\vspace{15pt}
\begin{nopagebreak}
\noindent{\fontsize{20pt}{10pt}\textbf{mlo} } \textit{second} (Ns)\\
\noindent {\tovian \fontsize{20pt}{10pt} \textbf{mlo} }\\
\noindent /ml{\textprimstress}o/\\
\noindent lit. ESS+two\\


\noindent History:

\vspace{-0pt}
\hspace{40pt}
\begin{tabular}{ccc}
\textit{0} & /pu-loe/&$\rightarrow$ & \textit{1000} & /puloe/&$\rightarrow$ & \textit{3200} & /ploe/&$\rightarrow$ & \textit{9500} & /plo/&$\rightarrow$ & \textit{11000} & /mlo/& \\
\end{tabular}

\vspace{20pt}\hline

\end{nopagebreak}
\filbreak



\vspace{15pt}
\begin{nopagebreak}
\noindent{\fontsize{20pt}{10pt}\textbf{lorlesh} } \textit{September} (N)\\
\noindent {\tovian \fontsize{20pt}{10pt} \textbf{lorlesh} }\\
\noindent /l{\textprimstress}orle{\textesh}/\\
\noindent lit. one+harvest\\
\noindent \textit{Messihold}\\


\noindent History:

\vspace{-0pt}
\hspace{40pt}
\begin{tabular}{ccc}
\textit{6000} & /lomllre{\textyogh}/&$\rightarrow$ & \textit{6500} & /lomlre{\textyogh}/&$\rightarrow$ & \textit{6503} & /lomrle{\textyogh}/&$\rightarrow$ & \textit{7600} & /lorle{\textyogh}/&$\rightarrow$ & \textit{12001} & /lorle{\textesh}/& \\
\end{tabular}

\vspace{20pt}\hline

\end{nopagebreak}
\filbreak



\vspace{15pt}
\begin{nopagebreak}
\noindent{\fontsize{20pt}{10pt}\textbf{pol} } \textit{seven/7} (N)\\
\noindent {\tovian \fontsize{20pt}{10pt} \textbf{pol} }\\
\noindent /p{\textprimstress}ol/\\


\noindent History:

\vspace{-0pt}
\hspace{40pt}
\begin{tabular}{ccc}
\textit{0} & /pol/& \\
\end{tabular}

\vspace{20pt}\hline

\end{nopagebreak}
\filbreak



\vspace{15pt}
\begin{nopagebreak}
\noindent{\fontsize{20pt}{10pt}\textbf{lan} } \textit{shadow/twilight} (N)\\
\noindent {\tovian \fontsize{20pt}{10pt} \textbf{lan} }\\
\noindent /l{\textprimstress}an/\\
\noindent lit. darkness+light\\


\noindent History:

\vspace{-0pt}
\hspace{40pt}
\begin{tabular}{ccc}
\textit{500} & /bralahana/&$\rightarrow$ & \textit{3200} & /brlahana/&$\rightarrow$ & \textit{5000} & /brlaana/&$\rightarrow$ & \textit{5500} & /brlana/&$\rightarrow$ & \textit{7501} & /rlana/&$\rightarrow$ & \textit{7600} & /lana/&$\rightarrow$ & \textit{9500} & /lan/& \\
\end{tabular}

\vspace{20pt}\hline

\end{nopagebreak}
\filbreak



\vspace{15pt}
\begin{nopagebreak}
\noindent{\fontsize{20pt}{10pt}\textbf{nal} } \textit{sharing} (N)\\
\noindent {\tovian \fontsize{20pt}{10pt} \textbf{nal} }\\
\noindent /n{\textprimstress}al/\\
\noindent lit. joining+possession\\


\noindent History:

\vspace{-0pt}
\hspace{40pt}
\begin{tabular}{ccc}
\textit{500} & /{\texttoptiebar{t\textbeltl}}unaelis/&$\rightarrow$ & \textit{3200} & /{\texttoptiebar{t\textbeltl}}naels/&$\rightarrow$ & \textit{6501} & /naels/&$\rightarrow$ & \textit{10000} & /nals/&$\rightarrow$ & \textit{13001} & /nal/& \\
\end{tabular}

\vspace{20pt}\hline

\end{nopagebreak}
\filbreak



\vspace{15pt}
\begin{nopagebreak}
\noindent{\fontsize{20pt}{10pt}\textbf{shith} } \textit{sharpness/alertness/wakefulness/cutting} (N)\\
\noindent {\tovian \fontsize{20pt}{10pt} \textbf{shith} }\\
\noindent /{\textesh}{\textprimstress}i{\texttheta}/\\


\noindent History:

\vspace{-0pt}
\hspace{40pt}
\begin{tabular}{ccc}
\textit{0} & /{\textyogh}i{\texttheta}/&$\rightarrow$ & \textit{12001} & /{\textesh}i{\texttheta}/& \\
\end{tabular}

\vspace{20pt}\hline

\end{nopagebreak}
\filbreak



\vspace{15pt}
\begin{nopagebreak}
\noindent{\fontsize{20pt}{10pt}\textbf{munan} } \textit{shelter/protection/covering} (N)\\
\noindent {\tovian \fontsize{20pt}{10pt} \textbf{munan} }\\
\noindent /m{\textprimstress}unan/\\


\noindent History:

\vspace{-0pt}
\hspace{40pt}
\begin{tabular}{ccc}
\textit{0} & /benani/&$\rightarrow$ & \textit{3200} & /bnani/&$\rightarrow$ & \textit{9500} & /bnan/&$\rightarrow$ & \textit{11000} & /mnan/&$\rightarrow$ & \textit{11996} & /munan/& \\
\end{tabular}

\vspace{20pt}\hline

\end{nopagebreak}
\filbreak



\vspace{15pt}
\begin{nopagebreak}
\noindent{\fontsize{20pt}{10pt}\textbf{nalor} } \textit{shine} (N)\\
\noindent {\tovian \fontsize{20pt}{10pt} \textbf{nalor} }\\
\noindent /n{\textprimstress}alor/\\
\noindent lit. light+action\\


\noindent History:

\vspace{-0pt}
\hspace{40pt}
\begin{tabular}{ccc}
\textit{5000} & /aanalore/&$\rightarrow$ & \textit{5500} & /analore/&$\rightarrow$ & \textit{7500} & /nalore/&$\rightarrow$ & \textit{9500} & /nalor/& \\
\end{tabular}

\vspace{20pt}\hline

\end{nopagebreak}
\filbreak



\vspace{15pt}
\begin{nopagebreak}
\noindent{\fontsize{20pt}{10pt}\textbf{fed} } \textit{showing} (N)\\
\noindent {\tovian \fontsize{20pt}{10pt} \textbf{fed} }\\
\noindent /f{\textprimstress}ed/\\


\noindent History:

\vspace{-0pt}
\hspace{40pt}
\begin{tabular}{ccc}
\textit{0} & /fe{\textsubbridge{t}}e/&$\rightarrow$ & \textit{2000} & /fede/&$\rightarrow$ & \textit{9500} & /fed/& \\
\end{tabular}

\vspace{20pt}\hline

\end{nopagebreak}
\filbreak



\vspace{15pt}
\begin{nopagebreak}
\noindent{\fontsize{20pt}{10pt}\textbf{menan} } \textit{silver} (N)\\
\noindent {\tovian \fontsize{20pt}{10pt} \textbf{menan} }\\
\noindent /m{\textprimstress}enan/\\
\noindent lit. metal+light\\


\noindent History:

\vspace{-0pt}
\hspace{40pt}
\begin{tabular}{ccc}
\textit{4000} & /{\dh}nhmernahana/&$\rightarrow$ & \textit{5000} & /{\dh}nhmernaana/&$\rightarrow$ & \textit{5500} & /{\dh}nhmernana/&$\rightarrow$ & \textit{6400} & /{\dh}hmernana/&$\rightarrow$ & \textit{7501} & /hmernana/&$\rightarrow$ & \textit{7600} & /hmenana/&$\rightarrow$ & \textit{9500} & /hmenan/&$\rightarrow$ & \textit{13002} & /menan/& \\
\end{tabular}

\vspace{20pt}\hline

\end{nopagebreak}
\filbreak



\vspace{15pt}
\begin{nopagebreak}
\noindent{\fontsize{20pt}{10pt}\textbf{tem} } \textit{six/6} (N)\\
\noindent {\tovian \fontsize{20pt}{10pt} \textbf{tem} }\\
\noindent /{\textsubbridge{t}}{\textprimstress}em/\\


\noindent History:

\vspace{-0pt}
\hspace{40pt}
\begin{tabular}{ccc}
\textit{0} & /{\textsubbridge{t}}ep/&$\rightarrow$ & \textit{11000} & /{\textsubbridge{t}}em/& \\
\end{tabular}

\vspace{20pt}\hline

\end{nopagebreak}
\filbreak



\vspace{15pt}
\begin{nopagebreak}
\noindent{\fontsize{20pt}{10pt}\textbf{lhyatev} } \textit{sixteen/16} (N)\\
\noindent {\tovian \fontsize{20pt}{10pt} \textbf{lhyatev} }\\
\noindent /{\textbeltl}y{\textprimstress}a{\textsubbridge{t}}ev/\\
\noindent lit. twelve+four\\


\noindent History:

\vspace{-0pt}
\hspace{40pt}
\begin{tabular}{ccc}
\textit{4000} & /{\textbeltl}ya{\textsubbridge{t}}ev/& \\
\end{tabular}

\vspace{20pt}\hline

\end{nopagebreak}
\filbreak



\vspace{15pt}
\begin{nopagebreak}
\noindent{\fontsize{20pt}{10pt}\textbf{rithil} } \textit{sky} (N)\\
\noindent {\tovian \fontsize{20pt}{10pt} \textbf{rithil} }\\
\noindent /r{\textprimstress}i{\texttheta}il/\\
\noindent lit. air+place\\


\noindent History:

\vspace{-0pt}
\hspace{40pt}
\begin{tabular}{ccc}
\textit{10} & /rai{\texttheta}{\texttheta}ile/&$\rightarrow$ & \textit{4502} & /rai{\texttheta}ile/&$\rightarrow$ & \textit{5500} & /ri{\texttheta}ile/&$\rightarrow$ & \textit{9500} & /ri{\texttheta}il/& \\
\end{tabular}

\vspace{20pt}\hline

\end{nopagebreak}
\filbreak



\vspace{15pt}
\begin{nopagebreak}
\noindent{\fontsize{20pt}{10pt}\textbf{nin} } \textit{small} (N)\\
\noindent {\tovian \fontsize{20pt}{10pt} \textbf{nin} }\\
\noindent /n{\textprimstress}in/\\


\noindent History:

\vspace{-0pt}
\hspace{40pt}
\begin{tabular}{ccc}
\textit{0} & /nina/&$\rightarrow$ & \textit{9500} & /nin/& \\
\end{tabular}

\vspace{20pt}\hline

\end{nopagebreak}
\filbreak



\vspace{15pt}
\begin{nopagebreak}
\noindent{\fontsize{20pt}{10pt}\textbf{rithor} } \textit{smell} (N)\\
\noindent {\tovian \fontsize{20pt}{10pt} \textbf{rithor} }\\
\noindent /r{\textprimstress}i{\texttheta}or/\\
\noindent lit. air+action\\


\noindent History:

\vspace{-0pt}
\hspace{40pt}
\begin{tabular}{ccc}
\textit{500} & /rai{\texttheta}lore/&$\rightarrow$ & \textit{5500} & /ri{\texttheta}lore/&$\rightarrow$ & \textit{9500} & /ri{\texttheta}lor/&$\rightarrow$ & \textit{12500} & /ri{\texttheta}or/& \\
\end{tabular}

\vspace{20pt}\hline

\end{nopagebreak}
\filbreak



\vspace{15pt}
\begin{nopagebreak}
\noindent{\fontsize{20pt}{10pt}\textbf{maon} } \textit{smile} (N)\\
\noindent {\tovian \fontsize{20pt}{10pt} \textbf{maon} }\\
\noindent /m{\textprimstress}aon/\\
\noindent lit. happy+mouth\\


\noindent History:

\vspace{-0pt}
\hspace{40pt}
\begin{tabular}{ccc}
\textit{400} & /ma{\textsubbridge{t}}{\textsubbridge{t}}aeolan/&$\rightarrow$ & \textit{3000} & /ma{\textsubbridge{t}}aeolan/&$\rightarrow$ & \textit{3200} & /m{\textsubbridge{t}}aeoln/&$\rightarrow$ & \textit{7600} & /m{\textsubbridge{t}}aeon/&$\rightarrow$ & \textit{8750} & /maeon/&$\rightarrow$ & \textit{10000} & /maon/& \\
\end{tabular}

\vspace{20pt}\hline

\end{nopagebreak}
\filbreak



\vspace{15pt}
\begin{nopagebreak}
\noindent{\fontsize{20pt}{10pt}\textbf{lheth} } \textit{softness/gentleness/smoothness} (N)\\
\noindent {\tovian \fontsize{20pt}{10pt} \textbf{lheth} }\\
\noindent /{\textbeltl}{\textprimstress}e{\texttheta}/\\


\noindent History:

\vspace{-0pt}
\hspace{40pt}
\begin{tabular}{ccc}
\textit{0} & /{\textbeltl}e{\texttheta}/& \\
\end{tabular}

\vspace{20pt}\hline

\end{nopagebreak}
\filbreak



\vspace{15pt}
\begin{nopagebreak}
\noindent{\fontsize{20pt}{10pt}\textbf{tesh} } \textit{sound} (N)\\
\noindent {\tovian \fontsize{20pt}{10pt} \textbf{tesh} }\\
\noindent /{\textsubbridge{t}}{\textprimstress}e{\textesh}/\\


\noindent History:

\vspace{-0pt}
\hspace{40pt}
\begin{tabular}{ccc}
\textit{0} & /{\textsubbridge{t}}e{\textsubbridge{t}}e{\textyogh}ah/&$\rightarrow$ & \textit{1000} & /{\textsubbridge{t}}{\textsubbridge{t}}e{\textyogh}ah/&$\rightarrow$ & \textit{3000} & /{\textsubbridge{t}}e{\textyogh}ah/&$\rightarrow$ & \textit{3200} & /{\textsubbridge{t}}e{\textyogh}h/&$\rightarrow$ & \textit{4502} & /{\textsubbridge{t}}e{\textyogh}/&$\rightarrow$ & \textit{12001} & /{\textsubbridge{t}}e{\textesh}/& \\
\end{tabular}

\vspace{20pt}\hline

\end{nopagebreak}
\filbreak



\vspace{15pt}
\begin{nopagebreak}
\noindent{\fontsize{20pt}{10pt}\textbf{rar} } \textit{south} (N)\\
\noindent {\tovian \fontsize{20pt}{10pt} \textbf{rar} }\\
\noindent /r{\textprimstress}ar/\\


\noindent History:

\vspace{-0pt}
\hspace{40pt}
\begin{tabular}{ccc}
\textit{0} & /rar/& \\
\end{tabular}

\vspace{20pt}\hline

\end{nopagebreak}
\filbreak



\vspace{15pt}
\begin{nopagebreak}
\noindent{\fontsize{20pt}{10pt}\textbf{fethr} } \textit{speaking} (N)\\
\noindent {\tovian \fontsize{20pt}{10pt} \textbf{fethr} }\\
\noindent /f{\textprimstress}e{\texttheta}r/\\
\noindent lit. showing+thought\\


\noindent History:

\vspace{-0pt}
\hspace{40pt}
\begin{tabular}{ccc}
\textit{100} & /fe{\textsubbridge{t}}ei{\texttheta}ir/&$\rightarrow$ & \textit{1000} & /f{\textsubbridge{t}}ei{\texttheta}ir/&$\rightarrow$ & \textit{3000} & /fei{\texttheta}ir/&$\rightarrow$ & \textit{3200} & /fei{\texttheta}r/&$\rightarrow$ & \textit{5500} & /fe{\texttheta}r/& \\
\end{tabular}

\vspace{20pt}\hline

\end{nopagebreak}
\filbreak



\vspace{15pt}
\begin{nopagebreak}
\noindent{\fontsize{20pt}{10pt}\textbf{fethr} } \textit{speech} (N)\\
\noindent {\tovian \fontsize{20pt}{10pt} \textbf{fethr} }\\
\noindent /f{\textprimstress}e{\texttheta}r/\\


\noindent History:

\vspace{-0pt}
\hspace{40pt}
\begin{tabular}{ccc}
\textit{0} & /fe{\dh}ra/&$\rightarrow$ & \textit{9500} & /fe{\dh}r/&$\rightarrow$ & \textit{12002} & /fe{\texttheta}r/& \\
\end{tabular}

\vspace{20pt}\hline

\end{nopagebreak}
\filbreak



\vspace{15pt}
\begin{nopagebreak}
\noindent{\fontsize{20pt}{10pt}\textbf{tlesh} } \textit{speed} (N)\\
\noindent {\tovian \fontsize{20pt}{10pt} \textbf{tlesh} }\\
\noindent /{\texttoptiebar{t\textbeltl}}{\textprimstress}e{\textesh}/\\


\noindent History:

\vspace{-0pt}
\hspace{40pt}
\begin{tabular}{ccc}
\textit{0} & /{\texttoptiebar{t\textbeltl}}e{\textyogh}/&$\rightarrow$ & \textit{12001} & /{\texttoptiebar{t\textbeltl}}e{\textesh}/& \\
\end{tabular}

\vspace{20pt}\hline

\end{nopagebreak}
\filbreak



\vspace{15pt}
\begin{nopagebreak}
\noindent{\fontsize{20pt}{10pt}\textbf{mlhelor} } \textit{Spring} (N)\\
\noindent {\tovian \fontsize{20pt}{10pt} \textbf{mlhelor} }\\
\noindent /m{\textbeltl}{\textprimstress}elor/\\
\noindent lit. present+plant\\


\noindent History:

\vspace{-0pt}
\hspace{40pt}
\begin{tabular}{ccc}
\textit{0} & /num{\textbeltl}ele{\textschwa}mellore/&$\rightarrow$ & \textit{2300} & /num{\textbeltl}elemellore/&$\rightarrow$ & \textit{3200} & /nm{\textbeltl}lemllore/&$\rightarrow$ & \textit{6500} & /nm{\textbeltl}lemlore/&$\rightarrow$ & \textit{7501} & /m{\textbeltl}lemlore/&$\rightarrow$ & \textit{7600} & /m{\textbeltl}lelore/&$\rightarrow$ & \textit{9500} & /m{\textbeltl}lelor/&$\rightarrow$ & \textit{12500} & /m{\textbeltl}elor/& \\
\end{tabular}

\vspace{20pt}\hline

\end{nopagebreak}
\filbreak



\vspace{15pt}
\begin{nopagebreak}
\noindent{\fontsize{20pt}{10pt}\textbf{ledul} } \textit{spring vegetable} (N)\\
\noindent {\tovian \fontsize{20pt}{10pt} \textbf{ledul} }\\
\noindent /l{\textprimstress}edul/\\
\noindent lit. plant+first\\


\noindent History:

\vspace{-0pt}
\hspace{40pt}
\begin{tabular}{ccc}
\textit{100} & /{\textschwa}mellorepu-lo/&$\rightarrow$ & \textit{1000} & /{\textschwa}mellorepulo/&$\rightarrow$ & \textit{2000} & /{\textschwa}mellorebulo/&$\rightarrow$ & \textit{2300} & /mellorebulo/&$\rightarrow$ & \textit{3200} & /mllrebulo/&$\rightarrow$ & \textit{6500} & /mlrebulo/&$\rightarrow$ & \textit{6503} & /mrlebulo/&$\rightarrow$ & \textit{7501} & /rlebulo/&$\rightarrow$ & \textit{7600} & /lebulo/&$\rightarrow$ & \textit{9500} & /lebul/&$\rightarrow$ & \textit{11001} & /ledul/& \\
\end{tabular}

\vspace{20pt}\hline

\end{nopagebreak}
\filbreak



\vspace{15pt}
\begin{nopagebreak}
\noindent{\fontsize{20pt}{10pt}\textbf{nan} } \textit{star} (N)\\
\noindent {\tovian \fontsize{20pt}{10pt} \textbf{nan} }\\
\noindent /n{\textprimstress}an/\\
\noindent lit. small+light\\


\noindent History:

\vspace{-0pt}
\hspace{40pt}
\begin{tabular}{ccc}
\textit{500} & /ninaahana/&$\rightarrow$ & \textit{3200} & /nnaahana/&$\rightarrow$ & \textit{5000} & /nnaaana/&$\rightarrow$ & \textit{5500} & /nnaana/&$\rightarrow$ & \textit{6500} & /naana/&$\rightarrow$ & \textit{9500} & /naan/&$\rightarrow$ & \textit{12005} & /nan/& \\
\end{tabular}

\vspace{20pt}\hline

\end{nopagebreak}
\filbreak



\vspace{15pt}
\begin{nopagebreak}
\noindent{\fontsize{20pt}{10pt}\textbf{ing} } \textit{stickiness} (N)\\
\noindent {\tovian \fontsize{20pt}{10pt} \textbf{ing} }\\
\noindent /{\textprimstress}i{\ng}/\\


\noindent History:

\vspace{-0pt}
\hspace{40pt}
\begin{tabular}{ccc}
\textit{0} & /inga/&$\rightarrow$ & \textit{6000} & /i{\ng}ga/&$\rightarrow$ & \textit{8750} & /i{\ng}a/&$\rightarrow$ & \textit{9500} & /i{\ng}/& \\
\end{tabular}

\vspace{20pt}\hline

\end{nopagebreak}
\filbreak



\vspace{15pt}
\begin{nopagebreak}
\noindent{\fontsize{20pt}{10pt}\textbf{lor} } \textit{stomach} (N)\\
\noindent {\tovian \fontsize{20pt}{10pt} \textbf{lor} }\\
\noindent /l{\textprimstress}or/\\


\noindent History:

\vspace{-0pt}
\hspace{40pt}
\begin{tabular}{ccc}
\textit{0} & /sildabore/&$\rightarrow$ & \textit{3200} & /sldbore/&$\rightarrow$ & \textit{4501} & /slbore/&$\rightarrow$ & \textit{7501} & /lbore/&$\rightarrow$ & \textit{8750} & /lore/&$\rightarrow$ & \textit{9500} & /lor/& \\
\end{tabular}

\vspace{20pt}\hline

\end{nopagebreak}
\filbreak



\vspace{15pt}
\begin{nopagebreak}
\noindent{\fontsize{20pt}{10pt}\textbf{men} } \textit{stone/rock/earth} (N)\\
\noindent {\tovian \fontsize{20pt}{10pt} \textbf{men} }\\
\noindent /m{\textprimstress}en/\\


\noindent History:

\vspace{-0pt}
\hspace{40pt}
\begin{tabular}{ccc}
\textit{0} & /meron/&$\rightarrow$ & \textit{3200} & /mern/&$\rightarrow$ & \textit{7600} & /men/& \\
\end{tabular}

\vspace{20pt}\hline

\end{nopagebreak}
\filbreak



\vspace{15pt}
\begin{nopagebreak}
\noindent{\fontsize{20pt}{10pt}\textbf{pal} } \textit{stopping/ceasing} (N)\\
\noindent {\tovian \fontsize{20pt}{10pt} \textbf{pal} }\\
\noindent /p{\textprimstress}al/\\


\noindent History:

\vspace{-0pt}
\hspace{40pt}
\begin{tabular}{ccc}
\textit{0} & /palor/&$\rightarrow$ & \textit{3200} & /palr/&$\rightarrow$ & \textit{6503} & /parl/&$\rightarrow$ & \textit{7600} & /pal/& \\
\end{tabular}

\vspace{20pt}\hline

\end{nopagebreak}
\filbreak



\vspace{15pt}
\begin{nopagebreak}
\noindent{\fontsize{20pt}{10pt}\textbf{lrith} } \textit{storm} (N)\\
\noindent {\tovian \fontsize{20pt}{10pt} \textbf{lrith} }\\
\noindent /lr{\textprimstress}i{\texttheta}/\\
\noindent lit. wind+water\\


\noindent History:

\vspace{-0pt}
\hspace{40pt}
\begin{tabular}{ccc}
\textit{50} & /kwilyi-rai{\texttheta}z{\textsubbridge{t}}el/&$\rightarrow$ & \textit{1000} & /kwilyirai{\texttheta}z{\textsubbridge{t}}el/&$\rightarrow$ & \textit{2000} & /kwilyirai{\texttheta}zdel/&$\rightarrow$ & \textit{3200} & /kwlyrai{\texttheta}zdl/&$\rightarrow$ & \textit{4500} & /kwlyrai{\texttheta}zl/&$\rightarrow$ & \textit{4502} & /kwlyrai{\texttheta}l/&$\rightarrow$ & \textit{5500} & /kwlyri{\texttheta}l/&$\rightarrow$ & \textit{7501} & /wlyri{\texttheta}l/&$\rightarrow$ & \textit{7600} & /lri{\texttheta}l/&$\rightarrow$ & \textit{12500} & /lri{\texttheta}/& \\
\end{tabular}

\vspace{20pt}\hline

\end{nopagebreak}
\filbreak



\vspace{15pt}
\begin{nopagebreak}
\noindent{\fontsize{20pt}{10pt}\textbf{wi-} } \textit{SUB/subjunctive mood marker} (MOOD)\\
\noindent {\tovian \fontsize{20pt}{10pt} \textbf{wi-} }\\
\noindent /w{\textprimstress}i/\\


\noindent History:

\vspace{-0pt}
\hspace{40pt}
\begin{tabular}{ccc}
\textit{0} & /wi-/& \\
\end{tabular}

\vspace{20pt}\hline

\end{nopagebreak}
\filbreak



\vspace{15pt}
\begin{nopagebreak}
\noindent{\fontsize{20pt}{10pt}\textbf{mlheag} } \textit{Summer} (N)\\
\noindent {\tovian \fontsize{20pt}{10pt} \textbf{mlheag} }\\
\noindent /m{\textbeltl}{\textprimstress}eag/\\
\noindent lit. present+heat\\


\noindent History:

\vspace{-0pt}
\hspace{40pt}
\begin{tabular}{ccc}
\textit{0} & /num{\textbeltl}eleasaka/&$\rightarrow$ & \textit{2000} & /num{\textbeltl}eleasaga/&$\rightarrow$ & \textit{3200} & /nm{\textbeltl}leasaga/&$\rightarrow$ & \textit{5100} & /nm{\textbeltl}leaaga/&$\rightarrow$ & \textit{5500} & /nm{\textbeltl}leaga/&$\rightarrow$ & \textit{7501} & /m{\textbeltl}leaga/&$\rightarrow$ & \textit{9500} & /m{\textbeltl}leag/&$\rightarrow$ & \textit{12500} & /m{\textbeltl}eag/& \\
\end{tabular}

\vspace{20pt}\hline

\end{nopagebreak}
\filbreak



\vspace{15pt}
\begin{nopagebreak}
\noindent{\fontsize{20pt}{10pt}\textbf{ef} } \textit{sun} (N)\\
\noindent {\tovian \fontsize{20pt}{10pt} \textbf{ef} }\\
\noindent /{\textprimstress}ef/\\


\noindent History:

\vspace{-0pt}
\hspace{40pt}
\begin{tabular}{ccc}
\textit{0} & /eifol/&$\rightarrow$ & \textit{3200} & /eifl/&$\rightarrow$ & \textit{5500} & /efl/&$\rightarrow$ & \textit{12500} & /ef/& \\
\end{tabular}

\vspace{20pt}\hline

\end{nopagebreak}
\filbreak



\vspace{15pt}
\begin{nopagebreak}
\noindent{\fontsize{20pt}{10pt}\textbf{ngam} } \textit{surpassing/exceeding/overwhelming} (N)\\
\noindent {\tovian \fontsize{20pt}{10pt} \textbf{ngam} }\\
\noindent /{\ng}{\textprimstress}am/\\


\noindent History:

\vspace{-0pt}
\hspace{40pt}
\begin{tabular}{ccc}
\textit{0} & /{\ng}o{\ng}abi/&$\rightarrow$ & \textit{3200} & /{\ng}{\ng}abi/&$\rightarrow$ & \textit{6500} & /{\ng}abi/&$\rightarrow$ & \textit{9500} & /{\ng}ab/&$\rightarrow$ & \textit{11000} & /{\ng}am/& \\
\end{tabular}

\vspace{20pt}\hline

\end{nopagebreak}
\filbreak



\vspace{15pt}
\begin{nopagebreak}
\noindent{\fontsize{20pt}{10pt}\textbf{hol} } \textit{sweet} (N)\\
\noindent {\tovian \fontsize{20pt}{10pt} \textbf{hol} }\\
\noindent /ħ{\textprimstress}ol/\\
\noindent lit. GEN+fruit\\


\noindent History:

\vspace{-0pt}
\hspace{40pt}
\begin{tabular}{ccc}
\textit{400} & /hi-{\textyogh}ola/&$\rightarrow$ & \textit{1000} & /hi{\textyogh}ola/&$\rightarrow$ & \textit{3200} & /h{\textyogh}ola/&$\rightarrow$ & \textit{4502} & /hola/&$\rightarrow$ & \textit{9500} & /hol/&$\rightarrow$ & \textit{14006} & /ħol/& \\
\end{tabular}

\vspace{20pt}\hline

\end{nopagebreak}
\filbreak



\vspace{15pt}
\begin{nopagebreak}
\noindent{\fontsize{20pt}{10pt}\textbf{thuv} } \textit{sword} (N)\\
\noindent {\tovian \fontsize{20pt}{10pt} \textbf{thuv} }\\
\noindent /{\texttheta}{\textprimstress}uv/\\
\noindent lit. tool+PUR+killing\\


\noindent History:

\vspace{-0pt}
\hspace{40pt}
\begin{tabular}{ccc}
\textit{50} & /{\dh}an{\textbeltl}u-vekh/&$\rightarrow$ & \textit{1000} & /{\dh}an{\textbeltl}uvekh/&$\rightarrow$ & \textit{3200} & /{\dh}n{\textbeltl}uvkh/&$\rightarrow$ & \textit{4500} & /{\dh}n{\textbeltl}uvh/&$\rightarrow$ & \textit{4502} & /{\dh}n{\textbeltl}uv/&$\rightarrow$ & \textit{6400} & /{\dh}{\textbeltl}uv/&$\rightarrow$ & \textit{12002} & /{\texttheta}{\textbeltl}uv/&$\rightarrow$ & \textit{13000} & /{\texttheta}uv/& \\
\end{tabular}

\vspace{20pt}\hline

\end{nopagebreak}
\filbreak



\vspace{15pt}
\begin{nopagebreak}
\noindent{\fontsize{20pt}{10pt}\textbf{tlun} } \textit{tar} (N)\\
\noindent {\tovian \fontsize{20pt}{10pt} \textbf{tlun} }\\
\noindent /{\texttoptiebar{t\textbeltl}}{\textprimstress}un/\\
\noindent lit. stickiness+joining\\


\noindent History:

\vspace{-0pt}
\hspace{40pt}
\begin{tabular}{ccc}
\textit{50} & /inga{\texttoptiebar{t\textbeltl}}una/&$\rightarrow$ & \textit{3200} & /ing{\texttoptiebar{t\textbeltl}}una/&$\rightarrow$ & \textit{6000} & /i{\ng}g{\texttoptiebar{t\textbeltl}}una/&$\rightarrow$ & \textit{7500} & /{\ng}g{\texttoptiebar{t\textbeltl}}una/&$\rightarrow$ & \textit{7501} & /g{\texttoptiebar{t\textbeltl}}una/&$\rightarrow$ & \textit{9500} & /g{\texttoptiebar{t\textbeltl}}un/&$\rightarrow$ & \textit{11501} & /{\texttoptiebar{t\textbeltl}}un/& \\
\end{tabular}

\vspace{20pt}\hline

\end{nopagebreak}
\filbreak



\vspace{15pt}
\begin{nopagebreak}
\noindent{\fontsize{20pt}{10pt}\textbf{lhad} } \textit{teaching} (N)\\
\noindent {\tovian \fontsize{20pt}{10pt} \textbf{lhad} }\\
\noindent /{\textbeltl}{\textprimstress}ad/\\
\noindent lit. knowledge+showing\\


\noindent History:

\vspace{-0pt}
\hspace{40pt}
\begin{tabular}{ccc}
\textit{0} & /{\textbeltl}arafe{\textsubbridge{t}}e/&$\rightarrow$ & \textit{2000} & /{\textbeltl}arafede/&$\rightarrow$ & \textit{3200} & /{\textbeltl}rafede/&$\rightarrow$ & \textit{9500} & /{\textbeltl}rafed/&$\rightarrow$ & \textit{12003} & /{\textbeltl}rad/&$\rightarrow$ & \textit{12500} & /{\textbeltl}ad/& \\
\end{tabular}

\vspace{20pt}\hline

\end{nopagebreak}
\filbreak



\vspace{15pt}
\begin{nopagebreak}
\noindent{\fontsize{20pt}{10pt}\textbf{tu-} } \textit{TEMP/temporal case marker} (CASE)\\
\noindent {\tovian \fontsize{20pt}{10pt} \textbf{tu-} }\\
\noindent /{\textsubbridge{t}}{\textprimstress}u/\\


\noindent History:

\vspace{-0pt}
\hspace{40pt}
\begin{tabular}{ccc}
\textit{0} & /{\textsubbridge{t}}u-/& \\
\end{tabular}

\vspace{20pt}\hline

\end{nopagebreak}
\filbreak



\vspace{15pt}
\begin{nopagebreak}
\noindent{\fontsize{20pt}{10pt}\textbf{tir} } \textit{ten/10} (N)\\
\noindent {\tovian \fontsize{20pt}{10pt} \textbf{tir} }\\
\noindent /{\textsubbridge{t}}{\textprimstress}ir/\\


\noindent History:

\vspace{-0pt}
\hspace{40pt}
\begin{tabular}{ccc}
\textit{0} & /{\textsubbridge{t}}ir/& \\
\end{tabular}

\vspace{20pt}\hline

\end{nopagebreak}
\filbreak



\vspace{15pt}
\begin{nopagebreak}
\noindent{\fontsize{20pt}{10pt}\textbf{lod} } \textit{them (inanimate)} (Ps)\\
\noindent {\tovian \fontsize{20pt}{10pt} \textbf{lod} }\\
\noindent /l{\textprimstress}od/\\


\noindent History:

\vspace{-0pt}
\hspace{40pt}
\begin{tabular}{ccc}
\textit{0} & /lo{\textsubbridge{t}}e/&$\rightarrow$ & \textit{2000} & /lode/&$\rightarrow$ & \textit{9500} & /lod/& \\
\end{tabular}

\vspace{20pt}\hline

\end{nopagebreak}
\filbreak



\vspace{15pt}
\begin{nopagebreak}
\noindent{\fontsize{20pt}{10pt}\textbf{t} } \textit{they (plural)} (Ps)\\
\noindent {\tovian \fontsize{20pt}{10pt} \textbf{t} }\\
\noindent /{\textsubbridge{t}}/\\


\noindent History:

\vspace{-0pt}
\hspace{40pt}
\begin{tabular}{ccc}
\textit{0} & /{\textsubbridge{t}}e/&$\rightarrow$ & \textit{12006} & /{\textsubbridge{t}}/& \\
\end{tabular}

\vspace{20pt}\hline

\end{nopagebreak}
\filbreak



\vspace{15pt}
\begin{nopagebreak}
\noindent{\fontsize{20pt}{10pt}\textbf{lhul} } \textit{thickness/denseness} (N)\\
\noindent {\tovian \fontsize{20pt}{10pt} \textbf{lhul} }\\
\noindent /{\textbeltl}{\textprimstress}ul/\\


\noindent History:

\vspace{-0pt}
\hspace{40pt}
\begin{tabular}{ccc}
\textit{0} & /{\textbeltl}unol/&$\rightarrow$ & \textit{3200} & /{\textbeltl}unl/&$\rightarrow$ & \textit{7600} & /{\textbeltl}ul/& \\
\end{tabular}

\vspace{20pt}\hline

\end{nopagebreak}
\filbreak



\vspace{15pt}
\begin{nopagebreak}
\noindent{\fontsize{20pt}{10pt}\textbf{l} } \textit{thing} (N)\\
\noindent {\tovian \fontsize{20pt}{10pt} \textbf{l} }\\
\noindent /l/\\


\noindent History:

\vspace{-0pt}
\hspace{40pt}
\begin{tabular}{ccc}
\textit{0} & /{\textsubbridge{t}}iel/&$\rightarrow$ & \textit{2201} & /{\textsubbridge{t}}iyel/&$\rightarrow$ & \textit{3200} & /{\textsubbridge{t}}iyl/&$\rightarrow$ & \textit{6300} & /{\textsubbridge{t}}yl/&$\rightarrow$ & \textit{7501} & /yl/&$\rightarrow$ & \textit{7600} & /l/& \\
\end{tabular}

\vspace{20pt}\hline

\end{nopagebreak}
\filbreak



\vspace{15pt}
\begin{nopagebreak}
\noindent{\fontsize{20pt}{10pt}\textbf{nalh} } \textit{thinking/believing/pondering} (N)\\
\noindent {\tovian \fontsize{20pt}{10pt} \textbf{nalh} }\\
\noindent /n{\textprimstress}a{\textbeltl}/\\


\noindent History:

\vspace{-0pt}
\hspace{40pt}
\begin{tabular}{ccc}
\textit{0} & /na{\textbeltl}/& \\
\end{tabular}

\vspace{20pt}\hline

\end{nopagebreak}
\filbreak



\vspace{15pt}
\begin{nopagebreak}
\noindent{\fontsize{20pt}{10pt}\textbf{lhyal} } \textit{thirteen/13} (N)\\
\noindent {\tovian \fontsize{20pt}{10pt} \textbf{lhyal} }\\
\noindent /{\textbeltl}y{\textprimstress}al/\\
\noindent lit. twelve+one\\


\noindent History:

\vspace{-0pt}
\hspace{40pt}
\begin{tabular}{ccc}
\textit{4000} & /{\textbeltl}yalo/&$\rightarrow$ & \textit{9500} & /{\textbeltl}yal/& \\
\end{tabular}

\vspace{20pt}\hline

\end{nopagebreak}
\filbreak



\vspace{15pt}
\begin{nopagebreak}
\noindent{\fontsize{20pt}{10pt}\textbf{femlh} } \textit{thirty six/36} (N)\\
\noindent {\tovian \fontsize{20pt}{10pt} \textbf{femlh} }\\
\noindent /f{\textprimstress}em{\textbeltl}/\\
\noindent lit. three+twelve\\


\noindent History:

\vspace{-0pt}
\hspace{40pt}
\begin{tabular}{ccc}
\textit{4000} & /fem{\textbeltl}ya/&$\rightarrow$ & \textit{9500} & /fem{\textbeltl}y/&$\rightarrow$ & \textit{13001} & /fem{\textbeltl}/& \\
\end{tabular}

\vspace{20pt}\hline

\end{nopagebreak}
\filbreak



\vspace{15pt}
\begin{nopagebreak}
\noindent{\fontsize{20pt}{10pt}\textbf{loelhyatem} } \textit{thirty/30} (N)\\
\noindent {\tovian \fontsize{20pt}{10pt} \textbf{loelhyatem} }\\
\noindent /loe{\textbeltl}y{\textprimstress}a{\textsubbridge{t}}em/\\
\noindent lit. twenty four+six\\


\noindent History:

\vspace{-0pt}
\hspace{40pt}
\begin{tabular}{ccc}
\textit{4000} & /loe{\textbeltl}ya{\textsubbridge{t}}ep/&$\rightarrow$ & \textit{11000} & /loe{\textbeltl}ya{\textsubbridge{t}}em/& \\
\end{tabular}

\vspace{20pt}\hline

\end{nopagebreak}
\filbreak



\vspace{15pt}
\begin{nopagebreak}
\noindent{\fontsize{20pt}{10pt}\textbf{ithr} } \textit{thought} (N)\\
\noindent {\tovian \fontsize{20pt}{10pt} \textbf{ithr} }\\
\noindent /{\textprimstress}i{\texttheta}r/\\


\noindent History:

\vspace{-0pt}
\hspace{40pt}
\begin{tabular}{ccc}
\textit{0} & /i{\texttheta}ir/&$\rightarrow$ & \textit{3200} & /i{\texttheta}r/& \\
\end{tabular}

\vspace{20pt}\hline

\end{nopagebreak}
\filbreak



\vspace{15pt}
\begin{nopagebreak}
\noindent{\fontsize{20pt}{10pt}\textbf{fem} } \textit{three/3} (N)\\
\noindent {\tovian \fontsize{20pt}{10pt} \textbf{fem} }\\
\noindent /f{\textprimstress}em/\\


\noindent History:

\vspace{-0pt}
\hspace{40pt}
\begin{tabular}{ccc}
\textit{0} & /fem/& \\
\end{tabular}

\vspace{20pt}\hline

\end{nopagebreak}
\filbreak



\vspace{15pt}
\begin{nopagebreak}
\noindent{\fontsize{20pt}{10pt}\textbf{ysha} } \textit{throwing} (N)\\
\noindent {\tovian \fontsize{20pt}{10pt} \textbf{ysha} }\\
\noindent /y{\textesh}{\textprimstress}a/\\


\noindent History:

\vspace{-0pt}
\hspace{40pt}
\begin{tabular}{ccc}
\textit{0} & /huy{\textyogh}a/&$\rightarrow$ & \textit{2100} & /hy{\textyogh}a/&$\rightarrow$ & \textit{7501} & /y{\textyogh}a/&$\rightarrow$ & \textit{12001} & /y{\textesh}a/& \\
\end{tabular}

\vspace{20pt}\hline

\end{nopagebreak}
\filbreak



\vspace{15pt}
\begin{nopagebreak}
\noindent{\fontsize{20pt}{10pt}\textbf{lrithesh} } \textit{thunder} (N)\\
\noindent {\tovian \fontsize{20pt}{10pt} \textbf{lrithesh} }\\
\noindent /lr{\textprimstress}i{\texttheta}e{\textesh}/\\
\noindent lit. storm+sound\\


\noindent History:

\vspace{-0pt}
\hspace{40pt}
\begin{tabular}{ccc}
\textit{100} & /kwilyi-rai{\texttheta}z{\textsubbridge{t}}el{\textsubbridge{t}}e{\textsubbridge{t}}e{\textyogh}ah/&$\rightarrow$ & \textit{1000} & /kwilyirai{\texttheta}z{\textsubbridge{t}}el{\textsubbridge{t}}{\textsubbridge{t}}e{\textyogh}ah/&$\rightarrow$ & \textit{2000} & /kwilyirai{\texttheta}zdel{\textsubbridge{t}}{\textsubbridge{t}}e{\textyogh}ah/&$\rightarrow$ & \textit{3000} & /kwilyirai{\texttheta}zdel{\textsubbridge{t}}e{\textyogh}ah/&$\rightarrow$ & \textit{3200} & /kwlyrai{\texttheta}zdl{\textsubbridge{t}}e{\textyogh}h/&$\rightarrow$ & \textit{4500} & /kwlyrai{\texttheta}zl{\textsubbridge{t}}e{\textyogh}h/&$\rightarrow$ & \textit{4501} & /kwlyrai{\texttheta}zle{\textyogh}h/&$\rightarrow$ & \textit{4502} & /kwlyrai{\texttheta}le{\textyogh}/&$\rightarrow$ & \textit{5500} & /kwlyri{\texttheta}le{\textyogh}/&$\rightarrow$ & \textit{7501} & /wlyri{\texttheta}le{\textyogh}/&$\rightarrow$ & \textit{7600} & /lri{\texttheta}le{\textyogh}/&$\rightarrow$ & \textit{12001} & /lri{\texttheta}le{\textesh}/&$\rightarrow$ & \textit{12500} & /lri{\texttheta}e{\textesh}/& \\
\end{tabular}

\vspace{20pt}\hline

\end{nopagebreak}
\filbreak



\vspace{15pt}
\begin{nopagebreak}
\noindent{\fontsize{20pt}{10pt}\textbf{nalor} } \textit{tie/bind/knot} (N)\\
\noindent {\tovian \fontsize{20pt}{10pt} \textbf{nalor} }\\
\noindent /n{\textprimstress}alor/\\
\noindent lit. joining+action\\


\noindent History:

\vspace{-0pt}
\hspace{40pt}
\begin{tabular}{ccc}
\textit{10} & /{\texttoptiebar{t\textbeltl}}unalore/&$\rightarrow$ & \textit{3200} & /{\texttoptiebar{t\textbeltl}}nalore/&$\rightarrow$ & \textit{6501} & /nalore/&$\rightarrow$ & \textit{9500} & /nalor/& \\
\end{tabular}

\vspace{20pt}\hline

\end{nopagebreak}
\filbreak



\vspace{15pt}
\begin{nopagebreak}
\noindent{\fontsize{20pt}{10pt}\textbf{ngothr} } \textit{time} (N)\\
\noindent {\tovian \fontsize{20pt}{10pt} \textbf{ngothr} }\\
\noindent /{\ng}{\textprimstress}o{\texttheta}r/\\


\noindent History:

\vspace{-0pt}
\hspace{40pt}
\begin{tabular}{ccc}
\textit{0} & /nko{\texttheta}ir/&$\rightarrow$ & \textit{2000} & /ngo{\texttheta}ir/&$\rightarrow$ & \textit{3200} & /ngo{\texttheta}r/&$\rightarrow$ & \textit{6000} & /{\ng}go{\texttheta}r/&$\rightarrow$ & \textit{8750} & /{\ng}o{\texttheta}r/& \\
\end{tabular}

\vspace{20pt}\hline

\end{nopagebreak}
\filbreak



\vspace{15pt}
\begin{nopagebreak}
\noindent{\fontsize{20pt}{10pt}\textbf{than} } \textit{tool} (N)\\
\noindent {\tovian \fontsize{20pt}{10pt} \textbf{than} }\\
\noindent /{\texttheta}{\textprimstress}an/\\


\noindent History:

\vspace{-0pt}
\hspace{40pt}
\begin{tabular}{ccc}
\textit{0} & /{\dh}an/&$\rightarrow$ & \textit{12002} & /{\texttheta}an/& \\
\end{tabular}

\vspace{20pt}\hline

\end{nopagebreak}
\filbreak



\vspace{15pt}
\begin{nopagebreak}
\noindent{\fontsize{20pt}{10pt}\textbf{gang} } \textit{tooth} (N)\\
\noindent {\tovian \fontsize{20pt}{10pt} \textbf{gang} }\\
\noindent /g{\textprimstress}a{\ng}/\\


\noindent History:

\vspace{-0pt}
\hspace{40pt}
\begin{tabular}{ccc}
\textit{0} & /ganga/&$\rightarrow$ & \textit{6000} & /ga{\ng}ga/&$\rightarrow$ & \textit{8750} & /ga{\ng}a/&$\rightarrow$ & \textit{9500} & /ga{\ng}/& \\
\end{tabular}

\vspace{20pt}\hline

\end{nopagebreak}
\filbreak



\vspace{15pt}
\begin{nopagebreak}
\noindent{\fontsize{20pt}{10pt}\textbf{thul} } \textit{top/height} (N)\\
\noindent {\tovian \fontsize{20pt}{10pt} \textbf{thul} }\\
\noindent /{\texttheta}{\textprimstress}ul/\\


\noindent History:

\vspace{-0pt}
\hspace{40pt}
\begin{tabular}{ccc}
\textit{0} & /{\dh}ul/&$\rightarrow$ & \textit{12002} & /{\texttheta}ul/& \\
\end{tabular}

\vspace{20pt}\hline

\end{nopagebreak}
\filbreak



\vspace{15pt}
\begin{nopagebreak}
\noindent{\fontsize{20pt}{10pt}\textbf{pel} } \textit{town/village} (N)\\
\noindent {\tovian \fontsize{20pt}{10pt} \textbf{pel} }\\
\noindent /p{\textprimstress}el/\\


\noindent History:

\vspace{-0pt}
\hspace{40pt}
\begin{tabular}{ccc}
\textit{0} & /pelor/&$\rightarrow$ & \textit{3200} & /pelr/&$\rightarrow$ & \textit{6503} & /perl/&$\rightarrow$ & \textit{7600} & /pel/& \\
\end{tabular}

\vspace{20pt}\hline

\end{nopagebreak}
\filbreak



\vspace{15pt}
\begin{nopagebreak}
\noindent{\fontsize{20pt}{10pt}\textbf{lheth} } \textit{trance/sleep} (N)\\
\noindent {\tovian \fontsize{20pt}{10pt} \textbf{lheth} }\\
\noindent /{\textbeltl}{\textprimstress}e{\texttheta}/\\


\noindent History:

\vspace{-0pt}
\hspace{40pt}
\begin{tabular}{ccc}
\textit{0} & /{\textbeltl}e{\dh}el/&$\rightarrow$ & \textit{3200} & /{\textbeltl}e{\dh}l/&$\rightarrow$ & \textit{12002} & /{\textbeltl}e{\texttheta}l/&$\rightarrow$ & \textit{12500} & /{\textbeltl}e{\texttheta}/& \\
\end{tabular}

\vspace{20pt}\hline

\end{nopagebreak}
\filbreak



\vspace{15pt}
\begin{nopagebreak}
\noindent{\fontsize{20pt}{10pt}\textbf{glar} } \textit{tree} (N)\\
\noindent {\tovian \fontsize{20pt}{10pt} \textbf{glar} }\\
\noindent /gl{\textprimstress}ar/\\


\noindent History:

\vspace{-0pt}
\hspace{40pt}
\begin{tabular}{ccc}
\textit{0} & /aklare/&$\rightarrow$ & \textit{2000} & /aglare/&$\rightarrow$ & \textit{7500} & /glare/&$\rightarrow$ & \textit{9500} & /glar/& \\
\end{tabular}

\vspace{20pt}\hline

\end{nopagebreak}
\filbreak



\vspace{15pt}
\begin{nopagebreak}
\noindent{\fontsize{20pt}{10pt}\textbf{lhya} } \textit{twelve/12} (N)\\
\noindent {\tovian \fontsize{20pt}{10pt} \textbf{lhya} }\\
\noindent /{\textbeltl}y{\textprimstress}a/\\


\noindent History:

\vspace{-0pt}
\hspace{40pt}
\begin{tabular}{ccc}
\textit{0} & /{\textbeltl}iya/&$\rightarrow$ & \textit{2100} & /{\textbeltl}ya/& \\
\end{tabular}

\vspace{20pt}\hline

\end{nopagebreak}
\filbreak



\vspace{15pt}
\begin{nopagebreak}
\noindent{\fontsize{20pt}{10pt}\textbf{loelhyal} } \textit{twenty five/25} (N)\\
\noindent {\tovian \fontsize{20pt}{10pt} \textbf{loelhyal} }\\
\noindent /lo{\textprimstress}e{\textbeltl}yal/\\
\noindent lit. twenty four+one\\


\noindent History:

\vspace{-0pt}
\hspace{40pt}
\begin{tabular}{ccc}
\textit{4000} & /loe{\textbeltl}yalo/&$\rightarrow$ & \textit{9500} & /loe{\textbeltl}yal/& \\
\end{tabular}

\vspace{20pt}\hline

\end{nopagebreak}
\filbreak



\vspace{15pt}
\begin{nopagebreak}
\noindent{\fontsize{20pt}{10pt}\textbf{loelh} } \textit{twenty four/24} (N)\\
\noindent {\tovian \fontsize{20pt}{10pt} \textbf{loelh} }\\
\noindent /l{\textprimstress}oe{\textbeltl}/\\
\noindent lit. two+twelve\\


\noindent History:

\vspace{-0pt}
\hspace{40pt}
\begin{tabular}{ccc}
\textit{4000} & /loe{\textbeltl}ya/&$\rightarrow$ & \textit{9500} & /loe{\textbeltl}y/&$\rightarrow$ & \textit{13001} & /loe{\textbeltl}/& \\
\end{tabular}

\vspace{20pt}\hline

\end{nopagebreak}
\filbreak



\vspace{15pt}
\begin{nopagebreak}
\noindent{\fontsize{20pt}{10pt}\textbf{loelhyalo} } \textit{twenty six/26} (N)\\
\noindent {\tovian \fontsize{20pt}{10pt} \textbf{loelhyalo} }\\
\noindent /loe{\textbeltl}y{\textprimstress}alo/\\
\noindent lit. twenty four+two\\


\noindent History:

\vspace{-0pt}
\hspace{40pt}
\begin{tabular}{ccc}
\textit{4000} & /loe{\textbeltl}yaloe/&$\rightarrow$ & \textit{9500} & /loe{\textbeltl}yalo/& \\
\end{tabular}

\vspace{20pt}\hline

\end{nopagebreak}
\filbreak



\vspace{15pt}
\begin{nopagebreak}
\noindent{\fontsize{20pt}{10pt}\textbf{lhyath} } \textit{twenty/20} (N)\\
\noindent {\tovian \fontsize{20pt}{10pt} \textbf{lhyath} }\\
\noindent /{\textbeltl}y{\textprimstress}a{\texttheta}/\\
\noindent lit. twelve+eight\\


\noindent History:

\vspace{-0pt}
\hspace{40pt}
\begin{tabular}{ccc}
\textit{4000} & /{\textbeltl}yaa{\texttheta}/&$\rightarrow$ & \textit{5500} & /{\textbeltl}ya{\texttheta}/& \\
\end{tabular}

\vspace{20pt}\hline

\end{nopagebreak}
\filbreak



\vspace{15pt}
\begin{nopagebreak}
\noindent{\fontsize{20pt}{10pt}\textbf{lo} } \textit{two/2} (Ns)\\
\noindent {\tovian \fontsize{20pt}{10pt} \textbf{lo} }\\
\noindent /l{\textprimstress}o/\\


\noindent History:

\vspace{-0pt}
\hspace{40pt}
\begin{tabular}{ccc}
\textit{0} & /lo/& \\
\end{tabular}

\vspace{20pt}\hline

\end{nopagebreak}
\filbreak



\vspace{15pt}
\begin{nopagebreak}
\noindent{\fontsize{20pt}{10pt}\textbf{hi-sulan} } \textit{undershine} (N)\\
\noindent {\tovian \fontsize{20pt}{10pt} \textbf{hi-sulan} }\\
\noindent /ħis{\textprimstress}ulan/\\
\noindent lit. GEN+bottom+light\\


\noindent History:

\vspace{-0pt}
\hspace{40pt}
\begin{tabular}{ccc}
\textit{14000} & /hi-sulan/&$\rightarrow$ & \textit{14006} & /ħi-sulan/& \\
\end{tabular}

\vspace{20pt}\hline

\end{nopagebreak}
\filbreak



\vspace{15pt}
\begin{nopagebreak}
\noindent{\fontsize{20pt}{10pt}\textbf{shim} } \textit{unknown/mystery} (N)\\
\noindent {\tovian \fontsize{20pt}{10pt} \textbf{shim} }\\
\noindent /{\textesh}{\textprimstress}im/\\


\noindent History:

\vspace{-0pt}
\hspace{40pt}
\begin{tabular}{ccc}
\textit{0} & /{\textesh}imi/&$\rightarrow$ & \textit{9500} & /{\textesh}im/& \\
\end{tabular}

\vspace{20pt}\hline

\end{nopagebreak}
\filbreak



\vspace{15pt}
\begin{nopagebreak}
\noindent{\fontsize{20pt}{10pt}\textbf{ther} } \textit{up} (N)\\
\noindent {\tovian \fontsize{20pt}{10pt} \textbf{ther} }\\
\noindent /{\texttheta}{\textprimstress}er/\\


\noindent History:

\vspace{-0pt}
\hspace{40pt}
\begin{tabular}{ccc}
\textit{0} & /{\texttheta}era/&$\rightarrow$ & \textit{9500} & /{\texttheta}er/& \\
\end{tabular}

\vspace{20pt}\hline

\end{nopagebreak}
\filbreak



\vspace{15pt}
\begin{nopagebreak}
\noindent{\fontsize{20pt}{10pt}\textbf{n} } \textit{us/we} (Ps)\\
\noindent {\tovian \fontsize{20pt}{10pt} \textbf{n} }\\
\noindent /n/\\


\noindent History:

\vspace{-0pt}
\hspace{40pt}
\begin{tabular}{ccc}
\textit{0} & /ne/&$\rightarrow$ & \textit{12006} & /n/& \\
\end{tabular}

\vspace{20pt}\hline

\end{nopagebreak}
\filbreak



\vspace{15pt}
\begin{nopagebreak}
\noindent{\fontsize{20pt}{10pt}\textbf{var} } \textit{value} (N)\\
\noindent {\tovian \fontsize{20pt}{10pt} \textbf{var} }\\
\noindent /v{\textprimstress}ar/\\


\noindent History:

\vspace{-0pt}
\hspace{40pt}
\begin{tabular}{ccc}
\textit{0} & /velare/&$\rightarrow$ & \textit{3200} & /vlare/&$\rightarrow$ & \textit{9500} & /vlar/&$\rightarrow$ & \textit{12500} & /var/& \\
\end{tabular}

\vspace{20pt}\hline

\end{nopagebreak}
\filbreak



\vspace{15pt}
\begin{nopagebreak}
\noindent{\fontsize{20pt}{10pt}\textbf{klumlor} } \textit{vine} (N)\\
\noindent {\tovian \fontsize{20pt}{10pt} \textbf{klumlor} }\\
\noindent /kl{\textprimstress}umlor/\\
\noindent lit. arm+plant\\


\noindent History:

\vspace{-0pt}
\hspace{40pt}
\begin{tabular}{ccc}
\textit{200} & /kelu{\textschwa}mellore/&$\rightarrow$ & \textit{2202} & /keluw{\textschwa}mellore/&$\rightarrow$ & \textit{2300} & /keluwmellore/&$\rightarrow$ & \textit{3200} & /kluwmllore/&$\rightarrow$ & \textit{6500} & /kluwmlore/&$\rightarrow$ & \textit{7600} & /klumlore/&$\rightarrow$ & \textit{9500} & /klumlor/& \\
\end{tabular}

\vspace{20pt}\hline

\end{nopagebreak}
\filbreak



\vspace{15pt}
\begin{nopagebreak}
\noindent{\fontsize{20pt}{10pt}\textbf{anshish} } \textit{waking dreambringer} (N)\\
\noindent {\tovian \fontsize{20pt}{10pt} \textbf{anshish} }\\
\noindent /{\textprimstress}an{\textesh}i{\textesh}/\\
\noindent lit. dream+COM+alertness\\


\noindent History:

\vspace{-0pt}
\hspace{40pt}
\begin{tabular}{ccc}
\textit{10} & /{\textbeltl}enayi-{\textyogh}i{\texttheta}/&$\rightarrow$ & \textit{1000} & /{\textbeltl}enayi{\textyogh}i{\texttheta}/&$\rightarrow$ & \textit{2100} & /{\textbeltl}enyi{\textyogh}i{\texttheta}/&$\rightarrow$ & \textit{3200} & /{\textbeltl}nyi{\textyogh}{\texttheta}/&$\rightarrow$ & \textit{4502} & /{\textbeltl}nyi{\textyogh}/&$\rightarrow$ & \textit{6000} & /{\textbeltl}nd{\textyogh}i{\textyogh}/&$\rightarrow$ & \textit{7501} & /nd{\textyogh}i{\textyogh}/&$\rightarrow$ & \textit{8750} & /n{\textyogh}i{\textyogh}/&$\rightarrow$ & \textit{8760} & /an{\textyogh}i{\textyogh}/&$\rightarrow$ & \textit{12001} & /an{\textesh}i{\textesh}/& \\
\end{tabular}

\vspace{20pt}\hline

\end{nopagebreak}
\filbreak



\vspace{15pt}
\begin{nopagebreak}
\noindent{\fontsize{20pt}{10pt}\textbf{nolh} } \textit{walking} (N)\\
\noindent {\tovian \fontsize{20pt}{10pt} \textbf{nolh} }\\
\noindent /n{\textprimstress}o{\textbeltl}/\\


\noindent History:

\vspace{-0pt}
\hspace{40pt}
\begin{tabular}{ccc}
\textit{0} & /no{\textbeltl}/& \\
\end{tabular}

\vspace{20pt}\hline

\end{nopagebreak}
\filbreak



\vspace{15pt}
\begin{nopagebreak}
\noindent{\fontsize{20pt}{10pt}\textbf{sor} } \textit{washing} (N)\\
\noindent {\tovian \fontsize{20pt}{10pt} \textbf{sor} }\\
\noindent /s{\textprimstress}or/\\
\noindent lit. water+action\\


\noindent History:

\vspace{-0pt}
\hspace{40pt}
\begin{tabular}{ccc}
\textit{0} & /z{\textsubbridge{t}}ellore/&$\rightarrow$ & \textit{2000} & /zdellore/&$\rightarrow$ & \textit{3200} & /zdllore/&$\rightarrow$ & \textit{4500} & /zllore/&$\rightarrow$ & \textit{6500} & /zlore/&$\rightarrow$ & \textit{9500} & /zlor/&$\rightarrow$ & \textit{12000} & /slor/&$\rightarrow$ & \textit{12500} & /sor/& \\
\end{tabular}

\vspace{20pt}\hline

\end{nopagebreak}
\filbreak



\vspace{15pt}
\begin{nopagebreak}
\noindent{\fontsize{20pt}{10pt}\textbf{sel} } \textit{water} (N)\\
\noindent {\tovian \fontsize{20pt}{10pt} \textbf{sel} }\\
\noindent /s{\textprimstress}el/\\


\noindent History:

\vspace{-0pt}
\hspace{40pt}
\begin{tabular}{ccc}
\textit{0} & /z{\textsubbridge{t}}el/&$\rightarrow$ & \textit{2000} & /zdel/&$\rightarrow$ & \textit{4500} & /zel/&$\rightarrow$ & \textit{12000} & /sel/& \\
\end{tabular}

\vspace{20pt}\hline

\end{nopagebreak}
\filbreak



\vspace{15pt}
\begin{nopagebreak}
\noindent{\fontsize{20pt}{10pt}\textbf{tirlhom} } \textit{week} (N)\\
\noindent {\tovian \fontsize{20pt}{10pt} \textbf{tirlhom} }\\
\noindent /{\textsubbridge{t}}{\textprimstress}ir{\textbeltl}om/\\
\noindent lit. ten+day\\


\noindent History:

\vspace{-0pt}
\hspace{40pt}
\begin{tabular}{ccc}
\textit{4000} & /{\textsubbridge{t}}ir{\textbeltl}onm/&$\rightarrow$ & \textit{7600} & /{\textsubbridge{t}}ir{\textbeltl}om/& \\
\end{tabular}

\vspace{20pt}\hline

\end{nopagebreak}
\filbreak



\vspace{15pt}
\begin{nopagebreak}
\noindent{\fontsize{20pt}{10pt}\textbf{tlem} } \textit{weight} (N)\\
\noindent {\tovian \fontsize{20pt}{10pt} \textbf{tlem} }\\
\noindent /{\texttoptiebar{t\textbeltl}}{\textprimstress}em/\\


\noindent History:

\vspace{-0pt}
\hspace{40pt}
\begin{tabular}{ccc}
\textit{0} & /a{\texttoptiebar{t\textbeltl}}emo/&$\rightarrow$ & \textit{7500} & /{\texttoptiebar{t\textbeltl}}emo/&$\rightarrow$ & \textit{9500} & /{\texttoptiebar{t\textbeltl}}em/& \\
\end{tabular}

\vspace{20pt}\hline

\end{nopagebreak}
\filbreak



\vspace{15pt}
\begin{nopagebreak}
\noindent{\fontsize{20pt}{10pt}\textbf{fil} } \textit{west} (N)\\
\noindent {\tovian \fontsize{20pt}{10pt} \textbf{fil} }\\
\noindent /f{\textprimstress}il/\\


\noindent History:

\vspace{-0pt}
\hspace{40pt}
\begin{tabular}{ccc}
\textit{0} & /fil/& \\
\end{tabular}

\vspace{20pt}\hline

\end{nopagebreak}
\filbreak



\vspace{15pt}
\begin{nopagebreak}
\noindent{\fontsize{20pt}{10pt}\textbf{lhan} } \textit{what} (INT)\\
\noindent {\tovian \fontsize{20pt}{10pt} \textbf{lhan} }\\
\noindent /{\textbeltl}{\textprimstress}an/\\


\noindent History:

\vspace{-0pt}
\hspace{40pt}
\begin{tabular}{ccc}
\textit{0} & /{\textbeltl}an/& \\
\end{tabular}

\vspace{20pt}\hline

\end{nopagebreak}
\filbreak



\vspace{15pt}
\begin{nopagebreak}
\noindent{\fontsize{20pt}{10pt}\textbf{lhan} } \textit{when} (INT)\\
\noindent {\tovian \fontsize{20pt}{10pt} \textbf{lhan} }\\
\noindent /{\textbeltl}{\textprimstress}an/\\


\noindent History:

\vspace{-0pt}
\hspace{40pt}
\begin{tabular}{ccc}
\textit{0} & /{\textbeltl}an/& \\
\end{tabular}

\vspace{20pt}\hline

\end{nopagebreak}
\filbreak



\vspace{15pt}
\begin{nopagebreak}
\noindent{\fontsize{20pt}{10pt}\textbf{lhan} } \textit{where} (INT)\\
\noindent {\tovian \fontsize{20pt}{10pt} \textbf{lhan} }\\
\noindent /{\textbeltl}{\textprimstress}an/\\


\noindent History:

\vspace{-0pt}
\hspace{40pt}
\begin{tabular}{ccc}
\textit{0} & /{\textbeltl}an/& \\
\end{tabular}

\vspace{20pt}\hline

\end{nopagebreak}
\filbreak



\vspace{15pt}
\begin{nopagebreak}
\noindent{\fontsize{20pt}{10pt}\textbf{lhan} } \textit{who} (INT)\\
\noindent {\tovian \fontsize{20pt}{10pt} \textbf{lhan} }\\
\noindent /{\textbeltl}{\textprimstress}an/\\


\noindent History:

\vspace{-0pt}
\hspace{40pt}
\begin{tabular}{ccc}
\textit{0} & /{\textbeltl}an/& \\
\end{tabular}

\vspace{20pt}\hline

\end{nopagebreak}
\filbreak



\vspace{15pt}
\begin{nopagebreak}
\noindent{\fontsize{20pt}{10pt}\textbf{lhan} } \textit{why} (INT)\\
\noindent {\tovian \fontsize{20pt}{10pt} \textbf{lhan} }\\
\noindent /{\textbeltl}{\textprimstress}an/\\


\noindent History:

\vspace{-0pt}
\hspace{40pt}
\begin{tabular}{ccc}
\textit{0} & /{\textbeltl}an/& \\
\end{tabular}

\vspace{20pt}\hline

\end{nopagebreak}
\filbreak



\vspace{15pt}
\begin{nopagebreak}
\noindent{\fontsize{20pt}{10pt}\textbf{shad} } \textit{width} (N)\\
\noindent {\tovian \fontsize{20pt}{10pt} \textbf{shad} }\\
\noindent /{\textesh}{\textprimstress}ad/\\


\noindent History:

\vspace{-0pt}
\hspace{40pt}
\begin{tabular}{ccc}
\textit{0} & /{\textyogh}ila{\textsubbridge{t}}e/&$\rightarrow$ & \textit{2000} & /{\textyogh}ilade/&$\rightarrow$ & \textit{3200} & /{\textyogh}lade/&$\rightarrow$ & \textit{9500} & /{\textyogh}lad/&$\rightarrow$ & \textit{12001} & /{\textesh}lad/&$\rightarrow$ & \textit{12500} & /{\textesh}ad/& \\
\end{tabular}

\vspace{20pt}\hline

\end{nopagebreak}
\filbreak



\vspace{15pt}
\begin{nopagebreak}
\noindent{\fontsize{20pt}{10pt}\textbf{lrith} } \textit{wind/breath} (N)\\
\noindent {\tovian \fontsize{20pt}{10pt} \textbf{lrith} }\\
\noindent /lr{\textprimstress}i{\texttheta}/\\
\noindent lit. movement+COM+air\\


\noindent History:

\vspace{-0pt}
\hspace{40pt}
\begin{tabular}{ccc}
\textit{10} & /kwilyi-rai{\texttheta}/&$\rightarrow$ & \textit{1000} & /kwilyirai{\texttheta}/&$\rightarrow$ & \textit{3200} & /kwlyrai{\texttheta}/&$\rightarrow$ & \textit{5500} & /kwlyri{\texttheta}/&$\rightarrow$ & \textit{7501} & /wlyri{\texttheta}/&$\rightarrow$ & \textit{7600} & /lri{\texttheta}/& \\
\end{tabular}

\vspace{20pt}\hline

\end{nopagebreak}
\filbreak



\vspace{15pt}
\begin{nopagebreak}
\noindent{\fontsize{20pt}{10pt}\textbf{mlhetush} } \textit{Winter} (N)\\
\noindent {\tovian \fontsize{20pt}{10pt} \textbf{mlhetush} }\\
\noindent /m{\textbeltl}{\textprimstress}e{\textsubbridge{t}}u{\textesh}/\\
\noindent lit. present+cold\\


\noindent History:

\vspace{-0pt}
\hspace{40pt}
\begin{tabular}{ccc}
\textit{0} & /num{\textbeltl}ele{\textsubbridge{t}}ku{\textesh}u/&$\rightarrow$ & \textit{3000} & /num{\textbeltl}ele{\textsubbridge{t}}u{\textesh}u/&$\rightarrow$ & \textit{3200} & /nm{\textbeltl}le{\textsubbridge{t}}u{\textesh}u/&$\rightarrow$ & \textit{7501} & /m{\textbeltl}le{\textsubbridge{t}}u{\textesh}u/&$\rightarrow$ & \textit{9500} & /m{\textbeltl}le{\textsubbridge{t}}u{\textesh}/&$\rightarrow$ & \textit{12500} & /m{\textbeltl}e{\textsubbridge{t}}u{\textesh}/& \\
\end{tabular}

\vspace{20pt}\hline

\end{nopagebreak}
\filbreak



\vspace{15pt}
\begin{nopagebreak}
\noindent{\fontsize{20pt}{10pt}\textbf{yeol} } \textit{wisdom/truth} (N)\\
\noindent {\tovian \fontsize{20pt}{10pt} \textbf{yeol} }\\
\noindent /y{\textprimstress}eol/\\


\noindent History:

\vspace{-0pt}
\hspace{40pt}
\begin{tabular}{ccc}
\textit{0} & /arayeol/&$\rightarrow$ & \textit{2100} & /aryeol/&$\rightarrow$ & \textit{7500} & /ryeol/&$\rightarrow$ & \textit{7600} & /yeol/& \\
\end{tabular}

\vspace{20pt}\hline

\end{nopagebreak}
\filbreak



\vspace{15pt}
\begin{nopagebreak}
\noindent{\fontsize{20pt}{10pt}\textbf{thraf} } \textit{wondering/curiosity} (N)\\
\noindent {\tovian \fontsize{20pt}{10pt} \textbf{thraf} }\\
\noindent /{\texttheta}r{\textprimstress}af/\\


\noindent History:

\vspace{-0pt}
\hspace{40pt}
\begin{tabular}{ccc}
\textit{0} & /{\texttheta}irafa/&$\rightarrow$ & \textit{3200} & /{\texttheta}rafa/&$\rightarrow$ & \textit{9500} & /{\texttheta}raf/& \\
\end{tabular}

\vspace{20pt}\hline

\end{nopagebreak}
\filbreak



\vspace{15pt}
\begin{nopagebreak}
\noindent{\fontsize{20pt}{10pt}\textbf{glar} } \textit{woods} (N)\\
\noindent {\tovian \fontsize{20pt}{10pt} \textbf{glar} }\\
\noindent /gl{\textprimstress}ar/\\


\noindent History:

\vspace{-0pt}
\hspace{40pt}
\begin{tabular}{ccc}
\textit{0} & /glar/& \\
\end{tabular}

\vspace{20pt}\hline

\end{nopagebreak}
\filbreak



\vspace{15pt}
\begin{nopagebreak}
\noindent{\fontsize{20pt}{10pt}\textbf{for} } \textit{writing/painting} (N)\\
\noindent {\tovian \fontsize{20pt}{10pt} \textbf{for} }\\
\noindent /f{\textprimstress}or/\\
\noindent lit. showing+action\\


\noindent History:

\vspace{-0pt}
\hspace{40pt}
\begin{tabular}{ccc}
\textit{10} & /fe{\textsubbridge{t}}elore/&$\rightarrow$ & \textit{1000} & /f{\textsubbridge{t}}elore/&$\rightarrow$ & \textit{3000} & /felore/&$\rightarrow$ & \textit{3200} & /flore/&$\rightarrow$ & \textit{9500} & /flor/&$\rightarrow$ & \textit{12500} & /for/& \\
\end{tabular}

\vspace{20pt}\hline

\end{nopagebreak}
\filbreak



\vspace{15pt}
\begin{nopagebreak}
\noindent{\fontsize{20pt}{10pt}\textbf{nalom} } \textit{year} (N)\\
\noindent {\tovian \fontsize{20pt}{10pt} \textbf{nalom} }\\
\noindent /n{\textprimstress}alom/\\
\noindent lit. small+cycle\\


\noindent History:

\vspace{-0pt}
\hspace{40pt}
\begin{tabular}{ccc}
\textit{10} & /ninaelakorim/&$\rightarrow$ & \textit{2000} & /ninaelagorim/&$\rightarrow$ & \textit{3200} & /nnaelgorm/&$\rightarrow$ & \textit{4501} & /nnaelorm/&$\rightarrow$ & \textit{6500} & /naelorm/&$\rightarrow$ & \textit{7600} & /naelom/&$\rightarrow$ & \textit{10000} & /nalom/& \\
\end{tabular}

\vspace{20pt}\hline

\end{nopagebreak}
\filbreak



\vspace{15pt}
\begin{nopagebreak}
\noindent{\fontsize{20pt}{10pt}\textbf{wa} } \textit{you} (P)\\
\noindent {\tovian \fontsize{20pt}{10pt} \textbf{wa} }\\
\noindent /w{\textprimstress}a/\\


\noindent History:

\vspace{-0pt}
\hspace{40pt}
\begin{tabular}{ccc}
\textit{0} & /wa/& \\
\end{tabular}

\vspace{20pt}\hline

\end{nopagebreak}
\filbreak



\vspace{15pt}
\begin{nopagebreak}
\noindent{\fontsize{20pt}{10pt}\textbf{w} } \textit{you (plural)} (Ps)\\
\noindent {\tovian \fontsize{20pt}{10pt} \textbf{w} }\\
\noindent /w/\\


\noindent History:

\vspace{-0pt}
\hspace{40pt}
\begin{tabular}{ccc}
\textit{0} & /we/&$\rightarrow$ & \textit{12006} & /w/& \\
\end{tabular}

\vspace{20pt}\hline

\end{nopagebreak}
\filbreak



\vspace{15pt}
\begin{nopagebreak}
\noindent{\fontsize{20pt}{10pt}\textbf{don} } \textit{young} (N)\\
\noindent {\tovian \fontsize{20pt}{10pt} \textbf{don} }\\
\noindent /d{\textprimstress}on/\\
\noindent lit. not+old\\


\noindent History:

\vspace{-0pt}
\hspace{40pt}
\begin{tabular}{ccc}
\textit{0} & /donan/&$\rightarrow$ & \textit{3200} & /donn/&$\rightarrow$ & \textit{6500} & /don/& \\
\end{tabular}

\vspace{20pt}\hline

\end{nopagebreak}
\filbreak



\onecolumn


\end{document}
